%%%%%%%%%%%%%%%%%%%%%%%%%%%%%%%%%%%%%%%%%%%%%%%%%%%%%%%%%%%%%%%%%
\chapter{Unit 6: Functions}
\label{chap:unit6}
%%%%%%%%%%%%%%%%%%%%%%%%%%%%%%%%%%%%%%%%%%%%%%%%%%%%%%%%%%%%%%%%%

%%===============================================================
\section{Functions}
\label{sec:functions}
%%===============================================================

\paragraph{Introduction}



Formally, a function f from a set X to a set Y is defined by a set G of

ordered pairs (x, y) such that $x \in X$, $y \in Y$, and every element of X is

the first component of exactly one ordered pair in G. In other words,

for every x in X, there is exactly one element y such that the ordered

pair (x, y) belongs to the set of pairs defining the function f. The set

G is called the graph of the function. Formally speaking, it may be

identified with the function, but this hides the usual interpretation of

a function as a process. Therefore, in common usage, the function is

generally distinguished from its graph.}} Whenever one quantity uniquely

determines the value of another quantity, we have a function. That is,

the set $X$ uniquely determines the set $Y$. You can think of a

\textit{function} as a kind of machine. You feed the machine raw materials, and

the machine changes the raw materials into a finished product.



\paragraph{Functions in everyday life}



Think about dropping a ball from a bridge. At each moment in time, the

ball is a height above the ground. The height of the ball is a function

of time. It was the job of physicists to come up with a formula for this

function. This type of function is called \textit{real-valued} since the

"finished product" is a number (or, more specifically, a real number).



Think about a wind storm. At different places, the wind can be blowing

in different directions with different intensities. The direction and

intensity of the wind can be thought of as a function of position. This

is a function of two real variables (a location is described by two

values - an $x$ and a $y$) which results in a vector (which is something

that can be used to hold a direction and an intensity). These functions

are studied in multivariable calculus (which is usually studied after a

one year college level calculus course). This a vector-valued function

of two real variables.



We will be looking at real-valued functions until studying multivariable

calculus. Think of a real-valued function as an \textit{input-output machine};

you give the function an input, and it gives you an output which is a

number (more specifically, a real number). For example, the squaring

function takes the input 4 and gives the output value 16. The same

squaring function takes the input -1 and gives the output value 1.



% [Image: This is an intuitive way to understand functions: a machine that makes

the input $x$ go through a transformation $f$ into the output

$f(x)$]



\paragraph{Notation}



Functions are used so much that there is a special notation for them.

The notation is somewhat ambiguous, so familiarity with it is important

in order to understand the intention of an equation or formula.



Though there are no strict rules for naming a function, it is standard

practice to use the letters $f$ , $g$ , and $h$ to denote functions, and

the variable $x$ to denote an independent variable. $y$ is used for both

dependent and independent variables.



When discussing or working with a function $f$ , it's important to know

not only the function, but also its independent variable $x$ . Thus,

when referring to a function $f$, you usually do not write $f$, but

instead $f(x)$ . The function is now referred to as "$f$ of $x$". The

name of the function is adjacent to the independent variable (in

parentheses). This is useful for indicating the value of the function at

a particular value of the independent variable. For instance, if



$$f(x)=7x+1$$ , and if we want to use the value of $f$ for $x$ equal to

$2$ , then we would substitute 2 for $x$ on both sides of the definition

above and write



$$f(2)=7(2)+1=14+1=15$$



This notation is more informative than leaving off the independent

variable and writing simply '$f$' , but can be ambiguous since the

parentheses next to $f$ can be misinterpreted as multiplication, $2f$.

To make sure nobody is too confused, follow this procedure:



1.  Define the function $f$ by equating it to some expression.

2.  Give a sentence like the following: "At $x=c$, the function $f$

    is\..."

3.  Evaluate the function.



\paragraph{Description}



There are many ways which people describe functions. In the examples

above, a verbal descriptions is given (the height of the ball above the

earth as a function of time). Here is a list of ways to describe

functions. The top three listed approaches to describing functions are

the most popular.



1.  A function is given a name (such as $f$) and a formula for the

    function is also given. For example, $f(x)=3x+2$ describes a

    function. We refer to the input as the \textbf{argument} of the function

    (or the \textbf{independent variable}), and to the output as the

    \textbf{value} of the function at the given argument.

2.  A function is described using an equation and two variables. One

    variable is for the input of the function and one is for the output

    of the function. The variable for the input is called the

    \textbf{independent variable}. The variable for the output is called the

    \textbf{dependent variable}. For example, $y=3x+2$ describes a function.

    The dependent variable appears by itself on the left hand side of

    equal sign.

3.  A verbal description of the function.



When a function is given a name (like in number 1 above), the name of

the function is usually a single letter of the alphabet (such as $f$ or

$g$). Some functions whose names are multiple letters (like the sine

function $y=\sin(x)$



\paragraph{Plugging a value into a function}



If we write $f(x)=3x+2$ , then we know that



\begin{itemize}
    \item The function $f$ is a function of $x$ .
    \item To evaluate the function at a certain number, replace the $x$ with
\end{itemize}
    that number.

\begin{itemize}
    \item Replacing $x$ with that number in the right side of the function
\end{itemize}
    will produce the function's output for that certain input.

\begin{itemize}
    \item In English, the definition of $f$ is interpreted, "Given a number,
\end{itemize}
    $f$ will return \textit{two more than the triple of that number}."



How would we know the value of the function $f$ at 3? We would have the

following three thoughts:



1.  $f(3)=3(3)+2$

2.  $3(3)+2=9+2$

3.  $9+2=11$



and we would write



$f(3)=3(3)+2=9+2=11$.



The value of $f$ at 3 is 11.



Note that $f(3)$ means the value of the dependent variable when $x$

takes on the value of 3. So we see that the number \textit{11} is the output of

the function when we give the number \textit{3} as the input. People often

summarize the work above by writing "the value of $f$ at three is

eleven", or simply "$f$ of three equals eleven".



\paragraph{Basic concepts of functions}



The formal definition of a function states that a function is actually a

\textit{mapping} that associates the elements of one set called the \textbf{domain}

of the function, $A$, with the elements of another set called the

\textit{range} of the function, $B$. For each value we select from the domain

of the function, there exists \textit{exactly one} corresponding element in the

range of the function. The definition of the function tells us which

element in the range corresponds to the element we picked from the

domain. An example of how is given below.



\paragraph{Example}



Let function $f(x)=5x+\frac{5x}{2}$ for all $x$. For what value of $x$

gives $f(x)=225\frac{1}{2}$ In mathematics, it is important to notice

any repetition. If something seems to repeat, keep a note of that in the

back of your mind somewhere.



Here, notice that $f(x)=5x+\frac{5x}{2}$ and $f(x)=225\frac{1}{2}$.

Because $f(x)$ is equal to two different things, it must be the case

that the other side of the equal sign to $f(x)$ is equal to the other.

This property is known as the transitive property and could thus make

the following equation below true:$$225\frac{1}{2}=5x+\frac{5x}{2}$$



Next, simplify --- make your life easier rather than harder. In this

instance, since $5x+\frac{5x}{2}$ has $x$ as a like-term, and the two

terms are fractions added to the other, put it over a common denominator

and simplify. Similar, since $225\frac{1}{2}$ is a mixed fraction,

$225\frac{1}{2}=225+\frac{1}{2}$.$$225\frac{1}{2}=5x+\frac{5x}{2}$$$$\Rightarrow\frac{225}{1}+\frac{1}{2}=\frac{5x}{1}+\frac{5x}{2}$$$$\Rightarrow\frac{450}{2}+\frac{1}{2}=\frac{10x}{2}+\frac{5x}{2}$$$$\Rightarrow\frac{451}{2}=\frac{15x}{2}$$



Multiply both sides by the reciprocal of the coefficient of $x$,

$\frac{2}{15}$ from both sides by multiplying by it.$$\frac{2}{15}\cdot\frac{451}{2}=x$$$$\frac{451}{15}=x$$ \textit{or} $x=\frac{451}{15}$.



The value of $x$ that makes $f(x)=225\frac{1}{2}$ is $x=\frac{451}{15}$.



Classically, the element picked from the domain is pictured as something

that is fed into the function and the corresponding element in the range

is pictured as the output. Since we "pick" the element in the domain

whose corresponding element in the range we want to find, we have

control over what element we pick and hence this element is also known

as the "independent variable". The element mapped in the range is

beyond our control and is "mapped to" by the function. This element is

hence also known as the "dependent variable", for it depends on which

independent variable we pick. Since the elementary idea of functions is

better understood from the classical viewpoint, we shall use it

hereafter. However, it is still important to remember the correct

definition of functions at all times.



\paragraph{Basic types of transformation}



To make it simple, for the function $f(x)$, all of the possible $x$

values constitute the domain, and all of the values $f(x)$ ($y$ on the

x-y plane) constitute the range. To put it in more formal terms, a

function $f$ is a mapping of some element $a\in A$, called the domain,

to exactly one element $b\in B$, called the range, such that $f:A\to B$.

The image below should help explain the modern definition of a function:



% [Image: $A$ is the domain of the function while $B$ is the range. This

transformation from set $A$ to $B$ is an example of one-to-one

function.]



1.  A function is considered \textbf{one-to-one} if an element $a\in A$ from

    domain $A$ of function $f$, leads to exactly one element $b\in B$

    from range $B$ of the function. By definition, since only one

    element $b$ is mapped by function $f$ from some element $a$,

    $f:A\to B$ implies that there exists only one element $b$ from the

    mapping. Therefore, there exists a one-to-one function because it

    complies with the definition of a function. This definition is

    similar to \textbf{\textit{Figure 1}}.

2.  A function is considered \textbf{many-to-one} if some elements

    $a\_{1},a\_{2},\ldots a\_{n}\in A$ from domain $A$ of function $f$ maps

    to exactly one element $b\in B$ from range $B$ of the function.

    Since some elements $\left( a\_{1},a\_{2},\ldots a\_{n} \right)\in A$

    must map onto exactly one element $b\in B$, $f:A\to B$ must be

    compliant with the definition of a function. Therefore, there exists

    a many-to-one function.

3.  A function is considered \textbf{one-to-many} if exactly one element

    $a\in A$ from domain $A$ of function $f$ maps to some elements

    $\left( b\_{1},b\_{2},\ldots b\_{n} \right)\in B$ from range $B$ of the

    function. If some element $a$ maps onto many distinct elements $b$,

    then $f:A\to B$ is non-functional since there exists many distinct

    elements $b$. Given many-to-one is non-compliant to the definition

    of a function, there exists no function that is one-to-many.



The modern definition describes a function sufficiently such that it

helps us determine whether some new type of "function" is indeed one.

Now that each case is defined above, you can now prove whether functions

are one of these function cases. Here is an example problem:



\paragraph{One-to-One Example}



\textbf{Given}: $f(x)=ax+b$, where $a$ and $b$ are constant and $a\ne0$.\textbf{Prove}: function $f$ is one-to-one.}} Notice that the only changing

element in the function $f$ is $x$. To prove a function is one-to-one by

applying the definition may be impossible because although two random

elements of domain set $A$ may not be many-to-one, there may be some

elements $\left( a\_{1},a\_{2},\ldots a\_{n} \right)\in A$ that may make

the function many-to-one. To account for this possibility, we must prove

that it is impossible to have some result like that.



Assume there exists two distinct elements $x\_{1}\ne x\_2$ that will

result in $f\left( x\_{1} \right) =f\left( x\_{2} \right)$. This would

make the function many-to-one. In consequence,$ax\_{1}+b=ax\_{2}+b$



Subtract $b$ from both sides of the equation.



$$\Rightarrow ax\_{1}=ax\_{2}$$



Subtract $ax\_{2}$ from both sides of the equation.



$$\Rightarrow ax\_{1}-ax\_{2}=0$$



Factor $a$ from both terms to the left of the equation.



$$\Rightarrow a\left( x\_{1}-x\_{2} \right) =0$$



Multiply $\frac{1}{a}$ to both sides of the equation.



$$\Rightarrow \left( x\_{1}-x\_{2} \right) =0$$



Add $x\_{2}$ to both sides of the equation.



$$\Rightarrow x\_{1}=x\_{2}$$



Notice that $x\_{1}=x\_{2}$. However, this is impossible because $x\_{1}$

and $x\_{2}$ are distinct. Ergo,

$f\left( x\_{1} \right)\ne f\left( x\_{2} \right)$. No two distinct inputs

can give the same output. Therefore, the function must be one-to-one.



\paragraph{Basic concepts}



There are a few more important ideas to learn from the modern definition of the function, and

it comes from knowing the difference between domain, range, and

codomain. We already discussed some of these, yet knowing about sets

adds a new definition for each of the following ideas. Therefore, let us

discuss these based on these new ideas. Let $A$ and $B$ be a set. If we

were to combine these two sets, it would be defined as the *cartesian

cross product* $A\times B$. The subset of this product is the function.

The below definitions are a little confusing. The best way to interpret

these is to draw an image. To the right of these definitions is the

image that corresponds to it.



![The domain is all the elements in set $A$ that can be mapped to the elements in set $B$. The range is those elements in set $B$ which are mapped with the domain. The codomain is all the elements in set $B$.] (https://upload.wikimedia.org/wikipedia/commons/thumb/0/0d/Function_Definition.svg/300px-Function_Definition.svg.png)



\paragraph{Definition of domain, range, and codomain of a function}



The \textbf{domain} is defined to be a set $A$ with all elements $a\in A$

mapping to at least one unique $b\in B$ .



The set of elements in $B$ is the \textbf{range} of the function mapping

$f:A\to B$ in the cartesian cross product, whereby the set of all

elements $a\in A$ maps to some element $b\in B$ .



The \textbf{codomain} is the set $B$, where it is not necessarily the case

that all elements $b\in B$ was mapped by some $a\in A$.



% [Image: The domain of the function is the interval from -1 to 1]  % [Image: The range of the function is the interval from 0 to 1]



Note that the codomain is not as important as the other two concepts.



Take $f(x)=\sqrt{1-x^2}$ for example:



Because of the square root, the content in the square root should be strictly not smaller than 0.



> $1-x^2\geq0$



> $\Leftrightarrow -1\leq x\leq 1$



Thus the domain is



> $x\in[-1,1]$



Correspondingly, the range will be



> $f(x)\in[0,1]$



\paragraph{Other types of transformation}



There is one more final aspect that needs to be defined. We already have

a good idea of what makes a mapping a function (e.g. it cannot be

one-to-many). However, three more definitions that you will often hear

will be necessary to talk about: \textit{injective}, \textit{surjective}, \textit{bijective}.



% [Image: The function mapping on the left is an example of an injective

function because it is one-to-one. However, it is not surjective because

the range and the codomain are not the

same.]



\begin{itemize}
    \item A function is \textbf{injective} if it is one-to-one.
    \item A function is \textbf{surjective} if it is "onto." That is, the mapping
\end{itemize}
    $f:A\to B$ has $B$ as the \textit{range} of the function, where the

    codomain and the range of the function are the same.

\begin{itemize}
    \item A function is \textbf{bijective} if it is both surjective and injective.
\end{itemize}

Again, the above definitions are often very confusing. Again, another

image is provided to the right of the bullet points. Along with that,

another example is also provided. Let us analyze the function

$g(x)=ax^{2}$.



\textbf{Given}: $g(x)=ax^{2}$, where $a$ is constant and $a\ne0$.\textbf{Prove}: function $g$ is non-surjective and non-injective.



Notice that the only changing element in the function $g$ is $x$. Let us check to see the conditions of this function.



\textbf{Is it injective}? Let $g(x)=y$, and solve for $x$. First, divide by

$a$.



$$y=ax^{2}$$



$$\Rightarrow \frac{y}{a}=x^{2}$$



Then take the square root of $x^2$. By definition,

$\sqrt{x^{2}}=\pm\sqrt{x^{2}}=\pm x$, so



$$\frac{y}{a}=x^{2}$$



$$\Rightarrow\pm\sqrt{\frac{y}{a}}=\pm x$$



Notice, however, what we learned from the above manipulation. As a

result of solving for $x$, we found that there are two solutions for

$x$. However, this resulted in two different values from $\frac{y}{a}$.

This means that for some individual $x$ that gives $y$, there are two

different inputs that result in the same value. Because

$f(x\_{1})=f(x\_{2})$ when $x\_{1}\ne x\_{2}$, this is by definition

non-injective.



\textbf{Is it surjective}? As a natural consequence of what we learned about

inputs, let us determine the range of the function. After all, the only

way to determine if something is surjective is to see if the range

applies to all real numbers. A good way to determine this is by finding

a pattern involving our domains. Let us say we input a negative number

for the domain: $f(-x)=a(-x)^{2}=a\cdot(-x)\cdot(-x)=ax^2$. This

seemingly trivial exercise tells us that negative numbers give us

non-negative numbers for our range. This is huge information! For

$x\in\mathbb{R}$, the function always results $y>0$ for our range. For

the set $B$, we have elements in that set that have no mappings from the

set $A$. As such, $y\le0$ is the codomain of set $B$. Therefore, this

function is non-surjective! 



\paragraph{Tests for equations}



\paragraph{The vertical line test}



The vertical line test is a systematic test to find out if an equation

involving $x$ and $y$ can serve as a function (with $x$ the independent

variable and $y$ the dependent variable). Simply graph the equation and

draw a vertical line through each point of the $x$-axis. If any vertical

line ever touches the graph at more than one point, then the equation is

not a function; if the line always touches at most one point of the

graph, then the equation is a function.



The circle, on the right, is not a function because the vertical line

intercepts two points on the graph when $x=1$.



% [Image: This is an example of an expression which

fails the vertical line

test.]



\paragraph{The horizontal line and the algebraic 1-1 test}



Similarly, the horizontal line test, though does not test if an equation

is a function, tests if a function is injective (one-to-one). If any

horizontal line ever touches the graph at more than one point, then the

function is not one-to-one; if the line always touches at most one point

on the graph, then the function is one-to-one.



The algebraic 1-1 test is the non-geometric way to see if a function is

one-to-one. The basic concept is that:



Assume there is a function $f$. If:



> $f(a)=f(b)$, and $a=b$, then



function $f$ is one-to-one.



Here is an example: prove that $f(x)=\frac{1-2x}{1+x}$ is injective.



Since the notation is the notation for a function, the equation is a

function. So we only need to prove that it is injective. Let $a$ and $b$

be the inputs of the function and that $f(a)=f(b)$. Thus,



$$\frac{1-2a}{1+a}=\frac{1-2b}{1+b}$$



$$\Leftrightarrow(1+b)(1-2a)=(1+a)(1-2b)$$



$$\Leftrightarrow1-2a+b-2ab=1-2b+a-2ab$$$\Leftrightarrow1-2a+b=1-2b+a$



$$\Leftrightarrow1-2a+3b=1+a$$



$$\Leftrightarrow1+3b=1+3a$$



$$\Leftrightarrow a=b$$



So, the result is $a=b$, proving that the function is injective.



Another example is proving that $g(x)=x^2$ is not injective.



Using the same method, one can find that $a=\pm b$, which is not $a=b$.

So, the function is not injective.



\paragraph{Remarks}



The following arise as a direct consequence of the definition of

functions:



1.  By definition, for each "input" a function returns only one

    "output", corresponding to that input. While the same output may

    correspond to more than one input, one input cannot correspond to

    more than one output. This is expressed graphically as the *vertical

    line test*: a line drawn parallel to the axis of the dependent

    variable (normally vertical) will intersect the graph of a function

    only once. However, a line drawn parallel to the axis of the

    independent variable (normally horizontal) may intersect the graph

    of a function as many times as it likes. Equivalently, this has an

    algebraic (or formula-based) interpretation. We can always say if

    $a=b$ , then $f(a)=f(b)$ , but if we only know that $f(a)=f(b)$ then

    we can't be sure that $a=b$ .

2.  Each function has a set of values, the function's \textit{domain}, which

    it can accept as input. Perhaps this set is all positive real

    numbers; perhaps it is the set {pork, mutton, beef}. This set must

    be implicitly/explicitly defined in the definition of the function.

    You cannot feed the function an element that isn't in the domain,

    as the function is not defined for that input element.

3.  Each function has a set of values, the function's \textit{range}, which it

    can output. This may be the set of real numbers. It may be the set

    of positive integers or even the set {0,1}. This set, too, must be

    implicitly/explicitly defined in the definition of the function.



Functions are an important foundation of mathematics. This modern

interpretation is a result of the hard work of the mathematicians of

history. It was not until recently that the definition of the relation

was introduced by René Descartes in \textit{Geometry} (1637). Nearly a century

later, the standard notation ($f(x)=y$) was first introduced by Leonhard

Euler in \textit{Introductio in Analysin Infinitorum} and *Institutiones

Calculi Differentialis*. The term function was also a new innovation

during Euler's time as well, being introduced Gottfried Wilhelm Leibniz

in a 1673 letter "to describe a quantity related to points of a curve,

such as a coordinate or curve's slope." Finally, the modern definition

of the function being the relation among sets was first introduced in

1908 by Godfrey Harold Hardy where there is a relation between two

variables $x$ and $y$ such that "to some values of $x$ at any rate

correspond values of $y$."



\paragraph{Important functions}



\paragraph{Polynomials}



Polynomial functions are the most common and most convenient functions

in the world of calculus. Their frequent appearances and their

applications on computer calculations have made them important. A

polynomial in a single indeterminate x can always be written (or

rewritten) in the form:



$$f(x)=a\_n x^n+a\_{n-1}x^{n-1}+a\_{n-2}x^{n-2}+\cdot\cdot\cdot+a\_2x^2+a\_1x+a\_0$$



To be more concise, it can also be written in the summation form:



$$\sum\_{k=0}^n a\_k x^k$$



\paragraph{Constant}



When $n=0$, the polynomial can be rewritten into the following function:



> $f(x)=C$, where $C$ is a constant.



The graph of this function is a horizontal line passing the point

$(0,C)$.



\paragraph{Linear}



When $n=1$, the polynomial can be rewritten into



> $f(x)=mx+b$, where $m\text{ and }b$ are constants.



The graph of this function is a straight line passing the point $(0,b)$

and $(-\frac{b}{m},0)$, and the slope of this function is $m$. 



% [Image: Two linear functions are shown on the image. One is

$f(x)=\frac{1}{2}x+2$ and the other is

$g(x)=-2x+5$]



\paragraph{Quadratic}



When $n=2$, the polynomial can be rewritten into



> $f(x)=ax^2+bx+c$, where $a,b,c$ are constants.



The graph of this function is a parabola, like the trajectory of a

basketball thrown into the hoop.



If there are questions about the quadratic formula and other basic

polynomial concepts, please review the respective chapters in Algebra.



% [Image: This is

the graph of a quadratic

function.]



\paragraph{Trigonometric}



Trigonometric functions are also very important because it can connect

algebra and geometry. Trigonometric functions are explained in detail

here due to its importance and difficulty.



\paragraph{Exponential and Logarithmic}



Exponential and logarithmic functions have a unique identity when

calculating the derivatives, so this is a great time to review those

functions. The exponential function is defined as:



$$f(x)=a^x$$, where $a$ is a constant. while the logarithmic function is

defined as:



$$f(x)=\log\_b x$$, where $b$ is a constant. A special number will be

frequently seen in those functions: the Euler's constant, also known as

the base of the natural logarithm. Notated as $e$, it is defined as

$e=\sum\_{n=0}^\infty \frac{1}{n!}$.



% [Image: The curve on the left is an

exponential function while the curve on the right is a logarithmic

one\|256x256px]



\paragraph{Signum}



The Signum (sign) function is simply defined as



$$\text{sgn}(x)= \begin{cases} -1 & \text{if} & x0 \end{cases}$$



\paragraph{Properties of functions}



Sometimes, a lot of function manipulations cannot be achieved without

some help from basic properties of functions.



\paragraph{Domain and range}



This concept is discussed above.



\paragraph{Growth}



Although it seems obvious to spot a function increasing or decreasing,

without the help of graphing software, we need a mathematical way to

spot whether the function is increasing or decreasing, or else our human

minds cannot immediately comprehend the huge amount of information.



Assume a function $f(x)$ with inputs $x\_1,x\_2$, and that $x\_1\in[a,b]$,

$x\_2\in[a,b]$, and $x\_2>x\_1$ at all times.



> If for all $x\_1$ and $x\_2$, $f(x\_2)-f(x\_1)>0$, then

>

> $f(x)$ is increasing in $[a,b]$

>

> If for all $x\_1$ and $x\_2$, $f(x\_2)-f(x\_1)

> $f(x)$ is decreasing in $[a,b]$



\textbf{Example:} In which intervals is $f(x)=\frac{1}{1-x^2}$ increasing?



Firstly, the domain is important. Because the denominator cannot be 0,

the domain for this function is



$x\in(-\infty,-1)\cup(-1,1)\cup(1,\infty)$



In $(-\infty,-1)$, the growth of the function is:



> Let $x\_1,x\_2\in[-\infty,-1]$ and $x\_2>x\_1$

>

> Thus,

>

> > $f(x\_2)-f(x\_1)=\frac{1}{1-x\_2^2}-\frac{1}{1-x\_1^2}=\frac{x\_2^2-x\_1^2}{(1-x\_2^2)(1-x\_1^2)}$

>

> $\because$ both $x\_1,x\_2\in[-\infty,-1]$

>

> $\therefore$ $(1-x\_2^2)(1-x\_1^2)>0$

>

> $\because$ $x\_2>x\_1$ and $x\_1

> $\therefore$ $x\_2^2-x\_1^2

> > So, $f(x\_2)-f(x\_1)=\frac{x\_2^2-x\_1^2}{(1-x\_2^2)(1-x\_1^2)}

> $f(x)$ is decreasing in $(-\infty,-1)$



In $(-1,1)$



> Let $x\_1,x\_2\in[-1,1]$ and $x\_2>x\_1$

>

> Thus,

>

> > $f(x\_2)-f(x\_1)=\frac{1}{1-x\_2^2}-\frac{1}{1-x\_1^2}=\frac{x\_2^2-x\_1^2}{(1-x\_2^2)(1-x\_1^2)}$

>

> $\because$ both $x\_1,x\_2\in[-\infty,-1]$

>

> $\thereforeMATH29(1-x\_2^2)(1-x\_1^2)>0$

>

> $\because x\_2>x\_1\text{ and }0

> $\therefore x\_2^2-x\_1^2>0$

>

> > So, $f(x\_2)-f(x\_1)=\frac{x\_2^2-x\_1^2}{(1-x\_2^2)(1-x\_1^2)}>0$

>

> $f(x)$ is increasing in $(1,\infty)$.



Therefore, the intervals in which the function is increasing are

$(0,1)\And(1,\infty)$.



$\blacksquare$



After learning derivatives, there will be more ways to determine the

growth of a function.



\paragraph{Parity}



The properties odd and even are associated with symmetry. While even

functions have a symmetry about the $y$-axis, odd functions are

symmetric about the origin. In mathematical terms:



> A function is even when $f(-x)=f(x)$

>

> A function is odd when $f(-x)=-f(x)$



\textbf{Example:} Prove that $f(x)=\frac{1}{1-x^2}$ is an even function.



$\because f(-x)=\frac{1}{1-(-x)^2}=\frac{1}{1-x^2}=f(x)$



$\therefore f(x)$ is an even function



\paragraph{Manipulating functions}



\paragraph{Addition, Subtraction, Multiplication and Division of functions}



For two real-valued functions, we can add the functions, multiply the

functions, raised to a power, etc. For example:



\begin{itemize}
    \item If we add the functions $y=3x+2$ and $y=x^2$ , we obtain
\end{itemize}
    $y=x^2+3x+2$ .



<!-- -->



\begin{itemize}
    \item If we subtract $y=3x+2$ from $y=x^2$ , we obtain $y=x^2-(3x+2)$ . We
\end{itemize}
    can also write this as $y=x^2-3x-2$ .



<!-- -->



\begin{itemize}
    \item If we multiply the function $y=3x+2$ and the function $y=x^2$ , we
\end{itemize}
    obtain $y=(3x+2)x^2$ . We can also write this as $y=3x^3+2x^2$ .



<!-- -->



\begin{itemize}
    \item If we divide the function $y=3x+2$ by the function $y=x^2$ , we
\end{itemize}
    obtain $y=\frac{3x+2}{x^2}$ .



If a math problem wants you to add two functions $f$ and $g$ , there are

two ways that the problem will likely be worded:



1.  If you are told that $f(x)=3x+2$ , that $g(x)=x^2$ , that

    $h(x)=f(x)+g(x)$ and asked about $h$ , then you are being asked to

    add two functions. Your answer would be $h(x)=x^2+3x+2$ .

2.  If you are told that $f(x)=3x+2$ , that $g(x)=x^2$ and you are asked

    about $f+g$ , then you are being asked to add two functions. The

    addition of $f$ and $g$ is called $f+g$ . Your answer would be

    $(f+g)(x)=x^2+3x+2$ .



Similar statements can be made for subtraction, multiplication and

division.



Let $f(x)=3x+2$ and: $g(x)=x^2$ . Let's add, subtract, multiply and

divide.



$$\begin{aligned}(f+g)(x) &= f(x)+g(x) \\\\ &= (3x+2)+(x^2) \\\\ &= x^2+3x+2 \end{aligned}$$ 



$$\begin{aligned} (f-g)(x) &= f(x)-g(x) \\\\ &= (3x+2)-(x^2) \\\\ &= -x^2+3x+2 \end{aligned}$$ 



$$\begin{aligned} (f\times g)(x) &= f(x)\times g(x) \\\\ &= (3x+2)\times(x^2) \\\\ &= 3x^3+2x^2 \end{aligned}$$ 



$$\begin{aligned} \left(\frac{f}{g}\right)(x) &= \frac{f(x)}{g(x)} \\\\ &= \frac{3x+2}{x^2} \\\\ &= \frac{3}{x}+\frac{2}{x^2} \end{aligned}$$ 



\paragraph{Composition of functions}



We begin with a fun (and not too complicated) application of composition

of functions before we talk about what composition of functions is.



Example: Dropping a ball If we drop a ball from a bridge which is 20

meters above the ground, then the height of our ball above the earth is

a function of time. The physicists tell us that if we measure time in

seconds and distance in meters, then the formula for height in terms of

time is $h=-4.9t^2+20$ . Suppose we are tracking the ball with a camera

and always want the ball to be in the center of our picture. Suppose we

have $\theta=f(h)$ The angle will depend upon the height of the ball

above the ground and the height above the ground depends upon time. So

the angle will depend upon time. This can be written as

$\theta=f(-4.9t^2+20)$ . We replace $h$ with what it is equal to. This

is the essence of composition.



Composition of functions is another way to combine functions which is

different from addition, subtraction, multiplication or division.



The value of a function $f$ depends upon the value of another variable

$x$ ; however, that variable could be equal to another function $g$ , so

its value depends on the value of a third variable. If this is the case,

then the first variable is a function $h$ of the third variable; this

function ($h$) is called the \textbf{composition} of the other two functions

($f$ and $g$).



Example: Composing two functions Let $f(x)=3x+2$ and: $g(x)=x^2$ . The

composition of $f$ with $g$ is read as either "f composed with g" or

"f of g of x."



Let



$$h(x)=f(g(x))$$ Then



$$\begin{aligned} h(x) &= f(g(x)) \\\\ &= f(x^2) \\\\ &= 3(x^2)+2 \\\\ &= 3x^2+2 \end{aligned}$$ 



Sometimes a math problem asks you compute $(f\circ g)(x)$ when they want

you to compute $f(g(x))$ ,



Here, $h$ is the composition of $f$ and $g$ and we write $h=f\circ g$ .

Note that composition is not commutative:



$f(g(x))=3x^2+2$ , and



$$\begin{aligned} g(f(x)) &= g(3x+2) \\\\ &= (3x+2)^2 \\\\ &= 9x^2+12x+4 \end{aligned}$$



so $f(g(x))\ne g(f(x))$ .



Composition of functions is very common, mainly because functions

themselves are common. For instance, squaring and sine are both

functions:



$$\text{square}(x)=x^2$$



$$\text{sine}(x)=\sin(x)$$



Thus, the expression $\sin^2(x)$ is a composition of functions:



$$\begin{aligned} \sin^2(x) &=\text{square}(\sin(x)) \\\\ &=\text{square}(\text{sine}(x)) \\\\ &= 9x^2+12x+4 \end{aligned}$$



(Note that this is \textit{not} the same as

$\text{sine}(\text{square}(x))=\sin(x^2)$.) Since the function sine

equals $\tfrac{1}{2}$ if $x=\frac{\pi}{6}$ ,



$$\text{square}\left(\sin\left((frac{\pi}{6}\right)\right)=\text{square}\left(\frac{1}{2}\right).$$



Since the function square equals $\frac{1}{4}$ if $x=\frac{\pi}{6}$ ,



$$\sin^2\left(\frac{\pi}{6}\right)=\text{square}\left(\sin\left(\frac{\pi}{6}\right)\right)=\text{square}\left(\frac{1}{2}\right)=\frac{1}{4}.$$



\paragraph{Transformations}



Transformations are a type of function manipulation that are very

common. They consist of multiplying, dividing, adding or subtracting

constants to either the input or the output. Multiplying by a constant

is called \textbf{dilation} and adding a constant is called \textbf{translation}.

Here are a few examples:



$f(2\times x)$ Dilation



$f(x+2)$ Translation



$2\times f(x)$ Dilation



$2+f(x)$ Translation



Translations and dilations can be either horizontal or vertical.

Examples of both vertical and horizontal translations can be seen at

right. The red graphs represent functions in their 'original' state,

the solid blue graphs have been translated (shifted) horizontally, and

the dashed graphs have been translated vertically.



% [Image: Examples of horizontal and vertical

translations]

% [Image: Examples of horizontal and vertical

dilations]



Dilations are demonstrated in a similar fashion. The function



$$f(2\times x)$$ has had its input doubled. One way to think about this

is that now any change in the input will be doubled. If I add one to

$x$, I add two to the input of $f$, so it will now change twice as

quickly. Thus, this is a horizontal dilation by $\frac{1}{2}$ because

the distance to the $y$-axis has been \textbf{halved}. A vertical dilation,

such as



$$2\times f(x)$$ is slightly more straightforward. In this case, you

double the output of the function. The output represents the distance

from the $x$-axis, so in effect, you have made the graph of the function

'taller'. Here are a few basic examples where $a$ is any positive

constant:



$$

\begin{array}{|c|c|c|c|}

\hline

\text{Original graph} & f(x) & \text{Rotation about origin} & -f(-x) \\\\ \hline

\text{Horizontal translation by} \hspace{0.1cm} \textit{a} \hspace{0.1cm} \text{units} \hspace{0.1cm} \textbf{left} & f(x+a) & \text{Horizontal translation by} \hspace{0.1cm} \textit{a} \hspace{0.1cm} \text{units} \hspace{0.1cm} \textbf{right} & f(x-a) \\\\ \hline

\text{Horizontal dilation by a factor of} \hspace{0.1cm} \textit{a} & f(x \times \frac{1}{a}) & \text{Vertical dilation by a factor of} \hspace{0.1cm} \textit{a} & a\times f(x) \\\\ \hline

\text{Vertical translation by} \hspace{0.1cm} \textit{a} \hspace{0.1cm} \text{units} \hspace{0.1cm} \textbf{down} & f(x)-a & \text{Vertical translation by} \hspace{0.1cm} \textit{a} \hspace{0.1cm} \text{units} \hspace{0.1cm} \textbf{up} & f(x)+a \\\\ \hline

\text{Reflection about x-axis} & -f(x) & \text{Reflection about y-axis} & f(-x) \\\\ \hline

\end{array}

$$



You can use this Markdown in your documents or applications that support Markdown formatting.



\paragraph{Inverse functions}



We call $g(x)$ the inverse function of $f(x)$ if, for all $x$ :



$$g(f(x))=f(g(x))=x$$ A function $f(x)$ has an inverse function if and

only if $f(x)$ is one-to-one. For example, the inverse of $f(x)=x+2$ is

$g(x)=x-2$ . The function $f(x)=\sqrt{1-x^2}$ has no inverse because it

is not injective. Similarly, the inverse functions of trigonometric

functions have to undergo transformations and limitations to be

considered as valid functions.



\paragraph{Notation}



The inverse function of $f$ is denoted as $f^{-1}(x)$ . Thus,

$f^{-1}(x)$ is defined as the function that follows this rule



$$f(f^{-1}(x))=f^{-1}(f(x))=x$$ To determine $f^{-1}(x)$ when given a

function $f$ , substitute $f^{-1}(x)$ for $x$ and substitute $x$ for

$f(x)$ . Then solve for $f^{-1}(x)$ , provided that it is also a

function.



\textbf{Example:} Given $f(x)=2x-7$ , find $f^{-1}(x)$ .



Substitute $f^{-1}(x)$ for $x$ and substitute $x$ for $f(x)$ . Then

solve for $f^{-1}(x)$ :



$$f(x)=2x-7$$



$$\Leftrightarrow x=2[f^{-1}(x)]-7$$



$$\Leftrightarrow x+7=2[f^{-1}(x)]$$



$$\Leftrightarrow\frac{x+7}{2}=f^{-1}(x)$$



To check your work, confirm that $f^{-1}(f(x))=x$ :



> $f^{-1}(f(x))=MATH51\frac{(2x-7)+7}{2}=\frac{2x}{2}=x$



If $f$ isn't one-to-one, then, as we said before, it doesn't have an

inverse. Then this method will fail.



\textbf{Example:} Given $f(x)=x^2$ , find $f^{-1}(x)$.



Substitute $f^{-1}(x)$ for $x$ and substitute $x$ for $f(x)$ . Then

solve for $f^{-1}(x)$ :



$$f(x)=x^2$$



$$x=(f^{-1}(x))^2$$



$$f^{-1}(x)=\pm\sqrt{x}$$



Since there are two possibilities for $f^{-1}(x)$ , it's not a

function. Thus $f(x)=x^2$ doesn't have an inverse. Of course, we could

also have found this out from the graph by applying the Horizontal Line

Test. It's useful, though, to have lots of ways to solve a problem,

since in a specific case some of them might be very difficult while

others might be easy. For example, we might only know an algebraic

expression for $f(x)$ but not a graph.



\paragraph{Resources}



\begin{itemize}
    \item \href{https://www.khanacademy.org/math/algebra/x2f8bb11595b61c86:functions/x2f8bb11595b61c86:evaluating-functions/v/what-is-a-function}{What is a
\end{itemize}
    Function?},

    Khan Academy

\begin{itemize}
    \item \href{https://www.khanacademy.org/math/algebra2/functions\_and\_graphs}{Function and their
\end{itemize}
    graphs},

    Khan Academy

\begin{itemize}
    \item \href{https://www.cliffsnotes.com/study-guides/calculus/precalculus/functions/relations-vs-functions}{Relations vs.
\end{itemize}
    Functions},

    Cliff's Notes

\begin{itemize}
    \item \href{https://courses.lumenlearning.com/ivytech-collegealgebra/chapter/determine-whether-a-relation-represents-a-function/}{Determining if a Relation is a
\end{itemize}
    Function},

    Lumen Learning

\begin{itemize}
    \item \href{https://courses.lumenlearning.com/waymakercollegealgebra/chapter/identify-functions-using-graphs/}{Identifying Functions Using
\end{itemize}
    Graphs},

    Lumen Learning



\paragraph{Licensing}



Content obtained and/or adapted from:



\begin{itemize}
    \item \href{https://en.wikibooks.org/wiki/Calculus/Functions}{Functions, Wikibooks:
\end{itemize}
    Calculus} under a

    CC BY-SA license

\subsection{Learning Objectives}

\begin{enumerate}
    \item \textbf{Define and Identify Functions:} Students will be able to define a function and identify whether a given relation between sets constitutes a function. Students should understand the formal definition of a function as a mapping from a domain to a range where each element of the domain maps to exactly one element in the range. They should also be able to use the vertical line test to determine if a graph represents a function.
    \item \textbf{Evaluate and Interpret Functions:} Students will be able to evaluate a function at specific points and interpret the results in context. Students should be proficient in substituting values into functions, calculating outputs, and interpreting the meaning of these outputs in relation to the function's context. For example, given $f(x) = 3x + 2$, they should be able to find $f(4)$ and explain what the result signifies.
    \item \textbf{Classify and Analyze Function Types:} Students will be able to classify functions as one-to-one, many-to-one, or non-functional, and analyze functions to determine if they are injective, surjective, or bijective. Students should understand the difference between injective (one-to-one), surjective (onto), and bijective functions. They should be able to apply these concepts to determine the type of function given an expression or graph. For instance, they should be able to prove that $ f(x) = 2x + 1 $ is injective and determine if $ g(x) = x^2 $ is surjective when mapping from $\mathbb{R}$ to $\mathbb{R}$.
\end{enumerate}

\subsection{Assessment Instrument}

\begin{enumerate}
    \item \textbf{Part 1: Define and Identify Functions}
    \item \textbf{Multiple Choice:} Which of the following relations represents a function?
    \item A. $\{(1, 2), (2, 3), (1, 4)\}$
    \item B. $\{(1, 2), (2, 2), (3, 4)\}$
    \item C. $\{(1, 2), (2, 2), (2, 3)\}$
    \item D. $\{(1, 2), (3, 4), (5, 6), (1, 7)\}$
    \item \textbf{True/False:} The vertical line test states that if a vertical line intersects a graph more than once, then the graph does not represent a function. (True/False)
    \item \textbf{Short Answer:} Define a function formally and explain how the vertical line test helps in determining if a graph represents a function. \textbf{Part 2: Evaluate and Interpret Functions}
    \item \textbf{Calculation:} Given the function $f(x) = 5x - 3$, find $f(7)$.
    \item \textbf{Contextual Interpretation:} For the function $h(x) = \frac{1}{x + 2}$, if $h(3) = \frac{1}{5}$, explain what this result tells us about the value of $x$ and the function's behavior.
    \item \textbf{Short Answer:} Consider the function $g(x) = x^2 + 2x + 1$. If $g(-1) = 0$, what does this value signify in terms of the function’s output? \textbf{Part 3: Classify and Analyze Function Types}
    \item \textbf{Multiple Choice:} Determine the type of the function $f(x) = 3x + 4$:
    \item A. One-to-One
    \item B. Many-to-One
    \item C. Non-Functional
    \item D. Not Enough Information
    \item \textbf{True/False:} The function $k(x) = x^2$ is surjective when mapping from $\mathbb{R}$ to $\mathbb{R}$. (True/False)
    \item \textbf{Short Answer:} Prove that the function $m(x) = -2x + 5$ is injective and analyze whether the function $n(x) = x^3$ is bijective when considering the domain and range as $\mathbb{R}$.
\end{enumerate}

%%===============================================================
\section{Graph of Functions}
\label{sec:graph-of-functions}
%%===============================================================

In mathematics, the \textbf{graph} of a function $f$ is the set of ordered

pairs $(x, y)$, where $f(x) = y.$ In the common case where $x$ and

$f(x)$ are real numbers, these pairs are Cartesian coordinates of points

in two-dimensional space and thus form a subset of this plane.



In the case of functions of two variables, that is functions whose

domain consists of pairs $(x, y),$ the graph usually refers to the set

of ordered triples $(x, y, z)$ where $f(x,y) = z,$ instead of the pairs

$((x, y), z)$ as in the definition above. This set is a subset of

three-dimensional space; for a continuous real-valued function of two

real variables, it is a surface.



A graph of a function is a special case of a relation.



In science, engineering, technology, finance, and other areas, graphs

are tools used for many purposes. In the simplest case one variable is

plotted as a function of another, typically using rectangular axes; see

Plot (graphics) for details.



In the modern foundations of mathematics, and, typically, in set theory,

a function is actually equal to its graph. However, it is often useful

to see functions as mappings, which consist not only of the relation

between input and output, but also which set is the domain, and which

set is the codomain. For example, to say that a function is onto

(surjective) or not the codomain should be taken into account. The graph

of a function on its own doesn't determine the codomain. It is common

to use both terms function and graph of a function since even if

considered the same object, they indicate viewing it from a different

perspective.



% [Image:  Graph of the function $f(x) = x^3 - 9x.$]



\paragraph{Definition}



Given a mapping $f : X \to Y,$ in other words a function $f$ together

with its domain $X$ and codomain $Y,$ the graph of the mapping is the

set $G(f) = \\{(x,f(x)) : x \in X\\},$



which is a subset of $X\times Y$. In the abstract definition of a

function, $G(f)$ is actually equal to $f.$



One can observe that, if, $f : \R^n \to \R^m,$ then the graph $G(f)$ is

a subset of $\R^{n+m}$ (strictly speaking it is $\R^n \times \R^m,$ but

one can embed it with the natural isomorphism).



\paragraph{Examples}



\paragraph{Functions of one variable}



The graph of the function $f : \\{1,2,3\\} \to \\{a,b,c,d\\}$ defined by



$$f(x)=

\begin{cases} a, & \text{if}x=1, \\\\ d, & \text{if }x=2, \\\\ c, & \text{if }x=3, \end{cases}$$ 

is the subset of the set

$\\{1,2,3\\} \times \\{a,b,c,d\\}$ 

$$G(f) = \\{ (1,a), (2,d), (3,c) \\}.$$



From the graph, the domain $\\{1,2,3\\}$ is recovered as the set of first

component of each pair in the graph

$\\{1,2,3\\} = \\{x : \text{there exists } y,\text{ such that }(x,y) \in G(f)\\}$.

Similarly, the range can be recovered as

$\\{a,c,d\\} = \\{y : \text{there exists }x,\text{ such that }(x,y)\in G(f)\\}$.

The codomain $\\{a,b,c,d\\}$, however, cannot be determined from the graph

alone.



The graph of the cubic polynomial on the real line 

$$f(x) = x^3 - 9x$$ 

is

$$\\{ (x, x^3 - 9x) : x \text{ is a real number} \\}.$$



If this set is plotted on a Cartesian plane, the result is a curve (see

figure).



\paragraph{Functions of two variables}



The graph of the trigonometric function $$f(x,y) = \sin(x^2)\cos(y^2)$$ is

$$\\{(x, y, \sin(x^2) \cos(y^2)) : x \text{ and } y \text{ are real numbers} \\}.$$



If this set is plotted on a three dimensional Cartesian coordinate

system, the result is a surface (see figure).



% [Image: Graph of the function

$f(x, y) = \sin\left(x^2\right) \cdot \cos\left(y^2\right).$]



Oftentimes it is helpful to show with the graph, the gradient of the

function and several level curves. The level curves can be mapped on the

function surface or can be projected on the bottom plane. The second

figure shows such a drawing of the graph of the function:

$f(x, y) = -(\cos(x^2) + \cos(y^2))^2.$



% [Image: Plot of the graph of

$f(x, y) = - \left(\cos\left(x^2\right) + \cos\left(y^2\right)\right)^2,$

also showing its gradient projected on the bottom

plane.]



\paragraph{Generalizations}



The graph of a function is contained in a Cartesian product of sets. An

$X$--$Y$ plane is a Cartesian product of two lines, called $X$ and $Y,$

while a cylinder is a cartesian product of a line and a circle, whose

height, radius, and angle assign precise locations of the points. Fibre

bundles are not Cartesian products, but appear to be up close. There is

a corresponding notion of a graph on a fibre bundle called a section.



The graph of a multifunction, say the multifunction

$\mathcal{R} : X \rightrightarrows Y,$ is the set $\operatorname{gr} \mathcal{R} \:= \left \\{ (x, y) \in X \times Y : y \in \mathcal{R}(x) \\} \right .$



\paragraph{Licensing}



Content obtained and/or adapted from:



\begin{itemize}
    \item \href{https://en.wikipedia.org/wiki/Graph\_of\_a\_function}{Graph of a function,
\end{itemize}
    Wikipedia} under

    a CC BY-SA license

\subsection{Learning Objectives}

\begin{enumerate}
    \item \textbf{Identify and Describe Graphs:} Students will identify and describe the graphs of functions with one and two variables, including their representations as subsets of two-dimensional and three-dimensional spaces.
    \item \textbf{Graph Interpretation:} Students will interpret the graph of a function to determine its domain, range, and the relationship between input and output values.
    \item \textbf{Function and Graph Relationship:} Students will explain the relationship between a function and its graph, understanding how the graph represents the function and how to use it to analyze function properties.
\end{enumerate}

\subsection{Assessment Instrument}

\begin{enumerate}
    \item \textbf{Part 1: Identify and Describe Graphs}
    \item \textbf{Multiple Choice:} Which of the following graphs represents a quadratic function?
    \item A. A straight line with a slope of 2.
    \item B. A parabola opening upwards.
    \item C. A circle centered at the origin.
    \item D. A hyperbola with asymptotes.
    \item \textbf{True/False:} The graph of a function $f(x, y) = x^2 + y^2$ represents a plane in three-dimensional space. (True/False)
    \item \textbf{Short Answer:} Describe the graph of the function $f(x) = \sin(x)$ and explain how it represents a periodic function in two-dimensional space. \textbf{Part 2: Graph Interpretation}
    \item \textbf{Calculation:} Given the graph of the function $f(x)$, determine the domain and range if the graph is defined for $x$ values from -2 to 2 and the $y$ values from -3 to 1.
    \item \textbf{Short Answer:} If a function's graph intersects the $x$-axis at $x = 4$ and $x = -1$, what does this tell you about the function's roots?
    \item \textbf{Multiple Choice:} For the function $g(x) = \frac{1}{x-1}$, what is the domain of the function based on its graph?
    \item A. All real numbers except $x = 1$
    \item B. All real numbers
    \item C. $x \geq 1$
    \item D. $x \leq 1$ \textbf{Part 3: Function and Graph Relationship}
    \item \textbf{True/False:} The graph of a linear function $y = mx + b$ always forms a straight line, where $m$ is the slope and $b$ is the y-intercept. (True/False)
    \item \textbf{Short Answer:} Explain how the graph of the function $h(x) = e^x$ reflects the properties of exponential growth.
    \item \textbf{Multiple Choice:} Which of the following statements best describes the relationship between the function $f(x) = x^2$ and its graph?
    \item A. The graph is a line with a positive slope.
    \item B. The graph is a parabola that opens downward.
    \item C. The graph is a parabola that opens upward.
    \item D. The graph is a circle centered at the origin.
\end{enumerate}

%%===============================================================
\section{Applications of Linear Functions}
\label{sec:applications-of-linear-functions}
%%===============================================================

\paragraph{What are Linear Equations?}



In the functions section we talk about how a function is like a box that

takes an independent input value and uses a rule defined mathematically

to create a unique output value. The value for the output is dependent

on the value that is put in the box. We call the values that are going

into the box the independent values or the domain. We call the values

coming out of the box the dependent values or the range.



Unless we specify differently, on Cartesian graphs the domain is the

real numbers. In the Cartesian Plane section we saw how running

different values through a function to identify the points on the

Cartesian plane by picking the first (x) value of the point the domain

and the second (y) value from the range. To restate this: by convention

the two variables for a function on the Cartesian Plane are x for the

domain, the independent variable, and y for the range, the dependent

variable. The variable y is the same as writing f(x). Mathematicians

recognize this equivalence but generally prefer to write y because its

shorter. Because a function definition has an input and an output it

must also contain an equal sign. The section in this book solving

equations showed the various operations we can perform on both sides of

an equal sign and still maintain the notion of equivalence. In this

section we plug different values into the independent variable and solve

to find the associated dependent variable. For instance if we start with

the equation:



$$y = x$$



$$\textit{We can add a -x to both sides to get the equation}$$



$$y -x = x - x$$



$$\textit{which we then simplify to}$$



$$y -x = 0$$



$$\textit{Or we can add a -y to both sides to get the equation}$$



$$y -y = x - y$$



$$\textit{which we then simplify to}$$



$$0 = x - y$$



$$\textit{Since $y -x = 0  \equiv 0 = x - y$ then:}$$



$$y -x = 0 = x - y$$



$$\textit{Using the transitive property}$$



$$y -x  = x - y$$



$$\textit{adding x + y to both sides gives us}$$



$$(y -x )+( y + x)  = (x - y) + (y + x)$$



$$\textit{using the associative property we change this to}$$



$$(y - y )+( x + x)  = (x - x) + (y + y)$$



$$\textit{which simplifies to}$$



$$2x  = 2y$$



$$\textit{And divide both sides by 2}$$



$$\frac{2x}{2}  = \frac{2y}{2}$$



$$\textit{To simply reverse the order and show}$$



$$x = y$$



We have not really proved anything mathematically above, but these

operations allow us to manipulate equations to get the dependent

variable by itself on one side of the equals sign. Then we can plug

numbers into the independent variable to discover the function values

for those numbers. Then we can draw these values as points on the

Cartesian plane and get a feel for what the function would look like if

we could see all the points defined by the function at once.



\paragraph{Horizontal Linear Equations}



Equations of the form $y = C\_2$ are linear functions of the general form

y = m x + b where slope m = 0 and the constant $C\_2$ \textbf{is} the

y-intercept b (in the general form). The graph of this zero-slope

function is a straight horizontal line, intercepting the y-axis at $C\_2$,

including zero and extends infinitely in the positive and negative

directions for all \textbf{R} values of x (see the following diagram).



% [Image: y = c, where c = 3]



The domain for such functions is \textbf{R} covering all real numbers (unless

otherwise specified), but the range is just the set { c }. The equation

y = 0 is the x-axis.



\paragraph{Vertical Linear Equations}



Equation $x = C\_1$, x is one single value $C\_1$ and y, being unrestricted,

is every \textbf{R} number. The graph of $x = C\_1$ is a straight vertical line

where $x = C\_1$, covering \textbf{all} positive, negative and zero values of y

(see the following diagram).



% [Image: x = c, where c = 3]



Its domain is set { $C\_1$ } and range is set \textbf{R} (unless otherwise

specified). $x = C\_1$ is technically not a function (there is more than

one value of y for each value of x), but it's a relation. Vertical

lines have \textbf{no} slope (m = divide by zero, undefined, plus and minus

infinity). These are the only types of linear equations of the general

form shown previously which are not linear functions. The equation x = 0

is \textbf{exactly} the y-axis. Lines with steepness approaching vertical

have very large-magnitude slopes but are still functions.



\paragraph{General Linear Equations}



We are going to start by looking at simple functions called **linear

equations**. When none of the instances of x and y in the algebraic

expression defining the function rule have exponents then all the

instances of x and y can be combined into just two occurrences. the

graph of the expression can be represented as a straight line. The

equation that expresses the function is considered a \textbf{linear equation}

with two variables. The following equation is a simple example of such a

linear equation:







$y - x = 2 \,$







Since y is the dependent variable it is standing in for the function. We

can re-write the expression as f(x) - x = 2. If we add an x to both

sides the equality property holds and we get the expression f(x) - x + x

= x + 2. Simplifying we get f(x) = x + 2. In the following table we'll

pick 3 values for x, and then calculate the dependent (y) values from

f(x).



$$

\begin{array}{|c|c|c|}

\hline

\text{x value} & \text{y value (abscissa)} & \text{Coordinates} (x,y) \\\\ \hline

-1 & 1 & (-1,1) \\\\ \hline

0 & 2 & (0,2) \\\\ \hline

1 & 3 & (1,3) \\\\ \hline

\end{array}

$$



where x and y are variables to be plotted in a two-dimensional Cartesian

coordinate graph as shown here:



% [Image: ]



% [Image: ]



This function is equivalent to the previous example of a linear

equation, y - x = 2. The arrows at each end of the line indicate that

the line extends infinitely in both directions. All linear functions of

a single input variable have or can be algebraically arranged to have

the general form:







$y = f(x) = m x + b \,$







where x and y are variables, f(x) is the function of x, m is a constant

called the \textbf{slope} of the line, and b is a constant which is the

ordinate of the \textbf{y-intercept} (i. e. the value of y where the function

line crosses the y-axis). The slope indicates the steepness of the line.

In the previous example where y = x + 2, the slope m = 1 and the

y-intercept ordinate b = 2. The y = mx+b form of a linear function is

called the \textbf{slope-intercept form}.



Unless a domain for x is otherwise stated, the domain for linear

functions will be assumed to be all real numbers and so the lines in

graphs of all linear functions extend infinitely in both directions.

Also in linear functions with all real number domains, the range of a

linear function will cover the entire set of real numbers for y, unless

the slope m = 0 and the function equals a constant. In such cases, the

range is simply the constant.



Conversely, if a function has the general form y = mx + b or if it can

be arranged to have that form, the function is linear. A linear equation

with two variables has or can be algebraically rearranged to have the

general form^1^:







$A x + B y = C \,$







where x and y represent the linear variables, and the letters A, B, and

C can represent any real constants, either positive or negative.

Conversely, an equation with two variables x and y having that general

form, or being able to be arranged in that form, would be linear as long

as A and B are not both equal to 0. In the preceding equation, capital

letters are to avoid confusion with other constants in this chapter and

for consistency with Reference 1.



If one divides the preceding equation by B (when B is not 0) and solves

for y, the following form can be obtained:







$y = (-A/B) x + (C/B) \,$







If one equates -A/B to the slope m and C/B to the y-intercept ordinate

b, it can be seen that the general form for a linear equation and the

slope-intercept form for a linear function are practically

interconvertible except for the fact that, in a linear function, the B

constant in the linear equation form cannot equal 0.



\paragraph{X and Y axis intercepts}



An axis intercept point is a point where the graph of a function,

relation, or equation intersects the X or Y axes. This section is about

finding out how a particular set of functions: linear functions cross

the axes.



We know that the domains of most lines are infinite because they are

defined at every value of X. The exception is lines that are defined as

$X=c$ where c is a number we choose when we write the function. By

definition this line is only defined for one value of X. Since the

domain maps onto more than one value for the range this is actually a

relationship and not a function. We've seen the graph of this

relationships is a vertical line that passes through the point (c,Y).



Lines with the equation X=C intersect the X axis once and the Y axis 0

or infinite times.



We can also restrict the range of a function by simply writing $Y=c$

where c is again any number we choose. The graph of this line is a

horizontal line that passes through the point (X,c). When Y=o this line

is the same as the X axis. when $c \ne 0$, then the line can never

intercept the X axis.



Lines with the equation Y=C intersect the Y axis once and the X axis 0

or infinite times.



When looking at Cartesian graphs and linear equations we run into a

mathematical axiom: "Two points determine a line.". We will see how

this axiom affects the slope-intercept definition of a line

$y= f(x) = m x + b$ in the next section. When two lines intersect they

intersect at a point. If a line is not horizontal or perpendicular it

will have to intersect the X and Y axes once, but only once.



In this page we are going to accept the statement"At most one line can

be drawn through any point not on a given line parallel to the given

line in a plane."



We've seen that for the equation Y=mX + b the Y intercept will always

be at b because that is where X=0.



Using Algebra we can subtract b from both sides: Y - b = mX



and multiply by $\frac {1}{m}$



$\frac {Y}{m} - \frac {b}{m} = X$



we can see that the X intercept is going to be $-\frac{b}{m}$.



An axis intercept may simply refer to the number value on the axis where

the intersection occurs. For brevity we may say the line has an X

intercept of 1 and a Y intercept of 2. After graphing just a few lines

you will be able to tell this line points down and runs through

quadrants II, I, and IV. With a little more practice you will be able to

know that the equation for the line is Y=-2x + 2. We will see that by

specifying the two points we are actually implying the slope of the

line. There is an exception to this rule. If we say a line crosses the

axes at 0 we know that the line will pass through 2 quadrants instead of

3, but we won't know which quadrants or how steep the line is. When we

look at slope in the next section we will see why the equations above

specify a point and a slope.



When you are trying to graph a linear equation finding the axes

intercepts is often the easiest way to go about doing it. To find the

x-intercept, set y = 0 and solve for x. To find the y-intercept, set x =

0 and solve for y. For most examples the intercepts are different

points, and a line can be drawn through the two intercepts. If both

intercepts are (0,0), then another point must be determined to graph the

line. If the equations is in the form x = c or y = c, the horizontal or

vertical lines are very simple to plot.



\paragraph{Slope}



\textbf{Algebra/Slope}



Slope is the change in the vertical distance of a line on a coordinate

plane over the change in horizontal difference. In other words, it is

the "rise" over the "run" or the steepness of a line. Slope is usually

represented by the symbol $m$ like in the equation $y = mx + b$, m the

coefficient of x represents the slope of the line.



Slope is computed by measuring the change in vertical distance divided

by the change in horizontal difference, i.e.:







$\ m = \frac{\Delta y}{\Delta x} = \frac{rise}{run}$







The Greek uppercase letter $\Delta$ represents change, in this case

change in the y-coordinates divided by the change in x-coordinates.



\textbf{Positive Slope/ Negative Slope}



If a line goes up from left to right, then the slope has to be positive.

For example, a slope of ¾ would have a "rise" of 3, or go up 3; and a

"run" of 4, or go right 4. Both numbers in the slope are either negative

or positive in order to have a positive slope.



If a line goes down from left to right, then the slope has to be

negative. For example, a slope of -3/4 would have a "rise" of -3, or go

down 3; and a "run" of 4, or go right 4. Only one number in the slope

can be negative for a line to have a negative slope.



\textbf{Other Types of Slope}



There are two special circumstances, no slope and slope of zero. A

horizontal line has a slope of 0 and a vertical line has an undefined

slope.



Horizontal lines have the form: $\ y = a$ ; where a is a constant, i.e.

$a \in R$. Vertical lines have the form: $\ x = a$ ; where a is a constant, i.e.

$a \in R$



\paragraph{Determining Slope}



To determine the slope you need some information. This can include two

(or more) coordinates, a parallel slope and a coordinate, a

perpendicular slope and a coordinate, or the y-intercept and slope.



For completely horizontal lines, the difference in y coordinates between

any two points is 0, so the slope m = 0, indicating no steepness in the

line at all. If the line extends between right-upper (+,+) and

left-lower ( -, -) directions, then the slope is positive. As the slope

increases, the line becomes steeper until the line is almost vertical

when the slope is very large. When the slope m = 1, the line is diagonal

with an angle halfway between the x and y axes. If the line extends

between left-upper (-,+) and right-lower (+, -) directions, then the

slope is negative. As the slope changes from 0 to very negative numbers,

the steepness in the opposite direction increases. Compare the slope ( m

) values in the following graph of functions y = 1 (where m = 0), y = (1/2) x + 1, y = x + 1, y = 2 x, y = -(1/2) x + 1, y

= -x + 1, and y = -2 x + 1. For all two-variable linear equations that

can be converted to linear functions, the same calculation applies to

slopes for those lines.



% [Image: ]



For the most part finding slope when given information is a simple

matter. Simply take the slope equation y=mx+b and replace the variable

with whatever information you know, and solve.



\paragraph{Two Coordinates}



To find the slope with two coordinates, you must first find the slope.

Use the standard equation $\frac{y\_2-y\_1}{x\_2-x\_1}$. Put that into the

equation as m, and replace x and y with x and y from one of the

coordinates. Solve for b. Put that into the equation and you're done.



Example: (1,4) (4,8)



$m = \frac{y\_2-y\_1}{x\_2-x\_1}$



$m = \frac{4-1}{4-8}$



$m = \frac{3}{-4}$



Plug that right in.



$y = \frac{3}{-4}x+b$



$4 = \frac{3}{-4}(1)+b$



$4 = \frac{3}{-4}+b$



$4-\frac{3}{-4} = b$



$\frac{-19}{4} = b$



Put that in the equation and you're done.



$y = \frac{3}{-4}x + \frac{-19}{4}$



\paragraph{Parallel Lines}



If you have to find the slope of a line(Let's say AB) which is parallel

to line(Let's say XY) then using the coordinates of line XY you can

find the coordinates of slope of line AB As, Slope of line of a line

parallel to another line is equal,i.e. Slope of AB = Slope of XY







AB and XY are PARALLEL Lines.Find the slope of AB.



Let line XY have the coordinates:-



$X(2,4)=(x\_1,y\_1)$ and $Y(3,6)=(x\_2,y\_2)$



Slope of XY



$m = (y\_2 - y\_1)/(x\_2-x\_1)$ $= (6-4)/(3-2)$ $= 2/1$ $\Rightarrow m = 2$



Slope of AB = Slope of XY, so Slope of AB = 2



\paragraph{Resources}



\begin{itemize}
    \item \href{https://courses.lumenlearning.com/beginalgebra/chapter/introduction-to-solving-linear-equations/}{Solving Linear
\end{itemize}
    Equations},

    Lumen Learning

\begin{itemize}
    \item \href{https://openstax.org/books/intermediate-algebra-2e/pages/2-1-use-a-general-strategy-to-solve-linear-equations}{General Strategy to Solving Linear
\end{itemize}
    Equations},

    OpenStax

\begin{itemize}
    \item \href{https://openstax.org/books/intermediate-algebra-2e/pages/2-3-solve-a-formula-for-a-specific-variable}{Solving Linear Equations for a Specific
\end{itemize}
    Variable},

    OpenStax

\begin{itemize}
    \item \href{https://www.youtube.com/watch?v=bsGyJ0GeRy8}{Solving Linear Equations: Example
\end{itemize}
    1}. Produced by TA

    Catherine Sporer, UTSA

\begin{itemize}
    \item \href{https://www.youtube.com/watch?v=nGZcfHQEdqo}{Solving Linear Equations: Example
\end{itemize}
    2}. Produced by TA

    Catherine Sporer, UTSA



\paragraph{Licensing}



Content obtained and/or adapted from:



\begin{itemize}
    \item \href{https://en.wikibooks.org/wiki/Algebra/Linear_Equations_and_Functions}{Linear Equations and Functions, Wikibooks:
\end{itemize}
    Algebra}

    under a CC BY-SA license

\begin{itemize}
    \item \href{https://en.wikibooks.org/wiki/Algebra/Intercepts}{Intercepts, Wikibooks:
\end{itemize}
    Algebra} under a

    CC BY-SA license

\begin{itemize}
    \item \href{https://en.wikibooks.org/wiki/Algebra/Slope}{Slope, Wikibooks:
\end{itemize}
    Algebra} under a CC

    BY-SA license

\subsection{Learning Objectives}

\begin{enumerate}
    \item \textbf{Identify and Classify Linear Equations:} Students will be able to identify linear equations in various forms, including horizontal ($y = C$), vertical ($x = C$), and general linear equations ($Ax + By = C$), and classify them based on their properties, such as slope and intercepts.
    \item \textbf{Graph Linear Equations:} Students will be able to graph linear equations on the Cartesian plane by determining and plotting the x-intercept and y-intercept, and by understanding how the slope affects the orientation of the line. This includes recognizing special cases such as horizontal and vertical lines.
    \item \textbf{Calculate and Interpret Slope:} Students will be able to calculate the slope of a line given two points or from the slope-intercept form of the equation. They will also interpret the meaning of the slope in the context of the graph, such as positive, negative, zero, and undefined slopes.
\end{enumerate}

\subsection{Assessment Instrument}

\begin{enumerate}
    \item \textbf{Part 1: Identify and Classify Linear Equations}
    \item \textbf{Multiple Choice:} Which of the following equations represents a vertical line?
    \item A. $y = 2x + 3$
    \item B. $x = -5$
    \item C. $y = -4$
    \item D. $2x + y = 7$
    \item \textbf{True/False:} The equation $4x - 3y = 12$ can be classified as a general linear equation and has a slope of $\frac{4}{3}$. (True/False)
    \item \textbf{Short Answer:} Identify the type of line represented by the equation $y = -2$ and describe its properties, including slope and intercepts. \textbf{Part 2: Graph Linear Equations}
    \item \textbf{Calculation:} Graph the equation $2x - 3y = 6$ by finding the x-intercept and y-intercept. Provide the coordinates of these intercepts.
    \item \textbf{True/False:} The equation $y = \frac{1}{2}x + 1$ represents a line with a slope of $\frac{1}{2}$ and a y-intercept of 1. (True/False)
    \item \textbf{Short Answer:} Describe how the slope affects the orientation of the line when graphing the equation $y = -3x + 2$. Include details on how the line will be oriented in the Cartesian plane. \textbf{Part 3: Calculate and Interpret Slope}
    \item \textbf{Calculation:} Given the points $(2, 3)$ and $(5, 7)$, calculate the slope of the line passing through these points.
    \item \textbf{Short Answer:} Interpret the meaning of a line with a slope of $0$. What does this indicate about the line's orientation and its relationship with the y-axis?
    \item \textbf{Multiple Choice:} Which of the following slopes corresponds to a line that is perpendicular to a line with a slope of $-2$?
    \item A. $2$
    \item B. $-\frac{1}{2}$
    \item C. $1$
    \item D. $-\frac{1}{2}$
\end{enumerate}

%%===============================================================
\section{Quadratic Functions and Applications}
\label{sec:quadratic-functions-and-applications}
%%===============================================================

In algebra, a \textbf{quadratic function}, a \textbf{quadratic polynomial}, a

\textbf{polynomial of degree 2}, or simply a \textbf{quadratic}, is a polynomial

function with one or more variables in which the highest-degree term is

of the second degree.



For example, a \textit{univariate} (single-variable) quadratic function has the

form



$$f(x)=ax^2+bx+c,\quad a \ne 0$$ in the single variable \textit{x}. The graph

of a univariate quadratic function is a parabola whose axis of symmetry

is parallel to the y-axis, as shown at right.



If the quadratic function is set equal to zero, then the result is a

quadratic equation. The solutions to the univariate equation are called

the roots of the univariate function.



The bivariate case in terms of variables \textit{x} and \textit{y} has the form



$$f(x,y) = a x^2 + by^2 + cx y+ d x+ ey + f \,$$



with at least one of \textit{a, b, c} not equal to zero, and an equation

setting this function equal to zero gives rise to a conic section (a

circle or other ellipse, a parabola, or a hyperbola).



A quadratic function in three variables \textit{x}, \textit{y,} and \textit{z} contains

exclusively terms \textit{$x^2$}, \textit{$y^2$}, \textit{$z^2$}, \textit{xy}, \textit{xz}, \textit{yz}, \textit{x}, \textit{y},

\textit{z}, and a constant:



$$f(x,y,z)=ax^2+by^2+cz^2+dxy+exz+fyz+gx+hy+iz +j,$$



with at least one of the coefficients \textit{a, b, c, d, e,} or \textit{f} of the

second-degree terms being non-zero.



In general there can be an arbitrarily large number of variables, in

which case the resulting surface of setting a quadratic function to zero

is called a quadric, but the highest degree term must be of degree 2,

such as \textit{$x^2$}, \textit{xy}, \textit{yz}, etc.



% [Image: A quadratic polynomial with two real roots (crossings of the \textit{x} axis)

and hence no complex roots. Some other quadratic polynomials have their

minimum above the \textit{x} axis, in which case there are no real roots and

two complex

roots.]



\paragraph{Etymology}



The adjective \textit{quadratic} comes from the Latin word \textit{quadratum}

("square"). A term like \textit{$x^2$} is

called a square in algebra because it is the area of a \textit{square} with

side x.



\paragraph{Terminology}



\paragraph{Coefficients}



The coefficients of a polynomial are often taken to be real or complex

numbers, but in fact, a polynomial may be defined over any ring.



\paragraph{Degree}



When using the term "quadratic polynomial", authors sometimes mean

"having degree exactly 2", and sometimes "having degree at most 2".

If the degree is less than 2, this may be called a "degenerate case".

Usually the context will establish which of the two is meant.



Sometimes the word "order" is used with the meaning of "degree",

e.g. a second-order polynomial.



\paragraph{Variables}



A quadratic polynomial may involve a single variable \textit{x} (the univariate

case), or multiple variables such as \textit{x}, \textit{y}, and \textit{z} (the multivariate

case).



\paragraph{The one-variable case}



Any single-variable quadratic polynomial may be written as



$$ax^2 + bx + c,\,$$ where \textit{x} is the variable, and \textit{a}, \textit{b}, and \textit{c}

represent the coefficients. In elementary algebra, such polynomials

often arise in the form of a quadratic equation $ax^2 + bx + c = 0$. The

solutions to this equation are called the roots of the quadratic

polynomial, and may be found through factorization, completing the

square, graphing, Newton's method, or through the use of the quadratic

formula. Each quadratic polynomial has an associated quadratic function,

whose graph is a parabola.



\paragraph{Bivariate case}



Any quadratic polynomial with two variables may be written as



$$f(x,y) = a x^2 + b y^2 + cxy + dx+ e y + f, \,$$ where \textit{x} and \textit{y}

are the variables and \textit{a}, \textit{b}, \textit{c}, \textit{d}, \textit{e}, and \textit{f} are the

coefficients. Such polynomials are fundamental to the study of conic

sections, which are characterized by equating the expression for \textit{f}

(\textit{x}, \textit{y}) to zero. Similarly, quadratic polynomials with three or more

variables correspond to quadric surfaces and hypersurfaces. In linear

algebra, quadratic polynomials can be generalized to the notion of a

quadratic form on a vector space.



\paragraph{Forms of a univariate quadratic function}



A univariate quadratic function can be expressed in three formats:



\begin{itemize}
    \item $f(x) = a x^2 + b x + c \,$ is called the \textbf{standard form},
    \item $f(x) = a(x - r\_1)(x - r\_2)\,$ is called the \textbf{factored form},
\end{itemize}
    where \textit{$r\_1$} and

    \textit{$r\_2$} are the roots of the

    quadratic function and the solutions of the corresponding quadratic

    equation.

\begin{itemize}
    \item $f(x) = a(x - h)^2 + k \,$ is called the \textbf{vertex form}, where
\end{itemize}
    \textit{h} and k are

    the x and y

    coordinates of the vertex, respectively.



The coefficient \textit{a} is the same value in all

three forms. To convert the \textbf{standard form} to \textbf{factored form}, one

needs only the quadratic formula to determine the two roots

\textit{$r\_1$} and \textit{$r\_2$}. To convert the **standard

form\textbf{ to }vertex form**, one needs a process called completing the

square. To convert the factored form (or vertex form) to standard form,

one needs to multiply, expand and/or distribute the factors.



\paragraph{Graph of the univariate function}



Regardless of the format, the graph of a univariate quadratic function

$f(x) = ax^2 + bx + c$ is a parabola (as shown at the right).

Equivalently, this is the graph of the bivariate quadratic equation

$y = ax^2 + bx + c$.



\begin{itemize}
    \item If \textit{a}, the parabola opens upwards.
    \item If \textit{a}, the parabola opens
\end{itemize}
    downwards.



The coefficient \textit{a} controls the degree of

curvature of the graph; a larger magnitude of

\textit{a} gives the graph a more closed (sharply

curved) appearance.



The coefficients \textit{b} and

\textit{a} together control the location of the axis

of symmetry of the parabola (also the

x-coordinate of the vertex and the \textit{h}

parameter in the vertex form) which is at



$$x = -\frac{b}{2a}.$$



The coefficient \textit{c} controls the height of the

parabola; more specifically, it is the height of the parabola where it

intercepts the y-axis.



% [Image: $f(x) = ax^2 |\_{a=\{0.1,0.3,1,3\}} $]



% [Image: $f(x) = x^2 + bx |\_{b=\{1,2,3,4\}} $]



% [Image: $f(x) = x^2 + bx |\_{b=\{-1,-2,-3,-4\}} $]



\paragraph{Vertex}



The \textbf{vertex} of a parabola is the place where it turns; hence, it is

also called the \textbf{turning point}. If the quadratic function is in

vertex form, the vertex is \textit{(h,k)}. Using

the method of completing the square, one can turn the standard form



$$f(x) = a x^2 + b x + c \,$$ into



$\begin{aligned}

f(x) \&= a x^2 + b x + c \\\\\&= a (x - h)^2 + k \\\\\&= a\left(x - \frac{-b}{2a}\right)^2 + \left(c - \frac{b^2}{4a}\right), \\\\    \end{aligned}$



so the vertex, \textit{(h,k)}, of the parabola in

standard form is



$\left(-\frac{b}{2a}, c - \frac{b^2}{4a}\right).$



If the quadratic function is in factored form



$$f(x) = a(x - r\_1)(x - r\_2) \,$$ the average of the two roots, i.e.,



$\frac{r\_1 + r\_2}{2} \,$



is the x-coordinate of the vertex, and hence

the vertex \textit{(h,k)} is



$\left(\frac{r\_1 + r\_2}{2}, f\left(\frac{r\_1 + r\_2}{2}\right)\right).$



The vertex is also the maximum point if

\textit{a}, or the minimum point if

\textit{a}.



The vertical line



$x=h=-\frac{b}{2a}$



that passes through the vertex is also the \textbf{axis of symmetry} of the

parabola.



\paragraph{Maximum and minimum points}



Using calculus, the vertex point, being a maximum or minimum of the

function, can be obtained by finding the roots of the derivative:



$$f(x)=ax^2+bx+c \quad \Rightarrow \quad f'(x)=2ax+b \,\.$$

x is a root of

\textit{f} '(\textit{x}) if \textit{f} '(\textit{x}) = 0 resulting in



$$x=-\frac{b}{2a}$$ with the corresponding function value



$$f(x) = a \left (-\frac{b}{2a} \right)^2+b \left (-\frac{b}{2a} \right)+c = c-\frac{b^2}{4a} \,\,$$

so again the vertex point coordinates,

\textit{(h,k)}, can be expressed as



$$\left (-\frac {b}{2a}, c-\frac {b^2}{4a} \right).$$



\paragraph{Roots of the univariate function}



% [Image: Graph of $y = ax^2 + bx + c$, where a and the discriminant $b^2 - 4ac$ are positive, with 
\begin{itemize}
    \item Roots and y-intercept in ``{=html} red ``{=html}
    \item Vertex and axis of symmetry in ``{=html} blue ``{=html}
    \item Focus and directrix in ``{=html} pink ``{=html} ]
\end{itemize}

% [Image: Visualisation of the complex roots of $y = ax^2 + bx + c$: the parabola is rotated 180° about

its vertex (``{=html} orange

``{=html}). Its x-intercepts are rotated 90° around their

mid-point, and the Cartesian plane is interpreted as the complex plane

(``{=html} green

``{=html})]



\paragraph{Exact roots}



The roots (or \textit{zeros}), \textit{$r\_1$} and

\textit{$r\_2$}, of the univariate quadratic

function



$\begin{aligned}

f(x) \&= ax^2+bx+c \\\\&= a(x-r\_1)(x-r\_2), \\\\ 

\end{aligned}$



are the values of x for which \textit{f}(\textit{x}) = 0.



When the coefficients \textit{a},

\textit{b}, and \textit{c}, are real

or complex, the roots are



$$r\_1=\frac{-b - \sqrt{b^2-4ac}}{2a},$$



$$r\_2=\frac{-b + \sqrt{b^2-4ac}}{2a}.$$



\paragraph{Upper bound on the magnitude of the roots}



The modulus of the roots of a quadratic $ax^2+bx+c\,$ can be no greater

than $\frac{\max(|a|, |b|, |c|)}{|a|}\times \phi,\,$ where $\phi$ is the

golden ratio $\frac{1+\sqrt{5}}{2}.$



\paragraph{The square root of a univariate quadratic function}



The square root of a univariate quadratic function gives rise to one of

the four conic sections, almost always either to an ellipse or to a

hyperbola.



If $a>0\,$ then the equation $y = \pm \sqrt{a x^2 + b x + c}$ describes a hyperbola, as can be seen by squaring both sides. The

directions of the axes of the hyperbola are determined by the ordinate of the minimum point of the corresponding parabola

$y\_p = a x^2 + b x + c \,$. If the ordinate is negative, then the hyperbola's major axis (through its vertices) is horizontal, while if

the ordinate is positive then the hyperbola's major axis is vertical.



If $a0 \,$ the function has a minimum if \textit{A}\>0, and a maximum

if \textit{A}\0 and a maximum if \textit{A}\<0;

its graph forms a parabolic cylinder.



\paragraph{Resources}



\begin{itemize}
    \item \href{https://mathresearch.utsa.edu/wikiFiles/MAT1053/Quadratic\_Functions/MAT1053\_M2.2Quadratic\_Functions.pdf}{Quadratic
\end{itemize}
    Functions},

    Book Chapters

\begin{itemize}
    \item \href{https://mathresearch.utsa.edu/wikiFiles/MAT1053/Quadratic\_Functions/MAT1053\_M2.2Quadratic\_FunctionsGN.pdf}{Guided
\end{itemize}
    Notes}



\paragraph{Licensing}



Content obtained and/or adapted from:



\begin{itemize}
    \item \href{https://en.wikipedia.org/wiki/Quadratic\_function}{Quadratic function,
\end{itemize}
    Wikipedia} under a

    CC BY-SA license

\subsection{Learning Objectives}

\begin{enumerate}
    \item \textbf{Understand and Apply Quadratic Function Forms:} Students will be able to identify and convert between the standard form, factored form, and vertex form of a univariate quadratic function. They will demonstrate their understanding by expressing a given quadratic function in each of these forms and interpreting the significance of the coefficients in terms of the graph of the function.
    \item \textbf{Analyze the Properties of Quadratic Functions:} Students will be able to analyze and describe the properties of a quadratic function, including the direction in which the parabola opens, the location of the vertex, and the axis of symmetry. They will use these properties to sketch the graph of the function and determine the maximum or minimum value of the function based on the coefficient $a$.
    \item \textbf{Solve and Interpret Quadratic Equations:} Students will be able to solve quadratic equations using various methods, such as factoring, completing the square, and applying the quadratic formula. They will demonstrate their ability to find the roots of quadratic equations and interpret the solutions in the context of the graph, including understanding the nature of the roots (real or complex) and their impact on the graph of the function.
\end{enumerate}

\subsection{Assessment Instrument}

\begin{enumerate}
    \item \textbf{Part 1: Understand and Apply Quadratic Function Forms}
    \item \textbf{Multiple Choice:} Which of the following represents the standard form of a quadratic function?
    \item A. $y = a(x - h)^2 + k$
    \item B. $y = a(x - r_1)(x - r_2)$
    \item C. $y = ax^2 + bx + c$
    \item D. $y = \frac{a}{x - h} + k$
    \item \textbf{True/False:} If a quadratic function is given in the form $y = 2(x - 1)(x + 3)$, then the equivalent standard form of the function is $y = 2x^2 + 4x - 6$. (True/False)
    \item \textbf{Short Answer:} Convert the quadratic function $f(x) = x^2 - 6x + 8$ into its vertex form. Describe the significance of the vertex form in terms of the graph of the function. \textbf{Part 2: Analyze the Properties of Quadratic Functions}
    \item \textbf{Calculation:} Given the quadratic function $f(x) = -3x^2 + 12x - 7$, determine the direction in which the parabola opens, the coordinates of the vertex, and the axis of symmetry.
    \item \textbf{Short Answer:} Explain how the coefficient $a$ in the quadratic function $f(x) = ax^2 + bx + c$ affects the direction in which the parabola opens. What does a positive $a$ indicate compared to a negative $a$?
    \item \textbf{Multiple Choice:} For the quadratic function $g(x) = (x - 4)^2 - 9$, what is the maximum or minimum value of the function, and where does it occur?
    \item A. Minimum value of $-9$ at $x = 4$
    \item B. Maximum value of $9$ at $x = -4$
    \item C. Minimum value of $-9$ at $x = -4$
    \item D. Maximum value of $9$ at $x = 4$ \textbf{Part 3: Solve and Interpret Quadratic Equations}
    \item \textbf{Calculation:} Solve the quadratic equation $x^2 - 5x + 6 = 0$ by factoring and find the roots. Describe how these roots can be interpreted graphically.
    \item \textbf{Short Answer:} Use the quadratic formula to solve the quadratic equation $2x^2 - 4x - 6 = 0$. Show your steps and discuss the nature of the roots.
    \item \textbf{True/False:} The quadratic equation $x^2 + 2x + 5 = 0$ has real roots. (True/False) If false, explain the nature of the roots and how this relates to the graph of the function.
\end{enumerate}

%%===============================================================
\section{Limits}
\label{sec:limits}
%%===============================================================

Limits, the first step into calculus, explain the complex nature of the

subject. It is used to define the process of derivation and integration.

It is also used in other circumstances to intuitively demonstrate the

process of "approaching".



\paragraph{Introduction}



\paragraph{Intuitive Look into Limits}



The limit is one of the greatest tools in the hands of any

mathematician. We will give the limit an approach. Because mathematics

came only due to approaches ... remember?!



We designate limit in the form:



$$\lim\_{x\to a}f(x)$$ This is read as "The limit of $f$ of $x$ as $x$

approaches $a$". This is an important thing to remember, it is basic

notation which is accepted by the world.



We'll take up later the question of how we can determine whether a

limit exists for $f(x)$ at $a$ and, if so, what it is. For now, we'll

look at it from an intuitive standpoint.



Let's say that the function that we're interested in is $f(x)=x^2$ ,

and that we're interested in its limit as $x$ approaches $2$. Using the

above notation, we can write the limit that we're interested in as

follows:



$$\lim\_{x\to2}x^2$$ One way to try to evaluate what this limit is would

be to choose values near 2, compute $f(x)$ for each, and see what

happens as they get closer to 2. There are two ways to approach values

near 2. One is to approach from below, and the other is to approach from

above:



$$

\begin{array}{|c|c|c|c|c|c|c|}

\hline

x & 1.7 & 1.8 & 1.9 & 1.95 & 1.99 & 1.999 \\\\ \hline

f(x)=x^2 & 2.89 & 3.24 & 3.61 & 3.8025 & 3.9601 & 3.996001 \\\\ \hline

\end{array}

$$                                                  



The table above is the case from below.



$$

\begin{array}{|c|c|c|c|c|c|c|}

\hline

x & 2.3 & 2.2 & 2.1 & 2.05 & 2.01 & 2.001 \\\\ \hline

f(x)=x^2 & 5.29 & 4.84 & 4.41 & 4.2025 & 4.0401 & 4.004001 \\\\ \hline

\end{array}

$$ 



The table above is the case from above.



We can see from the tables that as $x$ grows closer and closer to 2,

$f(x)$ seems to get closer and closer to 4, regardless of whether $x$

approaches 2 from above or from below. For this reason, we feel

reasonably confident that the limit of $x^2$ as $x$ approaches 2 is 4,

or, written in limit notation,



$$\lim\_{x\to2}x^2=4$$ We could have also just substituted 2 into $x^2$

and evaluated: $(2)^2=4$. However, this will not work with all limits.



Now let's look at another example. Suppose we're interested in the

behavior of the function $f(x)=\frac{1}{x-2}$ as $x$ approaches 2.

Here's the limit in limit notation:



$$\lim\_{x\to 2}\frac{1}{x-2}$$ Just as before, we can compute function

values as $x$ approaches 2 from below and from above. Here's a table,

approaching from below:



$$

\begin{array}{|c|c|c|c|c|c|c|}

\hline

x & 1.7 & 1.8 & 1.9 & 1.95 & 1.99 & 1.999 \\\\ \hline

f(x)=\frac{1}{x-2} & -3.333 & -5 & -10 & -20 & -100 & -1000 \\\\ \hline

\end{array}

$$ 



And here from above:



$$

\begin{array}{|c|c|c|c|c|c|c|}

\hline

x & 2.3 & 2.2 & 2.1 & 2.05 & 2.01 & 2.001 \\\\ \hline

f(x)=\frac{1}{x-2} & 3.333 & 5 & 10 & 20 & 100 & 1000 \\\\ \hline

\end{array}

$$ 



In this case, the function doesn't seem to be approaching a single

value as $x$ approaches 2, but instead becomes an extremely large

positive or negative number (depending on the direction of approach).

Well, one says such a limit does not exist because no finite number is

approached. This arises the concept of infinity: an undefined quantity

and the limit is also called infinite limit or limit without a bound.



Note that we cannot just substitute 2 into $\frac{1}{x-2}$ and evaluate

as we could with the first example, since we would be dividing by 0.



Both of these examples may seem trivial, but consider the following

function:



$$f(x)=\frac{x^2(x-2)}{x-2}$$ This function is the same as



$$

f(x)=

\begin{cases}

x^2 & \text{if} x \ne 2 \\\\ \text{undefined} & \text{if} x=2 

\end{cases}

$$



Note that these functions are really completely identical; not just

"almost the same," but actually, in terms of the definition of a

function, completely the same; they give exactly the same output for

every input.



In elementary algebra, a typical approach is to simply say that we can

cancel the term $(x-2)$ , and then we have the function $f(x)=x^2$.

However, that would be inaccurate; the function that we have now is not

really the same as the one we started with, because it is defined when

$x=2$ , and our original function was specifically not defined when

$x=2$. This may seem like a minor point, but from making this kind of

assumptions we can easily derive absurd results, such that $0=1$ (see

Mathematical Fallacy § All numbers equal all other numbers in Wikipedia

for a complete example). Even without calculus we can avoid this error

by stating that:



$$f(x)=\frac{x^2(x-2)}{x-2}=x^2\text{ when }x\ne2$$ In calculus, we can

introduce a more intuitive and also correct way of looking at this type

of function. What we want is to be able to say that, although the

function isn't defined when $x=2$, it works almost as if it was. It may

not get there, but it gets really, really close. For instance,

$f(1.99999)=3.99996$. The only question that we have is: what do we mean

by "close"?



\paragraph{Informal Definition of a Limit}



As the precise definition of a limit is a bit technical, it is easier to

start with an informal definition; we'll explain the formal definition

later.



We suppose that a function $f$ is defined for $x$ near $c$ (but we do

not require that it be defined when $x=c$).



> 

>

>    \textbf{Informal definition of a limit}

>    We call $L$ the \textbf{limit of $f(x)$ as $x$ approaches $c$} if

>     $f(x)$ becomes close to $L$ when $x$ is close (but not equal) to

>     $c$ , and if there is no other value $L'$ with the same property.

>    When this holds we write

>    $\lim\_{x\to c}f(x)=L$

>    or

>    $f(x)\to L\quad\text{as}\quad x\to c$



Notice that the definition of a limit is not concerned with the value of

$f(x)$ when $x=c$ (which may exist or may not). All we care about are

the values of $f(x)$ when $x$ is close to $c$ , on either the left or

the right (i.e. less or greater).



Limit can also be understood as: $x$ is infinitely approaching to $c$

but never equals to $c$, just like the function $f(x)=\frac{1}{x}$,

which infinitely approaches to $0$ but never equals $0$.



\paragraph{Basics}



\paragraph{Rules and Identities}



Now that we have defined, informally, what a limit is, we will list some

rules that are useful for working with and computing limits. You will be

able to prove all these once we formally define the fundamental concept

of the limit of a function.



First, the \textbf{constant rule} states that if $f(x)=b$ (that is, $f$ is

constant for all $x$) then the limit as $x$ approaches $c$ must be equal

to $b$. In other words



> 

>

>    \textbf{Constant Rule for Limits}

>    If $b$ and $c$ are constants then $\lim\_{x\to c}b=b$.



   \textbf{Example}: $\lim\_{x\to6}5=5$



Second, the \textbf{identity rule} states that if $f(x)=x$ (that is, $f$ just

gives back whatever number you put in) then the limit of $f$ as $x$

approaches $c$ is equal to $c$. That is,



> 

>

>    \textbf{Identity Rule for Limits}

>    If $c$ is a constant then $\lim\_{x\to c}x=c$.



   \textbf{Example}: $\lim\_{x\to6}x=6$



The next few rules tell us how, given the values of some limits, to

compute others.



> 

>

>    \textbf{Operational Identities for Limits}

>    Suppose that $\lim\_{x\to c}f(x)=L$ and $\lim\_{x\to c}g(x)=M$ and

>     that $k$ is constant. Then

>

> -   $\lim\_{x\to c}k\cdot f(x)=k\cdot\lim\_{x\to c}f(x)=k\cdot L$

> -   $\lim\_{x\to c}\Big[f(x)+g(x)\Big]=\lim\_{x\to c}f(x)+\lim\_{x\to c}g(x)=L+M$

> -   $\lim\_{x\to c}\Big[f(x)-g(x)\Big]=\lim\_{x\to c}f(x)-\lim\_{x\to c}g(x)=L-M$

> -   $\lim\_{x\to c}\Big[f(x)\cdot g(x)\Big]=\lim\_{x\to c}f(x)\cdot\lim\_{x\to c}g(x)=L\cdot M$

> -   $\lim\_{x\to c}\frac{f(x)}{g(x)}=\frac{\lim\limits\_{x\to c}f(x)}{\lim\limits\_{x\to c}g(x)}=\frac{L}{M}\quad\text{ provided }M\ne0$



Notice that in the last rule we need to require that $M$ is not equal to

0 (otherwise we would be dividing by zero which is an undefined

operation).



These rules are known as \textbf{identities}; they are the scalar product,

sum, difference, product, and quotient rules for limits. (A scalar is a

constant, and, when you multiply a function by a constant, we say that

you are performing \textbf{scalar multiplication}.)



Using these rules we can deduce another. Namely, using the rule for

products many times we get that



$$\lim\_{x\to c}f(x)^n=\left(\lim\_{x\to c}f(x)\right)^n=L^n$$ for a

positive integer $n$. This is called the \textbf{power rule}.



As a result, we can safely say that all limits for polynomial functions

can be deduced into several limits that satisfy the identity rule and

thus easier to compute.



Example 1



Find the limit $\lim\_{x\to 2}4x^3$.



We need to simplify the problem, since we have no rules about this

expression by itself. We know from the identity rule above that

$\lim\_{x\to2}x=2$. By the power rule,

$\lim\_{x\to2}x^3=\left(\lim\_{x\to2}x\right)^3=2^3=8$. Lastly, by the

scalar multiplication rule, we get

$\lim\_{x\to2}4x^3=4\cdot\lim\_{x\to2}x^3=4\cdot 8=32.$



$\blacksquare$



Example 2



Find the limit $\lim\_{x\to2}\Big[4x^3+5x+7\Big]$.



To do this informally, we split up the expression, once again, into its

components. As above, $\lim\_{x\to2}4x^3=32$.



Also $\lim\_{x\to2}5x=5\cdot\lim\_{x\to2}x=5\cdot 2=10$ and

$\lim\_{x\to2}7=7$. Adding these together gives



$$\lim\_{x\to2}\Big[4x^3+5x+7\Big]=\lim\_{x\to2}4x^3+\lim\_{x\to2}5x+\lim\_{x\to2}7=32+10+7=49.$$



$\blacksquare$



Example 3



Find the limit $\lim\_{x\to2}\frac{4x^3+5x+7}{(x-4)(x+10)}$.



From the previous example the limit of the numerator is

$\lim\_{x\to 2}\Big[4x^3+5x+7\Big]=49$. The limit of the denominator is



$$\lim\_{x\to2}(x-4)(x+10)=\lim\_{x\to2}\big[x-4\big]\cdot\lim\_{x\to2}\big[x+10\big]=(2-4)(2+10)=-24$$

As the limit of the denominator is not equal to zero we can divide. This

gives



$$\lim\_{x\to2}\frac{4x^3+5x+7}{(x-4)(x+10)}=-\frac{49}{24}.$$



$\blacksquare$



Example 4



Find the limit $\lim\_{x\to4}\frac{x^4-16x+7}{4x-5}$.



We apply the same process here as we did in the previous set of

examples;



$$\lim\_{x\to4}\frac{x^4-16x+7}{4x-5}=\frac{\lim\limits\_{x\to4}\Big[x^4-16x+7\Big]}{\lim\limits\_{x\to4}\Big[4x-5\Big]}=\frac{\lim\limits\_{x\to4}x^4-\lim\limits\_{x\to4}16x+\lim\limits\_{x\to4}7}{\lim\limits\_{x\to4}4x-\lim\limits\_{x\to4}5}$$.



We can evaluate each of these;

$\lim\_{x\to4}(x^4)=256\ ,\ \lim\_{x\to4}(16x)=64\ ,\ \lim\_{x\to4}(7)=7\ ,\ \lim\_{x\to4}(4x)=16\ ,\ \lim\_{x\to4}(5)=5.$

Thus, the answer is $\frac{199}{11}.$



$\blacksquare$



Example 5



Find the limit $\lim\_{x\to2}\frac{x^2-3x+2}{x-2}$.



In this example, evaluating the result directly will result in a

division by 0. While you can determine the answer experimentally, a

mathematical solution is possible as well.



First, the numerator is a polynomial that may be factored:

$\lim\_{x\to2}\frac{(x-2)(x-1)}{x-2}$



Now, you can divide both the numerator and denominator by $(x-2)$:

$\lim\_{x\to2}(x-1)=(2-1)=1$



Remember that the limit is a method to determine the approaching value

of a function instead of the value of the function itself. So, though

the function is undefined at $x=2$, the value of the function is

approaching to $1$ when $x \rightarrow 2$



$\blacksquare$



Example 6



Find the limit $\lim\_{x\to0}\left[\frac{1-\cos(x)}{x}\right]$.



To evaluate this seemingly complex limit, we will need to recall some

sine and cosine identities. We will also have to use two new facts. First, if $f(x)$ is a trigonometric function (that

is, one of sine, cosine, tangent, cotangent, secant or cosecant) and is

defined at $a$ , then $\lim\_{x\to a}f(x)=f(a)$.



Second, $\lim\_{x\to 0}\frac{\sin(x)}{x}=1$. This may be determined

experimentally, or by applying L'Hôpital's rule.



Method 1 (before learning L'Hôpital's rule):



To evaluate the limit, recognize that $1-\cos(x)$ can be multiplied by

$1+\cos(x)$ to obtain $1-\cos^2(x)$ which, by our trig identities, is

$\sin^2(x)$. So, multiply the top and bottom by $1+\cos(x)$. (This is

allowed because it is identical to multiplying by one.) This is a

standard trick for evaluating limits of fractions; multiply the

numerator and the denominator by a carefully chosen expression which

will make the expression simplify somehow. In this case, we should end

up with:



$\begin{aligned}

\lim\_{x \to 0}\left[\frac{1-\cos(x)}{x}\right]\&=\lim\_{x \to 0} \left[\frac{1-\cos(x)}{x} \cdot \frac{1}{1}\right] \\\\\&=\lim\_{x \to 0}\left[\frac{1-\cos(x)}{x} \cdot \frac{1+\cos(x)}{1+\cos(x)}\right] \\\\\&=\lim\_{x \to 0}\frac{(1-\cos(x))(1+\cos(x))}{x(1+\cos(x))} \\\\\&=\lim\_{x \to 0} \frac{1-\cos^2(x)}{x(1+\cos(x))} \\\\\&=\lim\_{x \to 0} \frac{\sin^2(x)}{x(1+\cos(x))} \\\\\&=\lim\_{x \to 0}\left[\frac{\sin(x)}{x} \cdot \frac{\sin(x)}{1+\cos(x)}\right]\\\\    \end{aligned}$



Our next step should be to break this up into

$\lim\_{x\to0}\left[\frac{\sin(x)}{x}\right]\cdot\lim\_{x\to0}\left[\frac{\sin(x)}{1+\cos(x)}\right]$

by the product rule. As mentioned above,

$\lim\_{x\to0}\left[\frac{\sin(x)}{x}\right]=1$.



Next,

$\lim\_{x\to0}\left[\frac{\sin(x)}{1+\cos(x)}\right]=\frac{\lim\limits\_{x\to0}\big[\sin(x)\big]}{\lim\limits\_{x\to0}\big[1+\cos(x)\big]}=\frac{0}{1+\cos(0)}=0$.



Thus, by multiplying these two results, we obtain 0.



$\blacksquare$



Method 2 (after learning L'Hôpital's rule):



We can simply apply L'Hôpital's rule:

$\lim\_{x \rightarrow 0}\frac{1-\cos x}{x}=\lim\_{x \rightarrow 0}\sin x=0$



It seems that L'Hôpital's rule can make a lot of limits like this very

convenient to be calculated. However, since it requires us to learn

derivatives, we cannot use it yet.



$\blacksquare$



We will now present an amazingly useful result, even though we cannot

prove it yet. We can find the limit at $c$ of any polynomial or rational

function, as long as that rational function is defined at $c$ (so we are

not dividing by 0). That is, $c$ must be in the domain of the function.



> 

>

>    \textbf{Limits of Polynomials and Rational functions}

>    If $f$ is a polynomial or rational function that is defined at $c$

>     then

>    $\lim\_{x\to c}f(x)=f(c)$

>     



We already learned this for trigonometric functions, so we see that it

is easy to find limits of polynomial, rational or trigonometric

functions wherever they are defined. In fact, this is true even for

combinations of these functions; thus, for example,

$\lim\_{x\to1}\bigg[\sin(x^2)+4\cos^3(3x-1)\bigg]=\sin(1^2)+4\cos^3(3\cdot 1-1)=\sin 1+4\cos^32$.



\paragraph{The Squeeze Theorem}



The Squeeze Theorem is very important in calculus, where it is typically

used to find the limit of a function by comparison with two other

functions whose limits are known.



It is called the Squeeze Theorem because it refers to a function $f$

whose values are squeezed between the values of two other functions $g$

and $h$ , both of which have the same limit $L$. If the value of $f$ is

trapped between the values of the two functions $g$ and $h$ , the values

of $f$ must also approach $L$.



Expressed more precisely:



> 

>

>    \textbf{Theorem: (Squeeze Theorem)}

>    Suppose that $g(x)\le f(x)\le h(x)$ holds for all $x$ in some open

>     interval containing $c$.

>

> <!-- -->

>

>    If $\lim\_{x\to c}g(x)=\lim\_{x\to c}h(x)=L$,

>    Then $\lim\_{x\to c}f(x)=L$.



% [Image: Graph showing $f$ (blue) being squeezed between $g$ (red) and $h$

(green)]



\textbf{Example}: Compute $\lim\_{x\to0}x\cdot\sin\left(\tfrac{1}{x}\right)$.



Since we know that



> $-1 \le \sin(\frac{1}{x}) \le 1$



Multiplying $x$ into the inequality yields



> $-|x| \le x\sin(\frac{1}{x}) \le |x|$



Now we apply the squeeze theorem



> $\lim\_{x\to0}-|x|\le\lim\_{x\to 0}x\sin(\frac{1}{x})\le\lim\_{x\to0}|x|$



Since $\lim\_{x\to0}-|x|=\lim\_{x\to0}|x|=0$,



$\lim\_{x\to0}x\cdot\sin\left(\tfrac{1}{x}\right)=0.$



$\blacksquare$



% [Image: Plot of $x\cdot\sin\left(\tfrac{1}{x}\right)$ for $-0.5\leq x \leq0.5$]



\paragraph{Finding Limits}



Now, we will discuss how, in practice, to find limits. First, if the function can be built out of rational, trigonometric, logarithmic, or exponential functions, then if a number $c$ is in the domain of the function, then the limit at $c$ is simply the value of the function at $c$:



> $\lim\_{x \to c} f(x)=f(c)$ when $f(x)$ can be built out of rational,

> trigonometric, logarithmic, or exponential functions and $c \in$ the

> Domain of $f(x)$



If $c$ is not in the domain of the function, then in many cases (as with

rational functions) the domain of the function includes all the points

near $c$, but not $c$ itself. An example would be if we wanted to find

$\lim\_{x\to0}\frac{x}{x}$ , where the domain includes all numbers

besides 0.



In that case, in order to find $\lim\_{x\to c}f(x)$ we want to find a

function $g(x)$ similar to $f(x)$ , except with the hole at $c$ filled

in. The limits of $f$ and $g$ will be the same, as can be seen from the

definition of a limit. By definition, the limit depends on $f(x)$ only

at the points where $x$ is close to $c$ but not equal to it, so the

limit at $c$ does not depend on the value of the function at $c$.

Therefore, if $\lim\_{x\to c}g(x)=L$ , $\lim\_{x\to c}f(x)=L$ also. And

since the domain of our new function $g$ includes $c$ , we can now

(assuming $g$ is still built out of rational, trigonometric, logarithmic

and exponential functions) just evaluate it at $c$ as before. Thus we

have $\lim\_{x\to c}f(x)=g(c)$.



In our example, this is easy; canceling the $x$'s gives $g(x)=1$, which

equals $f(x)=\frac{x}{x}$ at all points except 0. Thus, we have

$\lim\_{x\to0}\frac{x}{x}=\lim\_{x\to0}1=1$. In general, when computing

limits of rational functions, it's a good idea to look for common

factors in the numerator and denominator.



\paragraph{Specific DNE Situations}



Note that the limit might not exist at all (DNE means "does not

exist"). There are a number of ways in which this can occur:



\textbf{"Gap"}



   There is a gap (not just a single point) where the function is not defined. As an example, in

   $f(x)=\sqrt{x^2-16}$



$$\lim\_{x\to c}f(x)$$ does not exist when $-4\le c\le4$. There is no way

to "approach" the middle of the graph. Note that the function also has

no limit at the endpoints of the two curves generated (at $c=-4$ and

$c=4$). For the limit to exist, the point must be approachable from

\textit{both} the left and the right.



   Note also that there is no limit at a totally isolated point on a

    graph.



% [Image: $f(x)=\sqrt{x^2-16}$]



<!-- -->



"Jump": If the graph suddenly jumps to a different level, there is no limit at the point of the jump. For example, let $f(x)$ be the greatest integer $\le x$. Then, if $c$ is an integer, when $x$ approaches $c$ from the right $f(x)=c$ , while when $x$ approaches $c$ from the left $f(x)=c-1$. Thus $\lim\_{x\to c}f(x)$ will not exist.



Vertical asymptote: In $f(x)=\frac{1}{x^2}$

   the graph gets arbitrarily high as it approaches 0, so there is no

    limit. (In this case we sometimes say the limit is infinite; see the

    next section.)



% [Image: A graph of $\frac{1}{x^2}$ on the interval

$\[-2,2\]$.]



Infinite oscillation: These next two can be tricky to visualize. In this one, we mean that a graph continually rises above and falls below a horizontal line. In fact, it does this infinitely often as you approach a certain $x$-value. This often means that there is no limit, as the graph never approaches a particular value. However, if the height (and depth) of each oscillation diminishes as the graph approaches the $x$-value, so that the oscillations get arbitrarily smaller, then there might actually be a limit.

   The use of oscillation naturally calls to mind the trigonometric

    functions. An example of a trigonometric function that does not have

    a limit as $x$ approaches 0 is

   $f(x)=\sin\left(\tfrac{1}{x}\right)$

   As $x$ gets closer to 0 the function keeps oscillating between $-1$

    and 1. In fact, $\sin\left(\tfrac{1}{x}\right)$ oscillates an

    infinite number of times on the interval between 0 and any positive

    value of $x$. The sine function is equal to 0 whenever $x=k\pi$ ,

    where $k$ is a positive integer. Between every two integers $k$ ,

    $\sin(x)$ goes back and forth between 0 and $-1$ or 0 and 1. Hence,

    $\sin\left(\tfrac{1}{x}\right)=0$ for every $x=\tfrac{1}{k\pi}$. In

    between consecutive pairs of these values, $\tfrac{1}{k\pi}$ and

    $\tfrac{1}{(k+1)\pi}$ , $\sin\left(\tfrac{1}{x}\right)$ goes back

    and forth from 0, to either $-1$ or 1 and back to 0. We may also

    observe that there are an infinite number of such pairs, and they

    are all between 0 and $1/\pi$. There are a finite number of such

    pairs between any positive value of $x$ and $\tfrac{1}{\pi}$ , so

    there must be infinitely many between any positive value of $x$

    and 0. From our reasoning we may conclude that, as $x$ approaches 0

    from the right, the function $\sin\left(\tfrac{1}{x}\right)$ does

    not approach any specific value. Thus,

    $\lim\_{x\to0}\sin\left(\tfrac{1}{x}\right)$ does not exist.



% [Image: A graph of $\sin(\frac{1}{x})$ on the interval $(0,\frac{1}{\pi})$.]





\paragraph{Determining Limits}



The formal way to determine whether a limit exists is to find out

whether the value of the limit is the same when approaching from below

and above (see at the top of this chapter). The notation for the limit

approaching from below (in increasing order) is



> $\lim\_{x\to c^-}f(x)$, notice the negative sign on the constant $c$



The notation for the limit approaching from above (from decreasing

order) is



> $\lim\_{x\to c^+}f(x)$, notice the positive sign on the constant $c$



For example, let us find the limits of $f(x)=\frac{1}{x}$ when $x$ is

approaching $0$ in both directions. In other words, find

$\lim\_{x\to 0^-}\frac{1}{x}$ and $\lim\_{x\to 0^+}\frac{1}{x}$.



> Recall the table we made when we are trying to intuitively feel the

> limit. We can use that to help us. However, if familiar enough with

> reciprocal functions, we can simply determine the two values by

> imagining the graph of the function.

>

> The following table is when $x$ is approaching from below.



$$

\begin{array}{|c|c|c|c|c|c|c|}

\hline

x & -0.3 & -0.2 & -0.1 & -0.05 & -0.01 & -0.001 \\\\ \hline

f(x)=\frac{1}{x-2} & -3.333 & -5 & -10 & -20 & -100 & -1000 \\\\ \hline

\end{array}

$$     



> Thus, we found that when $x$ is approaching from below to $0$, the

> function approaches negative infinity. In mathematical terms:

>

> $\lim\_{x\to 0^-}\frac{1}{x}=-\infty$

>

> Now let's talk about the approach from above.



$$

\begin{array}{|c|c|c|c|c|c|c|}

\hline

x & 0.3 & 0.2 & 0.1 & 0.05 & 0.01 & 0.001 \\\\ \hline

f(x)=\frac{1}{x-2} & 3.333 & 5 & 10 & 20 & 100 & 1000 \\\\ \hline

\end{array}

$$ 



> We found that $\lim\_{x\to 0^+}\frac{1}{x}=\infty$



$\blacksquare$



The method of determining if limits exist is relatively intuitive. It

only requires some practice to be familiar with the process.



> 

>

>    \textbf{Determining Limits}

>    If $\lim\_{x\to c^-}f(x)=\lim\_{x\to c^+}f(x)=L$, then

>     $\lim\_{x\to c}f(x)=L$.

>    If $\lim\_{x\to c^-}f(x)\ne\lim\_{x\to c^+}f(x)=L$, then the limit

>     does not exist (DNE).



Let's use the same example: find $\lim\_{x\to 0}\frac{1}{x}$.



> Since we already found that $\lim\_{x\to 0^-}\frac{1}{x}=-\infty$ and

> $\lim\_{x\to 0^+}\frac{1}{x}=\infty$, and obviously,

>

> $\lim\_{x\to 0^-}\frac{1}{x}\ne\lim\_{x\to 0^+}\frac{1}{x}$

>

> We can say that $\lim\_{x\to 0}\frac{1}{x}$ does not exist.



$\blacksquare$



\paragraph{Infinity Situations}



Now consider the function



$$g(x)=\frac{1}{x^2}$$ What is the limit as $x$ approaches zero? The

value of $g(0)$ does not exist; it is not defined.



Notice, also, that we can make $g(x)$ as large as we like, by choosing a

small $x$ , as long as $x\ne0$. For example, to make $g(x)$ equal to

$10^{12}$ , we choose $x$ to be $10^{-6}$. Thus,

$\lim\_{x\to0}\frac{1}{x^2}$ does not exist.



However, we \textit{do} know something about what happens to $g(x)$ when $x$

gets close to 0 without reaching it. We want to say we can make $g(x)$

arbitrarily large (as large as we like) by taking $x$ to be sufficiently

close to 0, but not equal to 0. We express this symbolically as follows:



$$\lim\_{x\to0}g(x)=\lim\_{x\to0}\frac{1}{x^2}=\infty$$ Note that the

limit does not exist at $0$ ; for a limit, being $\infty$ is a special

kind of not existing. In general, we make the following definition.



> 

>

>    \textbf{Definition: Informal definition of a limit being $\pm\infty$}

>    We say the \textbf{limit of $f(x)$ as $x$ approaches $c$ is infinity}

>     if $f(x)$ becomes very big (as big as we like) when $x$ is close

>     (but not equal) to $c$.

>    In this case we write

>    $\lim\_{x\to c}f(x)=\infty$

>    or

>    $f(x)\to\infty\quad\text{as}\quad x\to c$.

>    Similarly, we say the **limit of $f(x)$ as $x$ approaches $c$ is

>     negative infinity** if $f(x)$ becomes very negative when $x$ is

>     close (but not equal) to $c$.

>    In this case we write

>    $\lim\_{x\to c}f(x)=-\infty$

>    or

>    $f(x)\to-\infty\quad \text{as} \quad x\to c$.



An example of the second half of the definition would be that

$\lim\_{x\to0}-\frac{1}{x^2}=-\infty$.



\paragraph{Applications of Limits}



To see the power of the concept of the limit, let's consider a moving

car. Suppose we have a car whose position is linear with respect to time

(that is, a graph plotting the position with respect to time will show a

straight line). We want to find the velocity. This is easy to do from

algebra; we just take the slope, and that's our velocity.



But unfortunately, things in the real world don't always travel in nice

straight lines. Cars speed up, slow down, and generally behave in ways

that make it difficult to calculate their velocities.



Now what we really want to do is to find the velocity at a given moment

(the instantaneous velocity). The trouble is that in order to find the

velocity we need two points, while at any given time, we only have one

point. We can, of course, always find the average speed of the car,

given two points in time, but we want to find the speed of the car at

one precise moment.



This is the basic trick of differential calculus, the first of the two

main subjects of this book. We take the average speed at two moments in

time, and then make those two moments in time closer and closer

together. We then see what the limit of the slope is as these two

moments in time are closer and closer, and say that this limit is the

slope at a single instant.



We will study this process in much greater depth later in the book.

First, however, we will need to study limits more carefully.



\paragraph{Resources}



\begin{itemize}
    \item \href{https://mathresearch.utsa.edu/wikiFiles/MAT1193/Limits/Presentation2\_IRC.pptx}{Limits}
\end{itemize}
    (Slides 33-41). PowerPoint file created by Professor Cynthia

    Roberts, UTSA.



\paragraph{Licensing}



Content obtained and/or adapted from:



\begin{itemize}
    \item \href{https://en.wikibooks.org/wiki/Calculus/Limits/An\_Introduction\_to\_Limits}{An Introduction to Limits, Wikibooks:
\end{itemize}
    Calculus}

    under a CC BY-SA license

\subsection{Learning Objectives}

\begin{enumerate}
    \item \textbf{Define and Interpret Limits:} Students will be able to define the concept of a limit in calculus and interpret its notation. They should demonstrate an understanding of the limit notation $\lim\_{x \to a} f(x)$ and explain its meaning in terms of a function approaching a particular value as $x$ approaches $a$.
    \item \textbf{Apply Limit Rules to Compute Limits:} Students will be able to apply basic limit rules and identities, such as the constant rule, identity rule, and operational identities (sum, difference, product, and quotient rules), to compute limits of polynomial and rational functions. They should be able to work through examples, simplify functions, and determine limits using these rules.
    \item \textbf{Evaluate Limits Using Tables and Graphical Methods:} Students will be able to evaluate limits of functions using tables of values and graphical methods. They should be able to analyze function behavior as $x$ approaches a particular value from both above and below, and determine whether the limit exists or is infinite. Additionally, students will understand the concept of limits where functions are undefined at certain points but approach a specific value.
\end{enumerate}

\subsection{Assessment Instrument}

\begin{enumerate}
    \item \textbf{Part 1: Define and Interpret Limits}
    \item \textbf{Multiple Choice:} What does the notation $\lim_{x \to 3} f(x) = 7$ signify?
    \item A. The function $f(x)$ is equal to 7 when $x = 3$.
    \item B. The function $f(x)$ approaches 7 as $x$ approaches 3.
    \item C. The value of $f(x)$ is undefined when $x = 3$.
    \item D. The function $f(x)$ is increasing as $x$ approaches 3.
    \item \textbf{True/False:} The limit $\lim_{x \to 2} \frac{x^2 - 4}{x - 2}$ is equal to 4. (True/False)
    \item \textbf{Short Answer:} Define what is meant by the limit of a function $f(x)$ as $x$ approaches $a$. Explain the significance of this concept in calculus. \textbf{Part 2: Apply Limit Rules to Compute Limits}
    \item \textbf{Calculation:} Compute the limit using the sum rule: $$ \lim_{x \to 2} (3x^2 + 5x - 4) $$
    \item \textbf{Short Answer:} Use the product rule to find the limit of the function $f(x) = (x + 1)(x - 2)$ as $x$ approaches 3. Show your work.
    \item \textbf{Multiple Choice:} Given the function $g(x) = \frac{2x^2 - 3x + 1}{x - 1}$, find the limit as $x$ approaches 1.
    \item A. 2
    \item B. -1
    \item C. 3
    \item D. The limit does not exist \textbf{Part 3: Evaluate Limits Using Tables and Graphical Methods}
    \item \textbf{Table of Values:} Create a table of values for the function $f(x) = \frac{1}{x}$ as $x$ approaches 0 from both positive and negative directions. Use this table to determine whether the limit exists and, if so, what it is.
    \item \textbf{Short Answer:} Explain how to use graphical methods to determine the limit of the function $h(x) = \frac{x^2 - 1}{x - 1}$ as $x$ approaches 1. What do you observe about the function’s behavior near $x = 1$?
    \item \textbf{True/False:} The limit $\lim_{x \to 0} \frac{\sin(x)}{x}$ can be evaluated as 1 using the graphical approach. (True/False) Explain why or why not, referring to the behavior of the function around $x = 0$.
\end{enumerate}

%%===============================================================
\section{One-Sided Limits}
\label{sec:one-sided-limits}
%%===============================================================

In calculus, a \textbf{one-sided limit} is either of the two limits of a

function \textit{f}(\textit{x}) of a real variable \textit{x} as \textit{x} approaches a specified

point either from the left or from the right.



% [Image: The function \textit{f}(\textit{x}) = \textit{$x^2$} + sign(\textit{x}) has a left limit of -1, a right limit of +1, and a function value of 0 at the point \textit{x} = 0.]



The limit as \textit{x} decreases in value approaching \textit{a} (\textit{x} approaches \textit{a}

or "from above") can be denoted:



$\lim\_{x \to a^+}f(x)\ \quad \text{or} \quad  \lim\_{x\ \downarrow\ a}\ f(x) \quad \text{or} \quad \lim\_{x \searrow a}\,f(x) \quad \text{or} \quad \lim\_{x \underset{>}{\to} a}f(x)$



The limit as \textit{x} increases in value approaching \textit{a} (\textit{x} approaches \textit{a}

or "from below") can be denoted:



$\lim\_{x \to a^-}f(x)\ \quad \text{or} \quad \lim\_{x\,\uparrow\,a}\, f(x) \quad \text{or} \quad \lim\_{x \nearrow a}\,f(x) \quad \text{or} \quad \lim\_{x \underset{<}{\to} a}f(x)$



In probability theory it is common to use the short notation:



$f(x-)$ for the left limit and $f(x+)$ for the right limit.



The two one-sided limits exist and are equal if the limit of \textit{f}(\textit{x}) as

\textit{x} approaches \textit{a} exists. In some cases in which the limit



$$\lim\_{x\to a} f(x)\,$$



does not exist, the two one-sided limits nonetheless exist.

Consequently, the limit as \textit{x} approaches \textit{a} is sometimes called a

"two-sided limit".



In some cases one of the two one-sided limits exists and the other does

not, and in some cases neither exists. The right-sided limit can be

rigorously defined as



$$\forall\varepsilon > 0\;\exists \delta >0 \;\forall x \in I \;(0 < x - a < \delta \Rightarrow |f(x) - L|<\varepsilon),$$



and the left-sided limit can be rigorously defined as



$$\forall\varepsilon > 0\;\exists \delta >0 \;\forall x \in I \;(0 < a - x < \delta \Rightarrow |f(x) - L|<\varepsilon),$$



where \textit{I} represents some interval that is within

the domain of \textit{f}.



\paragraph{Examples}

One example of a function with different one-sided limits is the

following (cf. picture):



$$\lim\_{x \to 0^+}{1 \over 1 + 2^{-1/x}} = 1,$$



whereas



$$\lim\_{x \to 0^-}{1 \over 1 + 2^{-1/x}} = 0.$$



% [Image: Plot of the function $1 / (1 + 2^{-1/x})$]



\paragraph{Relation to topological definition of limit}

The one-sided limit to a point \textit{p} corresponds to the general definition

of limit, with the domain of the function restricted to one side, by

either allowing that the function domain is a subset of the topological

space, or by considering a one-sided subspace, including \textit{p}.

Alternatively, one may consider the domain with a half-open interval

topology.



\paragraph{Abel's theorem}



A noteworthy theorem treating one-sided limits of certain power series

at the boundaries of their intervals of convergence is Abel's theorem.



\paragraph{Resources}



\begin{itemize}
    \item \href{https://www.khanacademy.org/math/ap-calculus-ab/ab-limits-new/ab-1-3/v/one-sided-limits-from-graphs}{One-sided limits from
\end{itemize}
    graphs},

    Khan Academy



\paragraph{Licensing}



Content obtained and/or adapted from:



\begin{itemize}
    \item \href{https://en.wikipedia.org/wiki/One-sided\_limit}{One-sided limit,
\end{itemize}
    Wikipedia} under a CC

    BY-SA license

\subsection{Learning Objectives}

\begin{enumerate}
    \item \textbf{Define and Distinguish Limits:} Students will be able to distinguish between left-sided and right-sided limits of a function as the variable approaches a specified point.
    \item \textbf{Compute One-Sided Limits:} Students will be able to compute and express one-sided limits using proper notation and definitions.
    \item \textbf{Analyze Functions:} Students will be able to analyze functions to determine when one-sided limits exist while the two-sided limit does not, and explain the implications of these scenarios.
\end{enumerate}

\subsection{Assessment Instrument}

\begin{enumerate}
    \item \textbf{Part 1: Define and Distinguish Limits}
    \item \textbf{Multiple Choice:} What does the notation $\lim_{x \to a^-} f(x)$ represent?
    \item A. The limit of $f(x)$ as $x$ approaches $a$ from the right.
    \item B. The limit of $f(x)$ as $x$ approaches $a$ from the left.
    \item C. The value of $f(x)$ at $x = a$.
    \item D. The limit of $f(x)$ as $x$ approaches infinity.
    \item \textbf{True/False:} If $\lim_{x \to 2^-} f(x) = 5$ and $\lim_{x \to 2^+} f(x) = 5$, then the two-sided limit $\lim_{x \to 2} f(x)$ exists and is equal to 5. (True/False)
    \item \textbf{Short Answer:} Define left-sided and right-sided limits in the context of a function $f(x)$ as $x$ approaches a specific point $a$. How do these concepts differ from the two-sided limit? \textbf{Part 2: Compute One-Sided Limits}
    \item \textbf{Calculation:} Compute the left-sided limit for the function $f(x) = \frac{1}{x - 3}$ as $x$ approaches 3 from the left. Express your answer using proper notation.
    \item \textbf{Short Answer:} Find the right-sided limit of the function $g(x) = x^2 - 4$ as $x$ approaches -2 from the right. Show your calculations and express your answer using proper notation.
    \item \textbf{Multiple Choice:} Determine the one-sided limit of the function $h(x) = \sqrt{x - 1}$ as $x$ approaches 1 from the right.
    \item A. 0
    \item B. 1
    \item C. Undefined
    \item D. -1 \textbf{Part 3: Analyze Functions}
    \item \textbf{True/False:} If a function $f(x)$ has a left-sided limit of 3 and a right-sided limit of 5 as $x$ approaches a point $a$, then the two-sided limit $\lim_{x \to a} f(x)$ exists. (True/False) Explain why or why not.
    \item \textbf{Short Answer:} Analyze the function $f(x) = \frac{1}{x - 2}$ and explain a scenario where the one-sided limits exist while the two-sided limit does not. Include the implications of this situation.
    \item \textbf{Multiple Choice:} For the piecewise function defined as: $$ f(x) = \begin{cases} x + 2 & \text{for } x < 1 \   3 & \text{for } x \geq 1 \end{cases} $$ What is the value of the left-sided limit and the right-sided limit as $x$ approaches 1?
    \item A. Left-sided limit = 3, Right-sided limit = 2
    \item B. Left-sided limit = 3, Right-sided limit = 3
    \item C. Left-sided limit = 2, Right-sided limit = 3
    \item D. Left-sided limit = 2, Right-sided limit = 2
\end{enumerate}

%%===============================================================
\section{Limits Involving Infinity}
\label{sec:limits-involving-infinity}
%%===============================================================

\paragraph{Limits involving infinity}



\paragraph{Limits at infinity}



Let $S\subseteq\mathbb{R}$, $x\in S$ and $f:S\mapsto\mathbb{R}$.



The limit of \textit{f} as \textit{x} approaches infinity is \textit{L}, denoted



$$\lim\_{x \to \infty}f(x) = L,$$



means that for all $\varepsilon > 0$, there exists \textit{c} such that

$|f(x) - L| < \varepsilon$ whenever \textit{x} \> \textit{c}. Or, symbolically:



$$\forall \varepsilon > 0 \; \exists c \; \forall x > c :\; |f(x) - L| < \varepsilon.$$



Similarly, the limit of \textit{f} as \textit{x} approaches negative infinity is \textit{L}, denoted



$$\lim\_{x \to -\infty}f(x) = L,$$



means that for all $\varepsilon > 0$ there exists \textit{c} such that

$|f(x) - L| < \varepsilon$ whenever \textit{x} < \textit{c}. Or, symbolically:



$$\forall \varepsilon > 0 \; \exists c \; \forall x < c :\; |f(x) - L| < \varepsilon.$$



% [Image: The limit of this function at infinity

exists.]



For example,



$$\lim\_{x \to -\infty}e^x = 0. \,$$



\paragraph{Infinite limits}



For a function whose values grow without bound, the function diverges

and the usual limit does not exist. However, in this case one may

introduce limits with infinite values. Let $S\subseteq\mathbb{R}$,

$x\in S$ and $f:S\mapsto\mathbb{R}$. The statement **the limit of \textit{f} as

\textit{x} approaches \textit{a} is infinity**, denoted



$$\lim\_{x \to a} f(x) = \infty,$$



means that for all $N > 0$ there exists $\delta > 0$ such that

$f(x) > N$ whenever $0 < |x - a| < \delta.$



These ideas can be combined in a natural way to produce definitions for

different combinations, such as



$$\lim\_{x \to \infty} f(x) = \infty, \lim\_{x \to a^+}f(x) = -\infty.$$



For example,



$$\lim\_{x \to 0^+} \ln x = -\infty.$$



Limits involving infinity are connected with the concept of asymptotes.



These notions of a limit attempt to provide a metric space

interpretation to limits at infinity. In fact, they are consistent with

the topological space definition of limit if



\begin{itemize}
    \item a neighborhood of -$\infty$ is defined to contain an interval \[-$\infty$, \textit{c})
\end{itemize}
    for some \textit{$c \in R$},

\begin{itemize}
    \item a neighborhood of $\infty$ is defined to contain an interval (\textit{c}, $\infty$\]
\end{itemize}
    where \textit{$c \in R$}, and

\begin{itemize}
    \item a neighborhood of \textit{$a \in R$} is defined in the normal way metric
\end{itemize}
    space \textit{R}.



In this case, \textit{$\overline{R}$} is a topological space and any function of the form \textit{f}: \textit{$X \rightarrow Y$} with

\textit{$X, Y \subseteq$} \textit{$\overline{R}$} is subject to the topological definition of a limit. Note that with this

topological definition, it is easy to define infinite limits at finite

points, which have not been defined above in the metric sense.



\paragraph{Alternative notation}



Many authors allow for the projectively extended real line to be used as

a way to include infinite values as well as extended real line. With

this notation, the extended real line is given as \textit{R} $\cup {-\infty, +\infty}$ and

the projectively extended real line is \textit{R} $\cup {\infty}$ where a neighborhood

of $\infty$ is a set of the form {\textit{x}: \|\textit{x}\| \> \textit{c}}. The advantage is that

one only needs three definitions for limits (left, right, and central)

to cover all the cases. As presented above, for a completely rigorous

account, we would need to consider 15 separate cases for each

combination of infinities (five directions: -$\infty$, left, central, right,

and +$\infty$; three bounds: -$\infty$, finite, or +$\infty$). There are also noteworthy

pitfalls. For example, when working with the extended real line,

$x^{-1}$ does not possess a central limit (which is normal):



$$\lim\_{x \to 0^{+}}{1\over x} = +\infty, \lim\_{x \to 0^{-}}{1\over x} = -\infty.$$



In contrast, when working with the projective real line, infinities

(much like 0) are unsigned, so, the central limit \textit{does} exist in that

context:



$$\lim\_{x \to 0^{+}}{1\over x} = \lim\_{x \to 0^{-}}{1\over x} = \lim\_{x \to 0}{1\over x} = \infty.$$





% [Image: Horizontal asymptote about

\textit{y} = 4]



In fact there are a plethora of conflicting formal systems in use. In

certain applications of numerical differentiation and integration, it

is, for example, convenient to have signed zeroes. A simple reason has

to do with the converse of $\lim\_{x \to 0^{-}}{x^{-1}} = -\infty$,

namely, it is convenient for $\lim\_{x \to -\infty}{x^{-1}} = -0$ to be

considered true. Such zeroes can be seen as an approximation to

infinitesimals.



\paragraph{Limits at infinity for rational functions}



There are three basic rules for evaluating limits at infinity for a

rational function \textit{f}(\textit{x}) = \textit{p}(\textit{x})/\textit{q}(\textit{x}): (where \textit{p} and \textit{q} are

polynomials):



\begin{itemize}
    \item If the degree of \textit{p} is greater than the degree of \textit{q}, then the
\end{itemize}
    limit is positive or negative infinity depending on the signs of the

    leading coefficients;

\begin{itemize}
    \item If the degree of \textit{p} and \textit{q} are equal, the limit is the leading
\end{itemize}
    coefficient of \textit{p} divided by the leading coefficient of \textit{q};

\begin{itemize}
    \item If the degree of \textit{p} is less than the degree of \textit{q}, the limit is 0.
\end{itemize}

If the limit at infinity exists, it represents a horizontal asymptote at

\textit{y} = \textit{L}. Polynomials do not have horizontal asymptotes; such

asymptotes may however occur with rational functions.



\paragraph{Resources}



\begin{itemize}
    \item \href{https://tutorial.math.lamar.edu/classes/calci/limitsatinfinityi.aspx}{Limits at Infinity Part
\end{itemize}
    I},

    Paul's Online Notes

\begin{itemize}
    \item \href{https://tutorial.math.lamar.edu/Classes/CalcI/LimitsAtInfinityII.aspx}{Limits at Infinity Part
\end{itemize}
    II},

    Paul's Online Notes

\begin{itemize}
    \item \href{https://www.khanacademy.org/math/ap-calculus-ab/ab-limits-new/ab-1-15/v/introduction-to-limits-at-infinity}{Introduction to limits at
\end{itemize}
    infinity},

    Khan Academy



\paragraph{Licensing}



Content obtained and/or adapted from:



\begin{itemize}
    \item \href{https://en.wikipedia.org/wiki/Limit\_of\_a\_function}{Limit of a function,
\end{itemize}
    Wikipedia} under

    a CC BY-SA license

\subsection{Learning Objectives}

\begin{enumerate}
    \item \textbf{Define and Compute Limits at Infinity}: Students will be able to apply the definition of limits at infinity to determine whether the limit of a function as $x$ approaches positive or negative infinity exists and find its value.
    \item \textbf{Analyze Infinite Limits} Students will be able to identify and interpret infinite limits where the function's values grow without bound, including limits approaching infinity or negative infinity.
    \item *Rules for Rational Functions:** Students will be able to use the rules for evaluating limits at infinity for rational functions to determine horizontal asymptotes and understand the behavior of the function at extreme values.
\end{enumerate}

\subsection{Assessment Instrument}

\begin{enumerate}
    \item \textbf{Part 1: Define and Compute Limits at Infinity}
    \item \textbf{Multiple Choice:} What is the limit of the function $f(x) = \frac{3x^2 - 2}{x^2 + 4}$ as $x$ approaches positive infinity?
    \item A. 3
    \item B. 0
    \item C. 1
    \item D. Negative infinity
    \item \textbf{True/False:} The limit of the function $g(x) = \frac{2x - 5}{x^2 + 1}$ as $x$ approaches negative infinity is 0. (True/False)
    \item \textbf{Short Answer:} Determine the limit of the function $h(x) = \sqrt{x^2 + 1}$ as $x$ approaches positive infinity. Explain the reasoning behind your answer. \textbf{Part 2: Analyze Infinite Limits}
    \item \textbf{Calculation:} Find the limit of the function $f(x) = \frac{1}{x - 3}$ as $x$ approaches 3 from the right and from the left. Discuss the behavior of the function as $x$ approaches 3.
    \item \textbf{Short Answer:} Describe the behavior of the function $g(x) = \frac{1}{x}$ as $x$ approaches positive infinity and negative infinity. What can be inferred about the limits?
    \item \textbf{Multiple Choice:} For the function $h(x) = x^2 + 2x + 1$, what is the behavior of $h(x)$ as $x$ approaches positive infinity?
    \item A. $h(x)$ approaches 1
    \item B. $h(x)$ approaches positive infinity
    \item C. $h(x)$ approaches negative infinity
    \item D. $h(x)$ approaches 2 \textbf{Part 3: Rules for Rational Functions}
    \item \textbf{Short Answer:} Use the rules for limits at infinity to find the horizontal asymptote of the function $f(x) = \frac{4x^3 - x + 2}{2x^3 + 5}$.
    \item \textbf{Multiple Choice:} Determine the horizontal asymptote of the rational function $g(x) = \frac{5x^2 + 3}{2x^2 - x + 1}$.
    \item A. $y = \frac{5}{2}$
    \item B. $y = 0$
    \item C. $y = \frac{5}{2}$
    \item D. $y = 1$
    \item \textbf{True/False:} The horizontal asymptote of the function $h(x) = \frac{3x^2 + 2x + 1}{x^2 - 4}$ is $y = 3$. (True/False) Explain the reasoning behind your answer.
\end{enumerate}

%%===============================================================
\section{Rates of Change}
\label{sec:rates-of-change}
%%===============================================================

In mathematics, a \textbf{rate} is the ratio between two related quantities

in different units.[1] If the denominator of the ratio is expressed as

a single unit of one of these quantities, and if it is assumed that this

quantity can be changed systematically (i.e., is an independent

variable), then the numerator of the ratio expresses the corresponding

\textit{rate of change} in the other (dependent) variable.



One common type of rate is "per unit of time", such as speed, heart

rate and flux. Ratios that have a non-time denominator include exchange

rates, literacy rates, and electric field (in volts per meter).



In describing the units of a rate, the word "per" is used to separate

the units of the two measurements used to calculate the rate (for

example a heart rate is expressed "beats per minute"). A rate defined

using two numbers of the same units (such as tax rates) or counts (such

as literacy rate) will result in a dimensionless quantity, which can be

expressed as a percentage (for example, the global literacy rate in 1998

was 80\%), fraction, or multiple.



Often \textit{rate} is a synonym of rhythm or frequency, a count per second

(i.e., hertz); e.g., radio frequencies, heart rates, or sample rates.



\paragraph{Introduction}



Rates and ratios often vary with time, location, particular element (or

subset) of a set of objects, etc. Thus they are often mathematical

functions.



A rate (or ratio) may often be thought of as an output-input ratio,

benefit-cost ratio, all considered in the broad sense. For example,

miles per hour in transportation is the output (or benefit) in terms of

miles of travel, which one gets from spending an hour (a cost in time)

of traveling (at this velocity).



A set of sequential indices i may be used to enumerate elements (or

subsets) of a set of ratios under study. For example, in finance, one

could define i by assigning consecutive integers to companies, to

political subdivisions (such as states), to different investments, etc.

The reason for using indices i, is so a set of ratios (i=0,N) can be

used in an equation so as to calculate a function of the rates such as

an average of a set of ratios. For example, the average velocity found

from the set of vi's mentioned above. Finding averages may involve

using weighted averages and possibly using the harmonic mean.



A ratio r=a/b has both a numerator "a" and a denominator "b". The

value of a and/or b may be a real number or integer. The inverse of a

ratio r is 1/r = b/a. A rate may be equivalently expressed as an inverse

of its value if the ratio of its units are also inverse. For example, 5

miles (mi) per kilowatt-hour (kWh) corresponds to 1/5 kWh/mi (or 200

Wh/mi).



Rates are relevant to many aspects of everyday life. For example: *How

fast are you driving?* The speed of car (often expressed in miles per

hour) is a rate. \textit{What interest does your savings account pay you?} The

amount of interest paid per year is a rate.



\paragraph{Rate of change}



Consider the case where the numerator $f$ of a rate is a function $f(a)$

where $a$ happens to be the denominator of the rate $\delta f/\delta a.$

A rate of change of $f$ with respect to $a$ (where $a$ is incremented by

$h$) can be formally defined in two ways:



$\begin{aligned}

    \mbox{Average rate of change} \&= \frac{f(a + h) - f(a)}{h} \\\    \mbox{Instantaneous rate of change} \&= \lim\_{h \to 0}\frac{f(a + h) - f(a)}{h}

    \end{aligned}$



where \textit{f(x)} is the function with respect to \textit{x} over the interval from

\textit{a} to \textit{a}+\textit{h}. An instantaneous rate of change is equivalent to a

derivative.



For example, the average speed of a car can be calculated using the

total distance travelled between two points, divided by the travel time,

while the instantaneous speed can be determined by viewing a

speedometer.



\paragraph{Resources}



\begin{itemize}
    \item \href{https://www.khanacademy.org/math/algebra/x2f8bb11595b61c86:functions/x2f8bb11595b61c86:average-rate-of-change/v/introduction-to-average-rate-of-change}{Introduction to average rate of
\end{itemize}
    change},

    Khan Academy



\paragraph{Licensing}



Content obtained and/or adapted from:



\begin{itemize}
    \item \href{https://en.wikipedia.org/wiki/Rate\_(mathematics}{Rate (mathematics),
\end{itemize}
    Wikipedia}) under a

    CC BY-SA license

\subsection{Learning Objectives}

\begin{enumerate}
    \item \textbf{Understand and Define Rates}: Students will be able to define a rate as the ratio between two related quantities with different units, and explain how rates can be represented with various units, including those involving time and non-time denominators.
    \item \textbf{Calculate and Interpret Rates of Change}: Students will be able to calculate both the average and instantaneous rates of change for a given function, using the formulas provided, and interpret these rates in practical contexts.
    \item \textbf{Apply Rates to Real-Life Scenarios}: Students will be able to apply the concept of rates to analyze real-life situations, such as speed or interest rates, and understand how these rates can be expressed and compared in different forms.
\end{enumerate}

\subsection{Assessment Instrument}

\begin{enumerate}
    \item \textbf{Part 1: Understand and Define Rates}
    \item \textbf{Multiple Choice:} Which of the following best describes a rate?
    \item A. A quantity with the same units in the numerator and denominator.
    \item B. A ratio that compares two quantities with the same units.
    \item C. A ratio that compares two quantities with different units.
    \item D. A constant value of a quantity over time.
    \item \textbf{True/False:} The rate of change of distance with respect to time, commonly known as speed, can be expressed in units such as miles per hour (mph) or meters per second (m/s). (True/False)
    \item \textbf{Short Answer:} Define the concept of a rate in your own words and provide an example where rates involve non-time denominators, such as price per unit or density. \textbf{Part 2: Calculate and Interpret Rates of Change}
    \item \textbf{Calculation:} For the function $f(x) = 3x^2 - 5x + 2$, calculate the average rate of change between $x = 1$ and $x = 4$. Show your work.
    \item \textbf{Short Answer:} Given the function $g(t) = 2t^3 - 4t + 1$, find the instantaneous rate of change at $t = 2$. Describe how you arrived at this rate.
    \item \textbf{Multiple Choice:} If the position of an object is given by the function $s(t) = 4t^2 + 3t$, what is the instantaneous rate of change of the position at $t = 3$?
    \item A. $15$
    \item B. $21$
    \item C. $27$
    \item D. $30$ \textbf{Part 3: Apply Rates to Real-Life Scenarios}
    \item \textbf{Short Answer:} A car travels 150 miles in 3 hours. Calculate the average speed of the car and express it in miles per hour. Explain the interpretation of this rate in a real-life context.
    \item \textbf{Multiple Choice:} If an investment grows from $1000 to $1200 over a year, what is the annual interest rate, expressed as a percentage? Assume simple interest for this calculation.
    \item A. 10\%
    \item B. 15\%
    \item C. 20\%
    \item D. 25\%
    \item \textbf{True/False:} The density of a material can be expressed as the ratio of mass to volume, such as grams per cubic centimeter (g/cm³). If two materials have the same density, they will have the same mass if they have the same volume. (True/False) Explain your reasoning.
\end{enumerate}

%%===============================================================
\section{Tangent Lines and Derivatives}
\label{sec:tangent-lines-and-derivatives}
%%===============================================================

In geometry, the \textbf{tangent line} (or simply \textbf{tangent}) to a plane

curve at a given point is the straight line that "just touches" the

curve at that point. Leibniz defined it as the line through a pair of

infinitely close points on the curve.[1] More precisely, a straight

line is said to be a tangent of a curve y = f(x) at a point x = c if the

line passes through the point (c, f(c)) on the curve and has slope

f'(c), where f' is the derivative of f. A similar definition applies

to space curves and curves in n-dimensional Euclidean space.



As it passes through the point where the tangent line and the curve

meet, called the \textbf{point of tangency}, the tangent line is "going in

the same direction" as the curve, and is thus the best straight-line

approximation to the curve at that point.



The tangent line to a point on a differentiable curve can also be

thought of as the graph of the affine function that best approximates

the original function at the given point.



Similarly, the \textbf{tangent plane} to a surface at a given point is the

plane that "just touches" the surface at that point. The concept of a

tangent is one of the most fundamental notions in differential geometry

and has been extensively generalized; see Tangent space.



The word "tangent" comes from the Latin tangere, "to touch".



% [Image: Tangent to a curve. The red line is tangential to the curve at the

point marked by a red

dot.]



\paragraph{Tangent line to a curve}



The intuitive notion that a tangent line "touches" a curve can be made

more explicit by considering the sequence of straight lines (secant

lines) passing through two points, A and B, those that lie on the

function curve. The tangent at A is the limit when point B approximates

or tends to A. The existence and uniqueness of the tangent line depends

on a certain type of mathematical smoothness, known as

"differentiability." For example, if two circular arcs meet at a sharp

point (a vertex) then there is no uniquely defined tangent at the vertex

because the limit of the progression of secant lines depends on the

direction in which "point B" approaches the vertex.



At most points, the tangent touches the curve without crossing it

(though it may, when continued, cross the curve at other places away

from the point of tangent). A point where the tangent (at this point)

crosses the curve is called an inflection point. Circles, parabolas,

hyperbolas and ellipses do not have any inflection point, but more

complicated curves do have, like the graph of a cubic function, which

has exactly one inflection point, or a sinusoid, which has two

inflection points per each period of the sine.



Conversely, it may happen that the curve lies entirely on one side of a

straight line passing through a point on it, and yet this straight line

is not a tangent line. This is the case, for example, for a line passing

through the vertex of a triangle and not intersecting it

otherwise---where the tangent line does not exist for the reasons

explained above. In convex geometry, such lines are called supporting

lines. 



% [Image: A tangent, a chord, and a secant to a

circle]



\paragraph{Analytical approach}



The geometrical idea of the tangent line as the limit of secant lines

serves as the motivation for analytical methods that are used to find

tangent lines explicitly. The question of finding the tangent line to a

graph, or the \textbf{tangent line problem,} was one of the central questions

leading to the development of calculus in the 17th century. In the

second book of his \textit{Geometry}, René Descartes said of the problem of

constructing the tangent to a curve, "And I dare say that this is not

only the most useful and most general problem in geometry that I know,

but even that I have ever desired to know".



\paragraph{Intuitive description}



Suppose that a curve is given as the graph of a function, \textit{y} =

\textit{f}(\textit{x}). To find the tangent line at the point \textit{p} = (\textit{a}, \textit{f}(\textit{a})),

consider another nearby point \textit{q} = (\textit{a} + \textit{h}, \textit{f}(\textit{a} + \textit{h})) on the

curve. The slope of the secant line passing through \textit{p} and \textit{q} is equal

to the difference quotient



$$\frac{f(a+h)-f(a)}{h}.$$



As the point \textit{q} approaches \textit{p}, which corresponds to making \textit{h} smaller

and smaller, the difference quotient should approach a certain limiting

value \textit{k}, which is the slope of the tangent line at the point \textit{p}. If

\textit{k} is known, the equation of the tangent line can be found in the

point-slope form:



$$y-f(a) = k(x-a).\,$$



\paragraph{More rigorous description}



To make the preceding reasoning rigorous, one has to explain what is

meant by the difference quotient approaching a certain limiting value

\textit{k}. The precise mathematical formulation was given by Cauchy in the

19th century and is based on the notion of limit. Suppose that the graph

does not have a break or a sharp edge at \textit{p} and it is neither plumb nor

too wiggly near \textit{p}. Then there is a unique value of \textit{k} such that, as

\textit{h} approaches 0, the difference quotient gets closer and closer to \textit{k},

and the distance between them becomes negligible compared with the size

of \textit{h}, if \textit{h} is small enough. This leads to the definition of the

slope of the tangent line to the graph as the limit of the difference

quotients for the function \textit{f}. This limit is the derivative of the

function \textit{f} at \textit{x} = \textit{a}, denoted \textit{f} '(\textit{a}). Using derivatives, the

equation of the tangent line can be stated as follows:



$$y=f(a)+f'(a)(x-a).\,$$



Calculus provides rules for computing the derivatives of functions that

are given by formulas, such as the power function, trigonometric

functions, exponential function, logarithm, and their various

combinations. Thus, equations of the tangents to graphs of all these

functions, as well as many others, can be found by the methods of

calculus.



% [Image: At each point, the moving line is always tangent to the curve.

Its slope is the derivative; green marks positive derivative, red marks

negative derivative and black marks zero derivative. The point (x,y) =

(0,1) where the tangent intersects the curve, is not a max, or a min,

but is a point of

inflection.]



\paragraph{How the method can fail}



Calculus also demonstrates that there are functions and points on their

graphs for which the limit determining the slope of the tangent line

does not exist. For these points the function \textit{f} is

\textit{non-differentiable}. There are two possible reasons for the method of

finding the tangents based on the limits and derivatives to fail: either

the geometric tangent exists, but it is a vertical line, which cannot be

given in the point-slope form since it does not have a slope, or the

graph exhibits one of three behaviors that precludes a geometric

tangent.



The graph \textit{$y = x^{\frac{1}{3}$} illustrates the first possibility: here the

difference quotient at \textit{a} = 0 is equal to \textit{$h^{\frac{1}{3}}$} \textit{/h} = \textit{$h^{\frac{-2}{3}$},

which becomes very large as \textit{h} approaches 0. This curve has a tangent

line at the origin that is vertical.



The graph \textit{$y = x^{\frac{2}{3}$} illustrates another possibility: this graph has

a cusp at the origin. This means that, when h approaches 0, the

difference quotient at a = 0 approaches plus or minus infinity depending

on the sign of x. Thus both branches of the curve are near to the half

vertical line for which y=0, but none is near to the negative part of

this line. Basically, there is no tangent at the origin in this case,

but in some context one may consider this line as a tangent, and even,

in algebraic geometry, as a double tangent.



The graph y = |x| of the absolute value function consists of two

straight lines with different slopes joined at the origin. As a point q

approaches the origin from the right, the secant line always has

slope 1. As a point q approaches the origin from the left, the secant

line always has slope -1. Therefore, there is no unique tangent to the

graph at the origin. Having two different (but finite) slopes is called

a corner.



Finally, since differentiability implies continuity, the contrapositive

states discontinuity implies non-differentiability. Any such jump or

point discontinuity will have no tangent line. This includes cases where

one slope approaches positive infinity while the other approaches

negative infinity, leading to an infinite jump discontinuity.



\paragraph{Equations}



When the curve is given by \textit{y} = \textit{f}(\textit{x}) then the slope of the tangent

is $\frac{dy}{dx},$ so by the point--slope formula the equation of the

tangent line at (\textit{X}, \textit{Y}) is



$$y-Y=\frac{dy}{dx}(X) \cdot (x-X)$$ 



where (\textit{x}, \textit{y}) are the coordinates of any point on the tangent line, and where the derivative

is evaluated at $x=X.$



When the curve is given by \textit{y} = \textit{f}(\textit{x}), the tangent line's equation

can also be found by using polynomial division to divide $f \, (x)$ by

$(x-X)^2$; if the remainder is denoted by $g(x),$ then the equation of

the tangent line is given by



$$y=g(x).$$



When the equation of the curve is given in the form \textit{f}(\textit{x}, \textit{y}) = 0

then the value of the slope can be found by implicit differentiation,

giving



$$\frac{dy}{dx}=-\frac{\frac{\partial f}{\partial x}}{\frac{\partial f}{\partial y}}.$$



The equation of the tangent line at a point (\textit{X},\textit{Y}) such that

\textit{f}(\textit{X},\textit{Y}) = 0 is then



$$\frac{\partial f}{\partial x}(X,Y) \cdot (x-X)+\frac{\partial f}{\partial y}(X,Y) \cdot (y-Y)=0.$$



This equation remains true if $\frac{\partial f}{\partial y}(X,Y) = 0$

but $\frac{\partial f}{\partial x}(X,Y) \neq 0$ (in this case the slope

of the tangent is infinite). If

$\frac{\partial f}{\partial y}(X,Y) =  \frac{\partial f}{\partial x}(X,Y)  =0,$

the tangent line is not defined and the point (\textit{X},\textit{Y}) is said to be

singular.



For algebraic curves, computations may be simplified somewhat by

converting to homogeneous coordinates. Specifically, let the homogeneous

equation of the curve be \textit{g}(\textit{x}, \textit{y}, \textit{z}) = 0 where \textit{g} is a

homogeneous function of degree \textit{n}. Then, if (\textit{X}, \textit{Y}, \textit{Z}) lies on the

curve, Euler's theorem implies



$$\frac{\partial g}{\partial x} \cdot X +\frac{\partial g}{\partial y} \cdot Y+\frac{\partial g}{\partial z} \cdot Z=ng(X, Y, Z)=0.$$



It follows that the homogeneous equation of the tangent line is



$$\frac{\partial g}{\partial x}(X,Y,Z) \cdot x+\frac{\partial g}{\partial y}(X,Y,Z) \cdot y+\frac{\partial g}{\partial z}(X,Y,Z) \cdot z=0.$$



The equation of the tangent line in Cartesian coordinates can be found

by setting \textit{z}=1 in this equation.



To apply this to algebraic curves, write \textit{f}(\textit{x}, \textit{y}) as



$$f=u\_n+u\_{n-1}+\dots+u\_1+u\_0\,$$ 



where each \textit{$u_r$} is the sum of all terms of degree \textit{r}. The homogeneous equation of the curve is then



$$g=u\_n+u\_{n-1}z+\dots+u\_1 z^{n-1}+u\_0 z^n=0.\,$$ 



Applying the equation above and setting \textit{z}=1 produces



$$\frac{\partial f}{\partial x}(X,Y) \cdot x + \frac{\partial f}{\partial y}(X,Y) \cdot y + \frac{\partial g}{\partial z}(X,Y,1) =0$$



as the equation of the tangent line. The equation in this form is often

simpler to use in practice since no further simplification is needed

after it is applied.



If the curve is given parametrically by



$$x=x(t),\quad y=y(t)$$ 



then the slope of the tangent is



$$\frac{dy}{dx}=\frac{\frac{dy}{dt}}{\frac{dx}{dt}}$$ 



giving the equation for the tangent line at $\, t=T, \, X=x(T), \, Y=y(T)$ as



$$\frac{dx}{dt}(T) \cdot (y-Y)=\frac{dy}{dt}(T) \cdot (x-X).$$



If $\frac{dx}{dt}(T)= \frac{dy}{dt}(T) =0,$ the tangent line is not

defined. However, it may occur that the tangent line exists and may be

computed from an implicit equation of the curve.



\paragraph{Normal line to a curve}



The line perpendicular to the tangent line to a curve at the point of

tangency is called the \textit{normal line} to the curve at that point. The

slopes of perpendicular lines have product -1, so if the equation of the

curve is \textit{y} = \textit{f}(\textit{x}) then slope of the normal line is



$$-\frac{1}{\frac{dy}{dx}}$$ 



and it follows that the equation of the normal line at (X, Y) is



$$(x-X)+\frac{dy}{dx}(y-Y)=0.$$ 



Similarly, if the equation of the curve has the form \textit{f}(\textit{x}, \textit{y}) = 0 then the equation of the normal line is

given by



$$\frac{\partial f}{\partial y}(x-X)-\frac{\partial f}{\partial x}(y-Y)=0.$$



If the curve is given parametrically by



$$x=x(t),\quad y=y(t)$$ 



then the equation of the normal line is



$$\frac{dx}{dt}(x-X)+\frac{dy}{dt}(y-Y)=0.$$



\paragraph{Angle between curves}



The angle between two curves at a point where they intersect is defined

as the angle between their tangent lines at that point. More

specifically, two curves are said to be tangent at a point if they have

the same tangent at a point, and orthogonal if their tangent lines are

orthogonal.



\paragraph{Multiple tangents at a point}



The formulas above fail when the point is a singular point. In this case

there may be two or more branches of the curve that pass through the

point, each branch having its own tangent line. When the point is the

origin, the equations of these lines can be found for algebraic curves

by factoring the equation formed by eliminating all but the lowest

degree terms from the original equation. Since any point can be made the

origin by a change of variables (or by translating the curve) this gives

a method for finding the tangent lines at any singular point.



For example, the equation of the limaçon trisectrix shown to the right is



$$(x^2+y^2-2ax)^2=a^2(x^2+y^2).\,$$



Expanding this and eliminating all but terms of degree 2 gives



$$a^2(3x^2-y^2)=0\,$$ 



which, when factored, becomes



$$y=\pm\sqrt{3}x.$$



So these are the equations of the two tangent lines through the origin.



When the curve is not self-crossing, the tangent at a reference point

may still not be uniquely defined because the curve is not

differentiable at that point although it is differentiable elsewhere. In

this case the left and right derivatives are defined as the limits of

the derivative as the point at which it is evaluated approaches the

reference point from respectively the left (lower values) or the right

(higher values). For example, the curve y = |x| is not differentiable

at x = 0: its left and right derivatives have respective slopes -1 and

1; the tangents at that point with those slopes are called the left and

right tangents.



Sometimes the slopes of the left and right tangent lines are equal, so

the tangent lines coincide. This is true, for example, for the curve \textit{$y = x^{\frac{2}{3}$}, for which both the left and right derivatives at \textit{x} = 0

are infinite; both the left and right tangent lines have equation \textit{x} = 0.



% [Image: The limaçon trisectrix: a curve with two tangents at the

origin.]



\paragraph{Licensing}



Content obtained and/or adapted from:



\begin{itemize}
    \item \href{https://en.wikipedia.org/wiki/Tangent}{Tangent, Wikipedia} under a
\end{itemize}
    CC BY-SA license

\subsection{Learning Objectives}

\begin{enumerate}
    \item \textbf{Define and Identify Tangent Lines and Planes}: Students will be able to define a tangent line to a curve and a tangent plane to a surface, and identify them based on their geometric properties and the point of tangency.
    \item \textbf{Calculate the Tangent Line Using Derivatives}: Students will be able to calculate the equation of the tangent line to a curve at a given point using derivatives and the point-slope form of a line, and apply these calculations to various functions.
    \item \textbf{Analyze Tangents in Different Contexts}: Students will be able to analyze and describe the behavior of tangent lines and planes in different scenarios, including identifying points of inflection, computing tangent lines for implicitly defined curves, and understanding situations with multiple tangents or non-differentiable points.
\end{enumerate}

\subsection{Assessment Instrument}

\begin{enumerate}
    \item \textbf{Part 1: Define and Identify Tangent Lines and Planes}
    \item \textbf{Multiple Choice:} What is the defining characteristic of a tangent line to a curve at a point $P$?
    \item A. It intersects the curve at exactly two points.
    \item B. It is parallel to the curve at point $P$.
    \item C. It touches the curve at $P$ and has the same slope as the curve at that point.
    \item D. It is perpendicular to the curve at point $P$.
    \item \textbf{True/False:} A tangent plane to a surface at a given point is the plane that intersects the surface only at that point and is perpendicular to the surface's normal vector at that point. (True/False)
    \item \textbf{Short Answer:} Define a tangent line to a curve at a point and a tangent plane to a surface at a point. Describe their geometric properties and how they relate to the points of tangency. \textbf{Part 2: Calculate the Tangent Line Using Derivatives}
    \item \textbf{Calculation:} Find the equation of the tangent line to the curve $f(x) = x^3 - 4x$ at the point where $x = 2$. Use the derivative to determine the slope and then write the equation in point-slope form.
    \item \textbf{Short Answer:} Given the function $g(x) = \sin(x)$, find the tangent line to the curve at $x = \frac{\pi}{4}$. Calculate the slope using the derivative and provide the equation of the tangent line.
    \item \textbf{Multiple Choice:} For the function $h(x) = e^x$, what is the slope of the tangent line at $x = 1$?
    \item A. $e$
    \item B. $e^1$
    \item C. $e^0$
    \item D. $e^2$ \textbf{Part 3: Analyze Tangents in Different Contexts}
    \item \textbf{Short Answer:} Describe how you would find the tangent line to an implicitly defined curve given by $x^2 + y^2 = 25$ at the point $(3, 4)$. Explain the steps and the role of implicit differentiation.
    \item \textbf{Multiple Choice:} For the function $f(x) = x^2 \sin(x)$, identify the behavior of the tangent line at the point where $x = 0$. What can be inferred about the tangent line at this point?
    \item A. The tangent line has a slope of 0.
    \item B. The tangent line is vertical.
    \item C. The function $f(x)$ has a point of inflection at $x = 0$.
    \item D. The tangent line is horizontal with a slope of 1.
    \item \textbf{True/False:} A function can have multiple tangent lines at a single point if the function is not differentiable at that point. (True/False) Explain your reasoning and provide an example or counterexample if applicable.
\end{enumerate}

%%===============================================================
\section{Techniques for Finding Derivatives}
\label{sec:techniques-for-finding-derivatives}
%%===============================================================

\paragraph{Techniques of Differentiation}



The definition of a derivative,

$f'(x)=\lim\_{{\Delta}x\rightarrow 0} \frac{f(x+{\Delta}x)-f(x)}{{\Delta}x}$

is not the only method of finding a derivative. This section provides

three different techniques to help find derivatives. These techniques

are often quicker than the definition of a derivative, depending on the

given equation. The last part of the section presents how to calculate

higher order derivatives.



\paragraph{The Power Rule}



The first technique allows you to find a derivative of a summation of

multiple terms.



This is known as \textbf{The Power Rule}.



Power Rule only applies in certain conditions.



\paragraph{Theorem}



If \textit{n} is a rational number, then the function

$f(x) = x^n$ is differentiable and



$$\frac {d}{dx} [x^n] = nx^{n-1}.$$



For f to be differentiable at x=0, n must be a number such that

$x^{n-1}$ is defined on an interval containing 0.



\paragraph{Proof}



$\begin{aligned}

\frac {d}{dx}[x^n]\&= \lim\_{{\Delta}x\rightarrow 0} \frac {(x + {\Delta}x)^n - x^n}{{\Delta}x} \\\\&= \lim\_{{\Delta}x\rightarrow 0} \frac{x^n + nx^{n-1} ({\Delta}x) + \frac {n(n-1)x^{n-2}}{2}({\Delta}x)^2 + ..... + ({\Delta}x)^n - x^n}{{\Delta}x} \\\\&= \lim\_{{\Delta}x\rightarrow 0} \left[nx^{n-1} + \frac{n(n-1)x^{n-2}}{2}({\Delta}x) + ..... + ({\Delta}x)^{n-1}\right] \\\\&= nx^{n-1} + 0 + ..... + 0 \\\\&= nx^{n-1} {\square}.

\end{aligned}$



\paragraph{Example}



The following is an example of the Power Rule:



$$f(x)=3x^2 + 2x + 1$$



Normally you would work through the definition of a derivative,



$$f'(x)=\lim\_{{\Delta}x\rightarrow 0} \frac{f(x+{\Delta}x)-f(x)}{{\Delta}x}$$



but instead you can apply the power rule to find your derivative.



First, take a look at the first term of the equation, $3x^2.$



To find $f'(x)$ first drop the power on the variable of this term down

to the front of the equation, $(2) * (3x^2).$



Next, subtract one from the power that was dropped down in front of the

equation, $(2) * (3x^1).$



Or, $(2)(3)x = 6x.$ So now we have the first part of the derivative,

$6x.$



Now, perform this operation to the second and the third terms of the

equation.



For the second part, $2x^1$



$(1) * (2x^1)$



$(1)(2)(x^0) = 2.$



This shows that the second part is $2.$



Now for the third part we notice that there is no variable, $x.$



There are two ways to look at this last part, either: one, there is no

variable so the part disappears, or two, since there is no variable, we

multiply the whole piece by 0, since that would be the power of the

variable.



Either way, there is no longer a third piece.



Now, if we put the two parts together, we have our answer



$$f'(x) = 6x + 2.$$



Here is another example of this method in action



$f(x) = 3x^3 - 4x^2 + 3x - 4$



$f'(x) = (3)(3x^{3-1}) - (2)(4x^{2-1}) + 3x^{1-1}$



$f'(x) = 9x^2 - 8x + 3.$



Here are some practice problems.



$f(x) = x^2 - 7x + 13$



$h(x) = 2x^3 + 2x - 3$



$g(x) = 4x^4 - 7x^3 + 2x^2 + 5x - 4$



\paragraph{The Product Rule}



The next technique is called \textbf{The Product Rule}.



As previously stated, not all derivatives can be found through the power

rule.



While it is a nice way to find many derivatives, other methods exist

depending on the equation you are currently solving.



The Power Rule can be seen as the sum of two separate functions, and

their respective derivative is the sum of the individual derivatives.



The product rule introduces how to solve the product of two functions.



\paragraph{Theorem}



The product of two differentiable functions $f(x)$ and $g(x)$ is itself

differentiable.



The derivative of $f(x)g(x)$ is $f(x) * g'(x) + g(x) * f'(x).$



\paragraph{Proof}



$\begin{aligned}

\frac{d}{dx} [f(x)g(x)] \&= \lim\_{{\Delta}x\rightarrow 0} \frac {f(x + {\Delta}x) \textit{ g(x + {\Delta}x) - f(x)g(x)}{{\Delta}x} \\\\&= \lim\_{{\Delta}x\rightarrow 0} \frac {f(x + {\Delta}x) } g(x + {\Delta}x)- f(x + {\Delta}x)g(x) + f(x + {\Delta}x)g(x) - f(x)g(x)}{{\Delta}x} \\\\&= \lim\_{{\Delta}x\rightarrow 0}\left [ f(x + {\Delta}x) \textit{ \frac {g(x + {\Delta}x) - g(x)}{{\Delta}x} + g(x) } \frac {f(x + {\Delta}x) - f(x)}{{\Delta}x}\right] \\\\&= \lim\_{{\Delta}x\rightarrow 0} \left[ f(x + {\Delta}x) \textit{ \frac {g(x + {\Delta}x) - g(x)}{{\Delta}x}\right] + \lim\_{{\Delta}x\rightarrow 0} \left[g(x) } \frac {f(x + {\Delta}x) - f(x)}{{\Delta}x}\right] \\\\&= \lim\_{{\Delta}x\rightarrow 0} f(x + {\Delta}x) \textit{ \lim\_{{\Delta}x\rightarrow 0} \frac{g(x + {\Delta}x) - g(x)}{{\Delta}x} + \lim\_{{\Delta}x\rightarrow 0} g(x) } \lim\_{{\Delta}x\rightarrow 0} \frac{f(x + {\Delta}x) - f(x)}{{\Delta}x} \\\\&= f(x)g'(x) + g(x)f'(x) {\square}.

\end{aligned}$



\paragraph{Example}



For example, take the equation $h(x)=(x^2 + 1)(x^2 - 2x - 1).$



To start, think of the derivative of each part. Let $f(x) = x^2 + 1$ and

let $g(x) = x^2 - 2x -1.$



$f'(x) = 2x$ and $g'(x) = 2x - 2.$



Now apply each part to the previous theorem.



$h'(x) = f(x) * g'(x) + g(x) * f'(x)$



$h'(x) = (x^2 + 1) * (2x-2) + (x^2 - 2x - 1) * (2x).$



Now all we need to do is simplify.



$h'(x) = (2x^3 + 2x -2x^2 -2) + (2x^3 - 4x^2 - 2x)$



$h'(x) = (2x^3 - 2x^2 + 2x - 2) + (2x^3 - 4x^2 - 2x)$



$h'(x) = 4x^3 - 6x^2 - 2.$



Now, if product rule doesn't seem to appeal to you, when possible, you

can just simplify the original equation.



If you multiple the parts by each other and then find the derivative,

you are simply doing the Power Rule.



Sometimes this may be easier, but in some cases you may be unable to

multiply the parts together.



Here is an example of changing a product rule.



$h(x) = (x^2 + 1)(x^2 - 2x - 1)$



$h(x) = x^4 + x^2 -2x^3 - 2x - x^2 - 1$



$h(x) = x^4 - 2x^3 - 2x - 1.$



Then using the power rule,



$h'(x) = 4x^3 - 6x^2 - 2.$



As you can see, both methods yield the same answer.



Also, product rule can be expanded beyond two pieces.



This would look like:



$\frac{d}{dx}[f(x)g(x)h(x)] = f'(x)g(x)h(x) + f(x)g'(x)h(x) + f(x)g(x)h'(x).$



For further work, here are some practice problems.



$h(x) = (x^2-2)(3x+1)$



$w(x) = (3x + 2)(2x^2 +4x)$



$p(y) = (y^2 -3y)(5y-4)(2y^2 + 1)$



\paragraph{The Quotient Rule}



The next technique is called \textbf{The Quotient Rule}.



Before we learned what to do with the product of two functions, quotient

rule teaches us what to do with a function divided by another function.



\paragraph{Theorem}



The quotient of two differentiable functions, $\frac {f(x)}{g(x)}$ is

itself differentiable for all$g(x) \ne 0.$



This derivative is given by the bottom function multiplied by the

derivative of the top function, minus the top function multiplied by the

derivative of the bottom function, all divided by the bottom function

squared.



Or, $\frac{d}{dx} \left[\frac{f(x)}{g(x)}\right] = \frac{g(x)f'(x) - f(x)g'(x)}{[g(x)]^2}.$



\paragraph{Proof}



$\begin{aligned}

\frac{d}{dx} \left[\frac{f(x)}{g(x)}\right] \&= \lim\_{{\Delta}x\rightarrow 0} \frac {\frac {f(x + {\Delta}x)}{g(x + {\Delta}x)} - \frac {f(x)}{g(x)}}{{\Delta}x} \\\\&= \lim\_{{\Delta}x\rightarrow 0} \frac {g(x)f(x + {\Delta}x) - f(x)g(x + {\Delta}x)}{{\Delta}xg(x)g(x + {\Delta}x)} \\\\&= \lim\_{{\Delta}x\rightarrow 0} \frac {g(x)f(x + {\Delta}x) - f(x)g(x) + f(x)g(x) - f(x)g(x + {\Delta}x)}{{\Delta}xg(x)g(x + {\Delta}x)} \\\\&= \frac {\lim\_{{\Delta}x\rightarrow 0} \frac {g(x)[f(x + {\Delta}x) - f(x)]}{{\Delta}x} - \lim\_{{\Delta}x\rightarrow 0} \frac {f(x)[g(x + {\Delta}x) - g(x)]}{{\Delta}x}}{\lim\_{{\Delta}x\rightarrow 0}  [g(x)g(x + {\Delta}x)]} \\\\&= \frac {g(x)[\lim\_{{\Delta}x\rightarrow 0} \frac {f(x + {\Delta}x) - f(x)}{{\Delta}x}] - f(x)[\lim\_{{\Delta}x\rightarrow 0} \frac {g(x + {\Delta}x) - g(x)}{{\Delta}x}]}{\lim\_{{\Delta}x\rightarrow 0}[g(x)g(x + {\Delta}x)]} \\\\&= \frac {g(x)f'(x) - f(x)g'(x)}{[g(x)]^2} {\square}.

\end{aligned}$



\paragraph{Example}



For an example, let



$h(x) = \frac {3x - 2}{x^2 + 1}$



and let



$f(x) = 3x - 2$



$g(x) = x^2 + 1.$



Start by finding the derivatives of each individual piece,



$f'(x) = 3$



$g'(x) = 2x.$



Now all we need do is plug each piece into the equation, so we have



$\begin{aligned}

h'(x) \&= \frac {(x^2 + 1)(3) - (3x - 2)(2x)}{(x^2 + 1)^2} \\\\&= \frac {3x^2 + 3 - (6x^2 - 4x)}{x^4 + 2x^2 + 1} \\\\&= \frac {-3x^2 + 4x + 3}{x^4 + 2x^2 + 1}.

\end{aligned}$



In many cases, the derivative will not be able to be simplified further

due to the division.



Sometimes, you may be able to perform some sort of division to turn a

quotient rule into a product rule.



If not, you may be able to set the bottom function to the power of (-1)

and then perform a product rule.



The only problem with this is that, in order to deal with the negative

power, the Chain Rule must be used, which is covered in another section.



Even so, if quotient rule doesn't appeal to you, there are ways to

change it.



Here are some practice problems for quotient rule:



$h(x) = \frac {x^2 + 2}{x}$



$z(x) = \frac {3x - 1}{x^2 + 2x + 1}$



$r(x) = \frac {2x^2 - 2x + 3}{3}$



\paragraph{Higher Order Derivatives}



The last thing to mention is \textbf{Higher Order Derivatives}.



It is possible to find a second, even third and fourth, derivative of a

function.



It is just finding the derivative of a derivative.



You denote higher order derivatives with two or more prime symbols, for

example,



$f''(x)$ or even, $g'''(x).$



Since the number of prime symbols can get out of hand quickly, many

mathematicians stop at three, in favor of using small numbers,



$f^4(x)$ or any number higher, $f^n(x).$



An example of a higher order derivative is:



$f(x) = x^4 + x^3 + 2x^2 - 3x + 1$



$f'(x) = 4x^3 + 3x^2 + 4x - 3$



$f''(x) = 12x^2 + 6x + 4$



$f'''(x) = 24x + 6$



$f^4(x) = 24.$



Some practice problems are:



$f(x) = 3x^3 + 2x^2 - 3$



$g(x) = (3x^2 - 1)(6x + 2)$



$r(x) = \frac {2x + 1}{3x}$



\paragraph{Licensing}



\begin{itemize}
    \item \href{https://en.wikibooks.org/wiki/High\_School\_Calculus/Techniques\_of\_Differentiation}{Techniques of Differentiation, Wikibooks: High School
\end{itemize}
    Calculus}

    under a CC BY-SA license

\subsection{Learning Objectives}

\begin{enumerate}
    \item \textbf{Differentiate Various Functions}: Students will be able to apply the Power Rule, Product Rule, and Quotient Rule to find the derivatives of polynomial, product, and quotient functions. This includes computing first and higher order derivatives accurately and simplifying the results as needed.
    \item \textbf{Verify and Interpret Derivatives}: Students will be able to verify the correctness of derivatives obtained through different rules by comparing them with derivatives computed using the definition of a derivative. They will also be able to interpret the significance of these derivatives in the context of the original functions, including understanding the implications of higher order derivatives for function behavior.
    \item \textbf{Notate and Simplify Higher Order Derivatives}: Students will be able to correctly notate higher order derivatives using both prime and superscript notation. They will also be able to compute and simplify higher order derivatives, analyzing the results to understand the function's curvature and motion.
\end{enumerate}

\subsection{Assessment Instrument}

\begin{enumerate}
    \item \textbf{Part 1: Differentiate Various Functions}
    \item \textbf{Calculation:} Find the derivative of the function $f(x) = 5x^4 - 3x^2 + 2x - 7$ using the Power Rule. Simplify your result.
    \item \textbf{Short Answer:} Compute the derivative of the function $g(x) = (3x^2 - x)(2x + 1)$ using the Product Rule. Show all steps and provide the simplified form of the derivative.
    \item \textbf{Multiple Choice:} Determine the derivative of the function $h(x) = \frac{x^3 - 2x}{x^2 + 1}$ using the Quotient Rule. Which of the following is the correct derivative?
    \item A. $\frac{(3x^2 - 2)(x^2 + 1) - (x^3 - 2x)(2x)}{(x^2 + 1)^2}$
    \item B. $\frac{(3x^2 - 2)(x^2 + 1) + (x^3 - 2x)(2x)}{(x^2 + 1)^2}$
    \item C. $\frac{(3x^2 - 2)(x^2 + 1) - (x^3 - 2x)(x^2)}{(x^2 + 1)^2}$
    \item D. $\frac{(3x^2 - 2)(x^2 + 1) + (x^3 - 2x)(x^2)}{(x^2 + 1)^2}$ \textbf{Part 2: Verify and Interpret Derivatives}
    \item \textbf{Short Answer:} Verify the derivative of the function $f(x) = x^3 - 4x + 1$ by using both the Power Rule and the definition of the derivative. Show your work and explain how you confirm the correctness of your derivative.
    \item \textbf{Multiple Choice:} For the function $g(x) = \cos(x) + \sin(x)$, what is the interpretation of its first derivative $g'(x)$ in terms of the behavior of the original function $g(x)$?
    \item A. The first derivative represents the concavity of $g(x)$.
    \item B. The first derivative represents the slope or rate of change of $g(x)$.
    \item C. The first derivative is the integral of $g(x)$.
    \item D. The first derivative shows where $g(x)$ is undefined.
    \item \textbf{True/False:} If a function $h(x)$ has a second derivative $h''(x)$ that is positive at a certain point, this indicates that the function $h(x)$ is concave up at that point. (True/False) Explain why or why not. \textbf{Part 3: Notate and Simplify Higher Order Derivatives}
    \item \textbf{Short Answer:} Notate and compute the second derivative of the function $f(x) = x^4 - 6x^2 + 3x$. Provide the result in both prime notation ($f''(x)$) and superscript notation ($f^{(2)}(x)$).
    \item \textbf{Multiple Choice:} For the function $g(x) = e^x \cdot \sin(x)$, what is the third derivative $g'''(x)$? Select the correct expression.
    \item A. $e^x \cdot (\sin(x) + 3 \cos(x))$
    \item B. $e^x \cdot (\sin(x) - 3 \cos(x))$
    \item C. $e^x \cdot (\sin(x) + 3 \cos(x) - e^x \sin(x))$
    \item D. $e^x \cdot (\sin(x) - \cos(x))$
    \item \textbf{True/False:} If the third derivative of a function is zero at a certain point, it implies that the function has a point of inflection at that point. (True/False) Explain the implications of a zero third derivative in relation to the function's curvature and motion.
\end{enumerate}

%%===============================================================
\section{Derivatives of Products and Quotients}
\label{sec:derivatives-of-products-and-quotients}
%%===============================================================

\paragraph{Product Rule}



Let $f(x)$ and $g(x)$ be differentiable functions. Then, the derivative

of their product, $f(x)g(x),$ is



$$\frac{d}{dx}(f(x)g(x)) = f'(x)g(x) + f(x)g'(x).$$



Examples:



\begin{itemize}
    \item $f(x) = 6x + 1,$ $g(x) = x^2.$ $f'(x) = 6$ and $g'(x) = 2x,$ so
\end{itemize}
    $\frac{d}{dx}(6x + 1)(x^2) = 6x^2 + (2x)(6x + 1).$

\begin{itemize}
    \item $f(x) = \sqrt{x},$ $g(x) = \sin{x}.$ $f'(x) = \frac{1}{2\sqrt{x}}$
\end{itemize}
    and $g'(x) = \cos{x},$ so

    $\frac{d}{dx}(\sqrt{x})(\sin{x}) = \left(\frac{1}{2\sqrt{x}}\right)(\cos{x}) + \sqrt{x}\sin{x}.$



\paragraph{Quotient Rule}



Let $f(x)$ and $g(x)$ be differentiable functions such that

$g(x) \neq 0.$ Then, the derivative of the quotient,

$\frac{f(x)}{g(x)},$ is



$$\frac{d}{dx}\left(\frac{f(x)}{g(x)}\right) = \frac{f'(x)g(x) - f(x)g'(x)}{(g(x))^2}.$$



Examples:



\begin{itemize}
    \item $f(x) = x + 1,$ $g(x) = x^3.$ $f'(x) = 1$ and $g'(x) = 3x^2,$ so
\end{itemize}
    $\frac{d}{dx}\left(\frac{x+1}{x^3}\right) = \frac{x^3 - (x+1)(3x^2)}{(x^3)^2} = \frac{x^3 - 3x^3 - 3x^2}{x^6} = \frac{-x^2(2x + 3)}{x^6} = -\frac{2x + 3}{x^4}.$



\begin{itemize}
    \item $f(x) = 5x,$ $g(x) = x^2 + 1.$ $f'(x) = 5$ and $g'(x) = 2x,$ so
\end{itemize}
    $\frac{d}{dx}\left(\frac{5x}{x^2 + 1}\right) = \frac{5(x^2+1) - 5x(2x)}{(x^2 + 1)^2} = \frac{5x^2+5 - 10x^2}{(x^2 + 1)^2} = \frac{5(1 - x^2)}{(x^2 + 1)^2}.$



\paragraph{Resources}



\begin{itemize}
    \item \href{http://www2.gcc.edu/dept/math/faculty/BancroftED/buscalc/chapter2/section2-4.php}{Product and Quotient
\end{itemize}
    Rules},

    Grover City College

\begin{itemize}
    \item \href{https://tutorial.math.lamar.edu/classes/calci/productquotientrule.aspx}{Product and Quotient
\end{itemize}
    Rule},

    Paul's Online Notes (Lamar University)

\subsection{Learning Objectives}

\begin{enumerate}
    \item \textbf{Apply the Product Rule}: Students will be able to apply the Product Rule to differentiate functions that are products of two differentiable functions.
    \item \textbf{Apply the Quotient Rule}: Students will be able to apply the Quotient Rule to differentiate functions that are quotients of two differentiable functions, given that the denominator is non-zero.
    \item \textbf{Simplify Derivatives Using Product and Quotient Rules}: Students will be able to simplify derivatives obtained from the Product Rule and Quotient Rule by combining like terms and reducing fractions.
\end{enumerate}

\subsection{Assessment Instrument}

\begin{enumerate}
    \item \textbf{Part 1: Apply the Product Rule}
    \item \textbf{Calculation:} Given the functions $u(x) = 2x^3$ and $v(x) = \sin(x)$, use the Product Rule to find the derivative of the function $f(x) = u(x) \cdot v(x)$. Show all steps and provide the simplified form of the derivative.
    \item \textbf{Multiple Choice:} What is the derivative of the function $f(x) = (x^2 + 1) \cdot \ln(x)$ using the Product Rule?
    \item A. $\frac{d}{dx} \left[(x^2 + 1) \cdot \ln(x)\right] = (2x \cdot \ln(x) + \frac{x^2 + 1}{x})$
    \item B. $\frac{d}{dx} \left[(x^2 + 1) \cdot \ln(x)\right] = (2x \cdot \ln(x) + \frac{x^2 + 1}{x}) \cdot \frac{1}{x}$
    \item C. $\frac{d}{dx} \left[(x^2 + 1) \cdot \ln(x)\right] = (2x \cdot \ln(x) + \frac{1}{x})$
    \item D. $\frac{d}{dx} \left[(x^2 + 1) \cdot \ln(x)\right] = (2x \cdot \ln(x) + \ln(x))$
    \item \textbf{Short Answer:} Compute the derivative of $f(x) = x^4 \cdot e^x$ using the Product Rule. Provide the result in its simplified form and explain each step. \textbf{Part 2: Apply the Quotient Rule}
    \item \textbf{Calculation:} Differentiate the function $f(x) = \frac{x^2 + 1}{x^3 - 2}$ using the Quotient Rule. Show your work and simplify the result.
    \item \textbf{True/False:} The derivative of the function $f(x) = \frac{3x^2 - 4}{2x + 1}$, using the Quotient Rule, is given by $\frac{(6x \cdot (2x + 1) - (3x^2 - 4) \cdot 2)}{(2x + 1)^2}$. (True/False) Explain your reasoning.
    \item \textbf{Short Answer:} Apply the Quotient Rule to find the derivative of the function $g(x) = \frac{\sin(x)}{x^2 + 1}$. Show all calculations and provide the simplified form of the derivative. \textbf{Part 3: Simplify Derivatives Using Product and Quotient Rules}
    \item \textbf{Calculation:} Simplify the derivative obtained from applying the Product Rule to the function $f(x) = x^2 \cdot \cos(x)$. Provide the simplified derivative and explain how you combined like terms.
    \item \textbf{Multiple Choice:} For the function $h(x) = \frac{x^3 \cdot e^x}{x + 1}$, what is the simplified form of the derivative using the Quotient Rule?
    \item A. $\frac{(3x^2 \cdot e^x \cdot (x + 1) - (x^3 \cdot e^x)}{(x + 1)^2}$
    \item B. $\frac{(x^3 \cdot e^x \cdot (1 + x) - (3x^2 \cdot e^x)}{(x + 1)^2}$
    \item C. $\frac{(x^3 \cdot e^x - 3x^2 \cdot e^x \cdot (x + 1))}{(x + 1)^2}$
    \item D. $\frac{(x^3 \cdot e^x - x^2 \cdot e^x \cdot (x + 1))}{(x + 1)^2}$
    \item \textbf{True/False:} The derivative of the function $f(x) = \frac{x^2 + 1}{\cos(x)}$ can be simplified to $\frac{\cos(x) \cdot (2x) + (x^2 + 1) \cdot \sin(x)}{\cos^2(x)}$. (True/False) Provide a brief explanation of how you simplified the derivative.
\end{enumerate}

%%===============================================================
\section{The Chain Rule}
\label{sec:the-chain-rule}
%%===============================================================

In calculus, the \textbf{chain rule} is a formula that expresses the

derivative of the composition of two differentiable functions

\textit{f} and \textit{g} in terms of the

derivatives \textit{f} and \textit{g}. More

precisely, if $h=f\circ g$ is the function such that $h(x)=f(g(x))$ for

every \textit{x}, then the chain rule is, in Lagrange's

notation,



$$h'(x) = f'(g(x)) g'(x).$$



or, equivalently,



$$h'=(f\circ g)'=(f'\circ g)\cdot g'.$$



The chain rule may also be expressed in Leibniz's notation. If a

variable \textit{z} depends on the variable

\textit{y}, which itself depends on the variable

\textit{x} (that is, \textit{y} and

\textit{z} are dependent variables), then

\textit{z} depends on \textit{x} as well,

via the intermediate variable \textit{y}. In this case,

the chain rule is expressed as



$$\frac{dz}{dx} = \frac{dz}{dy} \cdot \frac{dy}{dx},$$



and



$$\left.\frac{dz}{dx}\right|\_{x} = \left.\frac{dz}{dy}\right|\_{y(x)}

\cdot \left. \frac{dy}{dx}\right|\_{x} ,$$



for indicating at which points the derivatives have to be evaluated.



In integration, the counterpart to the chain rule is the substitution rule.



\paragraph{Intuitive explanation}



Intuitively, the chain rule states that knowing the instantaneous rate

of change of \textit{z} relative to

\textit{x} and that of \textit{x}

relative to \textit{x} allows one to calculate the

instantaneous rate of change of \textit{z} relative to

\textit{x} as the product of the two rates of change.



As put by George F. Simmons: "if a car travels twice as fast as a

bicycle and the bicycle is four times as fast as a walking man, then the

car travels 2 × 4 = 8 times as fast as the man."



The relationship between this example and the chain rule is as follows.

Let \textit{z}, \textit{y} and \textit{x} be the (variable) positions of the car, the

bicycle, and the walking man, respectively. The rate of change of

relative positions of the car and the bicycle is $\frac {dz}{dy}=2.$

Similarly, $\frac {dy}{dx}=4.$ So, the rate of change of the relative

positions of the car and the walking man is



$$\frac{dz}{dx}=\frac{dz}{dy}\cdot\frac{dy}{dx}=2\cdot 4=8.$$



The rate of change of positions is the ratio of the speeds, and the

speed is the derivative of the position with respect to the time; that

is,



$$\frac{dz}{dx}=\frac \frac{dz}{dt}\frac{dx}{dt},$$



or, equivalently,



$$\frac{dz}{dt}=\frac{dz}{dx}\cdot \frac{dx}{dt},$$



which is also an application of the chain rule.



\paragraph{Statement}



The simplest form of the chain rule is for real-valued functions of one

real variable. It states that if \textit{g} is a

function that is differentiable at a point \textit{c}

(i.e. the derivative \textit{g'(c)} exists) and

\textit{f} is a function that is differentiable at

\textit{g(c)}, then the composite function

$f\circ g$ is differentiable at \textit{c}, and the

derivative is



$$(f\circ g)'(c) = f'(g(c))\cdot g'(c).$$



The rule is sometimes abbreviated as



$$(f\circ g)' = (f'\circ g) \cdot g'.$$



If \textit{y=f(u)} and

\textit{u=g(x)}, then this abbreviated

form is written in Leibniz notation as:



$$\frac{dy}{dx} = \frac{dy}{du} \cdot \frac{du}{dx}.$$



The points where the derivatives are evaluated may also be stated

explicitly:



$$\left.\frac{dy}{dx}\right|\_{x=c} = \left.\frac{dy}{du}\right|\_{u = g(c)} \cdot \left.\frac{du}{dx}\right|\_{x=c}.$$



Carrying the same reasoning further, given \textit{n}

functions $f\_1, \ldots, f\_n\!$ with the composite function

$f\_1 \circ ( f\_2 \circ \cdots (f\_{n-1} \circ f\_n) )\!,$ if each function

$f\_i\!$ is differentiable at its immediate input, then the composite

function is also differentiable by the repeated application of Chain

Rule, where the derivative is (in Leibniz's notation):



$$\frac{df\_1}{dx} = \frac{df\_1}{df\_2}\frac{df\_2}{df\_3}\cdots\frac{df\_n}{dx}.$$



\paragraph{Proofs}



\paragraph{First proof}



One proof of the chain rule begins with the definition of the

derivative:



$$(f \circ g)'(a) = \lim\_{x \to a} \frac{f(g(x)) - f(g(a))}{x - a}.$$



Assume for the moment that $g(x)\!$ does not equal $g(a)$ for any

\textit{x} near \textit{a}. Then the

previous expression is equal to the product of two factors:



$$\lim\_{x \to a} \frac{f(g(x)) - f(g(a))}{g(x) - g(a)} \cdot \frac{g(x) - g(a)}{x - a}.$$



If $g$ oscillates near \textit{a}, then it might happen

that no matter how close one gets to \textit{a}, there is

always an even closer \textit{x} such that

\textit{g(x)=g(a)}. For example, this

happens near \textit{a=0} for the continuous

function \textit{g} defined by

\textit{g(x)=0} for

\textit{x=0} and

\textit{$g(x)=x^2 \sin(\frac{1}{x})$}

otherwise. Whenever this happens, the above expression is undefined

because it involves division by zero. To work around this, introduce a

function $Q$ as follows:



$$Q(y) = 

\begin{cases}

\frac{f(y) - f(g(a))}{y - g(a)}, & y \neq g(a), \\\\ f'(g(a)), & y = g(a).

\end{cases}$$ 



We will show that the difference quotient for \textit{f°g} is always equal to:



$$Q(g(x)) \cdot \frac{g(x) - g(a)}{x - a}.$$



Whenever \textit{g(x)} is not equal to \textit{g(a)}, this is clear because the factors of \textit{g(x) - g(a)} cancel. When \textit{g(x)} equals \textit{g(a)}, then the difference quotient for \textit{f ° g} is zero because \textit{f(g(x))} equals \textit{f(g(a))}, and the above product is zero because it equals \textit{f'(g(a))} times zero. So the above product is always equal to the difference quotient, and to show that the derivative of \textit{f ° g} at a exists and to determine its value, we need only show that the limit as \textit{x} goes to a of the above product exists and determine its value.



To do this, recall that the limit of a product exists if the limits of

its factors exist. When this happens, the limit of the product of these

two factors will equal the product of the limits of the factors. The two

factors are \textit{Q(g(x))} and

\textit{g(x) - g(a) / x - a}.

The latter is the difference quotient for \textit{g} at

\textit{a}, and because \textit{g} is

differentiable at \textit{a} by assumption, its limit as

\textit{x} tends to \textit{a} exists and

equals \textit{g'(a)}.



As for \textit{Q(g(x))}, notice that

\textit{Q} is defined wherever

\textbf{f} is. Furthermore, \textbf{f}

is differentiable at \textit{g(a)} by assumption,

so \textit{Q} is continuous at

\textit{g(a)}, by definition of the derivative.

The function \textit{g} is continuous at

\textit{a} because it is differentiable at

\textit{a}, and therefore

\textit{Q ° g} is continuous at

\textit{a}. So its limit as \textbf{x}

goes to \textbf{a} exists and equals

\textit{Q(g(a))}, which is

\textit{f'(g(a))}.



This shows that the limits of both factors exist and that they equal

\textit{f'(g(a))} and

\textit{g'(a)}, respectively. Therefore, the

derivative of \textit{f°g} at \textit{a} exists and

equals

\textit{f'(g(a))**g'(a)}.



\paragraph{Second proof}



Another way of proving the chain rule is to measure the error in the

linear approximation determined by the derivative. This proof has the

advantage that it generalizes to several variables. It relies on the

following equivalent definition of differentiability at a point: A

function \textit{g} is differentiable at \textit{a} if there exists a real number

\textit{g}'(\textit{a}) and a function \textit{e}(\textit{h}) that tends to zero as \textit{h} tends to

zero, and furthermore



$$g(a + h) - g(a) = g'(a) h + \varepsilon(h) h.$$ Here the left-hand

side represents the true difference between the value of \textit{g} at \textit{a} and

at \textit{a+h}, whereas the right-hand side

represents the approximation determined by the derivative plus an error

term.



In the situation of the chain rule, such a function \textit{e} exists because

\textit{g} is assumed to be differentiable at \textit{a}. Again by assumption, a

similar function also exists for \textit{f} at \textit{g}(\textit{a}). Calling this function

\textit{?}, we have



$$f(g(a) + k) - f(g(a)) = f'(g(a)) k + \eta(k) k.$$ The above definition

imposes no constraints on \textit{?}(0), even though it is assumed that

\textit{?}(\textit{k}) tends to zero as \textit{k} tends to zero. If we set

\textit{?}(0)=0, then \textit{?} is continuous at 0.



Proving the theorem requires studying the difference \textit{f(g(a + h)) - f(g(a))} as \textit{h} tends to zero. The first step is to substitute for \textit{g(a + h)} using the definition of differentiability of \textit{g} at \textit{a}:



$$f(g(a + h)) - f(g(a)) = f(g(a) + g'(a) h + \varepsilon(h) h) - f(g(a)).$$



The next step is to use the definition of differentiability of \textit{f} at \textit{g(a)}. This requires a term of the form \textit{f(g(a) + k)} for some \textit{k}. In the above equation, the correct \textit{k} varies with \textit{h}. Set \textit{$k\_h = g'(a) h + e(h) h$} and the right hand side becomes \textit{$f(g(a) + k\_h) - f(g(a))$}. Applying the definition of the derivative gives:



$$f(g(a) + k\_h) - f(g(a)) = f'(g(a)) k\_h + \eta(k\_h) k\_h.$$



To study the behavior of this expression as \textit{h} tends to zero, expand \textit{$k_h$}. After

regrouping the terms, the right-hand side becomes:



$$f'(g(a)) g'(a)h + [f'(g(a)) \varepsilon(h) + \eta(k\_h) g'(a) + \eta(k\_h) \varepsilon(h)] h.$$



Because \textit{e}(\textit{h}) and \textit{?}(\textit{k}\textasciitilde{}\textit{h}\textasciitilde{}) tend to zero as \textit{h} tends to zero,

the first two bracketed terms tend to zero as \textit{h} tends to zero.

Applying the same theorem on products of limits as in the first proof,

the third bracketed term also tends zero. Because the above expression

is equal to the difference \textit{f(g(a + h)) - f(g(a))},

by the definition of the derivative \textit{f°g}

is differentiable at \textit{a} and its derivative is \textit{f'(g(a)) g'(a)}.



The role of \textit{Q} in the first proof is played by \textit{?} in this proof. They

are related by the equation:



$$Q(y) = f'(g(a)) + \eta(y - g(a)).$$



The need to define \textit{Q} at \textit{g}(\textit{a}) is analogous to the need to define \textit{?} at zero.



\paragraph{Third proof}



Constantin Carathéodory's alternative definition of the

differentiability of a function can be used to give an elegant proof of

the chain rule.



Under this definition, a function \textit{f} is

differentiable at a point \textit{a} if and only if there

is a function \textit{q}, continuous at

\textit{a} and such that \textit{f(x) - f(a) = q(x)(x - a)}.

There is at most one such function, and if \textit{f} is

differentiable at \textit{a} then \textit{f '(a) = q(a)}.



Given the assumptions of the chain rule and the fact that differentiable

functions and compositions of continuous functions are continuous, we

have that there exist functions \textit{q}, continuous at \textit{g(a)}, and \textit{r}, continuous at \textit{a}, and such that,



$$f(g(x))-f(g(a))=q(g(x))(g(x)-g(a))$$



and



$$g(x)-g(a)=r(x)(x-a).$$



Therefore,



$$f(g(x))-f(g(a))=q(g(x))r(x)(x-a),$$



but the function given by \textit{h(x) = q(g(x))r(x)} is continuous at \textit{a}, and we get, for this \textit{a}



$$(f(g(a)))'=q(g(a))r(a)=f'(g(a))g'(a).$$



A similar approach works for continuously differentiable (vector-)functions of many variables. This

method of factoring also allows a unified approach to stronger forms of

differentiability, when the derivative is required to be Lipschitz

continuous, Hölder continuous, etc. Differentiation itself can be viewed

as the polynomial remainder theorem (the little Étienne Bézout|Bézout

theorem, or factor theorem), generalized to an appropriate class of functions.



\paragraph{Proof via infinitesimals}



If $y=f(x)$ and $x=g(t)$ then choosing infinitesimal $\Delta t\not=0$ we

compute the corresponding $\Delta x=g(t+\Delta t)-g(t)$ and then the

corresponding $\Delta y=f(x+\Delta x)-f(x),$ so that



$$\frac{\Delta y}{\Delta t}=\frac{\Delta y}{\Delta x} \frac{\Delta x}{\Delta t}$$



and applying the standard part we obtain



$$\frac{d y}{d t}=\frac{d y}{d x} \frac{dx}{dt}$$



which is the chain rule.



\paragraph{Resources}



\begin{itemize}
    \item \href{https://www.khanacademy.org/math/ap-calculus-ab/ab-differentiation-2-new/ab-3-1a/a/chain-rule-review}{Chain
\end{itemize}
    Rule},

    Khan Academy

\begin{itemize}
    \item \href{https://www.khanacademy.org/math/ap-calculus-ab/ab-differentiation-2-new/ab-3-1a/v/chain-rule-introduction}{Introduction to Chain
\end{itemize}
    Rule},

    Khan Academy

\begin{itemize}
    \item \href{https://www.youtube.com/watch?v=HaHsqDjWMLU}{Chain Rule For Finding
\end{itemize}
    Derivatives}, The

    Organic Chemistry Tutor



\paragraph{Licensing}



Content obtained and/or adapted from:



\begin{itemize}
    \item \href{https://en.wikipedia.org/wiki/Chain\_rule}{Chain rule, Wikipedia}
\end{itemize}
    under a CC BY-SA license

\subsection{Learning Objectives}

\begin{enumerate}
    \item \textbf{Apply the Chain Rule to Composite Functions}: Students will be able to apply the Chain Rule to find the derivative of a composite function $h(x) = f(g(x))$, using the formula $$h'(x) = f'(g(x)) \cdot g'(x).$$
    \item \textbf{Use Leibniz's Notation for the Chain Rule}: Students will be able to use Leibniz's notation to apply the Chain Rule in situations where variables depend on each other, using the formula $$\frac{dz}{dx} = \frac{dz}{dy} \cdot \frac{dy}{dx}.$$
    \item \textbf{Analyze and Prove the Chain Rule}: Students will be able to analyze and prove the Chain Rule using different methods, including definitions of differentiability, linear approximations, and infinitesimal calculations.
\end{enumerate}

\subsection{Assessment Instrument}

\begin{enumerate}
    \item \textbf{Part 1: Apply the Chain Rule to Composite Functions}
    \item \textbf{Calculation:} Given the composite function $h(x) = \sqrt{3x^2 + 2x}$, use the Chain Rule to find $h'(x)$. Show all steps and provide the simplified form of the derivative.
    \item \textbf{Multiple Choice:} If $f(x) = \ln(x^2 + 1)$ and $g(x) = 2x - 3$, what is the derivative of the composite function $h(x) = f(g(x))$ using the Chain Rule?
    \item A. $\frac{2}{(2x - 3)^2 + 1} \cdot 2$
    \item B. $\frac{1}{(2x - 3)^2 + 1} \cdot 2$
    \item C. $\frac{1}{(2x - 3)} \cdot 2$
    \item D. $\frac{2}{(2x - 3)} \cdot \frac{2}{x^2 + 1}$
    \item \textbf{Short Answer:} Find the derivative of the function $h(x) = \sin(2x^2 + 1)$ using the Chain Rule. Provide the derivative in its simplified form and explain each step. \textbf{Part 2: Use Leibniz's Notation for the Chain Rule}
    \item \textbf{Calculation:} Consider the functions $z = x^2 + y^2$ and $y = e^x$. Compute $\frac{dz}{dx}$ using Leibniz's notation, where $z$ is implicitly a function of $x$ through $y$. Provide all necessary steps.
    \item \textbf{Multiple Choice:} If $z$ depends on $y$ and $y$ depends on $x$, which of the following represents $\frac{dz}{dx}$ using Leibniz's notation?
    \item A. $\frac{dz}{dx} = \frac{dz}{dy} \cdot \frac{dx}{dy}$
    \item B. $\frac{dz}{dx} = \frac{dz}{dy} \cdot \frac{dy}{dx}$
    \item C. $\frac{dz}{dx} = \frac{dy}{dz} \cdot \frac{dx}{dy}$
    \item D. $\frac{dz}{dx} = \frac{dx}{dz} \cdot \frac{dy}{dx}$
    \item \textbf{Short Answer:} For the functions $z = \cos(y)$ and $y = x^2 + 1$, compute $\frac{dz}{dx}$ using Leibniz's notation. Show all calculations and explain the relationship between the variables. \textbf{Part 3: Analyze and Prove the Chain Rule}
    \item \textbf{Proof:} Prove the Chain Rule using the definition of the derivative. Start with the limit definition of the derivative and show how it leads to the formula $$h'(x) = f'(g(x)) \cdot g'(x).$$
    \item \textbf{Short Answer:} Explain how the Chain Rule can be derived using linear approximations. Provide a brief description of how linear approximations help in understanding the behavior of composite functions.
    \item \textbf{True/False:} The Chain Rule can be derived by considering the infinitesimal changes in $x$ and how they affect $g(x)$ and subsequently $f(g(x))$. (True/False) Provide a brief explanation of why or why not, including a description of how infinitesimal changes are used in the proof.
\end{enumerate}

%%===============================================================
\section{Derivatives of Exponential Functions}
\label{sec:derivatives-of-exponential-functions}
%%===============================================================

\paragraph{Logarithm Function}



We shall first look at the irrational number $e$ in order to show its

special properties when used with derivatives of exponential and

logarithm functions. The value of $e$ is approximately

$e\approx2.718282$ but it may also be calculated with the following

infinite limit:



$$e=\lim\_{n\to\infty}\left(1+\frac{1}{n}\right)^n$$



Now we find the derivative of $\ln(x)$ using the formal definition of

the derivative:



$$\frac{d}{dx}\ln(x)=\lim\_{h\to0}\frac{\ln(x+h)-\ln(x)}{h}=\lim\_{h\to0}\frac{1}{h}\ln\left(\frac{x+h}{x}\right)=\lim\_{h\to0}\ln\left(\frac{x+h}{x}\right)^\frac{1}{h}$$



Let $n=\frac{x}{h}$ . Note that as $n\to\infty$ , we get $h\to0$ . So we

can redefine our limit as:



$$\lim\_{n\to\infty}\ln\left(1+\frac{1}{n}\right)^\frac{n}{x}=\frac{1}{x}\ln\left(\lim\_{n\to\infty}\left(1+\frac{1}{n}\right)^n\right)=\frac{1}{x}\ln(e)=\frac{1}{x}$$



Here we could take the natural logarithm outside the limit because it

doesn't have anything to do with the limit (we could have chosen not

to). We then substituted the value of $e$ .



\textbf{Derivative of the Natural Logarithm}

$$\frac{d}{dx}\ln(x)=\frac{1}{x}$$         



If we wanted, we could go through that same process again for a

generalized base, but it is easier just to use properties of logs and

realize that:



$$\log\_a(x)=\frac{\ln(x)}{\ln(a)}$$



Since $\frac{1}{\ln(a)}$ is a constant, we can just take it outside of

the derivative:



$$\frac{d}{dx}\log\_a(x)=\frac{1}{\ln(a)}\cdot\frac{d}{dx}\ln(x)$$



Which leaves us with the generalized form of:



 \textbf{Derivative of the Logarithm}

 $$\frac{d}{dx}\log\_a(x)=\frac{1}{\ln(a)x}$$ 



An alternative approach to derivative of the logarithm refers to the

original expression of the logarithm as quadrature of the hyperbola \textit{y}

= 1/\textit{x} .



\paragraph{Exponential Function}



We shall take two different approaches to finding the derivative of

$\ln(e^x)$ . The first approach:



$$\frac{d}{dx}\ln(e^x)=\frac{d}{dx}x=1$$



The second approach:



$$\frac{d}{dx}\ln(e^x)=\frac{1}{e^x}\left(\frac{d}{dx}e^x\right)$$



Note that in the second approach we made some use of the chain rule. Thus:



$$\frac{1}{e^x}\left(\frac{d}{dx}e^x\right)=1$$



$$\frac{d}{dx}e^x=e^x$$



so that we have proved the following rule:



\textbf{Derivative of the exponential function}

 $$\frac{d}{dx}e^x=e^x$$               



Now that we have derived a specific case, let us extend things to the

general case. Assuming that $a$ is a positive real constant, we wish to

calculate:



$$\frac{d}{dx}a^x$$



One of the oldest tricks in mathematics is to break a problem down into

a form that we already know we can handle. Since we have already

determined the derivative of $e^x$ , we will attempt to rewrite $a^x$ in

that form.



Using that $e^{\ln(c)}=c$ and that $\ln(a^b)=b\cdot\ln(a)$ , we find

that:



$$a^x=e^{\ln(a)x}$$



Thus, we simply apply the chain rule:



$$\frac{d}{dx}e^{\ln(a)x}=e^{\ln(a)x}\cdot\frac{d}{dx}[\ln(a)x]=\ln(a)a^x$$



 \textbf{Derivative of the exponential function}

 $$\frac{d}{dx}a^x=\ln(a)a^x$$                





\paragraph{Logarithmic Differentiation}



We can use the properties of the logarithm, particularly the natural

log, to differentiate more difficult functions, such a products with

many terms, quotients of composed functions, or functions with variable

or function exponents. We do this by taking the natural logarithm of

both sides, re-arranging terms using the logarithm laws below, and then

differentiating both sides implicitly, before multiplying through by $y$.



 $$\log(a)+\log(b)=\log(ab)$$                     



 $$\log\left(\frac{a}{b}\right)=\log(a)-\log(b)$$ 



 $$\log(a^n)=n\log(a)$$                           



See the examples below.



\paragraph{Example 1}



We shall now prove the validity of the power rule using logarithmic

differentiation.



$$\frac{d}{dx}\ln(x^n)=n\frac{d}{dx}\ln(x)=nx^{-1}$$



$$\frac{d}{dx}\ln(x^n)=\frac{1}{x^n}\cdot\frac{d}{dx}x^n$$



Thus:



$$\frac{1}{x^n}\cdot\frac{d}{dx}x^n=nx^{-1}$$



$$\frac{d}{dx}x^n=nx^{n-1}$$



Example 2



Suppose we wished to differentiate   

                                                                                                    

$$y=\frac{(6x^2+9)^2}{\sqrt{3x^3-2}}$$



We take the natural logarithm of both sides

  

 $$\begin{aligned}    

 \ln(y)&=\ln\left(\frac{(6x^2+9)^2}{\sqrt{3x^3-2}}\right) \\\\       

 &=\ln\big((6x^2+9)^2\big)-\ln\big((3x^3-2)^{\frac12}\big) \\\\       

 &=2\ln(6x^2+9)-\frac{\ln(3x^3-2)}{2}                           

 \end{aligned}$$ 



Differentiating implicitly, recalling the chain rule 



 $$\begin{aligned}  

 \frac{1}{y}\cdot\frac                                                

 {dy}{dx}&=2\cdot\frac{12x}{6x^2+9}-\frac12\cdot\frac{9x^2}{3x^3-2} \\\\ 

 &=\frac{24x}{6x^2+9}-\frac{\frac{9}{2}x^2}{3x^3-2} \\\\                 

 &=\frac{24x(3x^3-2)-\frac92x^2(6x^2+9)}{(6x^2+9)(3x^3-2)}          

 \end{aligned}$$ 



Multiplying by $y$ , the original function

     

$$\frac{dy}{dx}=\frac{(6x^2+9)^2}{\sqrt{3x^3-2}}\cdot\frac{24x(3x^3-2)-\frac{9}{2}x^2(6x^2+9)}{(6x^2+9)(3x^3-2)}$$ 



Example 3



Let us differentiate a function



$$y=x^x$$ 



Taking the natural logarithm of left and right

                                                                      

 $$\begin{aligned}

\ln(y)&=\ln(x^x) \\\&=x\ln(x)

\end{aligned}$$ 



We then differentiate both sides, recalling the product and chain rules    

                                                                      

 $$\begin{aligned}

\frac{1}{y}\cdot\frac{dy}{dx}&=\ln(x)+x\frac{1}{x} \\\&=\ln(x)+1

\end{aligned}$$



 Multiplying by the original function $y$                             

                                                                     

 $$\frac{dy}{dx}=x^x(\ln(x)+1)$$                                     





Example 4



 Take a function

                                                              

 $$y=x^{6\cos(x)}$$



Then                                              

                                                                     

 $$\begin{aligned}

\ln(y)&=\ln(x^{6\cos(x)}) \\\&=6\cos(x)\ln(x)

\end{aligned}$$ 



We then differentiate

                                                

$$\frac{1}{y}\cdot\frac{dy}{dx}=-6\sin(x)\ln(x)+\frac{6\cos(x)}{x}$$



And finally multiply by $y$                                          

                                                                      

 $$\begin{aligned}                                                      

 \frac{dy}{dx}&=y\left(-6\sin(x)\ln(x)+\frac{6\cos(x)}{x}\right) \\\&=6x^{6\cos(x)}\left(\frac{\cos(x)}{x}-\sin(x)\ln(x)\right) 

 \end{aligned}$$                                                        



\paragraph{Resources}



\begin{itemize}
    \item \href{https://mathresearch.utsa.edu/wikiFiles/MAT1214/The\%20Derivative\%20of\%20the\%20Exponential\%20Function/MAT1214-3.9TheDerivativeOfTheExponentialFunctionPwPt.pptx}{Derivatives of the Exponential and Logarithmic
\end{itemize}
    Functions}

    PowerPoint file created by Dr. Sara Shirinkam, UTSA.



<!-- -->



\begin{itemize}
    \item \href{https://mathresearch.utsa.edu/wikiFiles/MAT1214/The\%20Derivative\%20of\%20the\%20Exponential\%20Function/MAT1214-3.9TheDerivativesOfExpAndLogFunctionsWS1.pdf}{Derivatives of the Exponential and Logarithmic Functions
\end{itemize}
    Worksheet}



\paragraph{Licensing}



Content obtained and/or adapted from:



\begin{itemize}
    \item \href{https://en.wikibooks.org/wiki/Calculus/Derivatives\_of\_Exponential\_and\_Logarithm\_Functions}{Derivatives of Exponential and Logarithm Functions, Wikibooks:
\end{itemize}
    Calculus}

    under a CC BY-SA license

\subsection{Learning Objectives}

\begin{enumerate}
    \item \textbf{Compute the Derivative of the Natural Logarithm:} Students will be able to derive and apply the formula for the derivative of the natural logarithm function.
    \item \textbf{Differentiate Exponential Functions and Their Generalizations:} Students will be able to find the derivatives of exponential functions, including those with base $ e $ and other positive real constants.
    \item \textbf{Apply Logarithmic Differentiation to Complex Functions:} Students will be able to use logarithmic differentiation to find the derivatives of more complex functions involving products, quotients, or variable exponents.
\end{enumerate}

\subsection{Assessment Instrument}

\begin{enumerate}
    \item \textbf{Part 1: Compute the Derivative of the Natural Logarithm}
    \item \textbf{Calculation:} Find the derivative of the function $f(x) = \ln(x^2 + 3x + 2)$. Provide the result in its simplified form and explain each step of your computation.
    \item \textbf{Multiple Choice:} Which of the following correctly represents the derivative of the natural logarithm function $f(x) = \ln(5x^3 + 4x - 7)$?
    \item A. $\frac{f'(x)}{5x^3 + 4x - 7}$
    \item B. $\frac{15x^2 + 4}{5x^3 + 4x - 7}$
    \item C. $\frac{15x^2 + 4}{5x^3 + 4x - 7}$
    \item D. $\frac{1}{5x^3 + 4x - 7} \cdot (15x^2 + 4)$
    \item \textbf{True/False:} The derivative of the natural logarithm function $f(x) = \ln(x)$ is given by $f'(x) = \frac{1}{x}$ for $x > 0$. (True/False) Provide a brief explanation for your answer. \textbf{Part 2: Differentiate Exponential Functions and Their Generalizations}
    \item \textbf{Calculation:} Compute the derivative of the function $f(x) = e^{3x^2 - 2x}$. Show all steps and provide the result in its simplified form.
    \item \textbf{Multiple Choice:} Which of the following is the derivative of the function $g(x) = 2^x$?
    \item A. $2^x \ln(2)$
    \item B. $2^x \cdot \frac{1}{\ln(2)}$
    \item C. $2^x \cdot \ln(x)$
    \item D. $2^x \cdot x$
    \item \textbf{Short Answer:} Find the derivative of the function $h(x) = e^{x^2} \cdot e^{-x}$ using properties of exponential functions. Show your work and provide the simplified result. \textbf{Part 3: Apply Logarithmic Differentiation to Complex Functions}
    \item \textbf{Calculation:} Use logarithmic differentiation to find the derivative of the function $f(x) = \frac{x^2 \cdot e^{x}}{\sqrt{x + 1}}$. Show all steps and provide the derivative in its simplified form.
    \item \textbf{Short Answer:} Differentiate the function $g(x) = (x^2 + 1)^x$ using logarithmic differentiation. Provide a clear explanation of the steps involved in applying logarithmic differentiation to this function.
    \item \textbf{Multiple Choice:} For the function $h(x) = \frac{\ln(x)}{x^2}$, which of the following is the correct derivative using logarithmic differentiation?
    \item A. $h'(x) = \frac{1 - 2 \ln(x)}{x^3}$
    \item B. $h'(x) = \frac{1 - \ln(x)}{x^2}$
    \item C. $h'(x) = \frac{1 - 2 \ln(x)}{x^2}$
    \item D. $h'(x) = \frac{1 - 2 \ln(x)}{x^2}$
\end{enumerate}

%%===============================================================
\section{Derivatives of Logarithmic Functions}
\label{sec:derivatives-of-logarithmic-functions}
%%===============================================================

\paragraph{Logarithm Function}



We shall first look at the irrational number $e$ in order to show its

special properties when used with derivatives of exponential and

logarithm functions. The value of $e$ is approximately

$e\approx2.718282$ but it may also be calculated with the following

infinite limit:



$$e=\lim\_{n\to\infty}\left(1+\frac{1}{n}\right)^n$$



Now we find the derivative of $\ln(x)$ using the formal definition of

the derivative:



$$\frac{d}{dx}\ln(x)=\lim\_{h\to0}\frac{\ln(x+h)-\ln(x)}{h}=\lim\_{h\to0}\frac{1}{h}\ln\left(\frac{x+h}{x}\right)=\lim\_{h\to0}\ln\left(\frac{x+h}{x}\right)^\frac{1}{h}$$



Let $n=\frac{x}{h}$ . Note that as $n\to\infty$ , we get $h\to0$ . So we

can redefine our limit as:



$$\lim\_{n\to\infty}\ln\left(1+\frac{1}{n}\right)^\frac{n}{x}=\frac{1}{x}\ln\left(\lim\_{n\to\infty}\left(1+\frac{1}{n}\right)^n\right)=\frac{1}{x}\ln(e)=\frac{1}{x}$$



Here we could take the natural logarithm outside the limit because it

doesn't have anything to do with the limit (we could have chosen not

to). We then substituted the value of $e$ .



\textbf{Derivative of the Natural Logarithm}

$$\frac{d}{dx}\ln(x)=\frac{1}{x}$$         



If we wanted, we could go through that same process again for a

generalized base, but it is easier just to use properties of logs and

realize that:



$$\log\_a(x)=\frac{\ln(x)}{\ln(a)}$$



Since $\frac{1}{\ln(a)}$ is a constant, we can just take it outside of

the derivative:



$$\frac{d}{dx}\log\_a(x)=\frac{1}{\ln(a)}\cdot\frac{d}{dx}\ln(x)$$



Which leaves us with the generalized form of:



 \textbf{Derivative of the Logarithm}

 $$\frac{d}{dx}\log\_a(x)=\frac{1}{\ln(a)x}$$ 



An alternative approach to derivative of the logarithm refers to the

original expression of the logarithm as quadrature of the hyperbola \textit{y}

= 1/\textit{x} .



\paragraph{Exponential Function}



We shall take two different approaches to finding the derivative of

$\ln(e^x)$ . The first approach:



$$\frac{d}{dx}\ln(e^x)=\frac{d}{dx}x=1$$



The second approach:



$$\frac{d}{dx}\ln(e^x)=\frac{1}{e^x}\left(\frac{d}{dx}e^x\right)$$



Note that in the second approach we made some use of the chain rule. Thus:



$$\frac{1}{e^x}\left(\frac{d}{dx}e^x\right)=1$$



$$\frac{d}{dx}e^x=e^x$$



so that we have proved the following rule:



\textbf{Derivative of the exponential function}

 $$\frac{d}{dx}e^x=e^x$$               



Now that we have derived a specific case, let us extend things to the

general case. Assuming that $a$ is a positive real constant, we wish to

calculate:



$$\frac{d}{dx}a^x$$



One of the oldest tricks in mathematics is to break a problem down into

a form that we already know we can handle. Since we have already

determined the derivative of $e^x$ , we will attempt to rewrite $a^x$ in

that form.



Using that $e^{\ln(c)}=c$ and that $\ln(a^b)=b\cdot\ln(a)$ , we find

that:



$$a^x=e^{\ln(a)x}$$



Thus, we simply apply the chain rule:



$$\frac{d}{dx}e^{\ln(a)x}=e^{\ln(a)x}\cdot\frac{d}{dx}[\ln(a)x]=\ln(a)a^x$$



 \textbf{Derivative of the exponential function}

 $$\frac{d}{dx}a^x=\ln(a)a^x$$                





\paragraph{Logarithmic Differentiation}



We can use the properties of the logarithm, particularly the natural

log, to differentiate more difficult functions, such a products with

many terms, quotients of composed functions, or functions with variable

or function exponents. We do this by taking the natural logarithm of

both sides, re-arranging terms using the logarithm laws below, and then

differentiating both sides implicitly, before multiplying through by $y$.



 $$\log(a)+\log(b)=\log(ab)$$                     



 $$\log\left(\frac{a}{b}\right)=\log(a)-\log(b)$$ 



 $$\log(a^n)=n\log(a)$$                           



See the examples below.



\paragraph{Example 1}



We shall now prove the validity of the power rule using logarithmic

differentiation.



$$\frac{d}{dx}\ln(x^n)=n\frac{d}{dx}\ln(x)=nx^{-1}$$



$$\frac{d}{dx}\ln(x^n)=\frac{1}{x^n}\cdot\frac{d}{dx}x^n$$



Thus:



$$\frac{1}{x^n}\cdot\frac{d}{dx}x^n=nx^{-1}$$



$$\frac{d}{dx}x^n=nx^{n-1}$$



Example 2



Suppose we wished to differentiate   

                                                                                                    

$$y=\frac{(6x^2+9)^2}{\sqrt{3x^3-2}}$$



We take the natural logarithm of both sides

  

 $$\begin{aligned}    

 \ln(y)&=\ln\left(\frac{(6x^2+9)^2}{\sqrt{3x^3-2}}\right) \\\\       

 &=\ln\big((6x^2+9)^2\big)-\ln\big((3x^3-2)^{\frac12}\big) \\\\       

 &=2\ln(6x^2+9)-\frac{\ln(3x^3-2)}{2}                           

 \end{aligned}$$ 



Differentiating implicitly, recalling the chain rule 



 $$\begin{aligned}  

 \frac{1}{y}\cdot\frac                                                

 {dy}{dx}&=2\cdot\frac{12x}{6x^2+9}-\frac12\cdot\frac{9x^2}{3x^3-2} \\\\ 

 &=\frac{24x}{6x^2+9}-\frac{\frac{9}{2}x^2}{3x^3-2} \\\\                 

 &=\frac{24x(3x^3-2)-\frac92x^2(6x^2+9)}{(6x^2+9)(3x^3-2)}          

 \end{aligned}$$ 



Multiplying by $y$ , the original function

     

$$\frac{dy}{dx}=\frac{(6x^2+9)^2}{\sqrt{3x^3-2}}\cdot\frac{24x(3x^3-2)-\frac{9}{2}x^2(6x^2+9)}{(6x^2+9)(3x^3-2)}$$ 



Example 3



Let us differentiate a function



$$y=x^x$$ 



Taking the natural logarithm of left and right

                                                                      

 $$\begin{aligned}

\ln(y)&=\ln(x^x) \\\&=x\ln(x)

\end{aligned}$$ 



We then differentiate both sides, recalling the product and chain rules    

                                                                      

 $$\begin{aligned}

\frac{1}{y}\cdot\frac{dy}{dx}&=\ln(x)+x\frac{1}{x} \\\&=\ln(x)+1

\end{aligned}$$



 Multiplying by the original function $y$                             

                                                                     

 $$\frac{dy}{dx}=x^x(\ln(x)+1)$$                                     





Example 4



 Take a function

                                                              

 $$y=x^{6\cos(x)}$$



Then                                              

                                                                     

 $$\begin{aligned}

\ln(y)&=\ln(x^{6\cos(x)}) \\\&=6\cos(x)\ln(x)

\end{aligned}$$ 



We then differentiate

                                                

$$\frac{1}{y}\cdot\frac{dy}{dx}=-6\sin(x)\ln(x)+\frac{6\cos(x)}{x}$$



And finally multiply by $y$                                          

                                                                      

 $$\begin{aligned}                                                      

 \frac{dy}{dx}&=y\left(-6\sin(x)\ln(x)+\frac{6\cos(x)}{x}\right) \\\&=6x^{6\cos(x)}\left(\frac{\cos(x)}{x}-\sin(x)\ln(x)\right) 

 \end{aligned}$$                                                        



\paragraph{Resources}



\begin{itemize}
    \item \href{https://mathresearch.utsa.edu/wikiFiles/MAT1214/The\%20Derivative\%20of\%20the\%20Exponential\%20Function/MAT1214-3.9TheDerivativeOfTheExponentialFunctionPwPt.pptx}{Derivatives of the Exponential and Logarithmic
\end{itemize}
    Functions}

    PowerPoint file created by Dr. Sara Shirinkam, UTSA.



<!-- -->



\begin{itemize}
    \item \href{https://mathresearch.utsa.edu/wikiFiles/MAT1214/The\%20Derivative\%20of\%20the\%20Exponential\%20Function/MAT1214-3.9TheDerivativesOfExpAndLogFunctionsWS1.pdf}{Derivatives of the Exponential and Logarithmic Functions
\end{itemize}
    Worksheet}



\paragraph{Licensing}



Content obtained and/or adapted from:



\begin{itemize}
    \item \href{https://en.wikibooks.org/wiki/Calculus/Derivatives\_of\_Exponential\_and\_Logarithm\_Functions}{Derivatives of Exponential and Logarithm Functions, Wikibooks:
\end{itemize}
    Calculus}

    under a CC BY-SA license

\subsection{Learning Objectives}

\begin{enumerate}
    \item \textbf{Compute the Derivative of the Natural Logarithm:} Students will be able to derive and apply the formula for the derivative of the natural logarithm function.
    \item \textbf{Differentiate Exponential Functions and Their Generalizations:} Students will be able to find the derivatives of exponential functions, including those with base $ e $ and other positive real constants.
    \item \textbf{Apply Logarithmic Differentiation to Complex Functions:} Students will be able to use logarithmic differentiation to find the derivatives of more complex functions involving products, quotients, or variable exponents.
\end{enumerate}

\subsection{Assessment Instrument}

\begin{enumerate}
    \item \textbf{Part 1: Compute the Derivative of the Natural Logarithm}
    \item \textbf{Calculation:} Find the derivative of the function $f(x) = \ln(x^2 + 3x + 2)$. Provide the result in its simplified form and explain each step of your computation.
    \item \textbf{Multiple Choice:} Which of the following correctly represents the derivative of the natural logarithm function $f(x) = \ln(5x^3 + 4x - 7)$?
    \item A. $\frac{f'(x)}{5x^3 + 4x - 7}$
    \item B. $\frac{15x^2 + 4}{5x^3 + 4x - 7}$
    \item C. $\frac{15x^2 + 4}{5x^3 + 4x - 7}$
    \item D. $\frac{1}{5x^3 + 4x - 7} \cdot (15x^2 + 4)$
    \item \textbf{True/False:} The derivative of the natural logarithm function $f(x) = \ln(x)$ is given by $f'(x) = \frac{1}{x}$ for $x > 0$. (True/False) Provide a brief explanation for your answer. \textbf{Part 2: Differentiate Exponential Functions and Their Generalizations}
    \item \textbf{Calculation:} Compute the derivative of the function $f(x) = e^{3x^2 - 2x}$. Show all steps and provide the result in its simplified form.
    \item \textbf{Multiple Choice:} Which of the following is the derivative of the function $g(x) = 2^x$?
    \item A. $2^x \ln(2)$
    \item B. $2^x \cdot \frac{1}{\ln(2)}$
    \item C. $2^x \cdot \ln(x)$
    \item D. $2^x \cdot x$
    \item \textbf{Short Answer:} Find the derivative of the function $h(x) = e^{x^2} \cdot e^{-x}$ using properties of exponential functions. Show your work and provide the simplified result. \textbf{Part 3: Apply Logarithmic Differentiation to Complex Functions}
    \item \textbf{Calculation:} Use logarithmic differentiation to find the derivative of the function $f(x) = \frac{x^2 \cdot e^{x}}{\sqrt{x + 1}}$. Show all steps and provide the derivative in its simplified form.
    \item \textbf{Short Answer:} Differentiate the function $g(x) = (x^2 + 1)^x$ using logarithmic differentiation. Provide a clear explanation of the steps involved in applying logarithmic differentiation to this function.
    \item \textbf{Multiple Choice:} For the function $h(x) = \frac{\ln(x)}{x^2}$, which of the following is the correct derivative using logarithmic differentiation?
    \item A. $h'(x) = \frac{1 - 2 \ln(x)}{x^3}$
    \item B. $h'(x) = \frac{1 - \ln(x)}{x^2}$
    \item C. $h'(x) = \frac{1 - 2 \ln(x)}{x^2}$
    \item D. $h'(x) = \frac{1 - 2 \ln(x)}{x^2}$
\end{enumerate}

%%===============================================================
\section{Derivatives and Graphs}
\label{sec:derivatives-and-graphs}
%%===============================================================

\paragraph{Application of Derivatives to Graphs}



Derivatives also help us to graph functions, by locating minimum,

maximum, and inflection points. They can also determine the interval on

which a function is convex or concave.



\textbf{Convex} - The graph is below the tangent lines. Example $-x^2,$ think

of a frown.



\textbf{Concave} - The graph is above the tangent lines. Example $x^2,$ think

of a smile.



\textbf{Inflection Point} - The point when a function changes from convex to

concave or vice versa. Example $x^3$ at x = 0.



\textbf{Maximum Point} - The highest point of a function on a convex

interval.



\textbf{Minimum Point} - The lowest point of a function on a concave

interval.



\textbf{Stationary Point} - A point where f'(c) = 0



\paragraph{Rules of Stationary Points}



\begin{itemize}
    \item If $f' \left ( c \right ) = 0$ and $f'' \left ( c \right ) 0,$ then
\end{itemize}
    c is a local minimum point of f(x). The graph of f(x) will be convex

    on the interval.

\begin{itemize}
    \item If $f' \left ( c \right ) = 0$ and $f'' \left ( c \right ) = 0$ and
\end{itemize}
    $f''' \left ( c \right ) \ne 0,$ then c is a local inflection point

    of f(x).



\paragraph{Locating And Evaluating Stationary Points}



\textit{The Function}



The original function is used to determine when the function crosses the

x and y axis's.



\textit{The First Derivative Test}



The first derivative is used to find all the minimum, maximum or

inflection points. At these points f'(x) = 0. Also if we make a sign

chart we can see at which intervals the function is increasing and

decreasing, but the second derivative is a better test of this.



\textit{The Second Derivative Test}



The second derivative is used to find the points when a function is

concave or when it is convex at these points f''(x) = 0. Then you need

to make a sign chart. The interval(s) that the sign of the second

derivative is positive the function is concave and if the sign is

negative the function is convex. When f'(a) is zero and on that

interval f''(x) is negative then the point \textit{a} is a maximum point,

before this point the function will be increasing afterwards the

function will be decreasing. When f'(a) is zero and on that interval

f''(x) is positive then the point \textit{a} is a minimum point, before this

the function will be decreasing afterwards the function will be

increasing.



\textit{The Third Derivative Test}



The third derivative test determines if a point is an inflection point.

If f''(c) = 0 and $f'''(c) \ne 0,$ then f(c) is an inflection point.



\paragraph{An Example}



Sketch the graph of $f(x) = x^3 -4x^2$



1) The function will tell us that:



$f(x) = x^2(x - 4)$

$0 = x^2(x - 4)$ x = 0 and x=4 are the x-intercepts.

$y = 0^2(0 - 4)$ y = 0 is the y-intercept.



2) Then we use the first derivative test.



$f'(x) = 3x^2 - 8x$

$0 = x(3x - 8)$ x = 0 and x=2.66 are the minimum or maximum or

    inflection points.



3) Then we use the second derivative test.



f''(x) = 6x - 8

0 = 6x -8 so x = $\frac{4}{3}.$

We now make a number line \\_\\_f'''(0) = -8\\_\\_\\_\\_\\_\\_

    1.333\\_\\_\\_\\_f'''(2) = 4.

Before 1.333 the function will be convex after 1.333 the function

    will be concave.



4) Then we use the third derivative test.



f'''(x) = 6

$f'''(\frac{4}{3}) = 6$ The point is an inflection point.



5) Now we gather all the information into a table. Also we need to find

the corresponding y values to all the x values that we obtained from

steps 1 to 3 using f(x).



\begin{array}



$$

\begin{array}{|c|c|}

\hline

\text{Interval or Point} & \text{$f(x)$ behaviour} \\\\ \hline

x < 0 & \text{Increasing} \\\\ \hline

(0,0) & \text{Maximum Point, x-intercept, and y-intercept} \\\\ \hline

0 < x < 1.33 & \text{Decreasing} \\\\ \hline

(1.33,-4.74) & \text{Inflection Point} \\\\ \hline

1.33 < x < 2.66 & \text{Decreasing} \\\\ \hline

(2.66,-9.48) & \text{Minimum Point} \\\\ \hline

2.66 < x < $\infty$ & \text{Increasing} \\\\ \hline

(4,0) & \text{x-intercept} \\\\ \hline

\end{array}

$$                      



6)Now we can draw our graph:



% [Image: ]



\paragraph{Resources}



\begin{itemize}
    \item \href{https://en.wikibooks.org/wiki/A-level\_Mathematics/OCR/C1/Differentiation}{Differentiation},
\end{itemize}
    Wikibooks: A-level Mathematics

\begin{itemize}
    \item \href{https://mathresearch.utsa.edu/wikiFiles/MAT1214/Derivatives\%20and\%20the\%20Shape\%20of\%20a\%20Graph\_/MAT1214-4.5DerivativesShapeOfGraphsPwPt.pptx}{Derivatives and the Shape of a
\end{itemize}
    Graph}

    PowerPoint file created by Dr. Sara Shirinkam, UTSA.



<!-- -->



\begin{itemize}
    \item \href{https://mathresearch.utsa.edu/wikiFiles/MAT1214/Derivatives\%20and\%20the\%20Shape\%20of\%20a\%20Graph\_/MAT1214-4.5DerivativesShapeOfGraphsWS1.pdf}{Derivatives and the Shape of a Graph
\end{itemize}
    Worksheet}



\href{https://youtu.be/rPMfVi-FiuE}{Cal1 Intervals of Increase and Decrease

Part1} video lecture by Instructor Beatty



\href{https://youtu.be/9ZP6VxofDaY}{Cal1 Intervals of Increase and Decrease

Part2} video lecture by Instructor Beatty



\href{https://youtu.be/6dQY7j5vGT4}{Cal1 Intervals of Concavity Part1} video

lecture by Instructor Beatty



\href{https://youtu.be/BCb8KYhgR3Y}{Cal1 Intervals of Increase and Decrease and

Concavity} video lecture by Instructor

Beatty



\href{https://youtu.be/lM8Ej8VqKkw}{Cal1 Intervals of Increase Decrease and Concavity

Part2} video lecture by Instructor Beatty



\paragraph{Licensing}



Content obtained and/or adapted from:



\begin{itemize}
    \item \href{https://en.wikibooks.org/wiki/A-level\_Mathematics/OCR/C1/Differentiation}{Differentiation, Wikibooks: A-level
\end{itemize}
    Mathematics/OCR/C1}

    under a CC BY-SA license

\subsection{Learning Objectives}

\begin{enumerate}
    \item \textbf{Identify and Classify Stationary Points:} Students will be able to identify stationary points of a function where the first derivative $f'(x) = 0$  and classify them as local maxima, local minima, or potential inflection points. They should use the first derivative test to determine whether these points are maxima or minima based on the sign changes in $f'(x)$.
    \item \textbf{Apply the Second Derivative Test for Concavity:} Students will be able to use the second derivative $f''(x)$ to determine the concavity of a function and classify points of inflection. They should be able to identify intervals where the function is convex (concave up) or concave (concave down) based on the sign of $f''(x)$.
    \item \textbf{Construct and Analyze Graphs Based on Derivative Information:} Students will be able to combine information from the first and second derivative tests to sketch the graph of a function, including identifying key features such as stationary points, intervals of increase/decrease, concavity, and inflection points. They should be able to compile this information into a table summarizing the behavior of the function and use it to draw an accurate graph.
\end{enumerate}

\subsection{Assessment Instrument}

\begin{enumerate}
    \item \textbf{Part 1: Identify and Classify Stationary Points}
    \item \textbf{Calculation:} Find the stationary points of the function $f(x) = x^3 - 6x^2 + 9x + 1$. Determine where $f'(x) = 0$ and classify each stationary point as a local maximum, local minimum, or potential inflection point using the first derivative test.
    \item \textbf{Multiple Choice:} For the function $f(x) = x^4 - 4x^3 + 6x^2$, which of the following correctly identifies the stationary points and their classification?
    \item A. Local maxima at $x = 1$ and $x = 2$, local minima at $x = 0$
    \item B. Local minima at $x = 1$ and $x = 2$, local maxima at $x = 0$
    \item C. Local maxima at $x = 0$, $x = 1$ and $x = 2$
    \item D. Local minima at $x = 1$, $x = 2$, and potential inflection points at $x = 0$
    \item \textbf{Short Answer:} Determine the stationary points of the function $f(x) = e^x - 2x$ and classify them using the first derivative test. Explain your classification process. \textbf{Part 2: Apply the Second Derivative Test for Concavity}
    \item \textbf{Calculation:} Use the second derivative test to determine the concavity of the function $f(x) = x^3 - 3x^2 + 2x$. Identify intervals of concavity and any points of inflection.
    \item \textbf{True/False:} If $f''(x) > 0$ on an interval, then the function $f(x)$ is concave up on that interval. (True/False) Provide a brief explanation of why this is true or false.
    \item \textbf{Short Answer:} For the function $f(x) = \ln(x) - x$, compute the second derivative and use it to determine where the function is concave up, concave down, and if there are any inflection points. Show all steps. \textbf{Part 3: Construct and Analyze Graphs Based on Derivative Information}
    \item \textbf{Calculation:} Sketch the graph of the function $f(x) = x^3 - 3x^2 + 2x$. Use the information from the first and second derivatives to identify and label stationary points, intervals of increase/decrease, concavity, and inflection points. Summarize this information in a table.
    \item \textbf{Short Answer:} Given the function $f(x) = x^4 - 4x^2$, analyze its behavior based on the first and second derivatives. Create a detailed summary table that includes information about stationary points, intervals of increase/decrease, concavity, and inflection points. Use this summary to sketch the function's graph.
    \item \textbf{Multiple Choice:} Which of the following statements accurately describes the graph of the function $f(x) = x^3 - 6x^2 + 9x$ based on the first and second derivative tests?
    \item A. The function has a local minimum at $x = 1$, local maximum at $x = 3$, and is concave up for all $x$.
    \item B. The function has a local maximum at $x = 1$, local minimum at $x = 3$, and is concave down for $x < 1$ and concave up for $x > 1$.
    \item C. The function has a local minimum at $x = 1$ and $x = 2$, and is concave up for all $x$.
    \item D. The function has a local maximum at $x = 2$, local minimum at $x = 1$, and is concave down for all $x$.
\end{enumerate}

%%===============================================================
\section{The Second Derivative}
\label{sec:the-second-derivative}
%%===============================================================

In calculus, the \textbf{second derivative}, or the **second order

derivative**, of a function \textit{f} is the

derivative of the derivative of \textit{f}. Roughly

speaking, the second derivative measures how the rate of change of a

quantity is itself changing; for example, the second derivative of the

position of an object with respect to time is the instantaneous

acceleration of the object, or the rate at which the velocity of the

object is changing with respect to time. In Leibniz notation:



$$\mathbf{a} =  \frac{d\mathbf{v}}{dt} = \frac{d^2\boldsymbol{x}}{dt^2},$$



where \textit{a} is acceleration, \textit{v} is velocity, \textit{t} is time, \textit{x} is

position, and d is the instantaneous "delta" or change. The last

expression $\tfrac{d^2\boldsymbol{x}}{dt^2}$ is the second derivative of

position (x) with respect to time.



On the graph of a function, the second derivative corresponds to the

curvature or concavity of the graph. The graph of a function with a

positive second derivative is upwardly concave, while the graph of a

function with a negative second derivative curves in the opposite way.



% [Image: The second derivative of a quadratic function is

constant.]





\paragraph{Second derivative power rule}



The power rule for the first derivative, if applied twice, will produce

the second derivative power rule as follows:



$$\frac{d^2}{dx^2}\left[x^n\right] = \frac{d}{dx}\frac{d}{dx}\left[x^n\right] = \frac{d}{dx}\left[nx^{n-1}\right] = n\frac{d}{dx}\left[x^{n-1}\right] = n(n - 1)x^{n-2}.$$



\paragraph{Notation}



The second derivative of a function $f(x)$ is usually denoted $f''(x).$ That is:



$$f'' = \left(f'\right)'$$



When using Leibniz's notation for derivatives, the second derivative of a dependent variable

\textit{x} with respect to an independent variable \textit{x} is written



$$\frac{d^2y}{dx^2}.$$ This notation is derived from the following

formula:



$$\frac{d^2y}{dx^2} \,=\, \frac{d}{dx}\left(\frac{dy}{dx}\right).$$



\paragraph{Alternative notation}



As the previous section notes, the standard Leibniz notation for the

second derivative is $\frac{d^2y}{dx^2}.$ However, this form is not

algebraically manipulable. That is, although it is formed looking like a

fraction of differentials, the fraction cannot be split apart into

pieces, the terms cannot be cancelled, etc. However, this limitation can

be remedied by using an alternative formula for the second derivative.

This one is derived from applying the quotient rule to the first

derivative. Doing this yields the formula:



$$y''(x) = \frac{d}{dx}\left(\frac{dy}{dx}\right) = \frac{d\left(\frac{dy}{dx}\right)}{dx} = \frac{d^2y}{dx^2} - \frac{dy}{dx}\frac{d^2x}{dx^2}$$



In this formula, $du$ represents the differential operator applied to

$u,$ i.e., $d(u),$ $d^2u$ represents applying the differential operator

twice, i.e., $d(d(u)),$ and $du^2$ refers to the square of the

differential operator applied to $u,$ i.e., $(d(u))^2.$



When written this way (and taking into account the meaning of the

notation given above), the terms of the second derivative can be freely

manipulated as any other algebraic term. For instance, the inverse

function formula for the second derivative can be deduced from algebraic

manipulations of the above formula, as well as the chain rule for the

second derivative. Whether making such a change to the notation is

sufficiently helpful to be worth the trouble is still under debate.



\paragraph{Example}



Given the function



$$f(x) = x^3,$$



the derivative of \textit{f} is the function



$$f^{\prime}(x) = 3x^2.$$



The second derivative of \textit{f} is the derivative of $f^{\prime},$ namely



$$f^{\prime\prime}(x) = 6x.$$



\paragraph{Relation to the graph}



\paragraph{Concavity}



The second derivative of a function \textit{f} can be

used to determine the \textbf{concavity} of the graph of

\textit{f}. A function whose second derivative is

positive will be concave up (also referred to as convex), meaning that

the tangent line will lie below the graph of the function. Similarly, a

function whose second derivative is negative will be concave down (also

simply called concave), and its tangent lines will lie above the graph

of the function.



\paragraph{Inflection points}



If the second derivative of a function changes sign, the graph of the

function will switch from concave down to concave up, or vice versa. A

point where this occurs is called an \textbf{inflection point}. Assuming the

second derivative is continuous, it must take a value of zero at any

inflection point, although not every point where the second derivative

is zero is necessarily a point of inflection.



\paragraph{Second derivative test}



The relation between the second derivative and the graph can be used to

test whether a stationary point for a function (i.e., a point where

$f'(x)=0$) is a local maximum or a local minimum. Specifically,



\begin{itemize}
    \item If $f^{\prime\prime}(x) < 0,$ then $f$ has a local maximum at $x.$
    \item If $f^{\prime\prime}(x) > 0,$ then $f$ has a local minimum at $x.$
    \item If $f^{\prime\prime}(x) = 0,$ the second derivative test says
\end{itemize}
    nothing about the point $x,$ a possible inflection point.



The reason the second derivative produces these results can be seen by

way of a real-world analogy. Consider a vehicle that at first is moving

forward at a great velocity, but with a negative acceleration. Clearly,

the position of the vehicle at the point where the velocity reaches zero

will be the maximum distance from the starting position -- after this

time, the velocity will become negative and the vehicle will reverse.

The same is true for the minimum, with a vehicle that at first has a

very negative velocity but positive acceleration.



% [Image: A plot of $f(x) = \sin(2x)$ from $-\pi/4$ to $5\pi/4.$ The tangent

line is blue where the curve is concave up, green where the curve is

concave down, and red at the inflection points (0, $\pi$/2, and

$\pi$).]



\paragraph{Limit}



It is possible to write a single limit for the second derivative:



$$f''(x) = \lim\_{h \to 0} \frac{f(x+h) - 2f(x) + f(x-h)}{h^2}.$$



The limit is called the second symmetric derivative. Note that the

second symmetric derivative may exist even when the (usual) second

derivative does not.



The expression on the right can be written as a difference quotient of

difference quotients:



$$\frac{f(x+h) - 2f(x) + f(x-h)}{h^2} = \frac{\frac{f(x+h) - f(x)}{h} - \frac{f(x) - f(x-h)}{h}}{h}.$$



This limit can be viewed as a continuous version of the second

difference for sequences.



However, the existence of the above limit does not mean that the

function $f$ has a second derivative. The limit above just gives a

possibility for calculating the second derivative---but does not provide

a definition. A counterexample is the sign function $\operatorname{sgn}(x),$ which is

defined as:



$$\operatorname{sgn}(x) = 

\begin{cases}

-1 & \text{if } x < 0, \\\\ 0 & \text{if } x = 0, \\\\ 1 & \text{if } x > 0. 

\end{cases}$$



The sign function is not continuous at zero, and therefore the second

derivative for $x=0$ does not exist. But the above limit exists for $x=0$:



$$\begin{aligned}

\lim\_{h \to 0} \frac{\operatorname{sgn}(0+h) - 2\operatorname{sgn}(0) + \operatorname{sgn}(0-h)}{h^2} &= \lim\_{h \to 0} \frac{\operatorname{sgn}(h) - 2\cdot 0 + \operatorname{sgn}(-h)}{h^2} \\\&= \lim\_{h \to 0} \frac{\operatorname{sgn}(h) + (-\operatorname{sgn}(h))}{h^2} = \lim\_{h \to 0} \frac{0}{h^2} = 0.

\end{aligned}$$



\paragraph{Quadratic approximation}



Just as the first derivative is related to linear approximations, the

second derivative is related to the best quadratic approximation for a

function \textit{f}. This is the quadratic function

whose first and second derivatives are the same as those of

\textit{f} at a given point. The formula for the best

quadratic approximation to a function \textit{f}

around the point \textit{x} = \textit{a} is



$$f(x) \approx f(a) + f'(a)(x-a) + \tfrac12 f''(a)(x-a)^2.$$



This quadratic approximation is the second-order Taylor polynomial for the function centered at \textit{x} = \textit{a}



\paragraph{Resources}



\begin{itemize}
    \item \href{https://www.youtube.com/watch?v=OhqNbQi9QPk}{Concavity, Inflection Points, and Second
\end{itemize}
    Derivative}, The

    Organic Chemistry Tutor



\paragraph{Licensing}



Content obtained and/or adapted from:



\begin{itemize}
    \item \href{https://en.wikipedia.org/wiki/Second\_derivative}{Second derivative,
\end{itemize}
    Wikipedia} under a

    CC BY-SA license

\subsection{Learning Objectives}

\begin{enumerate}
    \item \textbf{Define the Second Derivative:} Students will be able to define the second derivative of a function, recognize its notation in both Lagrange's and Leibniz's notations, and explain its significance in measuring the rate of change of the first derivative.
    \item \textbf{Apply the Second Derivative to Graph Analysis:} Students will be able to use the second derivative to determine the concavity of a function's graph, identify inflection points, and apply the second derivative test to classify stationary points as local maxima or minima.
    \item \textbf{Compute and Interpret the Second Derivative:} Students will be able to compute the second derivative of a given function using differentiation rules, and interpret its geometric and physical meaning, including its role in quadratic approximation and acceleration.
\end{enumerate}

\subsection{Assessment Instrument}

\begin{enumerate}
    \item \textbf{Part 1: Define the Second Derivative}
    \item \textbf{Multiple Choice:} Which of the following represents the second derivative of a function $f(x)$ in Lagrange’s notation?
    \item A. $f'(x)$
    \item B. $f''(x)$
    \item C. $f''(a)$
    \item D. $f^{(2)}(x)$
    \item \textbf{True/False:} The second derivative of a function provides information about the rate at which the first derivative is changing. (True/False)
    \item \textbf{Short Answer:} Explain in your own words the significance of the second derivative in relation to the first derivative of a function. \textbf{Part 2: Apply the Second Derivative to Graph Analysis}
    \item \textbf{Multiple Choice:} If the second derivative of a function $f(x)$ is positive on an interval, then the graph of $f(x)$ is:
    \item A. Concave up
    \item B. Concave down
    \item C. Linear
    \item D. Neither concave up nor concave down
    \item \textbf{Fill in the Blank:} An inflection point occurs where the second derivative changes its sign, meaning that the graph of the function changes from ___________ to ___________ or vice versa.
    \item \textbf{Short Answer:} Given the function $f(x) = x^3 - 3x^2 + 2x$, find the second derivative and use it to determine the nature of any inflection points. \textbf{Part 3: Compute and Interpret the Second Derivative}
    \item \textbf{Calculation:} Compute the second derivative of the function $f(x) = 4x^4 - 2x^3 + 6x - 1$.
    \item \textbf{Interpretation:} For the function $f(x) = e^{-x^2}$, interpret the meaning of the second derivative in terms of the function's concavity and its graphical shape.
    \item \textbf{Application:} A particle moves along a straight line with position given by $s(t) = 3t^4 - 4t^3 + t^2$. Calculate the second derivative of $s(t)$ and explain its physical significance with respect to the particle's acceleration.
\end{enumerate}

%%===============================================================
\section{Optimization Applications}
\label{sec:optimization-applications}
%%===============================================================

Optimization is one of the uses of calculus in the real world. Perhaps

we are a pizza parlor and wish to maximize profit. Perhaps we have a

flat piece of cardboard and we need to make a box with the greatest

volume. How does one go about this process?



This requires the use of maximums and minimums. We know that we find

maximums and minimums via derivatives. Therefore, one can conclude that

calculus will be a useful tool for maximizing or minimizing

(collectively known as "optimizing") a situation.



In general, an optimization problem has a constraint that changes how we

view the problem. The wording of the problem (whether subtle or not) can

also drastically change how we view the problem. The constraint is the

thing we will absolutely not change (such as the general shape, or the

amount that we are willing to pay for the entire operation, etc.). If we

can identify what we are going to change and what we are going to

require stay the same, we will be on our way to solve the problem.



\paragraph{How to Solve}



These general steps should be taken in order to complete an optimization

problem.



1.  Write out necessary formulas and other pieces of information given

    by the problem.

\begin{itemize}
    \item The problems should have a variable you control and a variable
\end{itemize}
        that you want to maximize/minimize.

\begin{itemize}
    \item The formulas you find may contain extra variables. Depending on
\end{itemize}
        how the question works out, they may be substituted out or can

        be ignored (which will be explained later).

2.  Combine the formulas together so that the variable you want to

    maximize/minimize is on one side of the equation and everything else

    on the other.

3.  Differentiate the formula. If your equation has multiple variables,

    pick any variable to differentiate as long as it is not the one you

    control (i.e. pick the variable that you could not get rid of from

    the formula).

\begin{itemize}
    \item Note that during differentiation, if you come across a variable
\end{itemize}
        that you have not picked, imagine it as a number and apply the

        necessary differentiation rule. Do not treat it as a variable in

        this case.

4.  Set the differentiated formula to equal 0 and solve for the variable

    you control.

5.  The value you get here is your answer. If you instead have another

    formula, that means that your answer depends on those other

    variables, which would usually be what the question asked for if you

    have such a situation that you have another variable to juggle to

    begin with.



The reason why this algorithm works comes from a few mathematical

theorems which you will probably not need to know when completing these

problems. Usually the problems given will be mathematically simple (in

other words, there are not a lot of cases to test). However, if you wish

to know, they work like this:



\begin{itemize}
    \item A derivative of 0 is either a global or local maximum or minimum.
\end{itemize}
    Usually the question will tend towards answering that question

    without much difficulty (like always positive numbers, for example)



\paragraph{Examples}



\paragraph{Volume Example}



A box manufacturer desires to create a closed box with a surface area of

100 inches squared and with a square base yet rectangular sides. What is

the maximum volume that can be formed by bending this material into a

box?



\begin{itemize}
    \item Write out known formulas and information
\end{itemize}

$A\_{base}=x^2$



$A\_{side}=x\cdot h$



$A\_{total}=2x^2+4x\cdot h=100$



$V=l\cdot w\cdot h=x^2\cdot h$



\begin{itemize}
    \item Write the variable $h$ in terms of $x$ in the volume equation.
\end{itemize}

$2x^2+4xh=100$



$x^2+2xh=50$



$2xh=50-x^2$



$h=\frac{50-x^{2}}{2x}$



$\begin{aligned}

 V \&=\left(x^{2}\right)\left(\frac{50-x^2}{2x}\right) \\\ \&=\frac{50x-x^3}{2}

\end{aligned}$



\begin{itemize}
    \item Find the derivative of the volume equation in order to maximize the
\end{itemize}
    volume



$\frac{dV}{dx}=\frac{50-3x^2}{2}$



\begin{itemize}
    \item Set $\frac{dV}{dx}=0$ and solve for $x$
\end{itemize}

$\frac{50-3x^2}{2}=0$



$50-3x^2=0$



$3x^2=50$



$x=\pm\frac{\sqrt{50}}{\sqrt3}$



\begin{itemize}
    \item Plug-in the $x$ value into the volume equation and simplify
\end{itemize}

$\begin{aligned}

 V \&=\frac{1}{2}\left[50\sqrt{\frac{50}{3}}-\left(\sqrt{\frac{50}{3}}\right)^3\right] \\\ \&=68.04138174...

\end{aligned}$



Answer: $V\_{max}=68.04138174...$



\paragraph{Volume Example II}



It is desired to make an open-top box of greatest possible volume from a

square piece of tin whose side is $\alpha$ by cutting equal squares out

of the corners and then folding up the tin to form the sides. What

should be the length of the side of a cut-out square?



% [Image: ]



If we call the side length of the cut-out squares $x,$ then each side of

the base of the folded box is $\alpha-2x,$ and the height is $x.$

Therefore, the volume function is

$V(x)=x(\alpha-2x)^2=x(\alpha^2-4\alpha x+4x^2)=\alpha^2x-4\alpha x^2+4x^3.$

We must optimize the volume by taking the derivative of the volume

function and setting it equal to 0. Since it does not change, $\alpha$

is treated as a constant, not a variable.



$$V(x)=\alpha^2x-4\alpha x^2+4x^3$$



$$\begin{aligned}

 0&=V'(x)=\alpha^2-8\alpha x+12x^2 \\\ &= 12x^{2}-8\alpha x+\alpha^{2}

\end{aligned}$$



We can now use the quadratic formula to solve for $x$:



$$x=\frac{-b\pm\sqrt{b^2-4ac}}{2a}$$



$\begin{aligned}

 x \&=\frac{-(-8\alpha)\pm\sqrt{(-8\alpha)^2-4(12)(\alpha^2)}}{2(12)} \\\ \&=\frac{8\alpha\pm\sqrt{64\alpha^2-48\alpha^2}}{24} \\\ \&=\frac{8\alpha\pm\sqrt{16\alpha^2}}{24} \\\ \&=\frac{8\alpha\pm 4\alpha}{24}

\end{aligned}$



$$x=\frac{\alpha}{6},\frac{\alpha}{2}$$



We reject $x=\frac{\alpha}{2},$ since it is a minimum (it results in the base length $\alpha-2x$ being

0, making the volume 0).



Therefore, the answer is $x=\frac{\alpha}{6}.$



\paragraph{Volume Example III}



A cylindrical can is required to hold $1.024\text{ L}$ of liquid.

Determine the measurements that will minimize the amount of material

required to construct the can of specified volume.



This is a classic problem for optimization. The problem is we do not

have any specified measurements. Luckily enough, we are given SI

units, so we can write what we are given accordingly,



\begin{itemize}
    \item Conversion factor: $1\text{ L}=1000\text{ cm}^{3}.$ This means that
\end{itemize}
    through dimensional analysis (or basic algebra):

    $1.024\text{ L}=1024\text{ cm}^{3}.$ Recall that this is the volume!

\begin{itemize}
    \item Volume of a cylinder: $V=\pi r^{2}h$
    \item Surface area of a cylinder: $S=2\pi rh+2\pi r^{2}$
\end{itemize}

Recall what we are trying to minimize: "the amount of material required

to construct the can of specified volume." This means the surface area

needs to be minimized! The unfortunate part is we have two variables (it

is not constant). Luckily, we can use the volume to write $h$ in terms

of $r$:



$$1024\text{ cm}^{3}=\pi r^{2}h$$



$$\Rightarrow \frac{1024\text{ cm}^{3}}{\pi {r\text{ cm}}^{2}}=\frac{1024}{\pi r^{2}}\text{ cm}=h$$



We can the substitute this information into the surface area formula and

minimize (set the first derivative to zero):



$$S(r)=2\pi r\left(\frac{1024}{\pi r^{2}}\right)+2\pi r^{2}.$$

 

Eliminate the $\pi$ and $r.$



$$\Rightarrow S(r)=2\cdot\frac{1024}{r}+2\pi r^{2}$$



Set the first derivative equal to zero.



$$\begin{aligned}

 \Rightarrow0=S'(r)&=\frac{d}{dr}\left(\frac{2048}{r}+2\pi r^{2}\right) \\\ &=\frac{d}{dr}\left(\frac{2048}{r}\right)+\frac{d}{dr}\left(2\pi r^{2}\right) \\\ &=\frac{d}{dr}\left(2048\cdot{r}^{-1}\right)+2\pi\frac{d}{dr}\left(r^{2}\right) \\\ &=2048\cdot\frac{d}{dr}\left({r}^{-1}\right)+2\pi\left(2r\right) \\\ &=2048\cdot(-1){r}^{-2}+4\pi r \\\ &=-\frac{2048}{{r}^{2}}+4\pi r \\\ &=\frac{4\pi r^{3}-2048}{{r}^{2}}

\end{aligned}$$



All that is left to do is find the critical values of the

operation. This ultimately means to let the numerator equal zero and

find what values of $r$ make the denominator equal to zero. Here is the

operation for the numerator:



$4\pi r^{3}-2048=0$$



$$\Leftrightarrow 4\pi r^{3}=2048$$



$$\Leftrightarrow r^{3}=\frac{2048}{4\pi}=\frac{512}{\pi}$$



$$\Leftrightarrow r=\sqrt[3]{\frac{512}{\pi}}=\frac{8}{\sqrt[3]{\pi}}\approx5.4623$$



Here is the operation for the denominator:



$${r}^{2}=0$$



$$r=\pm0$$



Let us note that we do not care about one of the critical

values being equal to zero. This problem is asking us to find the value

that makes the condition true but also minimized. While zero does

minimize the amount of material used to construct the can, we cannot

make a can if we do not show one. This solution would be interesting,

however, if we were asked about the graph of $S(r).$ Nevertheless, our

final answer may not necessarily be $r\approx5.4623\text{ cm}$ and

$h\approx10.9245\text{ cm}.$ We need to verify that the critical value

we found is the correct answer. That is where we substitute values for

$r\in\left(0,\,\frac{8}{\sqrt[3]{\pi}}\right)$ and

$r\in\left(\frac{8}{\sqrt[3]{\pi}},\,\infty\right)$ into the first

derivative:



For $r\in\left(0,\,\frac{8}{\sqrt[3]{\pi}}\right),$ let $r=5$: $S'(5)=\frac{4\pi(5)^{3}-2048}{(5)^{2}}0.$



Based on these results above, it is clear that

$r=\frac{8}{\sqrt[3]{\pi}}$ is a value that minimizes the function such

that it allows us to make a can. Keep in mind, if you do not show that

the value we found is a minimum, then this value could also be a

maximum. Imagine if that were the case, and we suggested the company to

produce a can that would maximize the amount of material used.

Potentially millions of dollars of material costs would offset the

profit from the optimum. Either way, we have found our answers:



Answer: the radius is $5.4623\text{ cm}$ and the height is $10.9245\text{ cm}.$



Sales Example



A small retailer can sell $n$ units of a product for a revenue of

$r(n)=8.1n$ and at a cost of $c(n)=n^3-7n^2+18n,$ with all amounts in

thousands. How many units does it need to sell to maximize its profit?



![](https://upload.wikimedia.org/wikipedia/commons/6/6a/Calculus_Graph-Finding_Maximum_Profit.png)



The retailer's profit is defined by the equation $p(n)=r(n)-c(n),$

which is the net revenue generated. The question asks for the maximum

amount of profit, or the maximum of the above equation. As previously

discussed, the maxima and minima of a graph are found when the slope of

the said graph is equal to zero. To find the slope, one finds the

derivative of the function of interest -- here, $p(n).$ By using the

subtraction rule $p'(n)=r'(n)-c'(n)$:



$$p(n)=r(n)-c(n)$$



$$p'(n)=\frac{d}{dn}[8.1n]-\frac{d}{dn}\left[n^3-7n^2+18n\right]=-3n^2+14n-9.9$$



Therefore, when $-3n^{2}+14n-9.9=0$ the profit will be maximized or

minimized. Use the quadratic formula to find the roots, giving

$n=0.869,3.798.$ To find which of these is the maximum and minimum the

function can be tested:



$$p(0.869)=-3.97321,p(3.798)=8.58802$$



Because we only consider the functions for all $n\ge 0$ (i.e., you can't have $n=-5$ units), the

only points that can be minima or maxima are those two listed above. To

show that $3.798$ is in fact a maximum (and that the function doesn't

remain constant past this point) check if the sign of $p'(n)$ changes at this point.



Here, it does, and for $n$ greater than $3.798,$ $P'(n)$ will continue

to decrease. This demonstrates that the firm will be maximizing its

profits at $n=3.798.$ The retailer selling $3.798$ units would return a

profit of $8,588.02.$



However, the answer is $n=3.798.$



Revenue Example



A widget firm has a linear demand function $D(q)=-\frac{1}{5}q+10$ where

$q$ kilograms of widgets (in hundreds) are shipped and sold to

customers. At what price should the widgets be sold such that the

revenue is maximized?



The revenue of a firm is calculated by multiplying the quantity sold to

the price. Since $D(q)$ is the price, and $q$ is the quantity sold,

$R(q)=q\cdot D(q).$ The question is asking us to find the price that

will maximize the quantity. This means finding the quantity is the first

goal. The final goal is the price found.



![Maximizing

revenue|490x490px](https://upload.wikimedia.org/wikipedia/commons/thumb/0/06/Maximizing\_revenue.png/490px-Maximizing\_revenue.png)



Our first goal is finding the quantity that will maximize the revenue.

The revenue function is maximized if the slope of the tangent line is

zero at a point and the sign of the slope around those points change

sign from positive to negative. This is achieved through taking the

derivative and applying the first derivative test. This is made easier

when the revenue function distributes the $q$:



$$\begin{aligned}

 R(q)\&=q\cdot D(q) \\\ \&=q\cdot\left(-\frac{1}{5}q+10\right) \\\ \&=-\frac{1}{5}q^{2}+10q

\end{aligned}$$



Let us now take the first derivative:



$$\frac{d}{dq}\left[R(q)\right]=\frac{d}{dq}\left[-\frac{1}{5}q^{2}+10q\right]$$



$$\begin{aligned}

 R'(q)\&=\frac{d}{dq}\left[-\frac{1}{5}q^{2}\right]+\frac{d}{dq}\left[10q\right] \\\ \&=-\frac{1}{5}\frac{d}{dq}\left[q^{2}\right]+10\frac{d}{dq}\left[q\right] \\\ \&=-\frac{1}{5}\left[2q\right]+10\left[1\right] \\\ \&=-\frac{2}{5}q+10

\end{aligned}$$



Set the first derivative equal to zero to find the critical values:



$$R'(q)=0=-\frac{2}{5}q+10$$



$$\Leftrightarrow-10=-\frac{2}{5}q$$



$$\Leftrightarrow25=q$$



Take note of the original revenue function.

Because the revenue function is a quadratic that has its $q^{2}$ term

multiplied by a factor less than 0, it is guaranteed there is only one

critical value, and it must be a maximum. The calculus-based approach

simply confirms this truth (using the first derivate test):



For any $q0.$ For any $q>25,$ $R'(q)\frac{1}{2},$

    $D'(x)>0.$ This implies the function before the critical value is

    decreasing, and the function after the critical value is increasing.

    This implies the function at the critical value is a local minimum.



As such, the $x$-value that will minimize the distance between $(x,y)$

and the origin is $x=\frac{1}{2}.$



However, you have not finished the question yet. You still have to find

the point of the function that the particle needs to stop at. This is

simply a matter of substitution into the function $f(x)$:



$$\begin{aligned}

 f\left(\frac{1}{2}\right)\&=\sqrt{4-\left(\frac{1}{2}\right)} \\\ \&=\sqrt{\frac{8}{2}-\frac{1}{2}} \\\ \&=\sqrt{\frac{7}{2}}

\end{aligned}$$



Most teachers in calculus classes don't mind if you leave

the answer as $\left(\frac{1}{2},\sqrt{\frac{7}{2}}\right).$ However,

for the minuscule number of teachers or testing organizations that care

about rationalizing the denominator, the work for $\sqrt{\frac{7}{2}}$

is done below:



$$\begin{aligned}

 \sqrt{\frac{7}{2}}\&=\frac{\sqrt{7}}{\sqrt{2}} \\\ \&=\frac{\sqrt{7}}{\sqrt{2}}\cdot\frac{\sqrt{2}}{\sqrt{2}} \\\ \&=\frac{\sqrt{7}\cdot\sqrt{2}}{2} \\\ \&=\frac{\sqrt{14}}{2} \\\\end{aligned}$$



This is the reason most teachers don't mind if the

denominator is left as a square root; it distracts from the calculus

(although slightly). This leaves us with the final answer:



The particle must leave the function at

$\left(\frac{1}{2},\frac{\sqrt{14}}{2}\right)$ to minimize the distance

required to get to the origin.$\blacksquare$



Non-standard way: implicit differentiation



Only the implicit differentiation will be shown, since any calculations

beyond that would result in the same procedures we outlined before.

Recall $D(x)=\sqrt{x^{2}-x+4}.$ Square both sides to return the function

seen below:



$$\left[D(x)\right]^{2}=x^{2}-x+4$$



Taking the derivative of the above function with respect to $x$ will require us to implicitly differentiate

(although this differentiation is neatly disguised). The calculations are demonstrated below:



$$\frac{d}{dx}\left[\left(D(x)\right)^{2}\right]=\frac{d}{dx}\left(x^{2}-x+4\right)$$



$$\begin{aligned}

 \Rightarrow2D(x)\frac{dD}{dx}\&=\frac{d}{dx}\left(x^{2}\right)-\frac{d}{dx}\left(x\right)+\frac{d}{dx}\left(4\right) \\\ \&=2x-1+0

\end{aligned}$$



$$2D(x)\frac{dD}{dx}=2x-1$$



$$\begin{aligned}

 \Leftrightarrow\frac{dD}{dx}\&=\frac{2x-1}{2D(x)} \\\ \&=\frac{2x-1}{2\sqrt{x^{2}-x+4}}

\end{aligned}$$



Notice how the calculation above resulted in the same first derivative when done the standard way, meaning the calculations

from herein will be the same. With this let us move on to the next problem.



Maximizing the Rectangular Area Formed by a Curved Shelf



Let there be a function $f(x)=15\log(x)$ that describes the curve of one

part of a stylized shelf. The curve ends $100$ centimeters horizontal of

the vertical wall. The manufacturer wants the front face of this portion

to have one large maximum rectangular indentation whereby the top-left

corner of the indentation forms a horizontal part that is parallel to

the shelf's bottom. How far from the left-hand side of the shelving

unit must the vertical wall formed in between be such that the condition

set by the manufacturer is true? A calculator may be used here.



This problem seems very weird. However, in actuality, this is a

disguised variation of another common optimization problem, the maximum

rectangular area formed by a function problem. The problem becomes much

more manageable when thinking about it in those terms. What the

manufacturer wants, after all, is a maximized indentation size. It is

not asking about the volume of that size, only the part where you can

find this is actually true.



The image shown on the right is actually the maximum area formed.

However, if the problem did not provide an image (and no numbers on the

axes), then one would not be approximating.



![The problem here is a disguised variation of another

problem.](https://upload.wikimedia.org/wikipedia/commons/thumb/6/60/Optimization_Problem_2.jpg/450px-Optimization_Problem_2.jpg)



Let $x$ be the distance from the left-hand side of the shelving unit,

and let the height of the rectangle formed be $y=f(x).$ Such a

construction means that $(100-x)$ is the length of the rectangle while

the height is $15\log(x).$ Therefore, the area function is



$$A(x)=15(100-x)\log(x)$$



We want to find the maximum of the function

given herein. Let us find the derivative of the area and set it equal to zero:



$$A(x)=15(100-x)\log(x)$$



$$\Rightarrow\frac{d}{dx}\left[A(x)\right]=\frac{d}{dx}\left[15(100-x)\log(x)\right]$$



$$\begin{aligned}

 \Rightarrow A'(x)\&=15\frac{d}{dx}\left((100-x)\log(x)\right) \hspace{1cm} \left(\text{Recall the product rule: }\frac{d}{dx}\left[f(x)\cdot g(x)\right]=f\cdot\frac{dg}{dx}+g\cdot\frac{df}{dx}\right) \\\ \&=15\left[(100-x)\cdot\frac{d}{dx}\left(\log(x)\right)+\log(x)\cdot\frac{d}{dx}\left(100-x\right)\right] \\\ \&=15\left[\frac{100-x}{x\ln(10)}-\log(x)\right] \\\ \&=15\left[\frac{100-x}{x\ln(10)}-\frac{\ln(x)}{\ln(10)}\right] \hspace{1cm} \left(\text{Multiply by $\frac{x}{x} to the right fraction.}\right) \\\ \&=15\left[\frac{100-x}{x\ln(10)}-\frac{x\ln(x)}{x\ln(10)}\right] \\\ \&=15\left[\frac{100-x-x\ln(x)}{x\ln(10)}\right]=0

\end{aligned}$$ 



However, after finding the first derivative, it is evident

that finding the value for $x$ in this specific situation is unknown to

us at this time because we do not know about the Lambert-W Function.

Therefore, a graphing calculator is acceptable to use now (typically,

many teachers want to see as much work as possible before allowing you

to use the calculator for solutions). If one does not have a graphing

calculator, use Desmos(https://www.desmos.com/calculator) or

Geogebra(https://www.geogebra.org/graphing?lang=en) or the many other

online graphing calculators.



The answer obtained by those websites is $x\approx 23.94721$ or an exact

answer of $x=e^{W(100e)-1}.$ We know this is also the maximum that we

are looking for because the derivative function is positive for

$xv\_{2}$ where $v\_{1}$ is the speed

of light traveling through the first material of index of refraction

$n\_{1}$ (and similar for $v\_{2}$).



-   $t\_{1}=\frac{\sqrt{\Lambda^{2}+\Upsilon^{2}}}{v\_{1}}$

-   $t\_{2}=\frac{\sqrt{(d-\Lambda)^{2}+\Gamma^{2}}}{v\_{2}}$



The time light took to travel through the two materials is given as a

function of $\Lambda,$ so



-   $T(\Lambda)=t\_{1}+t\_{2}=\frac{\sqrt{\Lambda^{2}+\Upsilon^{2}}}{v\_{1}}+\frac{\sqrt{(d-\Lambda)^{2}+\Gamma^{2}}}{v\_{2}}$



We now finally have an optimization problem. We want to minimize the

amount of time travelled. Set the first derivative of $T(\Lambda)$ equal

to zero. Remember, everything else besides $\Lambda$ is a constant.



$$T(\Lambda)=\frac{\sqrt{\Lambda^{2}+\Upsilon^{2}}}{v\_{1}}+\frac{\sqrt{(d-\Lambda)^{2}+\Gamma^{2}}}{v\_{2}}$$



$$\frac{d}{d\Lambda}\left[T(\Lambda)\right]=\frac{d}{d\Lambda}\left(\frac{\sqrt{\Lambda^{2}+\Upsilon^{2}}}{v\_{1}}+\frac{\sqrt{(d-\Lambda)^{2}+\Gamma^{2}}}{v\_{2}}\right)$$



$$\begin{aligned}

 0=T'(\Lambda)\&=\frac{d}{d\Lambda}\left(\frac{\sqrt{\Lambda^{2}+\Upsilon^{2}}}{v\_{1}}+\frac{\sqrt{(d-\Lambda)^{2}+\Gamma^{2}}}{v\_{2}}\right) \\\ \&=\frac{d}{d\Lambda}\left(\frac{\sqrt{\Lambda^{2}+\Upsilon^{2}}}{v\_{1}}\right)+\frac{d}{d\Lambda}\left(\frac{\sqrt{(d-\Lambda)^{2}+\Gamma^{2}}}{v\_{2}}\right) \\\ \&=\frac{1}{v\_{1}}\cdot\frac{d}{d\Lambda}\left(\sqrt{\Lambda^{2}+\Upsilon^{2}}\right)+\frac{1}{v\_{2}}\cdot\frac{d}{d\Lambda}\left(\sqrt{(d-\Lambda)^{2}+\Gamma^{2}}\right) \\\ \&=\frac{1}{v\_{1}}\cdot\left(\frac{2\Lambda}{2\sqrt{\Lambda^{2}+\Upsilon^{2}}}\right)+\frac{1}{v\_{2}}\cdot\left(\frac{2(d-\Lambda)\cdot(-1)}{2\sqrt{(d-\Lambda)^{2}+\Gamma^{2}}}\right) \\\ \&=\frac{1}{v\_{1}}\cdot\left(\frac{\Lambda}{\sqrt{\Lambda^{2}+\Upsilon^{2}}}\right)-\frac{1}{v\_{2}}\cdot\left(\frac{d-\Lambda}{\sqrt{(d-\Lambda)^{2}+\Gamma^{2}}}\right)

\end{aligned}$$



Recall what we are trying to solve for. We are not trying

to solve for $\Lambda$; we are trying to model the situation. If we did

have to solve for $\Lambda,$ it would be a nightmare (we may even leave

it as a non-trivial exercise for the reader), but nevertheless, we can

begin modeling the situation. First, move the negative term to the other side:



$$\frac{1}{v\_{1}}\cdot\left(\frac{\Lambda}{\sqrt{\Lambda^{2}+\Upsilon^{2}}}\right)-\frac{1}{v\_{2}}\cdot\left(\frac{d-\Lambda}{\sqrt{(d-\Lambda)^{2}+\Gamma^{2}}}\right)=0$$



$$\frac{1}{v\_{1}}\cdot\left(\frac{\Lambda}{\sqrt{\Lambda^{2}+\Upsilon^{2}}}\right)=\frac{1}{v\_{2}}\cdot\left(\frac{d-\Lambda}{\sqrt{(d-\Lambda)^{2}+\Gamma^{2}}}\right)$$



![The diagram relating to the angles involved as a result of applying

different geometric

properties.|430x430px](https://upload.wikimedia.org/wikipedia/commons/thumb/9/95/Snell%27s_Law_%28diagram_2%29.svg/430px-Snell%27s_Law_%28diagram_2%29.svg.png)



The next thing we should attempt to do is rewrite the expression so that

we clean-up those fractions (they are intimidating after all, and it

would not be practical to model).



We can use the trigonometric functions to help us here in this

department. Because the trigonometric functions are simply relations of

the ratio of a right triangle, this would be a very good solution.



The *normal* is defined as the line perpendicular a given "surface".

Because the distance $\Upsilon$ is parallel to the normal, the angle

generated by that is also $90^{\circ}.$ Because the horizontal distance

$d$ light travelled is parallel to the "surface," the incident angle

of the light, $\theta\_{1},$ is an angle between the distance $\Upsilon$

and the incident light vector $\mathbf{I},$ and the light vector

$\mathbf{I}$ is a transversal, by the corresponding angles theorem, the

angle between the distance $\Upsilon$ and the incident light vector

$\mathbf{I}$ is also $\theta\_{1}$ (see the upper-left corner of the

diagram for reference).



The result of this geometric derivation is simply this: we have found an

angle of part of a right triangle. As a result, we can write the

distance as a trigonometric equivalent



-   $\sin(\theta\_{1})=\frac{\Lambda}{\sqrt{\Lambda^{2}+\Upsilon^{2}}}$



We can use similar logic to show that



-   $\sin(\theta\_{2})=\frac{d-\Lambda}{\sqrt{(d-\Lambda)^{2}+\Gamma^{2}}}$



As a result, we can rewrite the equations involved into something much

more manageable and easy to calculate without measuring ratios:



$$\frac{1}{v\_{1}}\cdot\left(\frac{\Lambda}{\sqrt{\Lambda^{2}+\Upsilon^{2}}}\right)=\frac{1}{v\_{2}}\cdot\left(\frac{d-\Lambda}{\sqrt{(d-\Lambda)^{2}+\Gamma^{2}}}\right)$$



$$\Leftrightarrow\frac{1}{v\_{1}}\cdot\sin(\theta\_{1})=\frac{1}{v\_{2}}\sin(\theta\_{2})$$

We are technically not done with the problem. Although we may have

related the two different light vectors, we did not model the situation

in terms of index of refractions. This leaves us with the two final

steps. Recall the speed of light going through a certain material of

index of refraction $n$ is $v=\frac{c}{n}.$ This is equivalent to saying

$\frac{1}{v}=\frac{n}{c}.$ Since the speeds relate to the index of

refraction, we may rewrite the above equation as



$$\frac{n\_{1}\cdot\sin(\theta\_{1})}{c}=\frac{n\_{2}\cdot\sin(\theta\_{2})}{c}$$



Multiply both sides by $c$ and we finally found \textit{Snell's Law}:



Answer: $n\_{1}\sin(\theta\_{1})=n\_{2}\sin(\theta\_{2})$



Keep in mind, this result is still a derivative. However, we are allowed

to model real life scenarios in terms of differential equations.

Sometimes, it may be impossible to model situations without differential

equations. We are allowed to write our equation like this (physics

certainly does not care here). Finally, although this is not \textit{strictly}

an optimization problem, thinking about in that sense allowed us to

obtain a nice rule of the world (a \textit{law} in science). We proved this

behavior, so it seems light will be consistent according to this

equation. This in effect has allowed us to confirm in 2003 that negative

index of refraction exists (even through experimental observations) and

allowed us to discover properties of spin waves.



\paragraph{Licensing}



Content obtained and/or adapted from:



\begin{itemize}
    \item \href{https://en.wikibooks.org/wiki/Calculus/Optimization}{Optimization,
\end{itemize}
    Wikibooks}

    under a CC BY-SA license

\subsection{Learning Objectives}

\begin{enumerate}
    \item \textbf{Identifying Variables and Constraints in Optimization Problems}
    \item Students will be able to identify and articulate the constraint and variable to be optimized in a given optimization problem.
    \item \textbf{Applying Derivatives to Find Maximum and Minimum Values}
    \item Students will be able to apply the differentiation technique to derive and solve the formula representing the variable to be optimized in order to find its maximum or minimum value.
    \item \textbf{Utilizing Optimization Techniques for Real-World Applications}
    \item Students will be able to analyze the second derivative or conduct the first derivative test to confirm whether the critical points found in an optimization problem correspond to maximum or minimum values.
\end{enumerate}

\subsection{Assessment Instrument}

\begin{enumerate}
    \item \textbf{Part 1: Identifying Variables and Constraints in Optimization Problems}
    \item \textbf{Multiple Choice:} In the optimization problem of maximizing the area of a rectangle with a fixed perimeter, which of the following represents the constraint and the variable to be optimized?
    \item A. Constraint: Area, Variable: Perimeter
    \item B. Constraint: Perimeter, Variable: Area
    \item C. Constraint: Length and Width, Variable: Perimeter
    \item D. Constraint: Length and Width, Variable: Area
    \item \textbf{True/False:} In a problem involving maximizing profit subject to budget constraints, the budget constraint represents the variable to be optimized. (True/False)
    \item \textbf{Short Answer:} Describe how to identify the variable to be optimized and the constraints in the problem of minimizing the cost of manufacturing a cylindrical can with a fixed volume. \textbf{Part 2: Applying Derivatives to Find Maximum and Minimum Values}
    \item \textbf{Calculation:} Given the cost function $C(x) = 5x^2 + 10x + 15$ representing the cost of producing $x$ units of a product, find the value of $x$ that minimizes the cost.
    \item \textbf{Multiple Choice:} To find the maximum value of the function $f(x) = -2x^2 + 4x + 6$, which of the following steps is correct?
    \item A. Set $f'(x) = 0$ and solve for $x$ to find the critical points, then use $f''(x)$ to determine concavity.
    \item B. Use $f''(x) = 0$ to find the maximum point directly.
    \item C. Find $f(x)$ at $x = 0$ and $x = 1$ to compare values.
    \item D. Set $f(x) = 0$ and solve for $x$ to find maximum values.
    \item \textbf{Short Answer:} Apply differentiation to find the minimum value of the function $g(x) = x^3 - 6x^2 + 9x + 2$. Explain the process and provide the minimum value. \textbf{Part 3: Utilizing Optimization Techniques for Real-World Applications}
    \item \textbf{Calculation:} Given the revenue function $R(x) = -3x^2 + 24x$, where $x$ is the number of units sold, compute the second derivative and use it to determine if the critical point corresponds to a maximum or minimum revenue.
    \item \textbf{True/False:} If the first derivative of a function changes from positive to negative at a critical point, then the function has a local maximum at that point. (True/False)
    \item \textbf{Short Answer:} In a problem where you need to minimize the surface area of a box with a fixed volume, find the dimensions that minimize the surface area and explain how you used optimization techniques to reach your conclusion.
\end{enumerate}

%%===============================================================
\section{Implicit Differentiation}
\label{sec:implicit-differentiation}
%%===============================================================

Generally, you will encounter functions expressed in explicit form, that

is, in the form $y=f(x)$ . To find the derivative of $y$ with respect to

$x$ , you take the derivative with respect to $x$ of both sides of the

equation to get



$$\frac{dy}{dx}=\frac{d}{dx}[f(x)]=f'(x)$$



But suppose you have a relation of the form $f\bigl(x,y(x)\bigr)=g\bigl(x,y(x)\bigr)$ . In this

case, it may be inconvenient or even impossible to solve for $y$ as a

function of $x$ . A good example is the relation $y^2+2yx+3=5x$ . In

this case you can utilize \textbf{implicit differentiation} to find the

derivative. To do so, one takes the derivative of both sides of the

equation with respect to $x$ and solves for $y'$ . That is, form



$$\frac{d}{dx}\bigl[f\bigl(x,y(x)\bigr)\bigr]=\frac{d}{dx}\big[g\bigl(x,y(x)\bigr)\big]$$



and solve for $\frac{dy}{dx}$ . You need to employ the chain rule

whenever you take the derivative of a variable with respect to a

different variable. For example,



$$\frac{d}{dx}(y^3)=\frac{d}{dy}[y^3]\cdot\frac{dy}{dx}=3y^2\cdot y'$$



\paragraph{Implicit Differentiation and the Chain Rule}



To understand how implicit differentiation works and use it effectively

it is important to recognize that the key idea is simply the chain rule.

First let's recall the chain rule. Suppose we are given two

differentiable functions $f(x),g(x)$ and that we are interested in

computing the derivative of the function $f\bigl(g(x)\bigr)$ , the chain

rule states that:



$$\frac{d}{dx}\bigl[f\bigl(g(x)\bigr)\bigr]=f'\bigl(g(x)\bigr)\cdot g'(x)$$

That is, we take the derivative of $f$ as normal and then plug in $g$ ,

finally multiply the result by the derivative of $g$ .



Now suppose we want to differentiate a term like $y^2$ with respect to

$x$ where we are thinking of $y$ as a function of $x$ , so for the

remainder of this calculation let's write it as $y(x)$ instead of just

$y$ . The term $y^2$ is just the composition of $f(x)=x^2$ and $y(x)$ .

That is, $f\bigl(y(x)\bigr)=y(x)^2$ . Recalling that $f'(x)=2x$ then the

chain rule states that:



$$\frac{d}{dx}\bigl[f\bigl(y(x)\bigr)\bigr]=f'\bigl(y(x)\bigr)\cdot y'(x)=2y(x)y'(x)$$



Of course it is customary to think of $y$ as being a function of $x$

without always writing $y(x)$ , so this calculation usually is just

written as



$$\frac{d}{dx}(y^2)=2yy'$$ Don't be confused by the fact that we don't

yet know what $y'$ is, it is some function and often if we are

differentiating two quantities that are equal it becomes possible to

explicitly solve for $y'$ (as we will see in the examples below.) This

makes it a very powerful technique for taking derivatives.



\paragraph{Explicit Differentiation}



For example, suppose we are interested in the derivative of $y$ with

respect to $x$ , where $x,y$ are related by the equation



$$x^2+y^2=1$$



This equation represents a circle of radius 1 centered on the origin.

Note that $y$ is not a function of $x$ since it fails the vertical line

test ($y=\pm1$ when $x=0$ , for example).



To find $y'$ , first we can separate variables to get



$$y^2=1-x^2$$



Taking the square root of both sides we get two separate functions for

$y$ :



$$y=\pm\sqrt{1-x^2}$$



We can rewrite this as a fractional power:



$$y=\pm(1-x^2)^{\frac{1}{2}}$$



Using the chain rule we get,



$$y'=\pm\frac{(1-x^2)^{-\frac{1}{2}}\cdot(-2x)}{2}=\mp\frac{x}{\sqrt{1-x^2}}$$



And simplifying by substituting $y$ back into this equation gives



$$y'=-\frac{x}{y}$$



\paragraph{Implicit Differentiation}



Using the same equation



$$x^2+y^2=1$$



First, differentiate with respect to $x$ on both sides of the equation:



$$\frac{d}{dx}[x^2+y^2]=\frac{d}{dx}[1]$$



$$\frac{d}{dx}[x^2]+\frac{d}{dx}[y^2]=0$$



To differentiate the second term on the left hand side of the equation

(call it $f(y(x))=y^2$), use the chain rule:



$$\frac{df}{dx}=\frac{df}{dy}\cdot\frac{dy}{dx}=2y\cdot y'$$



So the equation becomes



$$2x+2yy'=0$$



Separate the variables:



$$2yy'=-2x$$



Divide both sides by $2y$ , and simplify to get the same result as

above:



$$y'=-\frac{2x}{2y}$$



$$y'=-\frac{x}{y}$$



\paragraph{Uses}



Implicit differentiation is useful when differentiating an equation that

cannot be explicitly differentiated because it is impossible to isolate

variables.



For example, consider the equation,



$$x^2+xy+y^2=16$$



Differentiate both sides of the equation (remember to use the product

rule on the term $xy$):



$$2x+y+xy'+2yy'=0$$



Isolate terms with $y'$:



$$xy'+2yy'=-2x-y$$



Factor out a $y'$ and divide both sides by the other term:



$$y'=\frac{-2x-y}{x+2y}$$



\paragraph{Example}



$$xy=1$$



can be solved as:



$$y=\frac{1}{x}$$



then differentiated:



$$\frac{dy}{dx}=-\frac{1}{x^2}$$



However, using implicit differentiation it can also be differentiated

like this:



$$\frac{d}{dx}[xy]=\frac{d}{dx}[1]$$



use the product rule:



$$x\frac{dy}{dx}+y=0$$



solve for $\frac{dy}{dx}$:



$$\frac{dy}{dx}=-\frac{y}{x}$$



Note that, if we substitute $y=\frac{1}{x}$ into

$\frac{dy}{dx}=-\frac{y}{x}$ , we end up with

$\frac{dy}{dx}=-\frac{1}{x^2}$ again.



\paragraph{Application: inverse trigonometric functions}



Arcsine, arccosine, arctangent. These are the functions that allow you

to determine the angle given the sine, cosine, or tangent of that angle.



First, let us start with the arcsine such that:



$$y=\arcsin(x)$$



To find $\frac{\mathrm{d}y}{\mathrm{d}x}$ we first need to break this

down into a form we can work with:



$$x=\sin(y)$$



Then we can take the derivative of that:



$$1=\cos(y)\cdot\frac{\mathrm{d}y}{\mathrm{d}x}$$



...and solve for $\frac{\mathrm{d}y}{\mathrm{d}x}$ : 



$$\frac{\mathrm{d}y}{\mathrm{d}x}=\frac{1}{\cos(y)}$$



% [Image: $y=\arcsin(x)$

gives us this unit

triangle.]



At this point we need to go back to the unit triangle. Since $y$ is the

angle and the opposite side is $\sin(y)=x$ , the adjacent side is

$\cos(y)=\sqrt{1-x^2}$ , and the hypotenuse is 1. Since we have

determined the value of $\cos(y)$ based on the unit triangle, we can

substitute it back in to the above equation and get:



 \textbf{Derivative of the Arcsine}                                    

$$\frac{\mathrm{d}}{\mathrm{d}x}\arcsin(x)=\frac{1}{\sqrt{1-x^2}}$$





We can use an identical procedure for the arccosine and arctangent:



 \textbf{Derivative of the Arccosine}                                   

 $$\frac{\mathrm{d}}{\mathrm{d}x}\ar                                

 ccos(x) = \frac{\mathrm{d}}{\mathrm{d}x}\left(\frac{\pi}{2}\right) -

  \frac{\mathrm{d}}{\mathrm{d}x}\arcsin(x) = -\frac{1}{\sqrt{1-x^2}}$$ 



 \textbf{Derivative of the Arctangent}                                 

 $$\frac{\mathrm{d}}{\mathrm{d}x}\arctan(x)=\frac{1}{1+x^2}$$       





\paragraph{Resources}



\begin{itemize}
    \item \href{https://mathresearch.utsa.edu/wikiFiles/MAT1214/Implicit\%20Differentiation/MAT1214\_3.8ImplicitDifferentiationPwPt.pptx}{Implicit
\end{itemize}
    Differentiation}

    PowerPoint file created by Dr. Sara Shirinkam, UTSA.



<!-- -->



\begin{itemize}
    \item \href{https://mathresearch.utsa.edu/wikiFiles/MAT1214/Implicit\%20Differentiation/MAT1214\_3.8ImplicitDifferentiationWS1.pdf}{Implicit Differentiation Worksheet
\end{itemize}
    1}



<!-- -->



\begin{itemize}
    \item \href{https://mathresearch.utsa.edu/wikiFiles/MAT1214/Implicit\%20Differentiation/MAT1214\_3.8ImplicitDifferentiationWS2.pdf}{Implicit Differentiation Worksheet
\end{itemize}
    2}



\paragraph{Licensing}



Content obtained and/or adapted from:



\begin{itemize}
    \item \href{https://en.wikibooks.org/wiki/Calculus/Implicit\_Differentiation}{Implicit Differentiation, Wikibooks:
\end{itemize}
    Calculus}

    under a CC BY-SA license

\subsection{Learning Objectives}

\begin{enumerate}
    \item \textbf{Perform Implicit Differentiation:} Students will be able to differentiate implicitly given equations where the dependent and independent variables are mixed, such as $x^2 + y^2 = 1$. They should be able to apply the chain rule and product rule appropriately to find $\frac{dy}{dx}$ without explicitly solving for $y$ as a function of $x$.
    \item \textbf{Apply the Chain Rule in Implicit Differentiation:} Students will be able to effectively use the chain rule in implicit differentiation. This involves differentiating terms like $y^2$ where $y$ is implicitly a function of $x$. They should be able to demonstrate understanding by correctly applying the chain rule, as shown in differentiating $y^2$ to obtain $2yy'$ and similar expressions for other terms.
    \item \textbf{Compute Derivatives of Inverse Trigonometric Functions:} Students will be able to find the derivatives of inverse trigonometric functions using implicit differentiation and trigonometric identities.
\end{enumerate}

\subsection{Assessment Instrument}

\begin{enumerate}
    \item \textbf{Part 1: Perform Implicit Differentiation}
    \item \textbf{Calculation:} Differentiate the equation $x^2 + y^2 = 25$ implicitly with respect to $x$. Find $\frac{dy}{dx}$.
    \item \textbf{Multiple Choice:} For the implicit equation $xy + e^y = x^2$, which of the following represents the correct differentiation of $e^y$?
    \item A. $e^y \frac{dy}{dx}$
    \item B. $e^y \cdot 1$
    \item C. $e^y \cdot \frac{d^2y}{dx^2}$
    \item D. $\frac{d}{dx}(e^y) = e^y \frac{dy}{dx}$
    \item \textbf{Short Answer:} Given the implicit function $x^2 y + y^3 = 7$, find $\frac{dy}{dx}$ using implicit differentiation. Describe the steps involved. \textbf{Part 2: Apply the Chain Rule in Implicit Differentiation}
    \item \textbf{Calculation:} Differentiate the implicit equation $x^2 y^2 + y^3 = 4$ with respect to $x$. Use the chain rule where necessary to find $\frac{dy}{dx}$.
    \item \textbf{Multiple Choice:} When differentiating the term $y^2$ implicitly with respect to $x$, which of the following is correct?
    \item A. $2y \frac{dy}{dx}$
    \item B. $2y$
    \item C. $\frac{d}{dx}(y^2) = 2y \cdot \frac{d^2y}{dx^2}$
    \item D. $\frac{d}{dx}(y^2) = 2 \cdot \frac{dy}{dx}$
    \item \textbf{Short Answer:} Differentiate the implicit function $x \sin(y) + \cos(y) = 3x^2$ with respect to $x$. Explain how the chain rule is applied to differentiate $\sin(y)$ and $\cos(y)$. \textbf{Part 3: Compute Derivatives of Inverse Trigonometric Functions}
    \item \textbf{Calculation:} Find the derivative of the function $f(x) = \arcsin(x)$ with respect to $x$ using implicit differentiation.
    \item \textbf{Multiple Choice:} Which of the following correctly represents the derivative of $\arctan(x)$ with respect to $x$?
    \item A. $\frac{1}{1+x^2}$
    \item B. $\frac{1}{x^2 + 1}$
    \item C. $\frac{1}{\sqrt{1 - x^2}}$
    \item D. $\frac{1}{\sqrt{1 + x^2}}$
    \item \textbf{Short Answer:} Use implicit differentiation to find the derivative of $y = \arccos(x)$ with respect to $x$. Describe the steps and apply the trigonometric identity as needed.
\end{enumerate}

%%===============================================================
\section{Related Rates}
\label{sec:related-rates}
%%===============================================================

\paragraph{Introduction}



One useful application of derivatives is as an aid in the calculation of

\textbf{related rates}. What is a related rate? In each case in the following

examples the related rate we are calculating is a derivative with

respect to some value. We compute this derivative from a rate at which

some other known quantity is changing. Given the rate at which something

is changing, we are asked to find the rate at which a value related to

the rate we are given is changing.



\paragraph{How to Solve}



These general steps should be taken in order to complete a related rates

problem.



1.  Write out any relevant formulas and information about the problem.

\begin{itemize}
    \item The problem should have a variable you "control" (i.e. have
\end{itemize}
        knowledge of the value and rate of) and a variable that you want

        to find the related rate.

\begin{itemize}
    \item Usually, related rates problem ask for a rate in respect to
\end{itemize}
        time. *Do not panic if your equations do not appear to have any

        relationship to time!* This will be handled later.

2.  Combine the formulas together so that the variable you want to find

    the related rate of is on one side of the equation and everything

    else is on the other side.

3.  Differentiate the formula with respect to time. Any other variable

    not a simple constant (such as $\pi$) should be differentiated as

    well. Be wary! Chain rule usually should be used.

4.  The other variables that you have differentiated should have been

    given in the question or should be calculated separately.

    Nevertheless, plug-in known information and simplify.

5.  The value you get here is your answer.



The steps to solve a related rates problem is strikingly similar to an

optimization problem, except that the main variable to find is not

assigned to be 0 (it is supposed to be found) and that the *extra

variables* in the optimization problem algorithm are actual variables in

this case and are treated as variables instead of constants when

differentiating.



\paragraph{Notation}



Newton's dot notation is used to show the derivative of a variable with

respect to time. That is, if $f$ is a quantity that depends on time,

then $\dot f=\frac{df}{dt}$ , where $t$ represents the time. This

notation is a useful abbreviation in situations where time derivatives

are often used, as is the case with related rates.



\paragraph{Examples}



A cone with a circular base is being filled with water. Find a formula

which will find the rate with which water is pumped---if the rate that

water was being filled and the measurements of the cone were known.



\textbf{Example 1:}

% [Image: ]



\begin{itemize}
    \item Write out any relevant formulas or pieces of information.
\end{itemize}

$V=\frac{1}{3}\pi r^2h$



\begin{itemize}
    \item Take the derivative of the equation above with respect to time.
\end{itemize}
    Remember to use the chain rule and the product rule.



$V=\frac{1}{3}\pi r^2h$



$\dot V=\frac{\pi}{3}\left(r^2\dot h+2rh\dot r\right)$



Answer: $\dot V=\frac{\pi}{3}\left(r^2\dot h+2rh\dot r\right)$



\textbf{Example 2:}



A spherical hot air balloon is being filled with air. The volume is

changing at a rate of 2 cubic feet per minute.



How is the radius changing with respect to time when the radius is equal

to 2 feet?



\begin{itemize}
    \item Write out any relevant formulas and pieces of information.
\end{itemize}

$V\_{sphere}=\frac{4}{3}\pi r^3$



$\dot V=2$



$r=2$



\begin{itemize}
    \item Take the derivative of both sides of the volume equation with
\end{itemize}
    respect to time.



$V=\frac{4}{3}\pi r^3$



$$\begin{aligned}

\dot V &=\frac{4}{3}3\pi r^2\dot r \\&=4\pi r^2\dot r

\end{aligned}$$



\begin{itemize}
    \item Solve for $\dot r$ .
\end{itemize}

$\dot r=\frac{\dot V}{4\pi r^2}$



\begin{itemize}
    \item Plug-in known information.
\end{itemize}

$\dot r=\frac{2}{16\pi}$



Answer: $\dot r=\frac{1}{8\pi}$ ft/min.



\textbf{Example 3:}



An airplane is attempting to drop a box onto a house. The house is 300

feet away in horizontal distance and 400 feet in vertical distance. The

rate of change of the horizontal distance with respect to time is the

same as the rate of change of the vertical distance with respect to

time. How is the distance between the box and the house changing with

respect to time at the moment? The rate of change in the horizontal

direction with respect to time is -50 feet per second.



*Note: Because the vertical distance is downward in nature, the rate of

change of y is negative. Similarly, the horizontal distance is

decreasing, therefore it is negative (it is getting closer and closer).*



The easiest way to describe the horizontal and vertical relationships of

the plane's motion is the Pythagorean Theorem.



\begin{itemize}
    \item Write out any relevant formulas and pieces of information.
\end{itemize}

$x^2+y^2=s^2$ (where \textit{s} is the distance between the plane and the

house)



$x=300$



$y=400$



$s=\sqrt{x^2+y^2}=\sqrt{300^2+400^2}=500$



$\dot x=\dot y=-50$



\begin{itemize}
    \item Take the derivative of both sides of the distance formula with
\end{itemize}
    respect to time.



$x^2+y^2=s^2$



$2x\dot x+2y\dot y=2s\dot s$



\begin{itemize}
    \item Solve for $\dot s.$
\end{itemize}

$\begin{aligned}

\dot s \&=\frac{2x\dot x+2y\dot y}{2s} \\\\&=\frac{x\dot x+y\dot y}{s}

\end{aligned}$



\begin{itemize}
    \item Plug-in known information
\end{itemize}

$\begin{aligned}

\dot s \&=\frac{(300)(-50)+(400)(-50)}{(500)} \\\\&=\frac{-35000}{500} \\\\&=-70$ ft/s

\end{aligned}$



Answer: $\dot s=-70$ ft/sec.



Example 4:



Sand falls onto a cone shaped pile at a rate of 10 cubic feet per

minute. The radius of the pile's base is always 1/2 of its altitude.

When the pile is 5 ft deep, how fast is the altitude of the pile

increasing?



\begin{itemize}
    \item Write down any relevant formulas and information.
\end{itemize}

$V=\frac{\pi}{3}r^2h$



$\dot V=10$



$r=\frac{h}{2}$



$h=5$



Substitute $r=\frac{h}{2}$ into the volume equation.



$\begin{aligned}

V \&=\frac{\pi}{3}r^2h \\\\&=\frac{\pi}{3}h\left(\frac{h^2}{4}\right) \\\\&=\frac{\pi}{12}h^3

\end{aligned}$



\begin{itemize}
    \item Take the derivative of the volume equation with respect to time.
\end{itemize}

$V=\frac{\pi}{12}h^3$



$\dot V=\frac{\pi}{4}h^2\dot h$



\begin{itemize}
    \item Solve for $\dot h$ .
\end{itemize}

$\dot h=\frac{4}{\pi h^2}\dot V$



\begin{itemize}
    \item Plug-in known information and simplify.
\end{itemize}

$\begin{aligned}

\dot h \&=\frac{4}{\pi(5)^2}10 \\\\&=\frac{8}{5\pi} ft/min

\end{aligned}$



Answer: $\dot h=\frac{8}{5\pi}$ ft/min.



\textbf{Example 5:}



A 10 ft long ladder is leaning against a vertical wall. The foot of the

ladder is being pulled away from the wall at a constant rate of 2

ft/sec. When the ladder is exactly 8 ft from the wall, how fast is the

top of the ladder sliding down the wall?



\begin{itemize}
    \item Write out any relevant formulas and information.
\end{itemize}

Use the Pythagorean Theorem to describe the motion of the ladder.



$x^2+y^2=l^2$ (where \textit{l} is the length of the ladder)



$l=10$



$\dot x=2$



$x=8$



$y=\sqrt{l^2-x^2}=\sqrt{100-64}=\sqrt{36}=6$



\begin{itemize}
    \item Take the derivative of the equation with respect to time.
\end{itemize}

$2x\dot x+2y\dot y=0$ ($l$ is constant so $\frac{dl^2}{dt}=0$ .)



\begin{itemize}
    \item Solve for $\dot y$ .
\end{itemize}

$2x\dot x+2y\dot y=0$



$2y\dot y=-2x\dot x$



$\dot y=-\frac{x}{y}\dot x$



\begin{itemize}
    \item Plug-in known information and simplify.
\end{itemize}

$\begin{aligned}

\dot y \&=\left(-\frac{8}{6}\right)(2) \\\\&=-\frac{8}{3} ft/sec

\end{aligned}$



Answer: $\dot y=-\frac{8}{3}$ ft/sec.





\paragraph{Exercises}



1.  A spherical balloon is inflated at a rate of $100 ft^3/min$ .

    Assuming the rate of inflation remains constant, how fast is the

    radius of the balloon increasing at the instant the radius is

    $4 ft$?





2.  Water is pumped from a cone shaped reservoir (the vertex is pointed

    down) $10 ft$ in diameter and $10 ft$ deep at a constant rate of

    $3 ft^3/min$ . How fast is the water level falling when the depth of

    the water is $6 ft$?





3.  A boat is pulled into a dock via a rope with one end attached to the

    bow of a boat and the other wound around a winch that is $2ft$ in

    diameter. If the winch turns at a constant rate of $2rpm,$ how fast

    is the boat moving toward the dock?





4.  At time $t=0$ a pump begins filling a cylindrical reservoir with

    radius 1 meter at a rate of $e^{-t}$ cubic meters per second. At

    what time is the liquid height increasing at 0.001 meters per

    second?



\paragraph{Solutions}



1.  $\frac{25}{16\pi} \frac{ft}{min}$

2.  $\frac{1}{3\pi} \frac{ft}{min}$

3.  $4\pi\frac{ft}{min}$

4.  $t=-\ln(0.001\pi)$



\paragraph{Resources}



\begin{itemize}
    \item \href{https://mathresearch.utsa.edu/wikiFiles/MAT1214/Related\%20Rates/MAT1214-4.1RelatedRatesPwPt.pptx}{Related
\end{itemize}
    Rates}

    PowerPoint file created by Dr. Sara Shirinkam, UTSA.



<!-- -->



\begin{itemize}
    \item \href{https://mathresearch.utsa.edu/wikiFiles/MAT1214/Related\%20Rates/MAT1214-4.1RelatedRatesWS1.pdf}{Related Rates
\end{itemize}
    Worksheet}



<!-- -->



\begin{itemize}
    \item \href{https://en.wikibooks.org/wiki/Calculus/Related\_Rates}{Related
\end{itemize}
    Rates},

    WikiBooks: Calculus



<!-- -->



\begin{itemize}
    \item \href{https://youtu.be/YRM-x3lZELs}{Cal1 Related Rates Part1} video
\end{itemize}
    lecture by Instructor Beatty, UTSA.



\paragraph{Licensing}



Content obtained and/or adapted from:



\begin{itemize}
    \item \href{https://en.wikibooks.org/wiki/Calculus/Related\_Rates}{Related Rates, WikiBooks:
\end{itemize}
    Calculus}

    under a CC BY-SA license

\subsection{Learning Objectives}

\begin{enumerate}
    \item \textbf{Identify and Set Up Related Rates Equations:} Students will be able to identify the relevant quantities and their rates of change in related rates problems. They should be able to set up the initial equation or formula that relates these quantities, ensuring that it incorporates all necessary variables and constants.
    \item \textbf{Differentiate and Apply Chain Rule:} Students will be able to differentiate the related rates equation with respect to time, applying the chain rule where necessary. They should be able to handle the differentiation of equations involving multiple variables and their rates of change.
    \item \textbf{Solve for the Desired Rate and Interpret Results:} Students will be able to solve for the unknown rate of change by substituting known values into the differentiated equation and simplifying. They should also be able to interpret the result in the context of the problem.
\end{enumerate}

\subsection{Assessment Instrument}

\begin{enumerate}
    \item \textbf{Part 1: Identify and Set Up Related Rates Equations}
    \item \textbf{Multiple Choice:} A ladder 10 feet long is leaning against a wall. If the base of the ladder is sliding away from the wall at a rate of 2 feet per second, which of the following equations correctly relates the rate at which the top of the ladder is sliding down the wall to the rate at which the base is moving away?
    \item A. $x^2 + y^2 = 100$
    \item B. $x^2 + y^2 = 10$
    \item C. $x \frac{dx}{dt} + y \frac{dy}{dt} = 0$
    \item D. $x \frac{dy}{dt} + y \frac{dx}{dt} = 100$
    \item \textbf{True/False:} In a problem involving a balloon rising vertically and a car moving horizontally, the rate of change of the distance between the balloon and the car can be modeled by the Pythagorean theorem if the balloon and car start at the same point. (True/False)
    \item \textbf{Short Answer:} A spherical balloon is being inflated, causing its radius to increase at a rate of 0.1 meters per second. Write the formula that relates the rate of change of the volume of the balloon to the rate of change of the radius. \textbf{Part 2: Differentiate and Apply Chain Rule}
    \item \textbf{Calculation:} Given the equation for the volume of a sphere $V = \frac{4}{3} \pi r^3$, where $r$ is the radius, differentiate this equation with respect to time $t$ to find the rate of change of volume $\frac{dV}{dt}$ in terms of $\frac{dr}{dt}$.
    \item \textbf{Multiple Choice:} For the related rates problem where a tank is being filled with water, and the height of water $h$ in the tank is increasing at a rate of $\frac{dh}{dt} = 3$ cm/min, which of the following correctly represents the rate at which the volume $V$ of the water in the tank is changing if the tank's cross-sectional area $A$ is constant?
    \item A. $\frac{dV}{dt} = A \frac{dh}{dt}$
    \item B. $\frac{dV}{dt} = \frac{dh}{dt} / A$
    \item C. $\frac{dV}{dt} = A \cdot \frac{d^2h}{dt^2}$
    \item D. $\frac{dV}{dt} = \frac{dA}{dt} \cdot \frac{dh}{dt}$
    \item \textbf{Short Answer:} Differentiate the equation $x^2 + y^2 = 25$ with respect to time $t$ and apply the chain rule to find the relationship between $\frac{dx}{dt}$ and $\frac{dy}{dt}$. \textbf{Part 3: Solve for the Desired Rate and Interpret Results}
    \item \textbf{Calculation:} In the related rates problem where a balloon is rising at a rate of 4 meters per second and a car is moving away from the base of the balloon at a rate of 3 meters per second, calculate the rate at which the distance between the car and the balloon is increasing when the car is 5 meters away from the point directly under the balloon, and the balloon is 12 meters high.
    \item \textbf{Multiple Choice:} If the radius of a sphere is increasing at a rate of 0.2 cm/min, and the rate of change of the volume is given by $\frac{dV}{dt} = 4 \pi r^2 \frac{dr}{dt}$, what is the rate of change of the volume when the radius is 10 cm?
    \item A. $80 \pi$ cm³/min
    \item B. $40 \pi$ cm³/min
    \item C. $20 \pi$ cm³/min
    \item D. $10 \pi$ cm³/min
    \item \textbf{Short Answer:} Given a cylindrical tank with a radius of 3 meters and a height that is increasing at a rate of 2 meters per hour, determine how fast the volume of the tank is increasing. Provide the formula used and interpret the result in the context of the tank's volume change.
\end{enumerate}

%%===============================================================
\section{Antiderivatives}
\label{sec:antiderivatives}
%%===============================================================

\paragraph{Definition}



Now recall that $F$ is said to be an antiderivative of \textit{f} if

$F'(x)=f(x)$ . However, $F$ is not the only antiderivative. We can add

any constant to $F$ without changing the derivative. With this, we

define the \textbf{indefinite integral} as follows:



$\int f(x)dx=F(x)+C$ where $F$ satisfies $F'(x)=f(x)$ and $C$ is any

constant.



The function $f(x)$ , the function being integrated, is known as the

\textbf{integrand}. Note that the indefinite integral yields a \textit{family} of

functions.



\textbf{Example}



Since the derivative of $x^4$ is $4x^3,$ the general antiderivative of

$4x^3$ is $x^4$ plus a constant. Thus,



$$\int 4x^3dx=x^4+C$$



\textbf{Example: Finding antiderivatives}



Let's take a look at $6x^2$ . How would we go about finding the

integral of this function? Recall the rule from differentiation that



$$\frac{d}{dx}x^n=nx^{n-1}$$



In our circumstance, we have:



$$\frac{d}{dx}x^3=3x^2$$



This is a start! We now know that the function we seek will have a power of 3 in it. How would we get the constant of 6? Well,



$$2\frac{d}{dx}x^3=2\times 3x^2=6x^2$$



Thus, we say that $2x^3$ is an antiderivative of $6x^2$ .



\paragraph{Indefinite integral identities}



\paragraph{Basic Properties of Indefinite Integrals}



\textbf{Constant Rule for indefinite integrals}



If $c$ is a constant then $\int c\cdot f(x)dx=c\int f(x)dx$



\textbf{Sum/Difference Rule for indefinite integrals}



$$\int\Big(f(x)+g(x)\Big)dx=\int f(x)dx+\int g(x)dx$$



$$\int\Big(f(x)-g(x)\Big)dx=\int f(x)dx-\int g(x)dx$$



\paragraph{Indefinite integrals of Polynomials}



Say we are given a function of the form, $f(x)=x^n$ , and would like to

determine the antiderivative of $f$ . Considering that



$$\frac{d}{dx}\frac{1}{n+1}x^{n+1}=x^n$$



we have the following rule for indefinite integrals:



\textbf{Power rule for indefinite integrals}

$\int x^ndx=\frac{x^{n+1}}{n+1}+C$ for all $n\ne -1$



\paragraph{Integral of the Inverse function}



To integrate $f(x)=\frac{1}{x}$ , we should first remember



$$\frac{d}{dx}\ln(x)=\frac{1}{x}$$



Therefore, since $\frac{1}{x}$ is the derivative of $\ln(x)$ we can

conclude that



$$\int\frac{dx}{x}=\ln|x|+C$$



Note that the polynomial integration rule does not apply when the

exponent is $-1$ . This technique of integration must be used instead.

Since the argument of the natural logarithm function must be positive

(on the real line), the absolute value signs are added around its

argument to ensure that the argument is positive.



\paragraph{Integral of the Exponential function}



Since



$$\frac{d}{dx}e^x=e^x$$



we see that $e^x$ is its own antiderivative. This allows us to find the integral of an exponential function: $\int e^xdx=e^x+C$



\paragraph{Integral of Sine and Cosine}



Recall that



$$\frac{d}{dx}\sin(x)=\cos(x)$$



$$\frac{d}{dx}\cos(x)=-\sin(x)$$



So $\sin(x)$ is an antiderivative of $\cos(x)$ and $-\cos(x)$ is an

antiderivative of $\sin(x)$ . Hence we get the following rules for

integrating $\sin(x)$ and $\cos(x)$



$$\int\cos(x)dx=\sin(x)+C$$



$$\int\sin(x)dx=-\cos(x)+C$$



We will find how to integrate more complicated trigonometric functions

in the chapter on integration techniques.



\textbf{Example}



Suppose we want to integrate the function $f(x)=x^4+1+2\sin(x)$ . An

application of the sum rule from above allows us to use the power rule

and our rule for integrating $\sin(x)$ as follows,



$$\begin{aligned}

\int f(x)dx &=\int\Big(x^4+1+2\sin(x)\Big)dx \\\&=\int x^4dx+\int 1\,dx+\int 2\sin(x)dx \\\&=\frac{x^5}{5}+x-2\cos(x)+C.

\end{aligned}$$



\paragraph{The Substitution Rule}



The substitution rule is a valuable asset in the toolbox of any

integration greasemonkey. It is essentially the chain rule (a

differentiation technique you should be familiar with) in reverse.

First, let's take a look at an example:



\paragraph{Preliminary Example}



Suppose we want to find $\int x\cos(x^2)dx$ . That is, we want to find a

function such that its derivative equals $x\cos(x^2)$ . Stated yet

another way, we want to find an antiderivative of $f(x)=x\cos(x^2)$ .

Since $\sin(x)$ differentiates to $\cos(x)$ , as a first guess we might

try the function $\sin(x^2)$ . But by the Chain Rule,



$$\frac{d}{dx}\sin(x^2)=\cos(x^2)\cdot\frac{d}{dx}x^2=\cos(x^2)\cdot 2x=2x\cos(x^2)$$



Which is almost what we want apart from the fact that there is an extra

factor of 2 in front. But this is easily dealt with because we can

divide by a constant (in this case 2). So,



$$\frac{d}{dx}\frac{\sin(x^2)}{2}=\frac{1}{2}\cdot\frac{d}{dx}\sin(x^2)=\frac{1}{2}\cdot 2\cos(x^2)x=x\cos(x^2)=f(x)$$



Thus, we have discovered a function, $F(x)=\frac{\sin(x^2)}{2},$ whose

derivative is $x\cos(x^2)$ . That is, $F$ is an antiderivative of

$f(x)=x\cos(x^2)$ . This gives us



$$\int x\cos(x^2)dx=\frac{\sin(x^2)}{2}+C$$



\paragraph{Generalization}



In fact, this technique will work for more general integrands. Suppose

$u$ is a differentiable function. Then to evaluate

$\int u'(x)\cos\bigl(u(x)\bigr)dx$ we just have to notice that by the

Chain Rule



$$\frac{d}{dx}\sin\bigl(u(x)\bigr)=\cos\bigl(u(x)\bigr)\frac{du}{dx}=u'(x)\cos\bigl(u(x)\bigr)$$



As long as $u'$ is continuous we have that



$$\int\cos\bigl(u(x)\bigr)u'(x)dx=\sin\bigl(u(x)\bigr)+C$$



Now the right hand side of this equation is just the integral of $\cos(u)$ but with

respect to $u$ . If we write $u$ instead of $u(x)$ this becomes

$\int\cos(u(x))u'(x)dx=\sin(u)+C=\int\cos(u)du$



So, for instance, if $u(x)=x^3$ we have worked out that



$$\int\bigl(\cos(x^3)\cdot 3x^2\bigr)dx=\sin(x^3)+C$$



\paragraph{General Substitution Rule}



Now there was nothing special about using the cosine function in the

discussion above, and it could be replaced by any other function. Doing

this gives us the substitution rule for indefinite integrals:



\textbf{Substitution rule for indefinite integrals}



Assume $u$ is differentiable with continuous derivative and that $f$ is continuous

on the range of $u$. Then $\int f\bigl(u(x)\bigr)\frac{du}{dx}dx=\int f(u)du$



Notice that it looks like you can "cancel" in the expression

$\frac{du}{dx}dx$ to leave just a $du$ . This does not really make any

sense because $\frac{du}{dx}$ is \textbf{not a fraction}. But it's a good

way to remember the substitution rule.



\paragraph{Examples}



The following example shows how powerful a technique substitution can

be. At first glance the following integral seems intractable, but after

a little simplification, it's possible to tackle using substitution.



\textbf{Example}



We will show that



$$\int\frac{dx}{(x^2+a^2)\sqrt{x^2+a^2}}=\frac{x}{a^2\sqrt{x^2+a^2}}+C$$



First, we re-write the integral:



$$\begin{aligned}

\int\frac{dx}{(x^2+a^2)\sqrt{x^2+a^2}}&=\int (x^2+a^2)^{-\frac32}dx \\\&=\int\left(x^2\left(1+\tfrac{a^2}{x^2}\right)\right)^{-\frac32} dx \\\&=\int x^{-3}\left(1+\tfrac{a^2}{x^2}\right)^{-\frac32} dx \\\&=\int \left(1+\tfrac{a^2}{x^2}\right)^{-\frac32} \left(x^{-3} dx\right)

\end{aligned}$$



Now we perform the following substitution:



$$u=1+\frac{a^2}{x^2}$$



$$\frac{du}{dx}=-2a^2x^{-3}\ \implies\ x^{-3}dx=-\frac{du}{2a^2}$$ 



Which yields:



$$\int\left(1+\tfrac{a^2}{x^2}\right)^{-\frac32} \left(x^{-3}dx\right)=$$



$$=\int u^{-\frac32}\left(-\frac{du}{2a^2}\right)$$



$$=-\frac{1}{2a^2}\int u^{-\frac32}du$$



$$=-\frac{1}{2a^2}\left(-\frac{2}{\sqrt u}\right)+C$$



$$=\frac{1}{a^2 \sqrt{1+\frac{a^2}{x^2}}}+C$$



$$=\left(\frac{x}{x}\right) \frac{1}{a^2 \sqrt{1+\frac{a^2}{x^2}}} + C$$



$$=\frac{x}{a^2 \sqrt{x^2+a^2}}+C$$



\paragraph{Integration by Parts}



Integration by parts is another powerful tool for integration. It was

mentioned above that one could consider integration by substitution as

an application of the chain rule in reverse. In a similar manner, one

may consider integration by parts as the product rule in reverse.



\paragraph{Preliminary Example}



\paragraph{General Integration by Parts}



\textbf{Integration by parts for indefinite integrals} ``{=html} Suppose

$f$ and $g$ are differentiable and their derivatives are continuous.

Then



$$\int f(x)g(x)dx=f(x)\int g(x)dx-\int\left(f'(x)\int g(x)dx\right)dx$$



it is also very important to notice that



$$\int f(x)g(x)dx=f(x)\int g(x)dx-\int\left(f'(x)\int g(x)dx\right)dx$$



is not equal to



$$\int f(x)g(x)dx=g(x)\int f(x)dx-\int\left(g'(x)\int f(x)dx\right)dx$$



to set the $f(x)$ and $g(x)$ we need to follow the rule called I.L.A.T.E.



ILATE defines the order in which we must set the $f(x)$

\begin{itemize}
    \item I for inverse trigonometric function
    \item L for log functions
    \item A for algebraic functions
    \item T for trigonometric functions
    \item E for exponential function
\end{itemize}

f(x) and g(x) must be in the order of ILATE or else your final answers

will not match with the main key



\paragraph{Examples}



\textbf{Example}



Find $\int x\cos(x)dx$



Here we let:



$u=x$, so that $du=dx$,



$dv=\cos(x)dx$, so that $v=\sin(x)$.



Then:



$$\begin{aligned}

\int x\cos(x)dx&=\int u\,dv \\\&=uv-\int v\,du \\\&=x\sin(x)-\int\sin(x)dx \\\&=x\sin(x)+\cos(x)+C

\end{aligned}$$



\textbf{Example}



Find $\int x^2e^xdx$



In this example we will have to use integration by parts twice.



Here we let



$u=x^2$, so that $du=2xdx$,



$dv=e^xdx$,  so that $v=e^x$.



Then:



$$\begin{aligned}

\int x^2e^xdx \\\&=\int u\,dv \\\&= uv - \int v \,du \\\&=x^2e^x-\int 2xe^xdx \\\&=x^2e^x-2\int xe^xdx

\end{aligned}$$



Now to calculate the last integral we use integration by parts again.

Let



$u=x$, so that $du=dx$,



$dv=e^xdx$, so that $v=e^x$



and integrating by parts gives



$$\int xe^xdx=xe^x-\int e^xdx=e^x(x-1)$$



So, finally we obtain



$$\int x^2e^xdx=x^2e^x-2e^x(x-1)+C=e^x(x^2-2x+2)+C$$



\textbf{Example}



Find $\int\ln(x)dx$



The trick here is to write this integral as



$\int\ln(x)\cdot 1\,dx$



Now let



$u=\ln(x)$ so $du=\frac{dx}{x}$,



$v=x$ so $dv=1\,dx$.



Then using integration by parts,



$$\begin{aligned}

\int\ln(x)dx &=x\ln(x)-\int\frac{x}{x}dx \\\&=x\ln(x)-\int 1\,dx \\\&=x\ln(x)-x+C \\\&=x\bigl(\ln(x)-1\bigr)+C

\end{aligned}$$



\textbf{Example}



Find $\int\arctan(x)dx$



Again the trick here is to write the integrand as

$\arctan(x)=\arctan(x)\cdot 1$ . Then let



$u=\arctan(x)$ so $du=\frac{dx}{1+x^2}$



$v=x$ so $dv=1\,dx$



so using integration by parts,



$$\begin{aligned}

\int\arctan(x)dx &=x\arctan(x)-\int\frac{x}{1+x^2}dx \\\&=x\arctan(x)-\tfrac12\ln(1+x^2)+C

\end{aligned}$$



\textbf{Example}



Find $\int e^x\cos(x)dx$



This example uses integration by parts twice. First let,



$u=e^x$ so $du=e^xdx$



$dv=\cos(x)dx$ so $v=\sin(x)$



so



$$\int e^x\cos(x)dx=e^x\sin(x)-\int e^x\sin(x)dx$$



Now, to evaluate the remaining integral, we use integration by parts

again, with



$u=e^x$ so $du=e^xdx$



$v=-\cos(x)$ so $dv=\sin(x)dx$



Then



$$\begin{aligned}

\int e^x\sin(x)dx&=-e^x\cos(x)-\int -e^x\cos(x)dx \\\&=-e^x\cos(x)+\int e^x\cos(x)dx

\end{aligned}$$



Putting these together, we have



$$\int e^x\cos(x)dx=e^x\sin(x)+e^x\cos(x)-\int e^x\cos(x)dx$$



Notice that the same integral shows up on both sides of this equation,

but with opposite signs. The integral does not cancel; it doubles when

we add the integral to both sides to get



$$2\int e^x\cos(x)dx=e^x\bigl(\sin(x)+\cos(x)\bigr)$$



$$\int e^x\cos(x)dx=\frac{e^x\bigl(\sin(x)+\cos(x)\bigr)}{2}$$



\paragraph{Resources}



\begin{itemize}
    \item \href{https://mathresearch.utsa.edu/wikiFiles/MAT1193/The\%20Definite\%20Integral/Presentation11\_Distance\_Definite\_Integral.pptx}{Distance Definite
\end{itemize}
    Integral}

    (Slides 23-38). PowerPoint file created by Professor Cynthia

    Roberts, UTSA.

\begin{itemize}
    \item \href{https://mathresearch.utsa.edu/wikiFiles/MAT1193/The\%20Definite\%20Integral/Presentation12\_DefiniteIntegral\%20\&\%20Antiderivatives.pptx}{Definite Integral \&
\end{itemize}
    Antiderivatives}.

    PowerPoint file created by Professor Cynthia Roberts, UTSA.

\begin{itemize}
    \item \href{https://mathresearch.utsa.edu/wikiFiles/MAT1214/Antiderivatives/MAT1214-4.10AntiderivativesPwPt.pptx}{Antiderivatives}
\end{itemize}
    PowerPoint file created by Dr. Sara Shirinkam, UTSA.



<!-- -->



\begin{itemize}
    \item \href{https://mathresearch.utsa.edu/wikiFiles/MAT1214/Antiderivatives/MAT1214-4.10AntiderivativesWS1.pdf}{Antiderivatives
\end{itemize}
    Worksheet}



\paragraph{Licensing}



Content obtained and/or adapted from:



\begin{itemize}
    \item \href{https://en.wikibooks.org/wiki/Calculus/Indefinite\_integral}{Indefinite integral, Wikibooks:
\end{itemize}
    Calculus}

    under a CC BY-SA license

\subsection{Learning Objectives}

\begin{enumerate}
    \item \textbf{Compute Antiderivatives:}
    \item Students will be able to find the general antiderivative of a given polynomial function using the power rule for indefinite integrals, including determining the appropriate constant of integration.
    \item \textbf{Apply Integration Techniques:}
    \item Students will apply fundamental integration techniques, including the constant rule, sum/difference rule, and integration of exponential, trigonometric, and inverse functions, to compute indefinite integrals of given functions.
    \item \textbf{Use Substitution and Integration by Parts:}
    \item Students will use the substitution rule and integration by parts to evaluate more complex integrals, demonstrating the ability to apply these techniques to integrals involving products of functions or compositions of functions.
\end{enumerate}

\subsection{Assessment Instrument}

\begin{enumerate}
    \item \textbf{Part 1: Compute Antiderivatives}
    \item \textbf{Calculation:} Find the general antiderivative of the polynomial function $f(x) = 3x^4 - 5x^2 + 2$. Be sure to include the constant of integration.
    \item \textbf{Multiple Choice:} What is the general antiderivative of the function $f(x) = 7x^3 - 4x + 6$?
    \item A. $\frac{7}{4}x^4 - 2x^2 + 6x + C$
    \item B. $7x^4 - 2x^2 + 6x + C$
    \item C. $7x^4 - 4x^2 + 6x + C$
    \item D. $\frac{7}{4}x^4 - 2x^2 + \frac{6}{2}x^2 + C$
    \item \textbf{Short Answer:} Determine the general antiderivative of $f(x) = -x^3 + 4x - 1$ and provide the constant of integration in your solution. \textbf{Part 2: Apply Integration Techniques}
    \item \textbf{Calculation:} Compute the indefinite integral of $f(x) = e^x \sin(x) \, dx$. Use the appropriate integration technique to solve the integral.
    \item \textbf{Multiple Choice:} Which of the following is the correct antiderivative of the function $f(x) = \cos(x) + \frac{1}{x}$?
    \item A. $\sin(x) + \ln|x| + C$
    \item B. $\sin(x) - \ln|x| + C$
    \item C. $\cos(x) + \ln|x| + C$
    \item D. $\sin(x) + \ln(x) + C$
    \item \textbf{Short Answer:} Find the antiderivative of the function $f(x) = 5x^2 - 3x + \frac{1}{x}$. Explain which integration rules you used. \textbf{Part 3: Use Substitution and Integration by Parts}
    \item \textbf{Calculation:} Evaluate the integral $\int x \cdot e^{x^2} \, dx$ using the substitution method. Show the substitution step and the final result.
    \item \textbf{Multiple Choice:} For the integral $\int x \cos(x) \, dx$, which integration technique should be used?
    \item A. Substitution
    \item B. Integration by parts
    \item C. Partial fractions
    \item D. Trigonometric substitution
    \item \textbf{Short Answer:} Use integration by parts to find the integral of $f(x) = x \ln(x) \, dx$. Show all steps and provide the final result.
\end{enumerate}

%%===============================================================
\section{Integration by Substitution}
\label{sec:integration-by-substitution}
%%===============================================================

\paragraph{Integration by Substitution}



There is a theorem that will help you with substitution for integration.

It is called \textbf{Change of Variables for Definite Integrals}.



what the theorem looks like is this



$$\int\_{a}^{b} f(x)\operatorname {d}x = \int\_{\alpha}^{\beta} f(g(u))g\prime (u)\operatorname {d}u$$



In order to get $\alpha$ you must plug \textbf{a}

into the function \textbf{g} and to get $\beta$ you must plug

\textbf{b} into the function \textbf{g}.



The tricky part is trying to identify what you want to make your

\textbf{u} to be. Some times substitution will not

be enough and you will have to use the rules for integration by parts.

That will be covered in a different section



\paragraph{Steps}



$\begin(aligned}

\int\limits\_{x=a}^{x=b}f(x)dx\&=\int\limits\_{x=a}^{x=b} f(x)\ \frac{du}{du}\ dx \hspace{1cm} \text{(1) i.e.} \hspace{1cm} \frac{du}{du}=1$ \\\\&=\int\limits\_{x=a}^{x=b}{\left(f(x)\ \frac{dx}{du}\right)\left(\frac{du}{dx}\right)}\ dx \hspace{1cm} \text{(2) i.e.} \hspace{1cm} \frac{dx}{du}\cdot\frac{du}{dx}=1 \\\\&=\int\limits\_{x=a}^{x=b}\left(f(x)\ \frac{dx}{du}\right)g'(x)\ dx \hspace{1cm} \text{(3) i.e.} \hspace{1cm} \frac{du}{dx}=g'(x) \\\\&=\int\limits\_{x=a}^{x=b}h(g(x))g'(x)dx \hspace{1cm} \text{(4) i.e. Now equate} \hspace{1cm} \left(f(x)\ \frac{dx}{du}\right) \hspace{1cm} \text{with} \hspace{1cm} h(g(x)) \\\\&=\int\limits\_{x=a}^{x=b}h(u)g'(x)dx \hspace{1cm} \text{(5) i.e.} \hspace{1cm} g(x)=u \\\\&=\int\limits\_{u=g(a)}^{u=g(b)}h(u)du \hspace{1cm} \text{(6) i.e.} du=\frac{du}{dx}dx=g'(x)dx \\\\&=\int\limits\_{u=c}^{u=d}h(u)du \hspace{1cm} \text{(7) i.e. We have achieved our desired result}

\end{aligned}$



\paragraph{Example 1}



$\int\_{0}^{2} x(x^2+1)^2 \operatorname {d}x$



Instead of making this a big polynomial we will just use the

substitution method.



\textbf{Step 1}



Identify your \textit{u}



Let $u = x^2+1$



\textbf{Step 2}



Identify $\operatorname {d}u$



$\operatorname {d}u = 2x\operatorname {d}x$



\textbf{Step 3}



Now we plug in our limits of integration to our \textit{u} to find our new limits of integration



When $x = 0, u =0^2 + 1 = 1$



and when $x = 2, u = 2^2 + 1 = 5$



Now our integration problem looks something like this



$$\frac {1}{2} \int\_{0}^{5} (x^2 + 1)^2 (2x)\operatorname {d}x$$





\textbf{Step 4}



write your new integration problem



When we plug in our \textit{u} it looks like



$$\frac {1}{2} \int\_{0}^{5} (u)^2 \operatorname {d}u$$



\textbf{Step 5}



Evaluate the Integral



$$\frac {1}{2} \left[\frac {1}{3} u^3 \right]\_{0}^{5}$$



$$\frac {1}{2} \left[\left(\frac {1}{3} * 5^3 \right) - \left(\frac {1}{3} * 0^3 \right)\right]$$



$$\frac {1}{2} \left[\frac {1}{3} * 125 \right]$$



$$\frac {1}{2} \left[\frac {125}{3}\right]$$



$$\frac {125}{6}$$



As you can see this all simplified fairly nice. Using substitution will be hard, for most people, at first. Once you get the hang of

doing this it should come to you faster and faster each time.



\paragraph{Example 2}



$$\int 3x^2(x^3+1)^5dx$$



we see that $3x^2$ is the derivative of $x^3+1$ . Letting



$$u=x^3+1$$



we have



$$\frac{du}{dx}=3x^2$$



or, in order to apply it to the integral,



$$du=3x^2dx$$



With this we may write



$$\int 3x^2(x^3+1)^5dx=\int u^5du=\frac{u^6}{6}+C=\frac{(x^3+1)^6}{6}+C$$



Note that it was not necessary that we had \textit{exactly}the derivative of $u$ in our integrand.

It would have been sufficient to have any constant multiple of the derivative.



For instance, to treat the integral



$$\int x^4\sin(x^5)dx$$



we may let $u=x^5$ . Then



$$du=5x^4dx$$ and so



$$\frac{du}{5}=x^4dx$$



the right-hand side of which is a factor of our integrand. Thus,



$$\int x^4\sin(x^5)dx=\int\frac{\sin(u)}{5}du=-\frac{\cos(u)}{5}+C=-\frac{\cos(x^5)}{5}+C$$



\paragraph{Resources}



\begin{itemize}
    \item \href{https://en.wikibooks.org/wiki/Calculus/Integration\_techniques/Recognizing\_Derivatives\_and\_the\_Substitution\_Rule}{Recognizing Derivatives and Substitution
\end{itemize}
    Rules},

    WikiBooks: Calculus

\begin{itemize}
    \item \href{https://en.wikibooks.org/wiki/High\_School\_Calculus/Integration\_by\_Substitution}{Integration by
\end{itemize}
    Substitution},

    WikiBooks: High School Calculus

\begin{itemize}
    \item \href{https://www.youtube.com/watch?v=uoCW8S-I9Es}{Example 1}. Produced
\end{itemize}
    by Professor Zachary Sharon, UTSA



<!-- -->



\begin{itemize}
    \item \href{https://www.youtube.com/watch?v=zqMxMtjbaBE}{Example 2}. Produced
\end{itemize}
    by TA Catherine Sporer, UTSA



``{=html}Indefinite Integrals Using

Substitution``{=html}



\begin{itemize}
    \item \href{https://www.youtube.com/watch?v=568hUU4beuk}{Indefinite Integration Using
\end{itemize}
    Substitution} by James

    Sousa, Math is Power 4U

\begin{itemize}
    \item \href{https://youtu.be/VfvJ5C9FjcA}{Integration by Substitution Part 1}
\end{itemize}
    by James Sousa, Math is Power 4U

\begin{itemize}
    \item \href{https://youtu.be/pmekyFVIorE}{Integration by Substitution Part 2}
\end{itemize}
    by James Sousa, Math is Power 4U

\begin{itemize}
    \item \href{https://www.youtube.com/watch?v=3MbdmYeYB\_I}{Ex 1: Indefinite Integration Using
\end{itemize}
    Substitution} by James

    Sousa, Math is Power 4U

\begin{itemize}
    \item \href{https://www.youtube.com/watch?v=xLtPL1Td5f4}{Ex 2: Indefinite Integration Using
\end{itemize}
    Substitution} by James

    Sousa, Math is Power 4U

\begin{itemize}
    \item \href{https://www.youtube.com/watch?v=Cdk-3H35TrY}{Ex 3: Indefinite Integration Using
\end{itemize}
    Substitution} by James

    Sousa, Math is Power 4U

\begin{itemize}
    \item \href{https://www.youtube.com/watch?v=hauyUjzRTos}{Ex 4: Indefinite Integration Using
\end{itemize}
    Substitution} by James

    Sousa, Math is Power 4U

\begin{itemize}
    \item \href{https://www.youtube.com/watch?v=902ta0Kshoc}{Ex 5: Indefinite Integration Using
\end{itemize}
    Substitution} by James

    Sousa, Math is Power 4U

\begin{itemize}
    \item \href{https://www.youtube.com/watch?v=FJHHQdQEecE}{Ex 6: Indefinite Integration Using
\end{itemize}
    Substitution} by James

    Sousa, Math is Power 4U

\begin{itemize}
    \item \href{https://www.youtube.com/watch?v=VY5r1-9apdc}{Ex 7: Indefinite Integration Using
\end{itemize}
    Substitution} by James

    Sousa, Math is Power 4U

\begin{itemize}
    \item \href{https://www.youtube.com/watch?v=UnND2S4IQSA}{Ex 8: Indefinite Integration Using
\end{itemize}
    Substitution} by James

    Sousa, Math is Power 4U

\begin{itemize}
    \item \href{https://www.youtube.com/watch?v=CVfe5vNy4y4}{Ex 9: Indefinite Integration Using
\end{itemize}
    Substitution} by James

    Sousa, Math is Power 4U



<!-- -->



\begin{itemize}
    \item \href{https://youtu.be/qclrs-1rpKI}{Integration using U-Substitution} by
\end{itemize}
    patrickJMT

\begin{itemize}
    \item \href{https://youtu.be/x06V9xuLdqg}{U-Substitution - More Complicated
\end{itemize}
    Examples} by patrickJMT



<!-- -->



\begin{itemize}
    \item \href{https://youtu.be/HHSepCzP56M}{U-Substitution Example 1} by Krista
\end{itemize}
    King

\begin{itemize}
    \item \href{https://youtu.be/iRIhO-w46L0}{U-Substitution Example 2} by Krista
\end{itemize}
    King

\begin{itemize}
    \item \href{https://youtu.be/w1FHS-XV64Q}{U-Substitution Example 3} by Krista
\end{itemize}
    King

\begin{itemize}
    \item \href{https://youtu.be/\_r6xMeO7aEw}{U-Substitution Example 4} by Krista
\end{itemize}
    King

\begin{itemize}
    \item \href{https://youtu.be/gIj3objG09Q}{U-Substitution Example 5} by Krista
\end{itemize}
    King

\begin{itemize}
    \item \href{https://youtu.be/TlUW6ZenNXY}{U-Substitution Example 6} by Krista
\end{itemize}
    King

\begin{itemize}
    \item \href{https://youtu.be/W1WVFJ8iMqQ}{U-Substitution Example 7} by Krista
\end{itemize}
    King



<!-- -->



\begin{itemize}
    \item \href{https://youtu.be/IAh00vU3FSY}{How To Integrate Using
\end{itemize}
    U-Substitution} by The Organic

    Chemistry Tutor



``{=html}Definite Integrals Using Substitution``{=html}



\begin{itemize}
    \item \href{https://www.youtube.com/watch?v=hjr27pv8pHQ}{Definite Integration Using
\end{itemize}
    Subsitution} by James

    Sousa, Math is Power 4U

\begin{itemize}
    \item \href{https://www.youtube.com/watch?v=jUwNl\_PDeDY}{Ex 1: Definite Integration Using Substitution - Change Limits of
\end{itemize}
    Integration?} by James

    Sousa, Math is Power 4U

\begin{itemize}
    \item \href{https://www.youtube.com/watch?v=n3DGt7cnS70}{Ex 2: Definite Integration Using Substitution - Change Limits of
\end{itemize}
    Integration?} by James

    Sousa, Math is Power 4U

\begin{itemize}
    \item \href{https://www.youtube.com/watch?v=8ofGv7TQL34}{Ex 1: Definite Integration Using
\end{itemize}
    Substitution} by James

    Sousa, Math is Power 4U

\begin{itemize}
    \item \href{https://www.youtube.com/watch?v=M-IUhe-iQx0}{Ex 2: Definite Integration Using
\end{itemize}
    Substitution} by James

    Sousa, Math is Power 4U



<!-- -->



\begin{itemize}
    \item \href{https://youtu.be/3XOzg2-DdTA}{Integration by U-Substitution, Definite
\end{itemize}
    Integral} by patrickJMT

\begin{itemize}
    \item \href{https://youtu.be/FJoyIAIC1Ag}{U-Substitution: When Do I Have to Change the Limits of Integration
\end{itemize}
    ?} by patrickJMT



<!-- -->



\begin{itemize}
    \item \href{https://youtu.be/0A2RlnutO8U}{U-Substitution Integration, Indefinite \& Definite
\end{itemize}
    Integral} by The Organic Chemistry

    Tutor



\paragraph{Licensing}



Content obtained and/or adapted from:



\begin{itemize}
    \item \href{https://en.wikibooks.org/wiki/Calculus/Integration\_techniques/Recognizing\_Derivatives\_and\_the\_Substitution\_Rule}{Recognizing Derivatives and Substitution Rules, WikiBooks:
\end{itemize}
    Calculus}

    under a CC BY-SA license



<!-- -->



\begin{itemize}
    \item \href{https://en.wikibooks.org/wiki/High\_School\_Calculus/Integration\_by\_Substitution}{Integration by Substitution, WikiBooks: High School
\end{itemize}
    Calculus}

    under a CC BY-SA license

\subsection{Learning Objectives}

\begin{enumerate}
    \item \textbf{Apply the Change of Variables Theorem}: Students will be able to use the Change of Variables theorem to transform a definite integral into a new variable $u$ and accurately adjust the limits of integration accordingly.
    \item \textbf{Perform Substitution for Integration}: Students will be able to identify appropriate substitutions for integrals, compute the derivative $\frac{du}{dx}$, and convert the integral into terms of $u$ to simplify and solve it.
    \item \textbf{Evaluate Definite Integrals Using Substitution}: Students will be able to evaluate definite integrals by substituting variables, simplifying the integrand, and adjusting the limits of integration, then compute the integral using standard integration techniques.
\end{enumerate}

\subsection{Assessment Instrument}

\begin{enumerate}
    \item \textbf{Part 1: Apply the Change of Variables Theorem}
    \item \textbf{Calculation:} Apply the Change of Variables theorem to transform the integral $\int_{0}^{1} (3x^2 + 2) \, dx$ using the substitution $u = x^3$. Find the new limits of integration and the transformed integral in terms of $u$.
    \item \textbf{Multiple Choice:} For the integral $$\int_{-2}^{2} (4 - x^2) \, dx$$, if you use the substitution $u = x^2$, which of the following is the correct change of limits for $u$?
    \item A. $u$ changes from $4$ to $0$
    \item B. $u$ changes from $0$ to $4$
    \item C. $u$ changes from $-4$ to $4$
    \item D. $u$ changes from $-2$ to $2$
    \item \textbf{Short Answer:} Use the Change of Variables theorem to find the transformed integral of $\int_{1}^{e} \frac{1}{x \ln(x)} \, dx$ with the substitution $u = \ln(x)$. Provide the new limits of integration and the transformed integral. \textbf{Part 2: Perform Substitution for Integration}
    \item \textbf{Calculation:} Evaluate the integral $\int \frac{2x}{x^2 + 1} \, dx$ by using the substitution $u = x^2 + 1$. Compute $\frac{du}{dx}$, substitute into the integral, and simplify.
    \item \textbf{Multiple Choice:} If you have the integral $\int \sin(3x) \, dx$, which substitution would simplify the integral most effectively?
    \item A. $u = \sin(3x)$
    \item B. $u = 3x$
    \item C. $u = \cos(3x)$
    \item D. $u = 3 \sin(x)$
    \item \textbf{Short Answer:} Use the substitution $u = 2x - 1$ to evaluate the integral $\int (2x - 1) e^{2x - 1} \, dx$. Compute $\frac{du}{dx}$, transform the integral, and solve. \textbf{Part 3: Evaluate Definite Integrals Using Substitution}
    \item \textbf{Calculation:} Evaluate the definite integral $\int_{0}^{\pi/2} \cos^2(x) \sin(x) \, dx$ using the substitution $u = \sin(x)$. Adjust the limits of integration accordingly and compute the integral.
    \item \textbf{Multiple Choice:} For the integral $\int_{1}^{4} \frac{1}{\sqrt{x - 1}} \, dx$, which substitution is appropriate to simplify the integral?
    \item A. $u = \sqrt{x - 1}$
    \item B. $u = x - 1$
    \item C. $u = \frac{1}{\sqrt{x - 1}}$
    \item D. $u = \frac{1}{x - 1}$
    \item \textbf{Short Answer:} Evaluate the definite integral $\int_{0}^{1} x e^{x^2} \, dx$ using the substitution $u = x^2$. Provide the adjusted limits of integration and the final result.
\end{enumerate}

%%===============================================================
\section{The Definite Integral}
\label{sec:the-definite-integral}
%%===============================================================

Suppose we are given a function and would like to determine the area

underneath its graph over an interval. We could guess, but how could we

figure out the exact area? Below, using a few clever ideas, we actually

\textit{define} such an area and show that by using what is called the

\textbf{definite integral} we can indeed determine the exact area underneath

a curve.



\paragraph{Introduction}



The rough idea of defining the area under the graph of $f$ is to

approximate this area with a finite number of rectangles. Since we can

easily work out the area of the rectangles, we get an estimate of the

area under the graph. If we use a larger number of smaller-sized

rectangles we expect greater accuracy with respect to the area under the

curve and hence a better approximation. Somehow, it seems that we could

use our old friend from differentiation, the limit, and "approach" an

infinite number of rectangles to get the exact area. Let's look at such

an idea more closely.



Suppose we have a function $f$ that is positive on the interval $[a,b]$

and we want to find the area $S$ under $f$ between $a$ and $b$ . Let's

pick an integer $n$ and divide the interval into $n$ subintervals of

equal width. As the interval $[a,b]$ has width $b-a,$ each subinterval has width $\Delta x=\frac{b-a}{n}$ .

We denote the endpoints of the subintervals by $x\_0,x\_1,\ldots,x\_n$ .

This gives us



$$x\_i=a+i\Delta x\text{ for }i=0,1,\dots,n$$



% [Image: Figure 1: Approximation of the area under the curve $f(x)$ from

$x=x\_0$ to

$x=x\_4.$]



Now for each $i=1,\ldots,n$ pick a \textit{sample point} $x\_i^*$ in the

interval $[x\_{i-1},x\_i]$ and consider the rectangle of height $f(x\_i^*)$

and width $\Delta x$. The area of this rectangle is $f(x\_i^*)\Delta x$ . By adding up the area of all the

rectangles for $i=1,\ldots,n$ we get that the area $S$ is approximated by



$$A\_n=f(x\_1^*)\Delta x+\dots+f(x\_n^*)\Delta x$$



% [Image: Figure 2: Rectangle approximating the area under the curve from $x\_2$

to $x\_3$ with sample point $x\_3^*$

.]



A more convenient way to write this is with summation notation.



$$A\_n=\sum\_{i=1}^n f(x\_i^*)\Delta x$$



For each number $n$ we get a different approximation. As $n$ gets larger the width of the rectangles

gets smaller which yields a better approximation. In the limit of $A\_n$ as $n$ tends

to infinity we get the area $S$ .



% [Image: Figure 3: Riemann sums with an increasing number of subdivisions

yielding better

approximations.]



\paragraph{Definition of the Definite Integral}



Suppose $f$ is a continuous function on $[a,b]$ and

$\Delta x=\frac{b-a}{n}$ . Then the \textit{definite integral} of $f$ between

$a$ and $b$ is



$$\int\limits\_a^b f(x)dx=\lim\_{n\to\infty}A\_n=\lim\_{n\to\infty}\sum\_{i=1}^n f(x\_i^*)\Delta x$$



where $x\_i^*$ are any sample points in the interval $[x\_{i-1},x\_i]$ and

$x\_k=a+k\cdot\Delta x$ for $k=0,\dots,n$ .



It is a fact that if $f$ is continuous on $[a,b]$ then this limit always

exists and does not depend on the choice of the points

$x\_i^*\in[x\_{i-1},x\_i]$ . For instance they may be evenly spaced, or

distributed ambiguously throughout the interval. The proof of this is

technical and is beyond the scope of this section.



\paragraph{Notation}



When considering the expression, $\int\limits\_a^b f(x)dx$ (read "the

integral from $a$ to $b$ of the $f$ of $x$ $dx$"), the function $f$ is

called the \textit{integrand} and the interval $[a,b]$ is the interval of

integration. Also $a$ is called the \textit{lower limit} and $b$ the *upper

limit* of integration.



One important feature of this definition is that we also allow functions

which take negative values. If $f(x)0$ and $a



\begin{itemize}
    \item \href{https://mathresearch.utsa.edu/wikiFiles/MAT1214/The\%20Definite\%20Integral/MAT1214-5.3TheDefiniteIntegralPwPt.pptx}{The Definite
\end{itemize}
    Integral}

    PowerPoint file created by Dr. Sara Shirinkam, UTSA.



<!-- -->



\begin{itemize}
    \item \href{https://mathresearch.utsa.edu/wikiFiles/MAT1214/The\%20Definite\%20Integral/MAT1214-5.3TheDefiniteIntegralNotes.pdf}{The Definite Integral
\end{itemize}
    Notes}

    created by Professor Eduardo Duenez, UTSA.



\paragraph{Licensing}



Content obtained and/or adapted from:



\begin{itemize}
    \item \href{https://en.wikibooks.org/wiki/Calculus/Definite\_integral}{Definite integral, Wikibooks:
\end{itemize}
    Calculus}

    under a CC BY-SA license

\subsection{Learning Objectives}

\begin{enumerate}
    \item \textbf{Understand the Definition of a Definite Integral:} Students will be able to explain the concept of the definite integral as the limit of Riemann sums and use this definition to determine the exact area under a curve for a given function over a specified interval.
    \item \textbf{Apply the Properties of Definite Integrals:} Students will be able to apply fundamental properties of definite integrals, such as the Constant Rule, Addition and Subtraction Rules, and Comparison Rule, to evaluate and manipulate integrals in various contexts.
    \item \textbf{Use Riemann Sums for Approximation:} Students will be able to compute approximations of definite integrals using left-handed and right-handed Riemann sums and understand how these approximations improve as the number of subdivisions increases.
\end{enumerate}

\subsection{Assessment Instrument}

\begin{enumerate}
    \item \textbf{Part 1: Understand the Definition of a Definite Integral}
    \item \textbf{Multiple Choice:} Which of the following best describes the definite integral of a function $f(x)$ from $a$ to $b$ as the limit of Riemann sums?
    \item A. The area under the curve of $f(x)$ from $a$ to $b$ divided by the width of the intervals.
    \item B. The limit of the sum of the areas of rectangles with width $\Delta x$ and height $f(x_i)$ as $\Delta x \to 0$.
    \item C. The sum of the values of the function at discrete points between $a$ and $b$.
    \item D. The difference between the function values at $a$ and $b$.
    \item \textbf{True/False:} The definite integral of a function $f(x)$ from $a$ to $b$ can be approximated using a finite number of rectangles, where the height of each rectangle is determined by the value of the function at the right endpoint of the interval. (True/False)
    \item \textbf{Short Answer:} Explain how the definite integral of $f(x) = x^2$ from $1$ to $3$ can be understood as the limit of Riemann sums. Describe the process and the concept of Riemann sums in this context. \textbf{Part 2: Apply the Properties of Definite Integrals}
    \item \textbf{Calculation:} Evaluate the integral $$\int_{0}^{2} (3x^2 - 4x + 1) \, dx$$ using the properties of definite integrals. Apply the Constant Rule and the Addition Rule to simplify the process.
    \item \textbf{Multiple Choice:} Given the integrals $$\int_{a}^{b} f(x) \, dx$$ and $$\int_{a}^{b} g(x) \, dx$$, which of the following represents the result of $$\int_{a}^{b} [f(x) + g(x)] \, dx$$?
    \item A. $$\int_{a}^{b} f(x) \, dx + \int_{a}^{b} g(x) \, dx$$
    \item B. $$\int_{a}^{b} f(x) \cdot g(x) \, dx$$
    \item C. $$\int_{a}^{b} f(x) \, dx - \int_{a}^{b} g(x) \, dx$$
    \item D. $$\int_{a}^{b} f(x) \, dx \cdot \int_{a}^{b} g(x) \, dx$$
    \item \textbf{Short Answer:} Use the Comparison Rule to determine whether the integral $$\int_{0}^{4} (2x - 1) \, dx$$ is greater than or less than $$\int_{0}^{4} x \, dx$$. Explain your reasoning using the properties of definite integrals. \textbf{Part 3: Use Riemann Sums for Approximation}
    \item \textbf{Calculation:} Approximate the integral $$\int_{1}^{3} (2x + 1) \, dx$$ using a left-handed Riemann sum with 4 equal subdivisions. Show your calculations and the result.
    \item \textbf{Multiple Choice:} Which of the following best describes the behavior of Riemann sums as the number of subdivisions increases?
    \item A. The approximation of the integral becomes less accurate.
    \item B. The approximation of the integral becomes more accurate.
    \item C. The approximation remains constant regardless of the number of subdivisions.
    \item D. The approximation becomes exactly equal to the integral when using only one subdivision.
    \item \textbf{Short Answer:} Calculate the right-handed Riemann sum for the integral $$\int_{0}^{2} x^2 \, dx$$ using 3 subdivisions. Provide the width of each interval, the function values at the right endpoints, and the final sum.
\end{enumerate}

%%===============================================================
\section{The Fundamental Theorem of Calculus}
\label{sec:the-fundamental-theorem-of-calculus}
%%===============================================================

The fundamental theorem of calculus is a critical portion of calculus

because it links the concept of a derivative to that of an integral. As

a result, we can use our knowledge of derivatives to find the area under

the curve, which is often quicker and simpler than using the definition

of the integral.



\paragraph{Mean Value Theorem for Integration}



We will need the following theorem in the discussion of the Fundamental

Theorem of Calculus.



Suppose $f(x)$ is continuous on $[a,b]$ . Then

$\frac{1}{b-a}\int\limits\_a^b f(x)dx=f(c)$ for some $c\in[a,b].$



\paragraph{Proof of the Mean Value Theorem for Integration}



$f(x)$ satisfies the requirements of the Extreme Value Theorem, so it

has a minimum $m$ and a maximum $M$ in $[a,b].$ Since



$$\int\limits\_a^b f(x)dx=\lim\_{n\to\infty}\sum\_{k=1}^n f(x\_k^*)\cdot\frac{b-a}{n}=\lim\_{n\to\infty}\frac{b-a}{n}\cdot\sum\_{k=1}^n f(x\_k^*)$$



and since



$$m\le f(x\_k^*)\le M$$



for all $x\_k^*\in[a,b]$



we have



$$\begin{aligned}

&\lim\_{n\to\infty}\frac{b-a}{n}\cdot\sum\_{k=1}^nm\le\lim\_{n\to\infty}\frac{b-a}{n}\cdot\sum\_{k=1}^n f(x\_k^*)\le\lim\_{n\to\infty}\frac{b-a}{n}\cdot\sum\_{k=1}^n M \\\&\lim\_{n\to\infty}mn\cdot\frac{b-a}{n}\le\int\limits\_a^b f(x)dx\le\lim\_{n\to\infty}Mn\cdot\frac{b-a}{n} \\\&\lim\_{n\to\infty}m(b-a)\le\int\limits\_a^b f(x)dx\le\lim\_{n\to\infty}M(b-a) \\\&m(b-a)\le\int\limits\_a^b f(x)dx\le M(b-a) \\\&m\le\frac{1}{b-a}\int\limits\_a^b f(x)dx\le M

\end{aligned}$$



Since $f$ is continuous, by the Intermediate Value Theorem there is some

$f(c)$ with $c\in[a,b]$ such that



$$\frac{1}{b-a}\int\limits\_a^b f(x)dx=f(c)$$



\paragraph{Fundamental Theorem of Calculus}



\paragraph{Statement of the Fundamental Theorem}



Suppose that $f$ is continuous on $[a,b]$ . We can define a function $F$

by



$$F(x)=\int\limits\_a^x f(t)dt\quad\text{for }x\in[a,b]$$



Suppose $f$ is continuous on $[a,b]$ and $F$ is defined by



$$F(x)=\int\limits\_a^x f(t)dt$$



Then $F$ is differentiable on $(a,b)$ and for all $x\in(a,b),$



$$F'(x)=f(x)$$



When we have such functions $F$ and $f$ where $F'(x)=f(x)$ for every $x$

in some interval $I$ we say that $F$ is the \textbf{antiderivative} of $f$ on

$I.$



Suppose that $f$ is continuous on $[a,b]$ and that $F$ is any

antiderivative of $f$ . Then



$$\int\limits\_a^b f(x)dx=F(b)-F(a)$$



Note: a minority of mathematicians refer to part one as two and part two

as one. All mathematicians refer to what is stated here as part 2 as The

Fundamental Theorem of Calculus.



\paragraph{Proofs}



\paragraph{Proof of Fundamental Theorem of Calculus Part I}



Suppose $x\in(a,b)$ . Pick $\Delta x$ so that $x+\Delta x\in(a,b)$. Then



$$F(x)=\int\limits\_a^x f(t)dt$$



and



$$F(x+\Delta x)=\int\limits\_a^{x+\Delta x}f(t)dt$$



% [Image: Figure 1]



Subtracting the two equations gives



$$F(x+\Delta x)-F(x)=\int\limits\_a^{x+\Delta x}f(t)dt-\int\limits\_a^x f(t)dt$$



Now



$$\int\limits\_a^{x+\Delta x}f(t)dt=\int\limits\_a^x f(t)dt+\int\limits\_x^{x+\Delta x}f(t)dt$$



so rearranging this we have



$$F(x+\Delta x)-F(x)=\int\limits\_x^{x+\Delta x}f(t)dt$$



According to the Mean Value Theorem for Integration, there exists a

$c\in[x,x+\Delta x]$ such that



$$\int\limits\_x^{x+\Delta x}f(t)dt=f(c)\cdot\Delta x$$



Notice that $c$ depends on $\Delta x$ . Anyway what we have shown is that



$$F(x+\Delta x)-F(x)=f(c)\cdot\Delta x$$



and dividing both sides by $\Delta x$ gives



$$\frac{F(x+\Delta x)-F(x)}{\Delta x}=f(c)$$



Take the limit as $\Delta x\to0$ we get the definition of the derivative

of $F$ at $x$ so we have



$$F'(x)=\lim\_{\Delta x\to0}\frac{F(x+\Delta x)-F(x)}{\Delta x}=\lim\_{\Delta x\to0}f(c)$$



To find the other limit, we will use the \textbf{squeeze theorem}.

$c\in[x,x+\Delta x]$ , so $x\le c\le x+\Delta x$ . Hence,



$$\lim\_{\Delta x\to0}\Big[x+\Delta x\Big]=x\quad\Rightarrow\quad\lim\_{\Delta x\to0}c=x$$



As $f$ is continuous we have



$$F'(x)=\lim\_{\Delta x\to0}f(c)=f\left(\lim\_{\Delta x\to0}c\right)=f(x)$$



which completes the proof.



$\blacksquare$



\paragraph{Proof of Fundamental Theorem of Calculus Part II}



Define $P(x)=\int\limits\_a^x f(t)dt$ . Then by the Fundamental Theorem

of Calculus part I we know that $P$ is differentiable on $(a,b)$ and for

all $x\in(a,b)$



$$P'(x)=f(x)$$



So $P$ is an antiderivative of $f$ . Since we were

assuming that $F$ was also an antiderivative for all $x\in(a,b)$ ,



$$\begin{aligned}

&P'(x)=F'(x) \\\&P'(x)-F'(x)=0 \\\&\Big(P(x)-F(x)\Big)'=0

\end{aligned}$$



Let $g(x)=P(x)-F(x)$ . The Mean Value Theorem applied to $g(x)$ on

$[a,\xi]$ with $a<\xi 

>

>  \textbf{Power Rule of Integration I}

>  $\int\limits\_a^b x^ndx=\frac{x^{n+1}}{n+1}\Bigg|\_a^b=\frac{b^{n+1}-a^{n+1}}{n+1}$

>     as long as $n\ne-1$ and $0\notin[a,b]$ .



Notice that we allow all values of $n$ , even negative or fractional. If

$n>0$ then this works even if $[a,b]$ includes $0$ .



> 

>

> \textbf{Power Rule of Integration II}

> $\int\limits\_a^b x^ndx=\frac{x^{n+1}}{n+1}\Bigg|\_a^b=\frac{b^{n+1}-a^{n+1}}{n+1}$

>     as long as $n>0$ .



Examples



\begin{itemize}
    \item To find $\int\limits\_1^2x^3dx$ we raise the power by 1 and have to
\end{itemize}
    divide by 4. So



$$\int\limits\_1^2x^3dx=\frac{x^4}{4}\Bigg|\_1^2=\frac{2^4}{4}-\frac{1^4}{4}=\frac{15}{4}$$



\begin{itemize}
    \item The power rule also works for negative powers. For instance
\end{itemize}

$$\int\limits\_1^3\frac{dx}{x^3}=\int\limits\_1^3x^{-3}dx=\frac{x^{-2}}{-2}\Bigg|\_1^3=\frac{1}{-2}\left(3^{-2}-1^{-2}\right)=-\frac12\left(\frac{1}{3^2}-1\right)=-\frac12\left(\frac19-1\right)=\frac12\cdot\frac89=\frac49$$



\begin{itemize}
    \item We can also use the power rule for fractional powers. For instance
\end{itemize}

$$\int\limits\_0^5\sqrt{x}dx=\int\limits\_0^5x^\frac12dx=\frac{x\sqrt{x}}{\frac32}\Bigg|\_0^5=\frac23\left(5^\frac32-0^\frac32\right)=\frac{10\sqrt5}{3}$$



\begin{itemize}
    \item Using linearity the power rule can also be thought of as applying to
\end{itemize}
    constants. For example,



$$=\int\limits\_3^{11}7dx=\int\limits\_3^{11}7x^0dx=7\int\limits\_3^{11}x^0dx=7x\Bigg|\_3^{11}=7(11-3)=56$$



\begin{itemize}
    \item Using the linearity rule we can now integrate any polynomial. For
\end{itemize}
    example



$$\int\limits\_0^3(3x^2+4x+2)dx=(x^3+2x^2+2x)\Bigg|\_0^3=3^3+2\cdot 3^2+2\cdot3-0=27+18+6=51$$



\paragraph{Exercises}



\paragraph{Questions}



1.  Evaluate $\int\limits\_0^1x^6dx$

2.  Evaluate $\int\limits\_1^2x^6dx$

3.  Evaluate $\int\limits\_0^2x^6dx$



\paragraph{Solutions}



1.  $\frac17=0.\overline{142857}$

2.  $\frac{127}{7}=18.\overline{142857}$

3.  $\frac{128}{7}=18.\overline{285714}$



\paragraph{Resources}



``{=html}The Fundamental Theorem of Calculus, Part

1``{=html}



\begin{itemize}
    \item \href{https://mathresearch.utsa.edu/wikiFiles/MAT1193/The\%20Fundamental\%20Theorem\%20of\%20Calculus/Presentation12\_DefiniteIntegral\%20\&\%20Antiderivatives.pptx}{Definite Integral \&
\end{itemize}
    Antiderivatives}

    (Slides 6\&7). PowerPoint file created by Professor Cynthia Roberts,

    UTSA.

\begin{itemize}
    \item \href{https://mathresearch.utsa.edu/wikiFiles/MAT1193/The\%20Fundamental\%20Theorem\%20of\%20Calculus/Presentation13\_FTC\%20Part\%20I.pptx}{Fundamental Theorem of Calculus Part
\end{itemize}
    1}.

    PowerPoint file created by Professor Cynthia Roberts, UTSA.

\begin{itemize}
    \item \href{https://mathresearch.utsa.edu/wikiFiles/MAT1214/The\%20Fundamental\%20Theorem\%20of\%20Calculus\_/MAT1214-TheFundamentalTheoremPwPt.pptx}{The Fundamental Theorem of
\end{itemize}
    Calculus}

    PowerPoint file created by Dr. Sara Shirinkam, UTSA.



<!-- -->



\begin{itemize}
    \item \href{https://youtu.be/PGmVvIglZx8}{Fundamental Theorem of Calculus Part
\end{itemize}
    1} by patrickJMT

\begin{itemize}
    \item \href{https://youtu.be/yrQ6PnE6D8w}{PART 1 OF THE DREADED FUNDAMENTAL THEOREM OF
\end{itemize}
    CALCULUS!} by Krista King

\begin{itemize}
    \item \href{https://youtu.be/aeB5BWY0RlE}{Fundamental Theorem of Calculus Part
\end{itemize}
    1} by The Organic Chemistry Tutor

\begin{itemize}
    \item \href{https://youtu.be/jyRdHbHeUuU}{The Second Fundamental Theorem of
\end{itemize}
    Calculus} by James Sousa, Math is

    Power 4U

\begin{itemize}
    \item \href{https://youtu.be/YE9jpfxEFYk}{Ex 1: The Second Fundamental Theorem of
\end{itemize}
    Calculus} by James Sousa, Math is

    Power 4U

\begin{itemize}
    \item \href{https://youtu.be/gJtkTRjqM6I}{Ex 2: The Second Fundamental Theorem of Calculus (Reverse
\end{itemize}
    Order)} by James Sousa, Math is Power

    4U

\begin{itemize}
    \item \href{https://youtu.be/X5ke8bUiiQM}{Ex 3: The Second Fundamental Theorem of
\end{itemize}
    Calculus} by James Sousa, Math is

    Power 4U

\begin{itemize}
    \item \href{https://youtu.be/ep52C-MUzDY}{Ex 4: The Second Fundamental Theorem of Calculus with Chain
\end{itemize}
    Rule} by James Sousa, Math is Power 4U

\begin{itemize}
    \item \href{https://youtu.be/8IrK5JcjMvM}{Ex 5: The Second Fundamental Theorem of Calculus with Chain
\end{itemize}
    Rule} by James Sousa, Math is Power 4U

\begin{itemize}
    \item \href{https://youtu.be/6p5NQwAeZRc}{Ex 6: Second Fundamental Theorem of Calculus with Chain
\end{itemize}
    Rule} by James Sousa, Math is Power 4U

\begin{itemize}
    \item \href{https://youtu.be/18qhfQ2IeLU}{Ex 7: Second Fundamental Theorem of Calculus with Chain
\end{itemize}
    Rule} by James Sousa, Math is Power 4U



``{=html}The Fundamental Theorem of Calculus, Part

2``{=html}



\begin{itemize}
    \item \href{https://youtu.be/7AteFyUCyZo}{The Fundamental Theorem of Calculus}
\end{itemize}
    by James Sousa, Math is Power 4U

\begin{itemize}
    \item \href{https://youtu.be/nHnZVFeQvNQ}{The Fundamental Theorem of Calculus Part
\end{itemize}
    2} by patrickJMT

\begin{itemize}
    \item \href{https://youtu.be/T-J7SkiE39Y}{PART 2 OF THE FUNDAMENTAL THEOREM OF
\end{itemize}
    CALCULUS!} by Krista King

\begin{itemize}
    \item \href{https://youtu.be/ns8N1UuXl4w}{The Fundamental Theorem of Calculus Part
\end{itemize}
    2} by The Organic Chemistry Tutor



\paragraph{Licensing}



Content obtained and/or adapted from:



\begin{itemize}
    \item \href{https://en.wikibooks.org/wiki/Calculus/Fundamental\_Theorem\_of\_Calculus}{Fundamental Theorem of Calculus, Wikibooks:
\end{itemize}
    Calculus}

    under a CC BY-SA license

\subsection{Learning Objectives}

\begin{enumerate}
    \item \textbf{Understand the Fundamental Theorem of Calculus}: Students will be able to state and explain the Fundamental Theorem of Calculus, including both its parts, and describe how it links the concepts of differentiation and integration.
    \item \textbf{Apply the Mean Value Theorem for Integration}: Students will be able to use the Mean Value Theorem for Integration to find a value $c$ within a given interval such that the average value of a continuous function over that interval equals the function's value at $c$.
    \item \textbf{Evaluate Definite Integrals Using Antiderivatives}: Students will be able to evaluate definite integrals by finding antiderivatives of the integrand and applying the Fundamental Theorem of Calculus to compute the integral as the difference between the antiderivative evaluated at the upper and lower bounds.
\end{enumerate}

\subsection{Assessment Instrument}

\begin{enumerate}
    \item \textbf{Part 1: Understand the Fundamental Theorem of Calculus}
    \item \textbf{Multiple Choice:} Which of the following statements correctly describes the Fundamental Theorem of Calculus?
    \item A. The Fundamental Theorem of Calculus states that the integral of a function can be approximated using Riemann sums.
    \item B. The Fundamental Theorem of Calculus links the concept of differentiation to that of integration by stating that differentiation and integration are inverse operations.
    \item C. The Fundamental Theorem of Calculus provides a method for solving differential equations.
    \item D. The Fundamental Theorem of Calculus states that the limit of a Riemann sum approaches the derivative of the function.
    \item \textbf{True/False:} The Fundamental Theorem of Calculus Part 1 asserts that if a function $F$ is continuous on $[a, b]$ and $f$ is its derivative on $(a, b)$, then $$\int_{a}^{b} f(x) \, dx = F(b) - F(a)$$. (True/False)
    \item \textbf{Short Answer:} Explain the relationship between differentiation and integration as outlined in the Fundamental Theorem of Calculus. Describe both parts of the theorem and how they connect these two operations. \textbf{Part 2: Apply the Mean Value Theorem for Integration}
    \item \textbf{Calculation:} Given the function $f(x) = x^2$ on the interval $[1, 4]$, use the Mean Value Theorem for Integration to find the value of $c$ in $[1, 4]$ such that the average value of $f(x)$ over this interval equals $f(c)$.
    \item \textbf{Multiple Choice:} The Mean Value Theorem for Integration guarantees that for a continuous function $f(x)$ on $[a, b]$, there exists a point $c$ in $[a, b]$ such that:
    \item A. $$f(c) = \frac{1}{b-a} \int_{a}^{b} f(x) \, dx$$
    \item B. $$f(c) = \int_{a}^{b} f(x) \, dx - f(a)$$
    \item C. $$f(c) = \frac{1}{b-a} \int_{a}^{b} f'(x) \, dx$$
    \item D. $$f(c) = \frac{1}{b-a} \int_{a}^{b} f(x) \, dx + f(b)$$
    \item \textbf{Short Answer:} Apply the Mean Value Theorem for Integration to the function $f(x) = \sin(x)$ over the interval $[0, \pi]$. Determine the value $c$ in this interval where the function's value equals its average value over the interval. \textbf{Part 3: Evaluate Definite Integrals Using Antiderivatives}
    \item \textbf{Calculation:} Evaluate the definite integral $$\int_{1}^{3} (4x^3 - 2x) \, dx$$ by finding the antiderivative of the integrand and applying the Fundamental Theorem of Calculus. Provide the steps and the final result.
    \item \textbf{Multiple Choice:} To compute the definite integral $$\int_{2}^{5} \frac{1}{x^2} \, dx$$, which antiderivative of $\frac{1}{x^2}$ is used in the Fundamental Theorem of Calculus?
    \item A. $-\frac{1}{x}$
    \item B. $\frac{1}{x}$
    \item C. $-\frac{1}{x^2}$
    \item D. $\frac{1}{x^2}$
    \item \textbf{Short Answer:} Find the definite integral $$\int_{0}^{\pi} \cos(x) \, dx$$ using the Fundamental Theorem of Calculus. Describe the antiderivative used and compute the integral by evaluating it at the upper and lower bounds.
\end{enumerate}

%%===============================================================
\section{Applications of Integrals}
\label{sec:applications-of-integrals}
%%===============================================================

\paragraph{Area between two curves}



Suppose we are given two functions $y\_1=f(x)$ and $y\_2=g(x)$ and we want

to find the area between them on the interval $[a,b]$ . Also assume that

$f(x)\ge g(x)$ for all $x$ on the interval $[a,b]$ . Begin by

partitioning the interval $[a,b]$ into $n$ equal subintervals each

having a length of $\Delta x=\frac{b-a}{n}$ . Next choose any point in

each subinterval, $x\_i^*$ . Now we can 'create' rectangles on each

interval. At the point $x\_i*$ , the height of each rectangle is

$f(x\_i^*)-g(x\_i^*)$ and the width is $\Delta x$ . Thus the area of each

rectangle is $\bigl[f(x\_i^*)-g(x\_i^*)\bigr]\Delta x$ . An

\textit{approximation} of the area, $A$ , between the two curves is



$$A:=\sum\_{i=1}^n \Big[f(x\_i^*)-g(x\_i^*)\Big]\Delta x.$$



Now we take the limit as $n$ approaches infinity and get



$$A=\lim\_{n\to\infty}\sum\_{i=1}^n \Big[f(x\_i^*)-g(x\_i^*)\Big]\Delta x$$



which gives the exact area. Recalling the definition of the definite

integral we notice that



$$A=\int\limits\_a^b \bigl(f(x)-g(x)\bigr)dx.$$



This formula of finding the area between two curves is sometimes known as applying integration

with respect to the \textit{x}-axis since the rectangles used to approximate

the area have their bases lying parallel to the \textit{x}-axis. It will be

most useful when the two functions are of the form $y\_1=f(x)$ and

$y\_2=g(x)$ . Sometimes however, one may find it simpler to integrate

with respect to the \textit{y}-axis. This occurs when integrating with respect

to the \textit{x}-axis would result in more than one integral to be evaluated.

These functions take the form $x\_1=f(y)$ and $x\_2=g(y)$ on the interval

$[c,d]$ . Note that $[c,d]$ are values of $y$ . The derivation of this

case is completely identical. Similar to before, we will assume that

$f(y)\ge g(y)$ for all $y$ on $[c,d]$ . Now, as before we can divide the

interval into $n$ subintervals and create rectangles to approximate the

area between $f(y)$ and $g(y)$ . It may be useful to picture each

rectangle having their 'width', $\Delta y$ , parallel to the \textit{y}-axis

and 'height', $f(y\_i^*)-g(y\_i^*)$ at the point $y\_i^*,$ parallel to

the \textit{x}-axis. Following from the work above we may reason that an

\textit{approximation} of the area, $A$ , between the two curves is



$$A:=\sum\_{i=1}^n \Big[f(y\_i^*)-g(y\_i^*)\Big]\Delta y.$$



As before, we take the limit as $n$ approaches infinity to arrive at



$$A=\lim\_{n\to\infty}\sum\_{i=1}^n \Big[f(y\_i^*)-g(y\_i^*)\Big]\Delta y,$$



which is nothing more than a definite integral, so



$$A=\int\limits\_c^d \bigl(f(y)-g(y)\bigr)dy.$$



Regardless of the form of the functions, we basically use the same formula.



% [Image: Closed\_path\_integral\_defined]



\paragraph{Volume}



If the function $A(x)$ is continuous on $[a,b]$ , then the volume $V\_S$

of the solid $S$ is given by:



$$V\_S=\int\limits\_a^b A(x)dx$$



\paragraph{Example 1: A right cylinder}



Now we will calculate the volume of a right cylinder using our new ideas

about how to calculate volume. Since we already know the formula for the

volume of a cylinder this will give us a "sanity check" that our

formulas make sense. First, we choose a dimension along which to

integrate. In this case, it will greatly simplify the calculations to

integrate along the height of the cylinder, so this is the direction we

will choose. Thus we will call the vertical direction $x$. Now

we find the function, $A(x)$ , which will describe the cross-sectional

area of our cylinder at a height of $x$ . The cross-sectional area of a

cylinder is simply a circle. Now simply recall that the area of a circle

is $\pi r^2$ , and so $A(x)=\pi r^2$ . Before performing the

computation, we must choose our bounds of integration. In this case, we

simply define $x=0$ to be the base of the cylinder, and so we will

integrate from $x=0$ to $x=h$ , where $h$ is the height of the cylinder.



% [Image: Figure

1]



Finally, we integrate:



$$\begin{aligned}

V\_{\mathrm{cylinder}}&=\int\limits\_a^b A(x)dx \\\&=\int\limits\_0^h\pi r^2dx \\\&=\pi r^2\int\limits\_0^hdx \\\&=\pi r^2x\bigg|\_{x=0}^h \\\&=\pi r^2(h-0) \\\&=\pi r^2h

\end{aligned}$$



This is exactly the familiar formula for the volume of a cylinder.



\paragraph{Example 2: A right circular cone}



For our next example we will look at an example where the cross

sectional area is not constant. Consider a right circular cone. Once

again the cross sections are simply circles. But now the radius varies

from the base of the cone to the tip. Once again we choose $x$ to be the

vertical direction, with the base at $x=0$ and the tip at $x=h$ , and we

will let $R$ denote the radius of the base. While we know the cross

sections are just circles we cannot calculate the area of the cross

sections unless we find some way to determine the radius of the circle

at height $x$ .



Luckily in this case it is possible to use some of what we know from

geometry. We can imagine cutting the cone perpendicular to the base

through some diameter of the circle all the way to the tip of the cone.

If we then look at the flat side we just created, we will see simply a

triangle, whose geometry we understand well. The right triangle from the

tip to the base at height $x$ is similar to the right triangle from the

tip to the base at height $h$ . This tells us that

$\frac{r}{h-x}=\frac{R}{h}$ . So that we see that the radius of the

circle at height $x$ is $r(x)=\frac{R}{h}(h-x)$ . Now using the familiar

formula for the area of a circle we see that

$A(x)=\pi\frac{R^2}{h^2}(h-x)^2$ .



% [Image: Figure 2: The cross-section of a right circular cone by a plane

perpendicular to the axis of the cone is a

circle.]



Now we are ready to integrate.



$$\begin{aligned}

V\_{\mathrm{cone}}&=\int\limits\_a^b A(x)dx \\\&=\int\limits\_0^h \pi\frac{R^2}{h^2}(h-x)^2dx \\\&=\pi\frac{R^2}{h^2}\int\limits\_0^h(h-x)^2dx

\end{aligned}$$



By u-substitution we may let $u=h-x$ , then $du=-dx$ and our integral becomes



$$\begin{aligned}

&&=\pi\frac{R^2}{h^2}\left(-\int\limits\_h^0 u^2du\right) \\\&&=\pi\frac{R^2}{h^2}\left(-\frac{u^3}{3}\bigg|\_h^0\right) \\\&&=\pi\frac{R^2}{h^2}\left(-0+\frac{h^3}{3}\right) \\\&&=\frac{\pi}{3}R^2h

\end{aligned}$$



% [Image: Figure 3: Cross-section of the right circular cone by a plane

perpendicular to the base and passing through the

tip.]



\paragraph{Example 3: A sphere}



In a similar fashion, we can use our definition to prove the well known

formula for the volume of a sphere. First, we must find our

cross-sectional area function, $A(x)$ . Consider a sphere of radius $R$

which is centered at the origin in $\R^3$ . If we again integrate

vertically then $x$ will vary from $-R$ to $R$ . In order to find the

area of a particular cross section it helps to draw a right triangle

whose points lie at the center of the sphere, the center of the circular

cross section, and at a point along the circumference of the cross

section. As shown in the diagram the side lengths of this triangle will

be $R$ , $|x|$ , and $r$ . Where $r$ is the radius of the circular cross

section. Then by the Pythagorean theorem $r=\sqrt{R^2-|x|^2}$ and find

that $A(x)=\pi(R^2-|x|^2)$ . It is slightly helpful to notice that

$|x|^2=x^2$ so we do not need to keep the absolute value.



% [Image: Figure 4: Determining the radius of the cross-section of the sphere at

a distance $|x|$ from the sphere's

center.]



So we have that



$$\begin{aligned}

V\_{\mathrm{sphere}}&=\int\limits\_a^b A(x)dx \\\&=\int\limits\_{-R}^R\pi(R^2-x^2)dx \\\&=\pi\int\limits\_{-R}^R R^2dx-\pi\int\limits\_{-R}^R x^2dx \\\&=\pi R^2x\Bigg|\_{-R}^R-\pi\frac{x^3}{3}\Bigg|\_{-R}^R \\\&=\pi R^2(R-(-R))-\pi\left(\frac{R^3}{3}-\frac{(-R)^3}{3}\right) \\\&=2\pi R^3-\frac{2\pi}{3}R^3=\frac{4\pi}{3}R^3

\end{aligned}$$



\paragraph{Arc length}



Suppose that $f'$ is continuous on $[a,b]$ . Then the length of the

curve given by $y=f(x)$ between $a$ and $b$ is given by



$$L=\int\limits\_a^b \sqrt{1+f'(x)^2}dx$$



And in Leibniz notation



$$L=\int\limits\_a^b \sqrt{1+\left(\tfrac{dy}{dx}\right)^2}dx$$



\textbf{Proof:} Consider $y\_{i+1}-y\_i=f(x\_{i+1})-f(x\_i)$ . By the Mean Value

Theorem there is a point $z\_i$ in $(x\_{i+1},x\_i)$ such that



$$y\_{i+1}-y\_i=f(x\_{i+1})-f(x\_i)=f'(z\_i)(x\_{i+1}-x\_i)$$



So



$$\begin{aligned}

\bigl|P\_iP\_{i+1}\bigr&=\sqrt{(x\_{i+1}-x\_i)^2+(y\_{i+1}-y\_i)^2} \\\&=\sqrt{(x\_{i+1}-x\_i)^2+f'(z\_i)^2(x\_{i+1}-x\_i)^2} \\\&=\sqrt{\bigl(1+f'(z\_i)^2\bigr)(x\_{i+1}-x\_i)^2} \\\&=\sqrt{1+f'(z\_i)^2}\Delta x

\end{aligned}$$



Putting this into the definition of the length of $C$ gives



$$L=\lim\_{n\to\infty}\sum\_{i=0}^{n-1}\sqrt{1+f'(z\_i)^2}\Delta x$$



Now this is the definition of the integral of the function

$g(x)=\sqrt{1+f'(x)^2}$ between $a$ and $b$ (notice that $g$ is

continuous because we are assuming that $f'$ is continuous). Hence



$$L=\int\limits\_a^b \sqrt{1+f'(x)^2}dx$$



as claimed.



\paragraph{Example}



Length of the curve $y=2x$ from $x=0$ to $x=1$}} As a sanity check of

our formula, let's calculate the length of the "curve" $y=2x$ from

$x=0$ to $x=1$ . First let's find the answer using the Pythagorean

Theorem.



$$P\_0=(0,0)$$



and



$$P\_1=(1,2)$$



so the length of the curve, $s$ , is



$$s=\sqrt{2^2+1^2}=\sqrt5$$



Now let's use the formula



$$s=\int\limits\_0^1 \sqrt{1+\left(\tfrac{d(2x)}{dx}\right)^2}\,dx=\int\limits\_0^1 \sqrt{1+2^2}\,dx=\sqrt5x\bigg|\_0^1=\sqrt5$$



\paragraph{Resources}



\begin{itemize}
    \item \href{https://www.youtube.com/watch?v=FE0eJGdWC04}{Business Calculus: Application of definite
\end{itemize}
    integral}, Rajendra

    Dahal (YouTube)

\begin{itemize}
    \item \href{https://tutorial.math.lamar.edu/classes/calci/intappsintro.aspx}{Applications of Integrals - Cal
\end{itemize}
    I},

    Paul's Online Notes

\begin{itemize}
    \item \href{https://tutorial.math.lamar.edu/classes/calcii/intappsintro.aspx}{Applications of Integrals - Cal
\end{itemize}
    II},

    Paul's Online Notes



\paragraph{Licensing}



Content obtained and/or adapted from:



\begin{itemize}
    \item \href{https://en.wikibooks.org/wiki/Calculus/Area}{Area, Wikibooks:
\end{itemize}
    Calculus} under a CC

    BY-SA license

\begin{itemize}
    \item \href{https://en.wikibooks.org/wiki/Calculus/Volume}{Volume, Wikibooks:
\end{itemize}
    Calculus} under a CC

    BY-SA license

\begin{itemize}
    \item \href{https://en.wikibooks.org/wiki/Calculus/Arc\_length}{Arc length, Wikibooks:
\end{itemize}
    Calculus} under a

    CC BY-SA license

\subsection{Learning Objectives}

\begin{enumerate}
    \item \textbf{Calculate the Area Between Two Curves Using Definite Integrals}
    \item Students will be able to determine the area between two curves $y_1 = f(x)$ and $y_2 = g(x)$ over a specified interval $[a, b]$ by evaluating the integral $\int_a^b (f(x) - g(x)) \, dx$ or $\int_c^d (f(y) - g(y)) \, dy$ as appropriate.
    \item \textbf{Apply Integration Techniques to Compute Volumes of Solids of Revolution}
    \item Students will be able to compute the volume of various solids of revolution, including cylinders, cones, and spheres, by applying the formula $V = \int_a^b A(x) \, dx$ where $A(x)$ represents the cross-sectional area of the solid.
    \item \textbf{Determine the Arc Length of a Curve Using Integration}
    \item Students will be able to find the length of a curve $y = f(x)$ on the interval $[a, b]$ using the arc length formula $L = \int_a^b \sqrt{1 + (f'(x))^2} \, dx$, and apply this method to practical examples.
\end{enumerate}

\subsection{Assessment Instrument}

\begin{enumerate}
    \item \textbf{Part 1: Calculate the Area Between Two Curves Using Definite Integrals}
    \item \textbf{Calculation:} Find the area between the curves $y\_1 = x^2$ and $y\_2 = x$ over the interval $[0, 1]$. Set up and evaluate the integral that represents this area.
    \item \textbf{Multiple Choice:} To find the area between the curves $y\_1 = 3x$ and $y\_2 = x^2$ over the interval $[1, 2]$, which integral would you use?
    \item A. $$\int_{1}^{2} (3x - x^2) \, dx$$
    \item B. $$\int_{1}^{2} (x^2 - 3x) \, dx$$
    \item C. $$\int_{1}^{2} (x^2 + 3x) \, dx$$
    \item D. $$\int_{1}^{2} (3x + x^2) \, dx$$
    \item \textbf{Short Answer:} Determine the area between the curves $y\_1 = \sqrt{x}$ and $y\_2 = 2x - 1$ from $x = 0$ to $x = 2$. Provide the steps and the final result. \textbf{Part 2: Apply Integration Techniques to Compute Volumes of Solids of Revolution}
    \item \textbf{Calculation:} Compute the volume of the solid obtained by rotating the region bounded by $y = x^2$ and the x-axis from $x = 0$ to $x = 1$ about the x-axis. Use the disk method to find the volume.
    \item \textbf{Multiple Choice:} To find the volume of a solid formed by rotating the region between $y = x$ and $y = 0$ from $x = 1$ to $x = 3$ about the y-axis, which formula should be used?
    \item A. $$V = \pi \int_{1}^{3} (x^2) \, dx$$
    \item B. $$V = \pi \int_{1}^{3} (x^2 - 0) \, dx$$
    \item C. $$V = \pi \int_{1}^{3} (x^2) \, dx$$
    \item D. $$V = \pi \int_{1}^{3} (x - 0)^2 \, dx$$
    \item \textbf{Short Answer:} Find the volume of the solid generated by revolving the region bounded by $y = 2x$ and $y = 4 - x^2$ about the x-axis. Use the washer method and provide the integral setup and final result. \textbf{Part 3: Determine the Arc Length of a Curve Using Integration}
    \item \textbf{Calculation:} Calculate the arc length of the curve $y = \ln(x)$ on the interval $[1, e]$ using the arc length formula. Set up and evaluate the integral to find the length.
    \item \textbf{Multiple Choice:} To find the arc length of the curve $y = x^3$ from $x = 0$ to $x = 2$, which formula should you use?
    \item A. $$L = \int_{0}^{2} \sqrt{1 + (3x^2)^2} \, dx$$
    \item B. $$L = \int_{0}^{2} \sqrt{1 + (3x)^2} \, dx$$
    \item C. $$L = \int_{0}^{2} \sqrt{1 + (x^3)^2} \, dx$$
    \item D. $$L = \int_{0}^{2} \sqrt{1 + (3x^2)^2} \, dx$$
    \item \textbf{Short Answer:} Determine the arc length of the curve $y = \sqrt{x}$ from $x = 1$ to $x = 4$. Use the arc length formula and provide all steps leading to the final answer.
\end{enumerate}

%%===============================================================
\section{Differential Equations}
\label{sec:differential-equations}
%%===============================================================

\textbf{Ordinary differential equations} involve equations containing:



\begin{itemize}
    \item variables
    \item functions
    \item their derivatives
\end{itemize}

and their solutions.



In studying integration, you \textit{already} have considered solutions to very

simple differential equations. For example, when you look to solving



$$\int f(x) \,dx=g(x)$$



for g(x), you are really solving the differential equation



$$g'(x) = f(x) \,$$



\paragraph{Notations and terminology}



The notations we use for solving differential equations will be crucial

in the ease of solubility for these equations.



This document will be using \textbf{three} notations primarily:



\begin{itemize}
    \item f' to denote the derivative of f
    \item D \textit{f} to denote the derivative of \textit{f}
    \item ${df \over dx}$ to denote the derivative of \textit{f} (for separable
\end{itemize}
    equations).



\paragraph{Terminology}



Consider the differential equation



$$3 f^{\prime \prime}(x)+5xf(x)=11$$



Since the equation's highest derivative is 2, we say that the

differential equation is of \textit{order} 2.



\paragraph{Some simple differential equations}



A key idea in solving differential equations will be that of

integration.



Let us consider the second order differential equation (remember that a

function acts on a value).



$$f''(x)=2 \,$$



How would we go about solving this? It tells us that on differentiating

twice, we obtain the constant 2 so, if we integrate twice, we should

obtain our result.



Integrating once first of all:



$$\int f''(x) \,dx = \int 2 \,dx$$



$$f'(x)=2x+C\_1 \,$$



We have transformed the apparently difficult second order differential

equation into a rather simpler one, viz.



$$f'(x)=2x+C\_1 \,$$



This equation tells us that if we differentiate a function once, we get

$2x+C\_1.$ If we integrate once more, we should find the solution.



$$\int f'(x) \,dx = \int 2x+C\_1 \,dx$$



$$f(x)=x^2+C\_1x+C\_2 \,$$



This is the \textit{solution} to the differential equation. We will get

$f''=2 \,$ for \textit{all} values of $C\_1$ and $C\_2.$



The values $C\_1$ and $C\_2$ are related to quantities known as *initial

conditions*.



Why are initial conditions useful? ODEs (ordinary differential

equations) are useful in modeling physical conditions. We may wish to

model a certain physical system which is initially at rest (so one

initial condition may be zero), or wound up to some point (so an initial

condition may be nonzero, say 5 for instance) and we may wish to see how

the system reacts under such an initial condition.



When we solve a system with given initial conditions, we substitute them

after our process of integration.



\paragraph{Example}



When we solved $f''(x)=2 \,$ say we had the initial conditions

$f'(0)=3 \,$ and $f(0)=2 \,.$ (Note, initial conditions need not occur

at f(0)).



After we integrate we make substitutions:



$$f'(0)=2(0)+C\_1 \,$$



$$3=C\_1 \,$$



$$\int f'(x) \,dx = \int 2x+3 \,dx$$



$$f(x)=x^2+3x+C\_2 \,$$



$$f(0)=0^2+3(0)+C\_2 \,$$



$$2=C\_2 \,$$



$$f(x)=x^2+3x+2 \,$$



Without initial conditions, the answer we obtain is known as the

\textit{general solution} or the solution to the \textit{family of equations}. With

them, our solution is known as a \textit{specific solution}.



\paragraph{Basic first order DEs}



In this section we will consider \textit{four} main types of differential

equations:



\begin{itemize}
    \item separable
    \item homogeneous
    \item linear
    \item exact
\end{itemize}

There are many other forms of differential equation, however, and these

will be dealt with in the next section



\paragraph{Separable equations}



A \textit{separable} equation is in the form (using dy/dx notation which will

serve us greatly here)



$${dy \over dx} = f(x)/g(y)$$



Previously we have only dealt with simple differential equations with

g(\textit{y})=1. How do we solve such a separable equation as above?



We group \textit{x} and \textit{dx} terms together, and \textit{y} and \textit{dy} terms together as

well.



$$g(y)\ dy = f(x)\ dx$$



Integrating both sides with respect to y on the left hand side and x on the right hand side:



$$\int g(y)\,dy=\int f(x)\,dx+C$$



we will obtain the solution.



\paragraph{Worked example}



Here is a worked example illustrating the process.



We are asked to solve



$${dy \over dx} = 3x^2y$$



Separating



$${dy \over y} = (3x^2)\,dx$$



Integrating



$$\int {dy \over y} = \int 3x^2\,dx$$



$$\ln{y}=x^3+C \,\!$$



$$y=e^{x^3+C}$$



Letting $k = e^C$ where k is a constant we obtain



$$y=ke^{x^3}$$



which is the general solution.



\paragraph{Verification}



This step does not need to be part of your work, but if you want to

check your solution, you can verify your answer by differentiation.



We obtained



$$y=ke^{x^3}$$



as the solution to



$${dy \over dx} = 3x^2y$$



Differentiating our solution with respect to x,



$${dy \over dx} = 3kx^2e^{x^3}$$



And since $y=ke^{x^3},$ we can write



$${dy \over dx} = 3x^2y$$



We see that we obtain our original differential equation, thus our work must be correct.



\paragraph{Homogeneous equations}



A \textit{homogeneous} equation is in the form



$${dy \over dx} = f(y/x)$$



This looks difficult as it stands, however we can utilize the

substitution



$$v = {y \over x}$$



so that we are now dealing with F(v) rather than F(y/x).



Now we can express y in terms of v, as \textit{y}=\textit{xv} and use the product rule.



The equation above then becomes, using the product rule



$${dy \over dx} = v+x{dv \over dx}$$



Then



$$v+x{dv \over dx} = f(v)$$



$$x{dv \over dx} = f(v)-v$$



$${dv \over dx} = {f(v)-v \over x}$$



which is a separable equation and can be solved as above.



However let's look at a worked equation to see how homogeneous

equations are solved.



\paragraph{Worked example}



We have the equation



$${dy \over dx} = {y^2 + x^2 \over yx}$$



This does not appear to be immediately separable, but let us expand to get



$${dy \over dx} = {y^2 \over yx} + {x^2 \over yx}$$



$${dy \over dx} = {x \over y} + {y \over x}$$



Substituting \textit{y}=\textit{xv} which is the same as substituting \textit{v}=\textit{y}/\textit{x}:



$${dy \over dx} = 1/v + v$$



Now



$$v+x{dv \over dx} = 1/v + v$$



Canceling \textit{v} from both sides



$$x{dv \over dx} = 1/v$$



Separating



$$v\, dv = dx/x$$



Integrating both sides



$${1 \over 2}v^2+C= \ln(x) \,$$



$${1 \over 2}\left({y \over x}\right)^2= \ln(x)-C$$



$$y^2 = 2x^2 \ln(x) - 2Cx^2 \,$$



$$y = x\sqrt{2 \ln(x) - 2C}$$



which is our desired solution.



\paragraph{Linear equations}



A linear first order differential equation is a differential equation in the form



$$a(x){dy \over dx} + b(x)y=c(x)$$



Multiplying or dividing this equation by any non-zero function of \textit{x}

makes no difference to its solutions so we could always divide by

\textit{a}(\textit{x}) to make the coefficient of the differential 1, but writing the

equation in this more general form may offer insights.



At first glance, it is not possible to integrate the left hand side, but

there is one special case. If \textit{b} happens to be the differential of \textit{a}

then we can write



$$a(x){dy \over dx} + b(x)y = a(x){dy \over dx} + y{da \over dx}

= {d \over dx}a(x)y$$



and integration is now straightforward.



Since we can freely multiply by any function, lets see if we can use

this freedom to write the left hand side in this special form.



We multiply the entire equation by an arbitrary, \textit{I}(\textit{x}), getting



$$aI{dy \over dx} + bIy=cI$$



then impose the condition



$$\frac{d}{dx}aI = bI$$



If this is satisfied the new left hand side will have the special form.

Note that multiplying \textit{I} by any constant will leave this condition

still satisfied.



Rearranging this condition gives



$$\frac{1}{I}\frac{dI}{dx} = \frac{b-\frac{da}{dx}}{a}$$



We can integrate this to get



$$\ln I(x) = \int \frac{b(z)}{a(z)}dz - \ln a(x) + c \quad

I(x)=\frac{k}{a(x)}e^{\int \frac{b(z)}{a(z)}dz}$$



We can set the constant \textit{k} to be 1, since this makes no difference.



Next we use \textit{I} on the original differential equation, getting



$$e^{\int \frac{b(z)}{a(z)}dz}{dy \over dx} + 

e^{\int \frac{b(z)}{a(z)}dz} \frac{b(x)}{a(x)}y

=e^{\int \frac{b(z)}{a(z)}dz}\frac{c(x)}{a(x)}$$



Because we've chosen \textit{I} to put the left hand side in the special form

we can rewrite this as



$${d \over dx}(ye^{\int \frac{b(z)}{a(z)}dz}) = 

e^{\int \frac{b(z)}{a(z)}dz}\frac{c(x)}{a(x)}$$



Integrating both sides and dividing by $e^I$ we obtain the final result



$$y = e^{-\int \frac{b(z)}{a(z)}dz}

\left(\int e^{\int \frac{b(z)}{a(z)}dz}\frac{c(x)}{a(x)}dx + C\right)$$



We call \textit{I} an \textit{integrating factor}. Similar techniques can be used on

some other calculus problems.



\paragraph{Example}



Consider



$$\frac{dy}{dx} + y \tan x = 1 \quad y(0)=0$$



First we calculate the integrating factor.



$$I=e^{\int \tan x dx} = e^ {\ln \sec x} = \sec x$$



Multiplying the equation by this gives



$$\sec x \frac{dy}{dx} + y \sec x \tan x = \sec x$$



or



$$\frac{d}{dx} y\sec x = \sec x$$



We can now integrate



$$y = \cos x \int\_0^x \sec z \, dz = \cos x \ln (\sec x + \tan x)$$



\paragraph{Exact equations}



An exact equation is in the form



$$f(x, y) dx + g(x, y) dy = 0$$



and, has the property that



$$D\_x f = D\_y g$$



(If the differential equation does not have this property then we can't

proceed any further).



As a result of this, if we have an exact equation then there exists a

function h(\textit{x}, \textit{y}) such that



$$D\_y h = f \text{and} D\_x h = g$$



So then the solutions are in the form



$$h(x, y) = c$$



by using the fact of the total differential. We can find then h(\textit{x},

\textit{y}) by integration



\paragraph{Basic second and higher order ODE's}



The generic solution of a \textit{$n^{\text{th}}$} order ODE will contain \textit{n} constants

of integration. To calculate them we need \textit{n} more equations. Most

often, we have either



boundary conditions, the values of \textit{y} and its derivatives take for two different values of \textit{x}



or



initial conditions, the values of \textit{y} and its first \textit{n-1} derivatives take for one particular value of \textit{x}.



\paragraph{Reducible ODE's}



1. If the independent variable, \textit{x}, does not occur in the differential

equation then its order can be lowered by one. This will reduce a second

order ODE to first order.



Consider the equation:



$$F\left(y,\frac{dy}{dx},\frac{d^2y}{dx^2}\right)=0$$



Define



$$u=\frac{dy}{dx}$$



Then



$$\frac{d^2y}{dx^2}=\frac{du}{dx}=\frac{du}{dy}\cdot\frac{dy}{dx}=\frac{du}{dy}\cdot u$$



Substitute these two expression into the equation and we get



$$F\left(y,u,\frac{du}{dy}\cdot u\right)=0$$



which is a first order ODE



\paragraph{Example}



Solve



$$1+2y^2\operatorname{D}^2y=0$$



if at \textit{x}=0,  \textit{y}=D\textit{y}=1



First, we make the substitution, getting



$$1+2y^2 u \frac{du}{dy}=0$$



This is a first order ODE. By rearranging terms we can separate the variables



$$udu=-\frac{dy}{2y^2}$$



Integrating this gives



$$u^2/2=c+1/2y$$ 



We know the values of \textit{y} and \textit{u} when \textit{x}=0 so we can find \textit{c}



$$c=u^2/2-1/2y=1^2/2-1/(2\cdot 1)=1/2-1/2=0$$



Next, we reverse the substitution



$$\frac{dy}{dx}^2=u^2=\frac{1}{y}$$



and take the square root



$$\frac{dy}{dx}=\pm \frac{1}{\sqrt{y}}$$ 



To find out which sign of the square root to keep, we use the initial condition, D\textit{y}=1 at \textit{x}=0,

again, and rule out the negative square root. We now have another separable first order ODE,



$$\frac{dy}{dx}=\frac{1}{\sqrt{y}}$$



Its solution is



$$\frac{2}{3}y^\frac{3}{2}= x+d$$



Since \textit{y}=1 when \textit{x}=0, \textit{d}=2/3, and



$$y=\left(1 + \frac{3x}{2} \right)^\frac{2}{3}$$



2. If the dependent variable, \textit{y}, does not occur in the differential

equation then it may also be reduced to a first order equation.



Consider the equation:



$$F\left(x,\frac{dy}{dx},\frac{d^2y}{dx^2}\right)=0$$



Define



$$u=\frac{dy}{dx}$$



Then



$$\frac{d^2y}{dx^2}=\frac{du}{dx}$$



Substitute these two expressions into the first equation and we get



$$F\left(x,u,\frac{du}{dx}\right)=0$$



which is a first order ODE



\paragraph{Linear ODEs}



An ODE of the form



$$\frac{d^ny}{dx^n}+a\_1(x)\frac{d^{n-1}y}{dx^{n-1}}+ ... +a\_n y=F(x)$$



is called \textbf{linear}. Such equations are much simpler to solve than

typical non-linear ODEs. Though only a few special cases can be solved

exactly in terms of elementary functions, there is much that can be said

about the solution of a generic linear ODE. A full account would be

beyond the scope of this book



If \textit{F(x)=0} for all \textit{x} the ODE is called \textbf{homogeneous}



Two useful properties of generic linear equations are



1.  Any linear combination of solutions of an homogeneous linear

    equation is also a solution.

2.  If we have a solution of a nonhomogeneous linear equation and we add

    any solution of the corresponding homogenous linear equation we get

    another solution of the nonhomogeneous linear equation



\paragraph{Variation of constants}



Suppose we have a linear ODE,



$$\frac{d^ny}{dx^n}+a\_1(x)\frac{d^{n-1}y}{dx^{n-1}}+ ... +a\_n y=0$$



and we know one solution, \textit{y=w(x)}



The other solutions can always be written as \textit{y=wz}. This substitution

in the ODE will give us terms involving every differential of \textit{z} upto

the \textit{$n^{\text{th}}$}, no higher, so we'll end up with an \textit{$n^{\text{th}}$} order linear

ODE for \textit{z}.



We know that \textit{z} is constant is one solution, so the ODE for \textit{z} must

not contain a \textit{z} term, which means it will effectively be an \textit{$n-1^{\text{th}}$}

order linear ODE. We will have reduced the order by one.



Lets see how this works in practice.



\paragraph{Example}



Consider



$$\frac{d^2y}{dx^2}+\frac{2}{x}\frac{dy}{dx}-\frac{6}{x^2}y=0$$



One solution of this is \textit{y=x^2^}, so substitute \textit{y=zx^2^} into this

equation.



$$\left( x^2\frac{d^2z}{dx^2}+4x\frac{dz}{dx}+2z\right) 

+\frac{2}{x} \left( x^2\frac{dz}{dx}+2xz \right) -\frac{6}{x^2}x^2 z=0$$



Rearrange and simplify.



$$x^2 D^2 z + 6xD z=0$$



This is first order for D\textit{z}. We can solve it to get



$$z=A x^{-5} \quad y=A x^{-3}$$



Since the equation is linear we can add this to any multiple of the

other solution to get the general solution,



$$y=A x^{-3} + B x^2$$



\paragraph{Linear homogeneous ODE's with constant coefficients}



Suppose we have a ODE



$$(D^n+a\_1 D^{n-1}+ ... + a\_{n-1}D+a\_0)y=0$$



we can take an inspired guess at a solution (motivate this)



$$y=e^{px}$$



For this function \textit{$D^ny=p^ny$} so the ODE becomes



$$(p^n+a\_1 p^{n-1}+ ... + a\_{n-1}p+a\_0)y=0$$



\textit{y=0} is a trivial solution of the ODE so we can discard it. We are then

left with the equation



$$p^n+a\_1 p^{n-1}+ ... + a\_{n-1}p+a\_0)=0$$



This is called the \textit{characteristic} equation of the ODE.



It can have up to \textit{n} roots, \textit{$p\_1$}, \textit{$p\_2$} ... \textit{$p\_n$}, each root giving us a

different solution of the ODE.



Because the ODE is linear, we can add all those solution together in any

linear combination to get a general solution



$$y=A\_1 e^{p\_1 x} +A\_2 e^{p\_2 x} + ... + A\_n e^{p\_n x}$$



To see how this works in practice we will look at the second order case.

Solving equations like this of higher order uses exactly the same

principles; only the algebra is more complex.



\paragraph{Second order}



If the ODE is second order,



$$D^2 y + bDy+cy=0$$ then the characteristic equation is a quadratic,



$$p^2+bp+c=0$$ with roots



$$p\_{\pm}=\frac{-b \pm \sqrt{b^2-4c}}{2}$$



What these roots are like depends on the sign of \textit{$b^2-4c$}, so we have

three cases to consider.



\textit{1) $b^2 > 4c$}



In this case we have two different real roots, so we can write down the

solution straight away.



$$y=A\_{+}e^{p\_{+}}+A\_{-}e^{p\_{-}}$$



\textit{2) $b^2 < 4c$}



In this case, both roots are imaginary. We could just put them directly

in the formula, but if we are interested in real solutions it is more

useful to write them another way.



Defining \textit{$k^2=4c-b^2$}, then the solution is



$$y=A\_{+}e^{ikx-\frac{bx}{2}}+A\_{-}e^{-ikx-\frac{bx}{2}}$$



For this to be real, the \textit{A}'s must be complex conjugates



$$A\_{\pm}=A e^{\pm ia}$$



Make this substitution and we can write,



$$y=A e^{-bx/2}\cos (kx+a)$$



If \textit{b} is positive, this is a damped oscillation.



\textit{3) $b^2 = 4c$}



In this case the characteristic equation only gives us one root,

\textit{p=-b/2}. We must use another method to find the other solution.



We'll use the method of variation of constants. The ODE we need to

solve is,



$$D^2 y -2pDy+p^2y=0$$



rewriting \textit{b} and \textit{c} in terms of the root. From

the characteristic equation we know one solution is $y=e^{px}$ so we

make the substitution $y=ze^{px},$ giving



$$(e^{px}D^2z+2pe^{px}Dz+p^2e^{px}z)-2p(e^{px}Dz+pe^{px}z)+p^2e^{px}z=0$$



This simplifies to *$D^2z=0*$, which is easily solved. We get



$$z=Ax+B \quad y=(Ax+B)e^{px}$$ 



so the second solution is the first multiplied by \textit{x}.



Higher order linear constant coefficient ODE's behave similarly: an

exponential for every real root of the characteristic and a exponent

multiplied by a trig factor for every complex conjugate pair, both being

multiplied by a polynomial if the root is repeated.



E.g., if the characteristic equation factors to



$$(p-1)^4(p-3)(p^2+1)^2=0$$



the general solution of the ODE will be



$$y=(A+Bx+Cx^2+Dx^3)e^x + Ee^{3x}+ F \cos (x+a) +Gx \cos(x+b)$$



The most difficult part is finding the roots of the characteristic

equation.



\paragraph{Linear nonhomogeneous ODEs with constant coefficients}



First, let's consider the ODE



$$Dy-y=x$$ 



a nonhomogeneous first order ODE which we know how to solve.



Using the integrating factor \textit{$e^{-x}$} we find



$$y=c e^{-x} +1 -x$$



This is the sum of a solution of the corresponding homogeneous equation,

and a polynomial.



Nonhomogeneous ODE's of higher order behave similarly.



If we have a single solution, \textit{$y\_p$} of the nonhomogeneous ODE, called a

\textit{particular} solution,



$$(D^n+a\_1 D^{n-1} + \cdots + a\_n)y=F(x)$$



then the general solution is \textit{$y=y\_p+y\_h$}, where \textit{$y\_h$} is the general solution of the homogeneous ODE.



Find \textit{$y\_p$} for an arbitrary \textit{F(x)} requires methods beyond the scope of

this chapter, but there are some special cases where finding \textit{$y\_p$} is straightforward.



Remember that in the first order problem \textit{$y\_p$} for a polynomial \textit{F(x)}

was itself a polynomial of the same order. We can extend this to higher orders.



\textit{Example:}



$$D^2y+y=x^3-x+1$$



Consider a particular solution



$$y\_p=b\_0+b\_1 x+b\_2 x^2 + x^3$$



Substitute for \textit{y} and collect coefficients



$$x^3 + b\_2 x^2 +(6+b\_1)x +(2b\_2+b\_0)=x^3-x+1$$ 



So \textit{$b\_2=0$}, \textit{$b\_1=-7$}, \textit{$b\_0=1$}, and the general solution is



$$y=a \sin x + b \cos x + 1 -7x + x^3$$



This works because all the derivatives of a polynomial are themselves polynomials.



Two other special cases are



$$F(x)=P\_n e^{kx} \quad y\_p(x)=Q\_n e^{kx}$$



$$F(x)=A\_n \sin kx +B\_n \cos kx \quad 

y\_p(x)=P\_n \sin kx +Q\_n \cos kx$$



where \textit{$P\_n$},\textit{$Q\_n$},\textit{$A\_n$}, and \textit{$B\_n$} are all polynomials of degree \textit{n}.



Making these substitutions will give a set of simultaneous linear

equations for the coefficients of the polynomials.



\paragraph{Licensing}



Content obtained and/or adapted from:



\begin{itemize}
    \item \href{https://en.wikibooks.org/wiki/Calculus/Ordinary\_differential\_equations}{Ordinary differential equations,
\end{itemize}
    Wikibooks}

    under a CC BY-SA license

\subsection{Learning Objectives}

\begin{enumerate}
    \item \textbf{Identify and Interpret Differential Equations:} Students will be able to identify the order and type of a given differential equation by examining its highest derivative and functional form.
    \item \textbf{Solve First-Order Differential Equations:} Students will be able to solve separable, homogeneous, linear, and exact first-order differential equations using appropriate techniques, such as separation of variables, substitutions, and integrating factors.
    \item \textbf{Apply Initial Conditions to Differential Equations:} Students will be able to apply initial conditions to obtain specific solutions from general solutions of ordinary differential equations, including both first-order and second-order cases.
\end{enumerate}

\subsection{Assessment Instrument}

\begin{enumerate}
    \item \textbf{Part 1: Identify and Interpret Differential Equations}
    \item \textbf{Multiple Choice:} Given the differential equation $$\frac{d^2y}{dx^2} + 3\frac{dy}{dx} + 2y = 0$$, identify the order and type of the differential equation.
    \item A. First-order linear
    \item B. Second-order linear
    \item C. First-order nonlinear
    \item D. Second-order nonlinear
    \item \textbf{True/False:} The differential equation $$y'' + y = 0$$ is an example of a second-order nonlinear differential equation. (True/False) Explain why or why not.
    \item \textbf{Short Answer:} Determine the order and type of the differential equation $$\frac{d^3y}{dx^3} - \frac{dy}{dx} = e^x$$. Describe the characteristics that led to your classification. \textbf{Part 2: Solve First-Order Differential Equations}
    \item \textbf{Calculation:} Solve the separable differential equation $$\frac{dy}{dx} = y \sin(x)$$. Provide the general solution, including the integration process.
    \item \textbf{Multiple Choice:} To solve the first-order linear differential equation $$\frac{dy}{dx} + 2y = e^x$$, which of the following methods is appropriate?
    \item A. Separation of variables
    \item B. Integrating factor
    \item C. Homogeneous equation method
    \item D. Substitution method
    \item \textbf{Short Answer:} Solve the exact differential equation $$\frac{\partial M}{\partial y} = \frac{\partial N}{\partial x}$$ where $$M(x, y) = x^2 + 3xy$$ and $$N(x, y) = 3x^2 + y$$. Explain the steps to find the general solution. \textbf{Part 3: Apply Initial Conditions to Differential Equations}
    \item \textbf{Calculation:} Given the general solution to the differential equation $$\frac{dy}{dx} = y^2$$ is $$y = \frac{1}{C - x}$$, find the specific solution that satisfies the initial condition $y(0) = 2$.
    \item \textbf{Multiple Choice:} If the general solution to a second-order differential equation is $$y(x) = C_1 e^{2x} + C_2 e^{-x}$$, what is the particular solution if the initial conditions are $y(0) = 1$ and $y'(0) = 0$?
    \item A. $$y(x) = e^{2x} - e^{-x}$$
    \item B. $$y(x) = 1 + e^{2x}$$
    \item C. $$y(x) = e^{2x} - 1$$
    \item D. $$y(x) = 2e^{2x} - e^{-x}$$
    \item \textbf{Short Answer:} Apply the initial condition $y(1) = 3$ to the general solution of the differential equation $$y'' - 4y = 0$$. First, find the general solution and then determine the specific solution that satisfies the initial condition.
\end{enumerate}
